% Generated by build_chapter.py from the GY1973 files. Do not edit by hand; edit the files and rebuild.
\chapter{The academic year 1973--74}\label{ch:1973}
\noindent\textbf{Scope and editorial basis.}
This chapter accounts for all 53 supplied scan pages, P000501--P000553, and
32 records: the 30 supplied catalogue entries and two meetings recovered
from the scans. Records 331 and 332 belong to 25 March 1974, although the
archive index dates them to 1975. They are retained separately, with the date
repair documented. The calendar divider 324 has no mathematical item; the
opening plan 302 includes Serre's substantive questions. Source order is
preserved even when record numbers run backwards. The Russian transcription
is checked against the scan where a formula or doubtful reading matters.
Catalogue page ranges and actual item locations are distinguished.

\medskip\noindent\textbf{Main mathematical lines.}
Manin's three lectures develop non-Archimedean integration, Mellin transforms,
periods of cusp forms and ordinary \(p\)-adic interpolation
\citep{manin1973periods}. The relevant
results largely precede these September meetings; their earlier seminar
occurrences remain visible in the problem index. Fuchs's extended exposition
connects nerves, homotopy fibres, Quillen's theorems A and B, exact categories
and the \(Q\)-construction \citep{quillen1973higher}. Several missing loops and degree shifts materially
change the notebook's literal statements. Kirillov's discussions connect
operator extensions, symplectic spinors and \(p\)-adic measures; Mishchenko's
Fredholm representations return to higher signatures.

The November talks range from modular lattices, buildings and finite-field
point counts to stable rank and derived categories. They require especially
careful distinctions between representability and classification, arbitrary
fields and algebraic extensions of finite fields, and real and complex
coefficient problems. In spring, periodic Schr\"odinger operators, finite-gap
KdV, dressing methods, modulation and reflectionless interpolation form a
connected line. The final lectures on wavefront sets and spectral asymptotics
make the analytic hypotheses and spectral normalization indispensable.

\medskip\noindent\textbf{Historical outcomes.}
The occurrence index contains 301 items. The classes describe historical
fate, while correctness and verification are assessed independently.
\begin{center}
\begin{tabular}{@{}lrrrrrr@{}}
\toprule
Class & A & B & C & D & E & F\\
\midrule
Items & 25 & 7 & 4 & 0 & 229 & 36\\
\bottomrule
\end{tabular}
\end{center}
\noindent
A denotes near-contemporary publication by the speaker or Gelfand; B, later
development by others; C, a precise unresolved remainder; D, an undeveloped
identifiable question; E, pre-existing results; F, insufficient identification.
These counts include definitions and methods, not just theorems. They are
not a measure of originality or a claim that every corrected reconstruction
was spoken in full.

Exact publication matches include Rudakov's representation theorems,
Gelfand--Ponomarev's lattice constructions, Novikov's periodic KdV paper and
Zakharov--Shabat's dressing scheme
\citep{rudakov1974irreducible,gelfandponomarev1974modular,novikov1974periodic,zakharovshabat1974scheme}.
A 1973 public announcement places
Mishchenko's relevant results before the October talks, despite their fuller
1974 publication \citep{mishchenko1973infinite}. Shabat's 2014 paper explicitly recalls the 1974 seminar
interpolation problem \citep{shabat2014interpolation}. The late spectral-gap discussion is identified with
Shnirelman's 1975 theorem using the minimum of two adjacent gaps for the
square-root Laplacian \citep{shnirelman1975multiplicity}. Publication, manuscript receipt and first checked
public circulation are recorded separately in the entries.

\medskip\noindent\textbf{Questions and critical verification.}
There are 17 definite problem-index entries. The four class-C
occurrences concern coherence of \(\mathrm{SL}_3(\mathbb Z)\), unconditional
coherence of \(\mathrm{SL}_2(\mathbb Z[1/p])\), transcendence of \(\zeta(3)\),
and the arithmetic nature of Euler's constant. Their status is dated
9 November 2026 and limited to the checked literature. A conditional
coherence theorem is not counted as an unconditional solution
\citep[Theorem~4.2, version~3]{brody2025commensurator}. No class-D
item is inferred from an unsuccessful search.

No C/D item justified a bounded Stage-2 research contribution in this run.
The verification appendix instead supplies elementary proofs needed to
assess the signed measure's Bernoulli moments, tag dependence, categorical
constructions and the complex conic's polynomial reduction. These are explicitly marked as proofs
supplied here, with no historical novelty claim. Independent exact-arithmetic
checks are included. The 93 correction records distinguish mistranslation,
transcription, ambiguity, provenance and mathematical errors faithfully
preserved in the notebook.

\medskip\noindent\textbf{Connections across years.}
The ordinary interpolation question reuses GP-1972-403 and its 1972--73
occurrences. Four-subspace classification is linked to 1970-198 and
1972-277 without identifying it with the separate rigidity problem.
The categorical lectures continue the plus-construction and nerve discussions
of 1971-250, 1971-255 and 1971-257; the derived-category discussion connects to
1971-266 and 1972-281. These are thematic links unless the supplied earlier
index supports identity of the quantified problem.

\GStalk{302}{1973-09-03}{I. M. Gelfand}{P000501-P000502}{Plan and Serre's finiteness and coherence questions}
A plan with a substantive Serre interlude. The thirteen agenda headings are: derived categories (Kirillov, Bernstein); several-variable Ahlfors--Nevanlinna theory and Grauert (Henkin, Shabat); Quillen's simplicial \(K\)-theory (Fuchs); Lichtenbaum's special-value programme (Manin); an uncertain Weil/Gersten-name entry; Smale's optimization of several functions; Brown--Douglas--Fillmore \(K\)-theory (Kirillov); Mishchenko's representations; Lusztig on higher signatures; an Atiyah--Bott--Patodi spectral-asymmetry entry (Shubin, names retained as written); Donovan--Freilich on four subspaces (Ponomarev); Manin; and Vershik--Gelfand--Graev's six variations on \(\mathrm{SL}_2\). Bare agenda labels remain in this summary instead of becoming invented theorems. Derived categories connect thematically to 1971-266.3--6 and 1972-281.4. Quillen's subject was already circulating before this meeting \citep{quillen1973higher}.

The interlude separates finite generation from finite presentation for two function-field arithmetic families, then defines coherence and asks arithmetic examples. This material ends on P000501 at the Manin heading. The supplied catalogue range P000501--P000502 overlaps Manin's lecture and is retained only as catalogue metadata.

\begin{GSitem}{1973-302.1}{F}{Lichtenbaum's special-value programme in the agenda}
\GSfield{Kind; source}{conjecture; \GSpage{P000501}}
\GSfield{The notes}{Plan item 4 announces \ru{Гипотеза Лихтенбаума} connecting \(\zeta\)-values at integers with higher \(K\)-theories; Manin is the prospective expositor.}
\GSfield{Statement}{No field, scheme, integer range, leading-term convention, regulator or torsion factor is specified. A programme relating special values and algebraic \(K\)-groups is identifiable, but a unique quantified conjecture is not.}
\GSfield{Fate}{\citet{lichtenbaum1973values} is a contemporary thematic match: Seattle conference 1972, publication 1973. The notes do not identify a particular conjecture of that paper.}
\GSfield{Present status}{As of 9 November 2026: The agenda sentence has no single solved/open truth value and is not a definite problem.}
\GSfield{Assessment}{literal statement: not determinable; publication match: thematic only; verification: publisher metadata checked; development: none; confidence: medium}
\end{GSitem}
\begin{GSitem}{1973-302.2}{F}{Unresolved name in the planned Weil discussion}
\GSfield{Kind; source}{remark; \GSpage{P000501}}
\GSfield{The notes}{Plan item 5 reads \ru{Вей: идеология Герстенцуса}; the surname is already marked doubtful.}
\GSfield{Statement}{No proposition can be recovered from this uncertain name and the word \textit{ideology}.}
\GSfield{Fate}{The scan was compared with both transcriptions. Searches for the supplied Cyrillic form and plausible transliterations did not establish a unique work. No identification with Gersten or Gerstenhaber is asserted.}
\GSfield{Present status}{As of 9 November 2026: Unidentified agenda label, not evidence of an undeveloped idea.}
\GSfield{Assessment}{literal statement: not determinable; publication match: none found; verification: unresolved lead; development: none; confidence: low}
\end{GSitem}
\begin{GSitem}{1973-302.3}{E}{Finite generation of the two function-field families}
\GSfield{Kind; source}{claim; \GSpage{P000501}}
\GSfield{The notes}{For finite \(k\), the notes assert finite generation of \(\mathrm{SL}_n(k[t])\), \(n\geq3\), and \(\mathrm{SL}_n(k[t,t^{-1}])\), \(n\geq2\).}
\GSfield{Statement}{For every finite field \(\mathbb F_q\), both families are finitely generated in the displayed ranges.}
\GSfield{Fate}{The pre-1973 result is credited to Behr 1969 and Harder's reduction theory in \citet[p.~312]{buxkohlwitzel2013rank}.}
\GSfield{Present status}{As of 9 November 2026: Established. Finite generation alone does not imply finite presentation.}
\GSfield{Assessment}{literal statement: correct; publication match: exact; verification: read the result; development: none; confidence: high}
\end{GSitem}
\begin{GSitem}{1973-302.4}{B}{Serre's finite-presentation question over a polynomial ring}
\GSfield{Kind; source}{question; \GSpage{P000501}}
\GSfield{The notes}{Are \(\mathrm{SL}_n(k[t])\), \(k\) finite and \(n\geq3\), finitely presented?}
\GSfield{Statement}{For every prime power \(q\) and integer \(n\geq3\), determine finite presentability of \(\mathrm{SL}_n(\mathbb F_q[t])\).}
\GSfield{Fate}{No at \(n=3\); yes at \(n\geq4\). The Rank Theorem gives finiteness length \(n-2\) \citep[p.~312]{buxkohlwitzel2013rank}. It credits the complete finite-presentation criterion to Behr 1998. General-rank preprint 2011, final 2013; these are not asserted to be first dates for each special case.}
\GSfield{Present status}{As of 9 November 2026: Completely answered, including small finite fields.}
\GSfield{Assessment}{literal statement: correct; publication match: exact; verification: read the result; development: none; confidence: high; problem \hyperref[prob:GP-1973-01]{GP-1973-01}}
\end{GSitem}
\begin{GSitem}{1973-302.5}{B}{Serre's finite-presentation question over a Laurent ring}
\GSfield{Kind; source}{question; \GSpage{P000501}}
\GSfield{The notes}{The same question is posed for \(\mathrm{SL}_n(k[t,t^{-1}])\), \(k\) finite, \(n\geq2\).}
\GSfield{Statement}{For every finite field \(\mathbb F_q\) and integer \(n\geq2\), determine finite presentability of \(\mathrm{SL}_n(\mathbb F_q[t,t^{-1}])\).}
\GSfield{Fate}{The rank sum is \(2(n-1)\), and finiteness length \(2n-3\): no at \(n=2\), yes at \(n\geq3\) \citep[p.~312]{buxkohlwitzel2013rank}. The complete finite-presentation criterion was already established by Behr 1998.}
\GSfield{Present status}{As of 9 November 2026: Completely answered. Replacing \(\mathbb F_q\) by \(\mathbb Z\) changes the problem.}
\GSfield{Assessment}{literal statement: correct; publication match: exact; verification: read the result; development: none; confidence: high; problem \hyperref[prob:GP-1973-02]{GP-1973-02}}
\end{GSitem}
\begin{GSitem}{1973-302.6}{E}{Coherence of a group}
\GSfield{Kind; source}{definition; \GSpage{P000501}}
\GSfield{The notes}{\ru{Когер. группа}: every subgroup with finitely many generators is finitely presented.}
\GSfield{Statement}{A group \(G\), not necessarily finitely generated, is coherent when every finitely generated subgroup \(H\leq G\) is finitely presented.}
\GSfield{Fate}{This is the group-theoretic property in Serre's questions \citep[pp.~734--735]{serre1974problems}. Requiring the ambient group to be finitely presented would improperly exclude the question about \(\mathrm{GL}_n(\mathbb Q)\).}
\GSfield{Present status}{As of 9 November 2026: Standard terminology; unrelated to sheaf or ring coherence.}
\GSfield{Assessment}{literal statement: correct; publication match: exact; verification: indirect; development: none; confidence: high}
\end{GSitem}
\begin{GSitem}{1973-302.7}{E}{Free groups are coherent}
\GSfield{Kind; source}{claim; \GSpage{P000501}}
\GSfield{The notes}{The notes say that free groups are coherent.}
\GSfield{Statement}{Every subgroup of a free group is free. A finitely generated such subgroup has finite rank and a finite presentation with no relators.}
\GSfield{Fate}{This is the classical Nielsen--Schreier implication; free coherent groups are recalled in \citet[section~4]{brody2025commensurator}.}
\GSfield{Present status}{As of 9 November 2026: Established, even if the ambient free group has infinite rank.}
\GSfield{Assessment}{literal statement: correct; publication match: exact; verification: read the result; development: none; confidence: high}
\end{GSitem}
\begin{GSitem}{1973-302.8}{E}{Nilpotent groups are coherent}
\GSfield{Kind; source}{claim; \GSpage{P000501}}
\GSfield{The notes}{The faint abbreviation following free is read as \ru{нильпот.}.}
\GSfield{Statement}{Every finitely generated subgroup of a nilpotent group is finitely generated nilpotent, hence polycyclic and finitely presented.}
\GSfield{Fate}{This is the classical polycyclic implication; see the coherent classes recalled in \citet[section~4]{brody2025commensurator}. The scan supports the transmitted abbreviation.}
\GSfield{Present status}{As of 9 November 2026: Established; torsion-freeness and finite generation of the ambient group are unnecessary.}
\GSfield{Assessment}{literal statement: correct; publication match: exact; verification: read the result; development: none; confidence: high}
\end{GSitem}
\begin{GSitem}{1973-302.9}{E}{Scott's coherence theorem for three-manifold groups}
\GSfield{Kind; source}{theorem; \GSpage{P000501}}
\GSfield{The notes}{The notebook explicitly cites Scott, JLMS(2)6(1973),437--440, for coherence of three-manifold fundamental groups.}
\GSfield{Statement}{For a connected second-countable three-manifold \(M\), with or without boundary, every finitely generated subgroup of \(\pi_1(M)\) is finitely presented: apply Scott's theorem to the corresponding connected covering.}
\GSfield{Fate}{\citet[section~2]{scott1973finitely} is the exact work, published May 1973, before this meeting. Its original introductory statement was read in the author's collected works.}
\GSfield{Present status}{As of 9 November 2026: Established. The covering-space argument supplies the subgroup quantifier.}
\GSfield{Assessment}{literal statement: correct; publication match: exact; verification: read the result; development: none; confidence: high}
\end{GSitem}
\begin{GSitem}{1973-302.10}{E}{The imaginary-quadratic arithmetic example}
\GSfield{Kind; source}{claim; \GSpage{P000501}}
\GSfield{The notes}{\(\mathrm{SL}_2(\Lambda)\) is called coherent, with \(\Lambda\) described as \ru{чисто мнимые элем. квадр. поля}.}
\GSfield{Statement}{For an imaginary quadratic number field \(K\), \(\mathrm{SL}_2(\mathcal O_K)\) is coherent, where \(\mathcal O_K\) is its ring of integers.}
\GSfield{Reconstruction}{Replacing the transmitted description by the ring of integers is a mathematical reconstruction, not a new reading of the ink.}
\GSfield{Fate}{This is the expected Bianchi-group consequence of \citet{scott1973finitely}: use a torsion-free finite-index subgroup and stability of coherence under finite extensions.}
\GSfield{Present status}{As of 9 November 2026: The reconstruction is established. Purely imaginary elements themselves are not a unital coefficient ring.}
\GSfield{Assessment}{literal statement: correct with missing hypotheses; publication match: partial; verification: indirect; development: none; confidence: medium}
\end{GSitem}
\begin{GSitem}{1973-302.11}{C}{Is SL3(Z) coherent?}
\GSfield{Kind; source}{question; \GSpage{P000501}}
\GSfield{The notes}{The notes explicitly ask whether \(\mathrm{SL}_3(\mathbb Z)\) is coherent.}
\GSfield{Statement}{Is every finitely generated subgroup \(H\leq\mathrm{SL}_3(\mathbb Z)\) finitely presented?}
\GSfield{Fate}{The question remains explicit and open in \citet[p.~5]{audibert2025zariski}. Incoherence in ranks \(n\geq4\) does not answer it.}
\GSfield{Present status}{As of 9 November 2026: No unconditional resolution was found in the dated search; the universal subgroup quantifier remains.}
\GSfield{Assessment}{literal statement: correct; publication match: exact; verification: read the result; development: none; confidence: high; problem \hyperref[prob:GP-1973-03]{GP-1973-03}}
\GSfield{Worth developing}{Focused work on a specified subgroup family is worthwhile, but this run found no bounded route to full coherence.}
\end{GSitem}
\begin{GSitem}{1973-302.12}{C}{Is SL2(Z[1/p]) coherent?}
\GSfield{Kind; source}{question; \GSpage{P000501}}
\GSfield{The notes}{The same line asks about \(\mathrm{SL}_2(\mathbb Z[1/p])\).}
\GSfield{Statement}{For every prime \(p\), is every finitely generated subgroup of \(\mathrm{SL}_2(\mathbb Z[1/p])\) finitely presented?}
\GSfield{Reconstruction}{The conventional assumption that \(p\) is prime is made explicit.}
\GSfield{Fate}{\citet[Theorem~4.2, version~3]{brody2025commensurator} proves a conditional positive result from Greenberg--Shalom. First preprint 2023; checked revision 7 November 2025. The hypothesis is essential to the attribution.}
\GSfield{Present status}{As of 9 November 2026: Open unconditionally in the checked literature. Inverting two distinct primes is a different, negatively answered problem.}
\GSfield{Assessment}{literal statement: correct with missing hypotheses; publication match: exact; verification: read the result; development: none; confidence: high; problem \hyperref[prob:GP-1973-04]{GP-1973-04}}
\GSfield{Worth developing}{A specified subgroup family may admit a bounded study; the required arithmeticity hypothesis and full coherence exceed the contribution justified here.}
\end{GSitem}
\begin{GSitem}{1973-302.13}{B}{Coherence of GLn(Q)}
\GSfield{Kind; source}{question; \GSpage{P000501}}
\GSfield{The notes}{The line ends \(\mathrm{GL}_n(\mathbb Q)?\), with no rank restriction.}
\GSfield{Statement}{For which integers \(n\geq1\) is the abstract discrete group \(\mathrm{GL}_n(\mathbb Q)\) coherent?}
\GSfield{Reconstruction}{The rank 1 exception is explicit. The thesis's displayed Baumslag--Solitar label is not copied: conjugation by \(\operatorname{diag}(p/q,q/p)\) scales a unipotent entry by \((p/q)^2\).}
\GSfield{Fate}{Yes at \(n=1\), no at \(n\geq2\). Baumslag's rational triangular obstruction is recorded with Serre's questions and discussed in \citet[p.~74]{brody2022trees}; block inclusion handles larger ranks.}
\GSfield{Present status}{As of 9 November 2026: Answered. Finitely generated subgroups of \(\mathbb Q^\times\) are finitely generated abelian. The printed resolution is traced to 1974; the exact date of Baumslag's intervening observation is unverified.}
\GSfield{Assessment}{literal statement: correct; publication match: exact; verification: indirect; development: none; confidence: medium; problem \hyperref[prob:GP-1973-05]{GP-1973-05}}
\end{GSitem}
\begin{GSitem}{1973-302.14}{F}{Bracketed relations among irreducibles}
\GSfield{Kind; source}{question; \GSpage{P000501}}
\GSfield{The notes}{The bracket contains the doubtful word \ru{неприводимыми}. It asks for finitely many relations among different such objects.}
\GSfield{Statement}{The objects, operation and quantifier over relations are unspecified; a precise question cannot be recovered.}
\GSfield{Fate}{Group-presentation and representation-theoretic readings remain possible after scan comparison. Neither warrants a publication match.}
\GSfield{Present status}{As of 9 November 2026: Unresolved textual question; excluded from the definite problem index.}
\GSfield{Assessment}{literal statement: not determinable; publication match: none found; verification: unresolved lead; development: none; confidence: low}
\end{GSitem}
\begin{GSitem}{1973-302.15}{F}{Which fundamental groups occur?}
\GSfield{Kind; source}{question; \GSpage{P000501}}
\GSfield{The notes}{The second bracket reads \ru{Какие бывают фунд. группы?}}
\GSfield{Statement}{No class of spaces, dimension, compactness condition or equivalence relation is supplied; these choices materially change the question.}
\GSfield{Fate}{The preceding three-manifold discussion is possible context, but does not determine the missing hypotheses.}
\GSfield{Present status}{As of 9 November 2026: Not identifiable as a definite problem from the surviving wording.}
\GSfield{Assessment}{literal statement: not determinable; publication match: none found; verification: unresolved lead; development: none; confidence: low}
\end{GSitem}

\GStalk{303}{1973-09-03}{Yu. I. Manin}{P000501-P000506}{p-adic measures: spaces, distributions and integrals}
The actual mathematical text runs from the Manin heading on P000501 to the next dated heading on P000505; the catalogue range includes neighbouring evidence. Manin motivates non-Archimedean integration by special zeta values, recalls Serre\textquoteright s classification of compact analytic manifolds, and compares invariant and bounded distributions. The linear-function counterexample explains why finitely additive distributions need a growth condition. Several constructions repeat records 298--299 of the preceding year. The precise moderate-growth theorem is \citet[Theorem~8.6]{manin1973periods}; gaps in the notebook are distinguished from open questions.

\begin{GSitem}{1973-303.1}{E}{Finding an integration theory for compact p-adic spaces}
\GSfield{Kind; source}{problem; \GSpage{P000501}}
\GSfield{The notes}{The heading asks for a reasonable class of \(p\)-adic measures on a compact analytic \(X/K\), with integrals of analytic maps \(X\to L\).}
\GSfield{Statement}{For a finite extension \(L/K/\mathbb Q_p\), a compact zero-dimensional Hausdorff space \(X\), and a bounded finitely additive \(L\)-valued function on clopen subsets, integration extends uniquely from locally constant functions to a continuous \(L\)-linear functional on \(C(X,L)\). Locally analytic integrands also admit larger distribution classes after imposing growth bounds.}
\GSfield{Reconstruction}{Compactness, completeness and boundedness are explicit sufficient hypotheses supplied here; the notebook asks a programme-level question, not uniqueness of an optimal measure class.}
\GSfield{Fate}{\citet[\S8.1 and Theorem~8.6]{manin1973periods}; the concrete answers, rather than an unspecified universal class, are the subject of this exposition.}
\GSfield{Present status}{As of 9 November 2026: established in the stated scope.}
\GSfield{Assessment}{literal statement: correct with missing hypotheses; publication match: exact; verification: read the result; development: none; confidence: high}
\end{GSitem}
\begin{GSitem}{1973-303.2}{F}{Adelic Mellin motivation with an unspecified quotient}
\GSfield{Kind; source}{remark; \GSpage{P000502}}
\GSfield{The notes}{The notes schematically write \(L(\chi)=\int_{G_A/G_K}\chi\,d\mu\) and identify \(G_K\) with principal adeles.}
\GSfield{Statement}{The record does not specify \(G\), the measure, convergence, or the automorphic test function; no single integral formula is asserted here beyond its role as motivation.}
\GSfield{Fate}{\citet[Introduction and \S1]{manin1976nonarchimedean} supplies a developed adelic framework, but does not identify every missing symbol in this line.}
\GSfield{Present status}{As of 9 November 2026: the intended adelic analogy is clear, but this formula is not an independently identifiable theorem.}
\GSfield{Assessment}{literal statement: not determinable; publication match: thematic only; verification: indirect; development: none; confidence: low}
\end{GSitem}
\begin{GSitem}{1973-303.3}{E}{Bernoulli values of the classical zeta function}
\GSfield{Kind; source}{theorem; \GSpage{P000502}}
\GSfield{The notes}{The displays compare \(\zeta(2m)\) and \(\zeta(1-2m)\) with Bernoulli numbers, with a question mark at the latter sign.}
\GSfield{Statement}{With \(t/(e^t-1)=\sum B_n t^n/n!\), for \(m\ge1\), \(\zeta(1-2m)=-B_{2m}/(2m)\) and \(\zeta(2m)=(-1)^{m+1}B_{2m}(2\pi)^{2m}/(2(2m)!)\).}
\GSfield{Reconstruction}{The Bernoulli convention and omitted constants are supplied; divergent power sums are not ordinary convergent sums.}
\GSfield{Fate}{Classical formulas, restated in \citet[\S1 and \S5]{nica2024irrationality}.}
\GSfield{Present status}{As of 9 November 2026: established in the stated scope.}
\GSfield{Assessment}{literal statement: correct with missing hypotheses; publication match: exact; verification: read the result; development: none; confidence: high}
\end{GSitem}
\begin{GSitem}{1973-303.4}{F}{Gelfand on divergent sums and spectral parity}
\GSfield{Kind; source}{remark; \GSpage{P000502}; I. M. Gelfand}
\GSfield{The notes}{The parenthetical intervention juxtaposes \(\sum n^3\), Bernoulli numbers and sums of powers of eigenvalues of \(-y^{\prime\prime}+p(x)y=\lambda y\), saying that integers and half-integers differ.}
\GSfield{Statement}{As an ordinary series \(\sum_{n\ge1}n^3\) diverges. Zeta regularization assigns \(\zeta(-3)=1/120\), but a spectral analogue requires an operator domain, boundary conditions and a specified regularization, none of which are recorded.}
\GSfield{Fate}{The Bernoulli identity is classical; \citet[\S1 and \S5]{nica2024irrationality} does not settle the unformulated spectral comparison.}
\GSfield{Present status}{As of 9 November 2026: the spectral-parity remark remains insufficiently specified for a problem-index entry.}
\GSfield{Assessment}{literal statement: not determinable; publication match: thematic only; verification: indirect; development: none; confidence: low}
\end{GSitem}
\begin{GSitem}{1973-303.5}{E}{Removing the Euler factor and interpolating zeta values}
\GSfield{Kind; source}{theorem; \GSpage{P000502}}
\GSfield{The notes}{The values \((1-p^{2m-1})B_{2m}/(2m)\) are called \(p\)-adically continuous and analytic.}
\GSfield{Statement}{For odd \(p\), interpolation of \(-(1-p^{n-1})B_n/n\) is performed on a fixed residue class of \(n\bmod(p-1)\), with the standard Teichmuller twist. The zeta branch is meromorphic with its pole at \(s=1\); there is no single untwisted analytic function on all of \(\mathbb Z_p\) realizing every even weight without a branch convention.}
\GSfield{Reconstruction}{The sign, branch restriction and pole are indispensable editorial qualifications.}
\GSfield{Fate}{\citet[\S0.1--0.3]{manin1976nonarchimedean} restates the Kubota--Leopoldt construction (1964).}
\GSfield{Present status}{As of 9 November 2026: established in the stated scope.}
\GSfield{Assessment}{literal statement: correct with missing hypotheses; publication match: exact; verification: read the result; development: none; confidence: high}
\end{GSitem}
\begin{GSitem}{1973-303.6}{F}{The proposed cross-prime analytic analogy}
\GSfield{Kind; source}{remark; \GSpage{P000502}}
\GSfield{The notes}{The notes suggest that \(\mathbb Q_p\to\mathbb Q_\ell\), \(p\ne\ell\), may be easier than \(\mathbb Q_p\to\mathbb Q_p\), but call it unknown.}
\GSfield{Statement}{Neither the function space nor the regularity condition nor a target theorem is stated. Continuous maps between these totally disconnected fields alone do not specify the intended arithmetic problem.}
\GSfield{Fate}{No exact problem was identified after reading the corresponding integration sections of \citet[\S8]{manin1973periods} and \citet[Introduction]{manin1976nonarchimedean}.}
\GSfield{Present status}{As of 9 November 2026: an exploratory analogy, not a precise unresolved conjecture.}
\GSfield{Assessment}{literal statement: not determinable; publication match: none found; verification: unresolved lead; development: none; confidence: low}
\end{GSitem}
\begin{GSitem}{1973-303.7}{E}{Serre classification of compact non-Archimedean manifolds}
\GSfield{Kind; source}{theorem; \GSpage{P000502}}
\GSfield{The notes}{Serre\textquoteright s theorem is \(X\simeq rA^n\), with \(rA^n\simeq r^{\prime}A^n\iff r\equiv r^{\prime}\pmod{q-1}\).}
\GSfield{Statement}{Let \(K\) be a non-Archimedean local field with valuation ring \(A\) and residue cardinality \(q\). For fixed \(n\ge1\), every nonempty compact \(K\)-analytic manifold of pure dimension \(n\) is analytically isomorphic to a finite disjoint union \(rA^n\); the isomorphism invariant is \(i(X)=r\bmod(q-1)\).}
\GSfield{Reconstruction}{The English translation says analytic set, but the Russian abbreviation and scan say analytic manifold. Pure positive dimension and nonemptiness are stated explicitly here.}
\GSfield{Fate}{\citet[Theorem~1 and corollary, p.~409]{serre1965classification}; received March 1964, published 1965.}
\GSfield{Present status}{As of 9 November 2026: established in the stated scope.}
\GSfield{Assessment}{literal statement: correct with missing hypotheses; publication match: exact; verification: read the result; development: none; confidence: high}
\end{GSitem}
\begin{GSitem}{1973-303.8}{E}{Computing the Serre invariant by reduction}
\GSfield{Kind; source}{method; \GSpage{P000503}}
\GSfield{The notes}{The notebook asks \ru{Как вычислить} \(r\), then gives \(r(X)=\#X_0(A/\mathfrak m)\) for a smooth scheme over \(A\).}
\GSfield{Statement}{If \(\mathcal X/A\) is smooth, separated and of finite type, with nonempty \(\mathcal X(A)\) of pure relative dimension \(n\ge1\), the inverse image of each residue point in \(\mathcal X(k)\) is an analytic \(n\)-ball. Thus \(i(\mathcal X(A))\equiv\#\mathcal X(k)\pmod{q-1}\).}
\GSfield{Reconstruction}{The integer number of residue discs is one representative of the invariant, not a canonical integer-valued invariant of the analytic manifold.}
\GSfield{Fate}{\citet[example (i), p.~410]{serre1965classification}.}
\GSfield{Present status}{As of 9 November 2026: established in the stated scope.}
\GSfield{Assessment}{literal statement: correct with missing hypotheses; publication match: exact; verification: read the result; development: none; confidence: high}
\end{GSitem}
\begin{GSitem}{1973-303.9}{E}{The projective line as two analytic balls}
\GSfield{Kind; source}{construction; \GSpage{P000503}}
\GSfield{The notes}{The notes count \(q+1\) residue points and obtain two balls; they write \(\mathbb P^1(K)=\mathbb Q_p\cup\{\infty\}\) and \(pA\cup A^{-1}\).}
\GSfield{Statement}{For any non-Archimedean local field \(K\), \(\mathbb P^1(K)=K\cup\{\infty\}\) is the disjoint union of \(\mathfrak m\) and \(\{x:|x|\ge1\}\cup\{\infty\}\); inversion identifies the latter with \(A\). Therefore it is two analytic balls and \(i=2\bmod(q-1)\).}
\GSfield{Reconstruction}{Replace Qp by K and pA by the maximal ideal for ramified extensions; the omitted point at infinity is included in the inverse-coordinate chart.}
\GSfield{Fate}{An explicit instance of \citet[Theorem~1]{serre1965classification}.}
\GSfield{Present status}{As of 9 November 2026: established in the stated scope.}
\GSfield{Assessment}{literal statement: false as written; publication match: exact; verification: read the result; development: none; confidence: high}
\end{GSitem}
\begin{GSitem}{1973-303.10}{E}{Compact Lie groups and supernatural order}
\GSfield{Kind; source}{claim; \GSpage{P000503}}
\GSfield{The notes}{For a compact \(p\)-adic Lie group the supernatural order is \(p^\infty N\), \((p,N)=1\), followed by \(N=r(X)\) and an \(SL_2\) example.}
\GSfield{Statement}{A positive-dimensional compact \(\mathbb Q_p\)-analytic group has an open pro-\(p\) subgroup and supernatural order \(p^\infty N\) with finite prime-to-\(p\) part \(N\). The analytic invariant is \(i(G)\equiv N\pmod{p-1}\). For \(SL_2(\mathbb Z_p)\), \(N=p^2-1\), so \(i(G)=0\); its representative among \(1,\ldots,p-1\) is \(p-1\).}
\GSfield{Reconstruction}{The invariant is a residue class. The notebook\textquoteright s displayed modulus p in the SL2 example is incompatible with the preceding theorem.}
\GSfield{Fate}{\citet[example (ii), p.~410]{serre1965classification} gives the compact-group application.}
\GSfield{Present status}{As of 9 November 2026: established in the stated scope.}
\GSfield{Assessment}{literal statement: false as written; publication match: exact; verification: read the result; development: none; confidence: high}
\end{GSitem}
\begin{GSitem}{1973-303.11}{E}{Distributions on clopen subsets}
\GSfield{Kind; source}{definition; \GSpage{P000503}}
\GSfield{The notes}{An \(L\)-measure is an \(L\)-valued finitely additive function on open-and-closed subsets of compact \(X\).}
\GSfield{Statement}{An \(L\)-valued distribution on a compact totally disconnected Hausdorff space assigns \(\mu(\varnothing)=0\) and \(\mu(U\sqcup V)=\mu(U)+\mu(V)\) for disjoint clopens. It is a bounded measure when \(\sup_U|\mu(U)|_L<\infty\).}
\GSfield{Reconstruction}{Modern terminology distinguishes a possibly unbounded distribution from a bounded measure.}
\GSfield{Fate}{\citet[\S8.1]{manin1973periods}.}
\GSfield{Present status}{As of 9 November 2026: established in the stated scope.}
\GSfield{Assessment}{literal statement: correct; publication match: exact; verification: read the result; development: none; confidence: high}
\end{GSitem}
\begin{GSitem}{1973-303.12}{E}{The residue-class tree and compatibility equation}
\GSfield{Kind; source}{construction; \GSpage{P000503}}
\GSfield{The notes}{The tree vertices are clopen sets and \(\mu(a+p^m)=\sum_{b\equiv a\ (p^m)}\mu(b+p^{m+1})\).}
\GSfield{Statement}{For \(X=\mathbb Z_p\), specifying numbers \(\mu_m(a)\in L\) for \(a\in\mathbb Z/p^m\mathbb Z\), \(m\ge0\), subject to \(\mu_m(a)=\sum_{j=0}^{p-1}\mu_{m+1}(a+jp^m)\), determines exactly one clopen distribution.}
\GSfield{Fate}{\citet[\S8.1, equation~(24)]{manin1973periods}.}
\GSfield{Present status}{As of 9 November 2026: established in the stated scope.}
\GSfield{Assessment}{literal statement: correct; publication match: exact; verification: read the result; development: none; confidence: high}
\end{GSitem}
\begin{GSitem}{1973-303.13}{E}{Invariant p-adic distribution is unbounded}
\GSfield{Kind; source}{claim; \GSpage{P000503}}
\GSfield{The notes}{Translation-invariant values are \(c p^{-m}\), whose \(p\)-adic absolute values tend to infinity when \(c\ne0\).}
\GSfield{Statement}{The unique translation-invariant distribution on \(\mathbb Z_p\) of total mass \(c\in L\) is \(\mu(a+p^m\mathbb Z_p)=cp^{-m}\). For \(L/\mathbb Q_p\) and \(c\ne0\), it is unbounded and cannot define a continuous functional on all \(C(\mathbb Z_p,L)\).}
\GSfield{Fate}{\citet[\S8.4]{manin1973periods}; already occurrence 1972-298.2.}
\GSfield{Present status}{As of 9 November 2026: established in the stated scope.}
\GSfield{Assessment}{literal statement: correct; publication match: exact; verification: read the result; development: none; confidence: high}
\end{GSitem}
\begin{GSitem}{1973-303.14}{E}{The bounded two-valued distribution}
\GSfield{Kind; source}{construction; \GSpage{P000503-P000504}}
\GSfield{The notes}{The values are \(+1\) below \(p^m/2\), \(-1\) above, for unit representatives; the next page asks to check additivity and attributes the measure to Mazur.}
\GSfield{Statement}{For odd \(p\), define on all standard residue classes \(0\le a<p^m\) the value \(+1\) if \(2a<p^m\), and \(-1\) if \(2a>p^m\). Child values sum to the parent value, so this is a bounded distribution; restrict to \(\mathbb Z_p^\times\) if desired.}
\GSfield{Reconstruction}{Odd p excludes an unassigned midpoint and makes 2 a unit. P000505 uses the opposite overall sign.}
\GSfield{Fate}{\citet[\S8.4, regularized distributions]{manin1973periods}; this is a signed normalization of the c=2 construction.}
\GSfield{Present status}{As of 9 November 2026: established in the stated scope.}
\GSfield{Assessment}{literal statement: correct with missing hypotheses; publication match: exact; verification: read the result; development: annotated; confidence: high}
\end{GSitem}
\begin{GSitem}{1973-303.15}{E}{Integration by arbitrary tagged clopen partitions}
\GSfield{Kind; source}{definition; \GSpage{P000504}}
\GSfield{The notes}{\(S(f,R)=\sum_i f(a_i)\mu(U_i)\), followed by a limit under refinement.}
\GSfield{Statement}{For a compact \(p\)-adic manifold, take finite disjoint clopen partitions and tags \(a_i\in U_i\); a Riemann integral is the tag-independent limit of these sums as the maximal chart diameter tends to zero, whenever that limit exists.}
\GSfield{Reconstruction}{Disjointness, shrinking mesh and tag independence make the limiting convention explicit.}
\GSfield{Fate}{\citet[\S8.5--8.6]{manin1973periods}.}
\GSfield{Present status}{As of 9 November 2026: established in the stated scope.}
\GSfield{Assessment}{literal statement: correct with missing hypotheses; publication match: exact; verification: read the result; development: none; confidence: high}
\end{GSitem}
\begin{GSitem}{1973-303.16}{E}{The linear function fails to integrate against the invariant distribution}
\GSfield{Kind; source}{claim; \GSpage{P000504}}
\GSfield{The notes}{Canonical sums for \(f(x)=x\), \(\mu(U)=p^{-m}\), tend to \(-1/2\); replacing the representative 0 by \(p^m\) changes the answer.}
\GSfield{Statement}{The canonical level-m sum is \((p^m-1)/2\). Replacing just the tag of \(p^m\mathbb Z_p\) by \(p^m\) adds exactly 1 at every level. The two limits are \(-1/2\) and \(1/2\); hence no tag-independent Riemann integral exists.}
\GSfield{Reconstruction}{The second limiting value is calculated explicitly here. This does not deny the separately defined Volkenborn integral with fixed canonical representatives.}
\GSfield{Fate}{This is the elementary obstruction discussed before \citet[Theorem~8.6]{manin1973periods}.}
\GSfield{Present status}{As of 9 November 2026: established in the stated scope.}
\GSfield{Assessment}{literal statement: correct; publication match: exact; verification: read the result; development: annotated; confidence: high}
\end{GSitem}
\begin{GSitem}{1973-303.17}{E}{Moderate growth suffices for Lipschitz integration}
\GSfield{Kind; source}{theorem; \GSpage{P000504}}
\GSfield{The notes}{If \(|\mu(U)/\mu_{\rm inv}(U)|\to0\) and \(f\) is Lipschitz of exponent 1, the integral exists.}
\GSfield{Statement}{Let \(L/\mathbb Q_p\) be finite and \(\mu\) a distribution on \(\mathbb Z_p\). If \(p^{-m}\max_{a\bmod p^m}|\mu(a+p^m\mathbb Z_p)|_p\to0\), every \(L\)-valued Lipschitz function has a tag-independent integral. Finite disjoint unions of charts give the one-dimensional compact-manifold version.}
\GSfield{Reconstruction}{The convergence is uniform over all residue balls; pointwise shrinking alone is not the stated sufficient uniform estimate.}
\GSfield{Fate}{\citet[Theorem~8.6]{manin1973periods}.}
\GSfield{Present status}{As of 9 November 2026: established in the stated scope.}
\GSfield{Assessment}{literal statement: correct with missing hypotheses; publication match: exact; verification: read the result; development: none; confidence: high}
\end{GSitem}

\GStalk{304}{1973-09-10}{Yu. I. Manin}{P000504-P000510}{Mellin transforms, modular forms and Hecke-period distributions}
The dated text begins on P000505 and ends with the tentative correction on P000508. After repeating the integration theorem, Manin calculates moments, describes analytic character space, and introduces cusp forms and Hecke operators. A period integral over a rational cusp almost satisfies the distribution equation; a quadratic-root correction is attempted. P000509 revisits and repairs it in the next talk. The relevant primary source is \citet[\S\S8--10]{manin1973periods}, with earlier seminar occurrences 1972-298--299.

\begin{GSitem}{1973-304.1}{E}{Repeated Lipschitz integration criterion}
\GSfield{Kind; source}{theorem; \GSpage{P000505}}
\GSfield{The notes}{The previous theorem is restated for a one-dimensional \(p\)-adic space and a finite extension \(K/\mathbb Q_p\).}
\GSfield{Statement}{For clopen distributions in finitely many \(\mathbb Z_p\)-charts, the uniform condition \(p^{-m}\max_U|\mu(U)|\to0\) guarantees convergence of all tagged sums for Lipschitz integrands.}
\GSfield{Reconstruction}{Uniformity and the finite-chart interpretation are supplied as in item 1973-303.17.}
\GSfield{Fate}{\citet[Theorem~8.6]{manin1973periods}.}
\GSfield{Present status}{As of 9 November 2026: established in the stated scope.}
\GSfield{Assessment}{literal statement: correct with missing hypotheses; publication match: exact; verification: read the result; development: none; confidence: high}
\end{GSitem}
\begin{GSitem}{1973-304.2}{E}{Moments of the sign distribution}
\GSfield{Kind; source}{claim; \GSpage{P000505}}
\GSfield{The notes}{The sign convention is now \(-1\) below the half threshold and \(+1\) above. The displayed odd moments use \((-1)^r(1-2^{-r})B_{r-1}/(2r)\).}
\GSfield{Statement}{For odd \(p\) with this sign convention and standard Bernoulli numbers, \[\int_{\mathbb Z_p}x^k\,d\mu=\frac{(4-2^{1-k})B_{k+1}}{k+1},\qquad k\ge0.\] In particular the total mass is \(-1\), the first moment is \(1/4\), and every positive even moment is zero.}
\GSfield{Reconstruction}{The Bernoulli convention is not defined on this page. The formula here follows from finite power sums, not by silently correcting B indices in the source; see the included exact-arithmetic check.}
\GSfield{Fate}{The underlying regularized distribution is in \citet[\S8.4]{manin1973periods}; the exact normalization displayed here is independently reconstructed from the notebook definition.}
\GSfield{Present status}{As of 9 November 2026: established in the stated scope.}
\GSfield{Assessment}{literal statement: false as written; publication match: partial; verification: read the result; development: annotated; confidence: high}
\end{GSitem}
\begin{GSitem}{1973-304.3}{E}{Restriction to units produces the Euler factor}
\GSfield{Kind; source}{theorem; \GSpage{P000505}}
\GSfield{The notes}{\(\int_{\mathbb Z_p^\times}=\int_{\mathbb Z_p}-\int_{p\mathbb Z_p}\), but the displayed exponent changes from \(2r-1\) to \(2m-1\).}
\GSfield{Statement}{The sign distribution satisfies \(\mu(pU)=\mu(U)\). Thus for every integer \(k\ge0\), \(\int_{\mathbb Z_p^\times}x^k\,d\mu=(1-p^k)\int_{\mathbb Z_p}x^k\,d\mu\). At \(k=2r-1\), the factor is \(1-p^{2r-1}\).}
\GSfield{Reconstruction}{The level variable m is replaced by the moment exponent r only in the modern statement.}
\GSfield{Fate}{\citet[\S8.4--8.7]{manin1973periods}; notebook-specific exponents are checked directly.}
\GSfield{Present status}{As of 9 November 2026: established in the stated scope.}
\GSfield{Assessment}{literal statement: false as written; publication match: exact; verification: read the result; development: annotated; confidence: high}
\end{GSitem}
\begin{GSitem}{1973-304.4}{E}{Continuous character space and its disc coordinates}
\GSfield{Kind; source}{construction; \GSpage{P000506}}
\GSfield{The notes}{\(\mathbb Z_p^\times=(\mathbb Z/p\mathbb Z)^\times\times(1+p)^{\mathbb Z_p}\), and \(\chi_1(1+p)=1+t\).}
\GSfield{Statement}{For odd \(p\) and finite \(K/\mathbb Q_p\), continuous characters of \(\mathbb Z_p^\times\) have a finite Teichmuller component and a component parametrized by \(t\in\mathfrak m_K\) via \(\chi_t((1+p)^a)=(1+t)^a\). For \(p=2\) use \(\{\pm1\}\times(1+4\mathbb Z_2)\) and generator 5.}
\GSfield{Reconstruction}{Continuity, the Teichmuller lift, and the separate p=2 case are made explicit.}
\GSfield{Fate}{\citet[\S8.3]{manin1973periods}; occurrence 1972-299.4 already records the p=2 qualification.}
\GSfield{Present status}{As of 9 November 2026: established in the stated scope.}
\GSfield{Assessment}{literal statement: correct with missing hypotheses; publication match: exact; verification: read the result; development: none; confidence: high}
\end{GSitem}
\begin{GSitem}{1973-304.5}{E}{Analyticity of the p-adic Mellin transform}
\GSfield{Kind; source}{theorem; \GSpage{P000506}}
\GSfield{The notes}{A moderately growing measure has \(\int\chi_{(t)}d\mu=F(t)\in K[[t]]\), convergent for \(|t|<1\).}
\GSfield{Statement}{On each component of character space, the uniform moderate-growth bound on the distribution makes its Mellin transform analytic on the open unit disc of the character parameter. For integer-power characters \(x\mapsto x^n\), fix \(n\bmod(p-1)\) when \(p\) is odd.}
\GSfield{Fate}{\citet[Theorem~8.7]{manin1973periods}.}
\GSfield{Present status}{As of 9 November 2026: established in the stated scope.}
\GSfield{Assessment}{literal statement: correct with missing hypotheses; publication match: exact; verification: read the result; development: none; confidence: high}
\end{GSitem}
\begin{GSitem}{1973-304.6}{E}{Holomorphic cusp forms for the full modular group}
\GSfield{Kind; source}{definition; \GSpage{P000506}}
\GSfield{The notes}{\(P_w\) consists of forms with invariant \(\varphi(z)(dz)^{(w+2)/2}\) and Fourier expansion starting at n=1.}
\GSfield{Statement}{For even \(k=w+2\ge2\), \(S_k(SL_2(\mathbb Z))\) consists of holomorphic functions on the upper half-plane satisfying \(f((az+b)/(cz+d))=(cz+d)^k f(z)\) and vanishing at the cusp; their Fourier expansion is \(\sum_{n\ge1}a_n e^{2\pi inz}\).}
\GSfield{Fate}{\citet[\S1.1--1.2]{manin1973periods}.}
\GSfield{Present status}{As of 9 November 2026: established in the stated scope.}
\GSfield{Assessment}{literal statement: correct; publication match: exact; verification: read the result; development: none; confidence: high}
\end{GSitem}
\begin{GSitem}{1973-304.7}{E}{The discriminant form and low-weight dimensions}
\GSfield{Kind; source}{claim; \GSpage{P000506}}
\GSfield{The notes}{The first nonzero space occurs at weight 12, with gaps up to 26 and a parenthetical one-dimensional claim; \(\Delta=q\prod(1-q^n)^{24}\) is displayed.}
\GSfield{Statement}{\(\Delta\in S_{12}(SL_2(\mathbb Z))\), \(S_{12}=\mathbb C\Delta\). The low-weight cusp spaces are zero for even k<12 and for k=14; they have dimension 1 at k=16,18,20,22,26, and dimension 2 at k=24.}
\GSfield{Reconstruction}{The notebook shorthand must not be read as dimension one at every weight from 12 through26; weight24 is the first two-dimensional example.}
\GSfield{Fate}{\citet[\S1.2 and the weight tables in \S7]{manin1973periods}.}
\GSfield{Present status}{As of 9 November 2026: established in the stated scope.}
\GSfield{Assessment}{literal statement: correct with missing hypotheses; publication match: exact; verification: read the result; development: none; confidence: high}
\end{GSitem}
\begin{GSitem}{1973-304.8}{E}{Hecke eigenbasis and algebraic Fourier coefficients}
\GSfield{Kind; source}{theorem; \GSpage{P000506-P000507}}
\GSfield{The notes}{There is a basis with algebraic-integer coefficients whose Dirichlet series have Euler products.}
\GSfield{Statement}{The commuting Hecke operators on full-level \(S_k\) admit normalized simultaneous eigenforms \(a_1=1\), with coefficients in number fields and algebraic integers; for such a form, \(L(f,s)=\prod_\ell(1-a_\ell\ell^{-s}+\ell^{k-1-2s})^{-1}\) initially in a right half-plane.}
\GSfield{Reconstruction}{Normalized eigenforms, not every arbitrary algebraic basis, give the Euler product. The obscured exponent is resolved by the next clear display on P000507.}
\GSfield{Fate}{\citet[\S1.2 and \S3]{manin1973periods}.}
\GSfield{Present status}{As of 9 November 2026: established in the stated scope.}
\GSfield{Assessment}{literal statement: correct with missing hypotheses; publication match: exact; verification: read the result; development: none; confidence: high}
\end{GSitem}
\begin{GSitem}{1973-304.9}{E}{The prime Hecke operator and its coefficient recurrence}
\GSfield{Kind; source}{construction; \GSpage{P000507}}
\GSfield{The notes}{\((\Phi|T_p)(z)=p^{w+1}\Phi(pz)+p^{-1}\sum_{a=0}^{p-1}\Phi((z+a)/p)\).}
\GSfield{Statement}{For \(k=w+2\), this operator preserves \(S_k\) and sends coefficients to \(a_{pn}+p^{k-1}a_{n/p}\), with \(a_{n/p}=0\) unless p divides n. An eigenform obeys \(a_{pn}+p^{k-1}a_{n/p}=a_pa_n\), giving the local quadratic Euler factor.}
\GSfield{Fate}{\citet[\S3, formula~(11)]{manin1973periods}.}
\GSfield{Present status}{As of 9 November 2026: established in the stated scope.}
\GSfield{Assessment}{literal statement: correct; publication match: exact; verification: read the result; development: none; confidence: high}
\end{GSitem}
\begin{GSitem}{1973-304.10}{E}{Rational-cusp period integrals}
\GSfield{Kind; source}{construction; \GSpage{P000507}}
\GSfield{The notes}{\(P(x)=\int_x^{i\infty}\Phi(z)dz\) for rational x; it converges at both cusps and \(P(x+1)=P(x)\). The note contrasts irrational endpoints.}
\GSfield{Statement}{For a full-level cusp form, the path integral from a rational cusp x to infinity converges, interpreted with cusp approach inside the upper half-plane, and is periodic in x. No uniform divergence claim for every irrational endpoint is needed.}
\GSfield{Reconstruction}{The unqualified phrase about irrational points is retained as a heuristic, not a theorem with universal quantifier.}
\GSfield{Fate}{\citet[\S3.2 and Proposition~9.2]{manin1973periods}.}
\GSfield{Present status}{As of 9 November 2026: established in the stated scope.}
\GSfield{Assessment}{literal statement: correct with missing hypotheses; publication match: exact; verification: read the result; development: none; confidence: high}
\end{GSitem}
\begin{GSitem}{1973-304.11}{E}{The period distribution relation}
\GSfield{Kind; source}{theorem; \GSpage{P000507}}
\GSfield{The notes}{\(\lambda_pP(x)=p^wP(px)+\sum_{a=0}^{p-1}P((x+a)/p)\).}
\GSfield{Statement}{If \(T_p\Phi=a_p\Phi\), changing variables in the convergent cusp integrals gives \(\sum_{j=0}^{p-1}P((x+j)/p)=a_pP(x)-p^{k-2}P(px)\) for rational x.}
\GSfield{Fate}{\citet[Proposition~9.2(d)]{manin1973periods}.}
\GSfield{Present status}{As of 9 November 2026: established in the stated scope.}
\GSfield{Assessment}{literal statement: correct; publication match: exact; verification: read the result; development: none; confidence: high}
\end{GSitem}
\begin{GSitem}{1973-304.12}{E}{The tentative correction of the period half-measure}
\GSfield{Kind; source}{construction; \GSpage{P000508}}
\GSfield{The notes}{The page labelled correction has \(\rho^{-m}P(a/p^m)-p^w\rho^{-(m+m)}P(a/p^{m+1})\), with a manuscript question mark.}
\GSfield{Statement}{For \(\rho^2-a_p\rho+p^{k-1}=0\), a compatible distribution at level n is \(\rho^{-n}P(a/p^n)-p^{k-2}\rho^{-n-1}P(a/p^{n-1})\). The next talk writes an overall \(\rho\)-multiple of this formula.}
\GSfield{Reconstruction}{Both the second exponent and the direction of the denominator shift are repaired, using the explicit corrected repetition on P000509, not merely the printed formula.}
\GSfield{Fate}{\citet[Lemma~9.4 and formula~(39)]{manin1973periods}.}
\GSfield{Present status}{As of 9 November 2026: established in the stated scope.}
\GSfield{Assessment}{literal statement: false as written; publication match: partial; verification: read the result; development: annotated; confidence: high; problem \hyperref[prob:GP-1972-403]{GP-1972-403}}
\end{GSitem}

\GStalk{305}{1973-09-17}{Yu. I. Manin}{P000508-P000513}{Periods of cusp forms, p-adic Hecke functions and special zeta values}
The dated continuation occupies P000509--P000511. It proves finite generation of rational-cusp periods by continued fractions, separates their two parity periods, and transfers Hecke operations to period coordinates. The resulting algebraic-valued distributions lead to primitive-character L-values. A closing discussion connects zeta values with higher K-groups and explicitly notes the unproved transcendence of zeta(3). The principal identification is \citet[Theorem~1.2, Proposition~9.2 and Theorem~10.2]{manin1973periods}; prior-year occurrences 1972-299.10--14 prevent treating this repeated exposition as an orphan construction.

\begin{GSitem}{1973-305.1}{E}{The corrected quadratic-root distribution}
\GSfield{Kind; source}{construction; \GSpage{P000509}}
\GSfield{The notes}{For a root \(\rho\) of \(\rho^2-\lambda_p\rho+p^{w+1}=0\), the notes assign \(\rho^{-m}P(a/p^{m+1})-p^w\rho^{-m-1}P(a/p^m)\) to a class modulo \(p^{m+1}\).}
\GSfield{Statement}{For every n≥1, let \(\nu_n(a)=\rho^{1-n}P(a/p^n)-p^{k-2}\rho^{-n}P(a/p^{n-1})\). The Hecke distribution relation and the quadratic equation imply \(\sum_{j=0}^{p-1}\nu_{n+1}(a+jp^n)=\nu_n(a)\). This gives a complex-valued finitely additive distribution before period normalization.}
\GSfield{Reconstruction}{\(n=m+1\) and \(k=w+2\) reconcile indexing. A nonzero root is automatic since \(p^{k-1}\ne0\).}
\GSfield{Fate}{\citet[Lemma~9.4 and formula~(39)]{manin1973periods}, up to the stated overall factor \(\rho\).}
\GSfield{Present status}{As of 9 November 2026: established in the stated scope.}
\GSfield{Assessment}{literal statement: correct; publication match: exact; verification: read the result; development: none; confidence: high; problem \hyperref[prob:GP-1972-403]{GP-1972-403}}
\end{GSitem}
\begin{GSitem}{1973-305.2}{E}{Periods and critical L-values}
\GSfield{Kind; source}{definition; \GSpage{P000509}}
\GSfield{The notes}{\(\tau_j(\Phi)=\int_0^{i\infty}\Phi(z)z^jdz\), and \(\tau_j=j!i^{j+1}L_\Phi(j+1)/(2\pi)^{j+1}\).}
\GSfield{Statement}{For a weight \(k=w+2\) cusp form and \(0\le j\le k-2\), its period is \(\tau_j=i^{j+1}\Gamma(j+1)(2\pi)^{-j-1}L(f,j+1)\). The critical integers are \(1,\ldots,k-1\). The completed L-function satisfies its functional equation about \(s=k/2\).}
\GSfield{Reconstruction}{The notebook\textquoteright s phrase critical strip means the critical integer range, not the analytic critical strip of a Dirichlet series.}
\GSfield{Fate}{\citet[\S1.1--1.2]{manin1973periods}.}
\GSfield{Present status}{As of 9 November 2026: established in the stated scope.}
\GSfield{Assessment}{literal statement: correct; publication match: exact; verification: read the result; development: none; confidence: high}
\end{GSitem}
\begin{GSitem}{1973-305.3}{E}{Continued fractions and integral modular-symbol reduction}
\GSfield{Kind; source}{method; \GSpage{P000509}}
\GSfield{The notes}{Consecutive convergents give matrices \(g_j\in SL_2(\mathbb Z)\); the path from 0 to b/d is cut into transformed standard paths.}
\GSfield{Statement}{For rational b/d and weight k, Euclidean-algorithm edges join adjacent cusps of determinant ±1. Changing variables under the corresponding modular matrices expresses \(P(b/d)\) as \(\sum_{j=0}^{k-2}A_j(b/d)\tau_j\) with \(A_j\in\mathbb Z\), after consistent path orientations.}
\GSfield{Reconstruction}{The initial edge and its orientation repair the first sum\textquoteright s lower index 0, which otherwise invokes an undefined convergent -1.}
\GSfield{Fate}{\citet[Proposition~3.2 and Proposition~9.2(c)]{manin1973periods}.}
\GSfield{Present status}{As of 9 November 2026: established in the stated scope.}
\GSfield{Assessment}{literal statement: correct with missing hypotheses; publication match: exact; verification: read the result; development: none; confidence: high}
\end{GSitem}
\begin{GSitem}{1973-305.4}{E}{Parity and reality of periods}
\GSfield{Kind; source}{claim; \GSpage{P000509}}
\GSfield{The notes}{Even-index periods are purely imaginary and odd-index periods real.}
\GSfield{Statement}{If the cusp form has real Fourier coefficients, then \(f(iy)\in\mathbb R\) and \(\tau_j=i^{j+1}\int_0^\infty f(iy)y^jdy\); thus even j gives an imaginary period and odd j a real period.}
\GSfield{Reconstruction}{The real-coefficient hypothesis is necessary and appears on the following page.}
\GSfield{Fate}{\citet[Proposition~9.2, proof]{manin1973periods}.}
\GSfield{Present status}{As of 9 November 2026: established in the stated scope.}
\GSfield{Assessment}{literal statement: correct with missing hypotheses; publication match: exact; verification: read the result; development: none; confidence: high}
\end{GSitem}
\begin{GSitem}{1973-305.5}{E}{Algebraicity after two period normalizations}
\GSfield{Kind; source}{theorem; \GSpage{P000510}}
\GSfield{The notes}{There exist \(\omega^+,\omega^-\) such that the even and odd period ratios lie in the field of Fourier coefficients.}
\GSfield{Statement}{For a nonzero normalized simultaneous Hecke eigenform \(f\in S_k(SL_2(\mathbb Z))\) and its coefficient field E, there are nonzero complex numbers \(\Omega^+,\Omega^-\) for which \(\tau_j/\Omega^+\in E\) for even j and \(\tau_j/\Omega^-\in E\) for odd j.}
\GSfield{Reconstruction}{Normalization, eigenform and nonzero-period hypotheses are explicit. The theorem does not prove that the common factors themselves are transcendental.}
\GSfield{Fate}{\citet[Theorem~1.2]{manin1973periods}; submitted in 1973 after the earlier-year seminar exposition.}
\GSfield{Present status}{As of 9 November 2026: established in the stated scope.}
\GSfield{Assessment}{literal statement: correct with missing hypotheses; publication match: exact; verification: read the result; development: none; confidence: high; problem \hyperref[prob:GP-1972-403]{GP-1972-403}}
\end{GSitem}
\begin{GSitem}{1973-305.6}{E}{Eichler-Shimura period relations}
\GSfield{Kind; source}{theorem; \GSpage{P000510}}
\GSfield{The notes}{The order-two and order-three modular generators give two families of homogeneous period equations, including \(\tau_j+(-1)^j\tau_{w-j}=0\).}
\GSfield{Statement}{For even w, the period polynomial \(r_f(X)=\int_0^{i\infty}f(z)(z-X)^w dz\) satisfies \(r_f|(1+S)=0\) and \(r_f|(1+U+U^2)=0\), using the right weight-w polynomial action of the order-two and order-three elements of PSL2(Z). Expanding gives the notebook\textquoteright s linear relations.}
\GSfield{Reconstruction}{Polynomial-action conventions make the signs unambiguous; they are not a new conjecture.}
\GSfield{Fate}{\citet[\S2, equations~(4)--(5)]{manin1973periods}.}
\GSfield{Present status}{As of 9 November 2026: established in the stated scope.}
\GSfield{Assessment}{literal statement: correct; publication match: exact; verification: read the result; development: none; confidence: high}
\end{GSitem}
\begin{GSitem}{1973-305.7}{E}{Odd period coordinates and the one-dimensional even excess}
\GSfield{Kind; source}{theorem; \GSpage{P000510}}
\GSfield{The notes}{The real cusp-form space maps isomorphically onto \(S^-\), and into \(S^+\) with codimension one.}
\GSfield{Statement}{For even \(k\ge4\), the odd period-polynomial solution space has the same real dimension as the space of real-coefficient cusp forms, and the even solution space contains their period image with one-dimensional Eisenstein/coboundary excess. The period maps have the stated injectivity and isomorphism properties.}
\GSfield{Reconstruction}{The standard period-polynomial conventions and weight range are supplied; this assertion is not about arbitrary cohomology of all modular forms.}
\GSfield{Fate}{\citet[\S2, Proposition~2.3]{manin1973periods}.}
\GSfield{Present status}{As of 9 November 2026: established in the stated scope.}
\GSfield{Assessment}{literal statement: correct with missing hypotheses; publication match: exact; verification: read the result; development: none; confidence: high}
\end{GSitem}
\begin{GSitem}{1973-305.8}{E}{Hecke operators on period coordinates}
\GSfield{Kind; source}{construction; \GSpage{P000510}}
\GSfield{The notes}{A commutative diagram transfers \(T_p\) to \(T_p^-\), with arithmetic matrix entries. The rewritten analytic formula has a minus sign before the average.}
\GSfield{Statement}{The Hecke action on the odd period image is induced by an integral matrix in the standard period coordinates; the period map intertwines it with the usual plus-sign Hecke operator on cusp forms.}
\GSfield{Reconstruction}{The minus sign in the analytic formula is a notebook error, contradicted by its own P000507 formula and Fourier recurrence.}
\GSfield{Fate}{\citet[Proposition~3.2]{manin1973periods}.}
\GSfield{Present status}{As of 9 November 2026: established in the stated scope.}
\GSfield{Assessment}{literal statement: false as written; publication match: exact; verification: read the result; development: annotated; confidence: high}
\end{GSitem}
\begin{GSitem}{1973-305.9}{E}{Ramanujan congruence modulo 691}
\GSfield{Kind; source}{theorem; \GSpage{P000510}}
\GSfield{The notes}{The discriminant form\textquoteright s coefficients satisfy \(\tau(n)\equiv\sum_{d\mid n}d^{11}\pmod{691}\).}
\GSfield{Statement}{For \(\Delta(q)=q\prod_{m\ge1}(1-q^m)^{24}=\sum_{n\ge1}\tau(n)q^n\), the congruence \(\tau(n)\equiv\sigma_{11}(n)\pmod{691}\) holds for every n≥1.}
\GSfield{Reconstruction}{The statement was already known; a new proof mechanism is not a new open problem.}
\GSfield{Fate}{\citet[\S1.3 and \S7]{manin1973periods} supplies a period-based proof of the older congruence.}
\GSfield{Present status}{As of 9 November 2026: established in the stated scope.}
\GSfield{Assessment}{literal statement: correct; publication match: exact; verification: read the result; development: none; confidence: high}
\end{GSitem}
\begin{GSitem}{1973-305.10}{E}{Algebraic normalized period measures}
\GSfield{Kind; source}{construction; \GSpage{P000511}}
\GSfield{The notes}{\(P^+=i\operatorname{Im}P/\omega^+\) and \(P^-=\operatorname{Re}P/\omega^-\) are used in the corrected measure formula.}
\GSfield{Statement}{Choose parity periods as above. The rational-cusp values of the normalized even and odd components lie in a fixed number field with bounded denominators. After fixing an embedding into \(\overline{\mathbb Q}_p\) and a finite completion containing \(\rho\), the compatible distribution is p-adic valued.}
\GSfield{Reconstruction}{The choice of complex parity convention, number-field embedding and completion is indispensable.}
\GSfield{Fate}{\citet[Proposition~9.2(c), \S9.3--9.5]{manin1973periods}.}
\GSfield{Present status}{As of 9 November 2026: established in the stated scope.}
\GSfield{Assessment}{literal statement: correct with missing hypotheses; publication match: exact; verification: read the result; development: none; confidence: high; problem \hyperref[prob:GP-1972-403]{GP-1972-403}}
\end{GSitem}
\begin{GSitem}{1973-305.11}{E}{The ordinary root gives a bounded measure}
\GSfield{Kind; source}{theorem; \GSpage{P000511}}
\GSfield{The notes}{If \(\lambda_p\) is coprime to p, one root gives a bounded measure and the other a growing one.}
\GSfield{Statement}{Normalize \(v_p(p)=1\). If \(a_p\) is a p-adic unit, the roots of \(X^2-a_pX+p^{k-1}\) have valuations 0 and k−1. The unit root in the bounded-denominator distribution formula gives a bounded measure. The other root has a potentially growing upper bound; boundedness cannot be excluded merely from that bound.}
\GSfield{Reconstruction}{The root valuations and positive conclusion are exact. The assertion that every nonunit-root distribution must grow needs a noncancellation hypothesis.}
\GSfield{Fate}{\citet[\S9.5]{manin1973periods}.}
\GSfield{Present status}{As of 9 November 2026: established in the stated scope.}
\GSfield{Assessment}{literal statement: correct with missing hypotheses; publication match: exact; verification: read the result; development: none; confidence: high; problem \hyperref[prob:GP-1972-403]{GP-1972-403}}
\end{GSitem}
\begin{GSitem}{1973-305.12}{F}{The unresolved relation between a\_p and p}
\GSfield{Kind; source}{remark; \GSpage{P000511}}
\GSfield{The notes}{\ru{Неизв. связь} \(\lambda_p\) \ru{и} \(p\), after the ordinary-prime discussion.}
\GSfield{Statement}{The record does not specify whether it asks for a fixed-prime criterion, infinitely many ordinary primes, their density, or a uniform bound. These are different arithmetic questions.}
\GSfield{Fate}{\citet[\S9.5]{manin1973periods} discusses growth, but does not supply missing quantifiers for this remark.}
\GSfield{Present status}{As of 9 November 2026: no definite question can be indexed without strengthening the notebook.}
\GSfield{Assessment}{literal statement: not determinable; publication match: thematic only; verification: indirect; development: none; confidence: low}
\end{GSitem}
\begin{GSitem}{1973-305.13}{E}{Primitive-character interpolation of the critical value at one}
\GSfield{Kind; source}{theorem; \GSpage{P000511}}
\GSfield{The notes}{The notebook writes \(\int\chi\,d\mu_f^\pm=\cdots L_f(1,\chi)\), leaving constants as dots.}
\GSfield{Statement}{For a normalized full-level eigenform, choose compatible periods, a p-adic embedding and a root \(\rho\). For primitive Dirichlet characters of p-power conductor, the distribution\textquoteright s finite character sums equal the period-normalized twisted critical value at s=1, with the conductor power of \(\rho\), Gauss sum, inverse-character convention and parity fixed by the definition. Under moderate growth this is the restriction of an analytic Mellin transform.}
\GSfield{Reconstruction}{No numerical prefactor is attributed to the notebook: its ellipses prevent a literal normalization match. Locally constant character sums exist even without a bounded measure.}
\GSfield{Fate}{\citet[Theorem~10.2, formula~(40)]{manin1973periods} supplies the complete normalization. The notes alone do not determine their dots.}
\GSfield{Present status}{As of 9 November 2026: established in the stated scope.}
\GSfield{Assessment}{literal statement: correct with missing hypotheses; publication match: partial; verification: read the result; development: none; confidence: high; problem \hyperref[prob:GP-1972-403]{GP-1972-403}}
\end{GSitem}
\begin{GSitem}{1973-305.14}{F}{The Lichtenbaum regulator interpretation of zeta(3)}
\GSfield{Kind; source}{conjecture; \GSpage{P000511}}
\GSfield{The notes}{\(\zeta(3)=2\pi^2|K_4\mathbb Z|e_5(\alpha_5)\), where \(\alpha_5\) generates the free part of \(K_5\mathbb Z\) and \(e_5\) is described only as a certain real homomorphism.}
\GSfield{Statement}{The exact coefficient cannot be evaluated until a normalization and sign of \(e_5\) are given. A precise weaker modern statement is that the appropriate nonzero real regulator on \(K_5(\mathbb Z)/\mathrm{torsion}\) is related to \(\zeta(3)/\pi^2\) by a nonzero rational factor; it does not imply transcendence of either side.}
\GSfield{Reconstruction}{The weaker Borel comparison is identified separately. It is not claimed to verify the exact notebook coefficient 2, nor to solve all modern special-value conjectures.}
\GSfield{Fate}{\citet[\S4]{lichtenbaum1972values} is the historical conjectural setting; \citet[\S6.4]{borel1977zeta} establishes the regulator comparison up to rational factors, with its normalization explicitly defined.}
\GSfield{Present status}{As of 9 November 2026: the regulator framework is established, but the exact handwritten equality remains undernormalized.}
\GSfield{Assessment}{literal statement: not determinable; publication match: partial; verification: read the result; development: none; confidence: medium}
\end{GSitem}
\begin{GSitem}{1973-305.15}{C}{Transcendence of zeta(3)}
\GSfield{Kind; source}{question; \GSpage{P000511}}
\GSfield{The notes}{\ru{Не док., что} \(\sum_{n\ge1}n^{-3}\) \ru{трансценд.}}
\GSfield{Statement}{Is \(\zeta(3)=\sum_{n\ge1}n^{-3}\) transcendental over \(\mathbb Q\)?}
\GSfield{Fate}{Ap\'ery proved irrationality in 1978; \citet[\S1, p.~2; manuscript dated 26 October 2024]{nica2024irrationality} explicitly separates that theorem from the still-unknown transcendence question. A recent preprint, \citet[Theorem~1.6; v2 dated 15 October 2026]{lai2026linear}, reports linear independence of \(1,\zeta(3),\zeta(4)\). This is a strictly weaker arithmetic assertion and is not a proof of transcendence; its theorem statement and introduction, rather than the complete proof, were checked.}
\GSfield{Present status}{As of 9 November 2026: irrationality is proved, while no accepted transcendence proof was identified in the literature search completed on this date.}
\GSfield{Assessment}{literal statement: correct; publication match: exact; verification: read the result; development: none; confidence: high; problem \hyperref[prob:GP-1973-10]{GP-1973-10}}
\GSfield{Worth developing}{No bounded development is justified here: the precise transcendence problem remains far beyond the elementary estimates in these notes.}
\end{GSitem}
\begin{GSitem}{1973-305.16}{C}{The arithmetic nature of Euler's constant}
\GSfield{Kind; source}{question; \GSpage{P000511}}
\GSfield{The notes}{Immediately after the transcendence sentence: \ru{Пост. Эйлера тоже не док.}}
\GSfield{Statement}{Reading the carried-over predicate literally gives the question whether \(\gamma=\lim_{n\to\infty}(\sum_{j=1}^n1/j-\log n)\) is transcendental. The notebook does not separately distinguish irrationality from transcendence.}
\GSfield{Reconstruction}{The antecedent of “also” is identified, but the weak notebook wording is preserved; no claim of a newly posed precise irrationality conjecture is made.}
\GSfield{Fate}{\citet[Conjectures~1.0.1--1.0.2 and \S3]{lagarias2013euler} surveys the arithmetic problem; even irrationality remains unproved in the accepted results checked.}
\GSfield{Present status}{As of 9 November 2026: the checked literature gives partial arithmetic criteria, not a proof of irrationality or transcendence; no accepted later proof was identified.}
\GSfield{Assessment}{literal statement: correct with missing hypotheses; publication match: partial; verification: read the result; development: none; confidence: medium; problem \hyperref[prob:GP-1973-11]{GP-1973-11}}
\GSfield{Worth developing}{No bounded attack is justified by this one-line remark; it is useful as a historical occurrence of a familiar arithmetic question.}
\end{GSitem}

\GStalk{307}{1973-09-24}{I. M. Gelfand}{P000511-P000514}{Helton--Howe traces and a non-Archimedean analogy}
The 24 September opening on P000512 precedes the explicit Kirillov heading, so record 307 is placed before 306 in source order. The block names Pincus, Douglas, Coburn, an uncertain surname, Singer, Helton and Howe; it states the polynomial lemma and the trace formula, then suggests normed rings as a non-Archimedean substitute for differential operators. Attribution to Gelfand follows the catalogue, since the page itself labels the session rather than this speaker. The algebraic lemma is checked in \citet[\S4]{heltonhowe1975traces}; the trace formula circulated at the April 1973 Halifax conference.

\begin{GSitem}{1973-307.1}{E}{Wallach factorization through the Poisson bracket}
\GSfield{Kind; source}{theorem; \GSpage{P000512}}
\GSfield{The notes}{For a bilinear form \(B\) on polynomials, \(B(P_1(Q),P_2(Q))=0\) implies \(B(R,S)=L(\{R,S\})\). The inserted phrase specifies two variables, while the next line says any number.}
\GSfield{Statement}{Let \(B:\mathbb C[x,y]\times\mathbb C[x,y]\to\mathbb C\) be bilinear and suppose \(B(p(q),r(q))=0\) for all \(p,r\in\mathbb C[t]\) and \(q\in\mathbb C[x,y]\). There is a unique linear functional \(L\) on \(\mathbb C[x,y]\) such that \(B(f,g)=L(f_xg_y-f_yg_x)\).}
\GSfield{Reconstruction}{The polynomial ring and characteristic zero are fixed explicitly. The two-variable conclusion cannot be read as a theorem with an unspecified Poisson structure in an arbitrary number of variables.}
\GSfield{Fate}{The algebraic form of Wallach\textquoteright s lemma is proved in \citet[\S4, algebraic preliminaries and Proposition~4.2]{heltonhowe1975traces}, extending the 1973 work. This is an exposition of the Helton--Howe programme, not a new theorem by the seminar speaker.}
\GSfield{Present status}{As of 9 November 2026: established in the stated scope.}
\GSfield{Assessment}{literal statement: correct with missing hypotheses; publication match: exact; verification: read the result; development: none; confidence: high}
\end{GSitem}
\begin{GSitem}{1973-307.2}{E}{The Helton--Howe trace measure}
\GSfield{Kind; source}{theorem; \GSpage{P000512}}
\GSfield{The notes}{Boundedness is unstated: self-adjoint \(U,V\) commute modulo nuclear operators and \(\operatorname{Tr}[R(U,V),S(U,V)]=\iint\{R,S\}\,d\sigma\). The order of the factors is said to be immaterial.}
\GSfield{Statement}{For bounded self-adjoint operators \(U,V\) on a complex separable Hilbert space with \([U,V]\) trace class, there is a finite, compactly supported, purely imaginary Borel measure \(\sigma\) on \(\mathbb R^2\) such that the displayed formula holds for all complex polynomials in two variables. Any fixed ordering convention gives the same trace of the commutator; it need not give the same operator.}
\GSfield{Reconstruction}{Bounded operators and a complex measure remove domain and missing-factor ambiguities. Nuclear here means trace class.}
\GSfield{Fate}{\citet[main theorem, formula~\((1^{\prime})\)]{heltonhowe1973integral}, presented at Halifax in April 1973 and printed in October. The ordering assertion is checked in \citet[Theorem~I and preceding paragraph]{heltonhowe1975traces}. The exact measure normalization is restated in \citet[Theorem~3.3]{sugahara2026heltonhowe}; its newer Toeplitz applications are not needed for the 1973 assertion.}
\GSfield{Present status}{As of 9 November 2026: established in the stated scope.}
\GSfield{Assessment}{literal statement: correct with missing hypotheses; publication match: exact; verification: read the result; development: none; confidence: high}
\end{GSitem}
\begin{GSitem}{1973-307.3}{F}{Normed rings in a proposed p-adic operator theory}
\GSfield{Kind; source}{method; \GSpage{P000512}}
\GSfield{The notes}{After citing Douglas\textquoteright s book, the notes propose suitably chosen normed rings in the p-adic case in place of differential operators.}
\GSfield{Statement}{No ring, module category, norm, class of operators or theorem to be recovered is specified. The strongest supported statement is a methodological suggestion linking the preceding operator theory to Manin\textquoteright s integration talks.}
\GSfield{Fate}{The operator-algebra context is documented by \citet[Introduction]{heltonhowe1975traces}; comparison with \citet[Introduction]{manin1976nonarchimedean} does not identify a precise common construction.}
\GSfield{Present status}{As of 9 November 2026: an unformulated programme; no mathematical problem is manufactured from it.}
\GSfield{Assessment}{literal statement: not determinable; publication match: thematic only; verification: unresolved lead; development: none; confidence: low}
\end{GSitem}

\GStalk{306}{1973-09-24}{A. A. Kirillov}{P000511-P000513}{Brown--Douglas--Fillmore extensions and analytic K-homology}
The explicit Kirillov heading is on the lower part of P000512; the discussion ends before the Fuchs heading on P000513. It presents extensions by compact operators, their Calkin-algebra description, the emerging homology theory and index pairings, and applications to normal operators. The catalogue lists this record before Gelfand, but the notebook order is the reverse. \citet[pp.~699--702]{douglas1980extensions} retrospectively identifies which results were in the April 1973 Halifax paper and which received fuller treatment in 1977. A garbled degree formula and the unspecified affine-group example are not silently completed.

\begin{GSitem}{1973-306.1}{E}{C*-extensions by the compact operators}
\GSfield{Kind; source}{definition; \GSpage{P000512}}
\GSfield{The notes}{\(0\to\mathcal K\to\mathcal A\to C(X)\to0\), called an extension of Banach algebras; \(X\) is compact and \(\mathcal K\) acts on a separable Hilbert space.}
\GSfield{Statement}{For nonempty compact metrizable \(X\) and an infinite-dimensional separable complex Hilbert space \(H\), consider essential C*-algebra extensions of \(C(X)\) by \(\mathcal K(H)\), with *-homomorphisms and the usual equivalence modulo addition of split extensions. In this category the resulting \(\operatorname{Ext}(X)\) is an abelian group.}
\GSfield{Reconstruction}{C*-structure, metrizability, essentiality and stable equivalence replace the misleading unrestricted Banach-algebra wording.}
\GSfield{Fate}{\citet{browndouglasfillmore1973unitary}; the definitions and chronology are checked in the authors\textquoteright s retrospective \citet[pp.~699--701]{douglas1980extensions}.}
\GSfield{Present status}{As of 9 November 2026: established in the stated scope.}
\GSfield{Assessment}{literal statement: correct with missing hypotheses; publication match: exact; verification: read the result; development: none; confidence: high}
\end{GSitem}
\begin{GSitem}{1973-306.2}{E}{Busby maps, splitting and addition}
\GSfield{Kind; source}{construction; \GSpage{P000512}}
\GSfield{The notes}{Embeddings \(C(X)\hookrightarrow\mathcal L/\mathcal K\) are taken up to unitary equivalence; those lifting to \(\mathcal L\) are trivial. Addition is \(\tau_1\oplus\tau_2\) on \(H_1\oplus H_2\cong H\).}
\GSfield{Statement}{An essential extension determines an injective *-homomorphism \(\tau:C(X)\to\mathcal Q(H)=\mathcal B(H)/\mathcal K(H)\), and its pullback recovers the extension. A *-homomorphic lift defines a split class. Direct sum, followed by a unitary identification of infinite-dimensional separable Hilbert spaces, gives the addition on Ext; stable equivalence makes the choice immaterial.}
\GSfield{Fate}{\citet[p.~700]{douglas1980extensions} gives precisely these equivalent constructions and attributes the commutative case to the 1973 work.}
\GSfield{Present status}{As of 9 November 2026: established in the stated scope.}
\GSfield{Assessment}{literal statement: correct with missing hypotheses; publication match: exact; verification: read the result; development: none; confidence: high}
\end{GSitem}
\begin{GSitem}{1973-306.3}{E}{Suspensions and the homological interpretation}
\GSfield{Kind; source}{theorem; \GSpage{P000512}}
\GSfield{The notes}{\(\operatorname{Ext}_{1-n}(X)=\operatorname{Ext}(S^nX)\); the main result is said to reduce Atiyah\textquoteright s \(\operatorname{Ell}(X)\) to Ext.}
\GSfield{Statement}{For based finite CW complexes, reduced groups \(\widetilde E_{1-n}(X)=\operatorname{Ext}(\Sigma^nX)\) extend, with periodicity, to reduced complex K-homology. In particular \(\operatorname{Ext}(X)\cong K_1(X)\), whereas even Fredholm cycles represent \(K_0(X)\). The relation with Atiyah\textquoteright s Ell uses the even group and an equivalence relation on cycles; it is not literally \(\operatorname{Ell}(X)=\operatorname{Ext}_1(X)\).}
\GSfield{Reconstruction}{Reduced suspensions, the even/odd degree distinction, and the cycle quotient are supplied. The result was announced in the BDF programme before this exposition; the 1977 proof date is not its first circulation.}
\GSfield{Fate}{\citet[p.~702]{douglas1980extensions} distinguishes the two degrees and the solution of Atiyah\textquoteright s problem; the full treatment is \citet{browndouglasfillmore1977extensions}.}
\GSfield{Present status}{As of 9 November 2026: established in the stated scope.}
\GSfield{Assessment}{literal statement: correct with missing hypotheses; publication match: exact; verification: read the result; development: none; confidence: high}
\end{GSitem}
\begin{GSitem}{1973-306.4}{E}{Exact sequence for a closed subcomplex}
\GSfield{Kind; source}{theorem; \GSpage{P000513}}
\GSfield{The notes}{The sequence begins \(\operatorname{Ext}_1(A)\to\operatorname{Ext}_1(X)\to\operatorname{Ext}_1(X/A)\to\operatorname{Ext}_0(A)\).}
\GSfield{Statement}{For a based finite CW pair \((X,A)\) with \(A\ne\varnothing\) a subcomplex, the reduced BDF homology groups have the cofiber long exact sequence \(\cdots\to\widetilde E_1(A)\to\widetilde E_1(X)\to\widetilde E_1(X/A)\to\widetilde E_0(A)\to\cdots\).}
\GSfield{Reconstruction}{The finite CW-pair scope avoids treating an arbitrary inclusion as a cofiber sequence without qualifications.}
\GSfield{Fate}{The generalized homology theorem in \citet[p.~702]{douglas1980extensions} supplies the exactness framework; \citet{browndouglasfillmore1977extensions} is the full original treatment.}
\GSfield{Present status}{As of 9 November 2026: established in the stated scope.}
\GSfield{Assessment}{literal statement: correct with missing hypotheses; publication match: exact; verification: read the result; development: none; confidence: high}
\end{GSitem}
\begin{GSitem}{1973-306.5}{E}{The family index pairing}
\GSfield{Kind; source}{construction; \GSpage{P000513}}
\GSfield{The notes}{\(K^1(X\times Y)\times\operatorname{Ext}_1(X)\to K^0(Y)\).}
\GSfield{Statement}{For finite CW complexes \(X,Y\), an odd K-cohomology class over \(X\times Y\) and an extension class in \(K_1(X)\) give a class in \(K^0(Y)\): apply the Busby map to a matrix representative, lift to a continuous Fredholm family, and take its index bundle. For \(Y=\mathrm{pt}\) this is the integer Fredholm-index pairing.}
\GSfield{Reconstruction}{Finite complexes, stabilized representatives and norm-continuous local lifts are needed for the family index; the compressed first line in the notebook does not fix degrees.}
\GSfield{Fate}{The scalar pairing and the K-homology identification are checked in \citet[pp.~702--703]{douglas1980extensions}; the family refinement is a reconstruction of the explicit notebook formula, not a separately located theorem in the accessible 1973 text.}
\GSfield{Present status}{As of 9 November 2026: established in the stated scope.}
\GSfield{Assessment}{literal statement: correct with missing hypotheses; publication match: partial; verification: indirect; development: none; confidence: high}
\end{GSitem}
\begin{GSitem}{1973-306.6}{E}{The garbled relation between Ext and degree-zero K-homology}
\GSfield{Kind; source}{claim; \GSpage{P000513}}
\GSfield{The notes}{The boxed formula is transcribed \(\operatorname{Ext}_1\oplus\mathbb Z=K_0\) for simplicial complexes; the operation is overwritten, but the subscript 1 is visible in the scan.}
\GSfield{Statement}{With the suspension convention on P000512, for a connected based finite CW complex \(X\) the correct degree-zero formula is \(\operatorname{Ext}(\Sigma X)\oplus\mathbb Z\cong K_0(X)\). It cannot be \(\operatorname{Ext}(X)\oplus\mathbb Z\): at \(X=S^1\) that would incorrectly identify \(\mathbb Z^2\) with \(\mathbb Z\).}
\GSfield{Reconstruction}{This is a degree correction, not a proposed decipherment of the overwritten operation. The added summand is the unreduced basepoint summand.}
\GSfield{Fate}{\citet[p.~702]{douglas1980extensions} explicitly identifies \(E_0(X)\) with \(K_0(X)\).}
\GSfield{Present status}{As of 9 November 2026: established in the stated scope.}
\GSfield{Assessment}{literal statement: false as written; publication match: partial; verification: read the result; development: annotated; confidence: high}
\end{GSitem}
\begin{GSitem}{1973-306.7}{F}{The affine-group extension and two indices}
\GSfield{Kind; source}{question; \GSpage{P000513}}
\GSfield{The notes}{\(0\to\mathcal K\to C^*(ax+b)\to C(\infty)\to0\), \(\operatorname{Ext}(\infty)=\mathbb Z+\mathbb Z\); the notes ask for the pair for a specific operator.}
\GSfield{Statement}{If the infinity-shaped symbol denotes the figure-eight graph \(S^1\vee S^1\), its Ext group is \(\mathbb Z^2\), with two index coordinates. But the group version of \(ax+b\), its representation, the precise ideal and the operator whose indices are requested are not recorded, so no numerical pair is determined.}
\GSfield{Reconstruction}{The figure-eight interpretation is conditional. It is not substituted for the source symbol as an established transcription correction.}
\GSfield{Fate}{The planar index classification in \citet[pp.~701--702]{douglas1980extensions} explains the graph interpretation. It does not identify this particular affine-group extension. The scan confirms an infinity-shaped symbol, not a written definition of the space.}
\GSfield{Present status}{As of 9 November 2026: the numerical question is underdetermined; no problem-index entry or claim of a remaining research problem is justified.}
\GSfield{Assessment}{literal statement: not determinable; publication match: thematic only; verification: unresolved lead; development: none; confidence: low}
\end{GSitem}
\begin{GSitem}{1973-306.8}{E}{Extensions of the circle are classified by the index}
\GSfield{Kind; source}{theorem; \GSpage{P000513}}
\GSfield{The notes}{For \(X=\mathbb T\), the image \(T\) of the generator satisfies \(TT^*=1\bmod\mathcal K\); the answer is the index.}
\GSfield{Statement}{A unital Busby map \(C(S^1)\to\mathcal Q(H)\) sends the coordinate function to a unitary in the Calkin algebra. Every lift \(T\) satisfies both \(TT^*-1\in\mathcal K\) and \(T^*T-1\in\mathcal K\), is Fredholm, and determines the extension by its index. Thus \(\operatorname{Ext}(S^1)\cong\mathbb Z\), with sign depending on the choice of generator.}
\GSfield{Reconstruction}{The star homomorphism supplies the second unitary relation. A single displayed relation permits a coisometry whose kernel has infinite dimension. Such an operator is not Fredholm.}
\GSfield{Fate}{\citet[pp.~701--702]{douglas1980extensions}; compare the equivalent normal-operator classification in \citet[Theorem~1.1 and following paragraph]{friisrordam2001approximation}.}
\GSfield{Present status}{As of 9 November 2026: established in the stated scope.}
\GSfield{Assessment}{literal statement: correct with missing hypotheses; publication match: exact; verification: read the result; development: none; confidence: high}
\end{GSitem}
\begin{GSitem}{1973-306.9}{E}{Normal operators plus compact operators form a closed set}
\GSfield{Kind; source}{theorem; \GSpage{P000513}}
\GSfield{The notes}{\(\{N+K:N\text{ normal},K\text{ compact}\}\) is norm closed.}
\GSfield{Statement}{On an infinite-dimensional separable complex Hilbert space, the set of compact perturbations of bounded normal operators is closed in operator norm. Equivalently, an essentially normal operator belongs to it exactly when \(\operatorname{ind}(T-\lambda)=0\) for every \(\lambda\notin\sigma_{\rm ess}(T)\).}
\GSfield{Fate}{The 1973 BDF result is reproved in \citet[Theorem~1.1 and \S2, proof of Theorem~1.1]{friisrordam2001approximation}; the primary article identifies BDF Corollary~11.2 as the original source.}
\GSfield{Present status}{As of 9 November 2026: established in the stated scope.}
\GSfield{Assessment}{literal statement: correct; publication match: exact; verification: read the result; development: none; confidence: high}
\end{GSitem}
\begin{GSitem}{1973-306.10}{F}{The asserted failure of closure for pairs}
\GSfield{Kind; source}{claim; \GSpage{P000513}}
\GSfield{The notes}{The circled insertion says that pairs \((N_1+K_1,N_2+K_2)\), with commuting normal \(N_1,N_2\), are not closed among pairs of this kind.}
\GSfield{Statement}{One possible reading is nonclosure in \(\mathcal B(H)^2\) of compact perturbations of commuting bounded normal pairs. The notebook gives no example and its phrase ``all pairs of this kind leaves the ambient class unclear. This is kept separate from the verified one-operator theorem.}
\GSfield{Fate}{Exact-phrase and BDF-title searches, the available 1973 metadata, \citet[pp.~699--704]{douglas1980extensions}, and \citet[\S2]{friisrordam2001approximation} did not locate the asserted pair counterexample. Neither the one-variable theorem nor general Ext theory alone verifies it.}
\GSfield{Present status}{As of 9 November 2026: a recorded claim whose precise ambient hypotheses and primary counterexample remain unresolved in this audit; it is not classified as an unsolved research problem.}
\GSfield{Assessment}{literal statement: not determinable; publication match: none found; verification: unresolved lead; development: none; confidence: low}
\end{GSitem}

\GStalk{308}{1973-09-24}{D. B. Fuchs}{P000512-P000518}{Quillen I: simplicial objects, classifying spaces and the Q-construction}
The Fuchs heading starts on P000513 after the BDF discussion; this first installment ends on P000516. It gives a detailed introduction to simplicial sets and their realization, the nerve of a category and the classifying space of a group, before introducing exact categories and \(K_i(M)=\pi_{i+1}BQM\). The explicitly named source is \citet{quillen1973higher}, prepared in January--February 1973 after the Seattle 1972 conference. The scan faithfully contains two incorrect pedagogical sentences about realization, which are corrected separately from the valid construction.

\begin{GSitem}{1973-308.1}{E}{Simplicial objects and the simplex category}
\GSfield{Kind; source}{definition; \GSpage{P000513-P000514}}
\GSfield{The notes}{A simplicial object is a contravariant functor from Ord, whose morphisms are nondecreasing maps of \(\{0,\ldots,p\}\) to \(\{0,\ldots,q\}\).}
\GSfield{Statement}{Let \(\Delta\) have objects \([n]=\{0,\ldots,n\}\), \(n\ge0\), and order-preserving maps. A simplicial object in a category \(\mathcal S\) is a functor \(\Delta^{\mathrm{op}}\to\mathcal S\). The elementary injections omit one vertex; the elementary surjections identify adjacent vertices.}
\GSfield{Fate}{\citet[\S1, p.~105]{segal1968classifying}; \citet[\S1]{quillen1973higher}.}
\GSfield{Present status}{As of 9 November 2026: established in the stated scope.}
\GSfield{Assessment}{literal statement: correct; publication match: exact; verification: read the result; development: none; confidence: high}
\end{GSitem}
\begin{GSitem}{1973-308.2}{E}{The face and degeneracy presentation}
\GSfield{Kind; source}{claim; \GSpage{P000514}}
\GSfield{The notes}{The two ladders have \(q+1\) face and \(q+1\) degeneracy arrows; the relations are said to be easy to write down.}
\GSfield{Statement}{The structure consists of \(d_i:S_n\to S_{n-1}\) and \(s_i:S_n\to S_{n+1}\), \(0\le i\le n\), subject to \(d_id_j=d_{j-1}d_i\) for \(i<j\), \(s_is_j=s_{j+1}s_i\) for \(i\le j\), and \(d_is_j=s_{j-1}d_i\) for \(i<j\), \(d_js_j=d_{j+1}s_j=1\), \(d_is_j=s_jd_{i-1}\) for \(i>j+1\).}
\GSfield{Reconstruction}{The relations omitted by the notebook are supplied explicitly.}
\GSfield{Fate}{This expands the functorial definition in \citet[\S1]{segal1968classifying}; no additional result is attributed to the speaker.}
\GSfield{Present status}{As of 9 November 2026: established in the stated scope.}
\GSfield{Assessment}{literal statement: correct; publication match: exact; verification: read the result; development: none; confidence: high}
\end{GSitem}
\begin{GSitem}{1973-308.3}{E}{The singular simplicial set}
\GSfield{Kind; source}{construction; \GSpage{P000514}}
\GSfield{The notes}{\(S_q\) is the set of continuous maps \(\Delta^q\to X\); the face maps restrict and the degeneracy maps precompose with affine projections.}
\GSfield{Statement}{For a topological space \(X\), put \(\operatorname{Sing}(X)_q=C(\Delta^q,X)\), considered as a set. Precomposition with affine maps induced by \(\Delta\)-morphisms gives the simplicial structure. This is different from topologizing these function spaces as a simplicial space.}
\GSfield{Fate}{The construction is an instance of \citet[\S1]{segal1968classifying}; the distinction between sets and spaces is explicit there.}
\GSfield{Present status}{As of 9 November 2026: established in the stated scope.}
\GSfield{Assessment}{literal statement: correct; publication match: exact; verification: read the result; development: none; confidence: high}
\end{GSitem}
\begin{GSitem}{1973-308.4}{E}{Geometric realization and its CW cells}
\GSfield{Kind; source}{construction; \GSpage{P000515}}
\GSfield{The notes}{The notebook glues copies of \(\Delta^q\) for the nondegenerate simplices and calls the result a cell space.}
\GSfield{Statement}{For a simplicial set \(S\), \(|S|=(\coprod_{q\ge0}S_q\times\Delta^q)/\sim\), where \((S(\alpha)s,t)\sim(s,\alpha_*t)\). This is a CW complex with one open q-cell for each nondegenerate q-simplex. Attaching maps may identify boundary points and may map a face through a lower-dimensional degeneracy.}
\GSfield{Reconstruction}{The all-simplices quotient avoids the notebook\textquoteright s incorrect assumption that every face of a nondegenerate simplex stays nondegenerate.}
\GSfield{Fate}{\citet[\S1, p.~105]{segal1968classifying}; \citet[\S1, realization preliminaries]{quillen1973higher}.}
\GSfield{Present status}{As of 9 November 2026: established in the stated scope.}
\GSfield{Assessment}{literal statement: correct with missing hypotheses; publication match: exact; verification: read the result; development: none; confidence: high}
\end{GSitem}
\begin{GSitem}{1973-308.5}{E}{Faces of nondegenerate simplices can be degenerate}
\GSfield{Kind; source}{claim; \GSpage{P000515}}
\GSfield{The notes}{The handwritten sentence says that arrows from right to left take nondegenerate simplices to nondegenerate simplices.}
\GSfield{Statement}{This is false for general simplicial sets. In the nerve of the group \(C_2=\{1,g\}\), the 2-simplex \((g,g)\) is nondegenerate, but its middle face is \(g^2=1\), the degenerate 1-simplex. Thus the page\textquoteright s own forthcoming group example supplies a counterexample.}
\GSfield{Reconstruction}{The counterexample is editorial; the erroneous sentence is visibly in the scan and is not blamed on the transcription.}
\GSfield{Fate}{The nerve and its degeneracies are defined in \citet[\S2--3]{segal1968classifying}. \ours\ The counterexample uses these definitions directly and was independently checked against the face and degeneracy formulas.}
\GSfield{Present status}{As of 9 November 2026: established in the stated scope.}
\GSfield{Assessment}{literal statement: false as written; publication match: partial; verification: read the result; development: annotated; confidence: high}
\end{GSitem}
\begin{GSitem}{1973-308.6}{E}{Realizing a simplicial space retains a topology}
\GSfield{Kind; source}{claim; \GSpage{P000515}}
\GSfield{The notes}{The note says a topological structure no longer arises for simplicial topological spaces, because the pieces have product topologies.}
\GSfield{Statement}{The same coend construction for a simplicial topological space has a quotient topology. What need not hold without additional conditions is the previous description as a CW complex with discrete cells indexed by simplices. If every level is a suitable compactly generated space, realization is taken in that category.}
\GSfield{Reconstruction}{Replacing ``topological structure by ``the preceding cell structure is a mathematical emendation, not a reading supplied by the manuscript.}
\GSfield{Fate}{\citet[\S1, p.~106]{segal1968classifying} explicitly defines realization of simplicial spaces.}
\GSfield{Present status}{As of 9 November 2026: established in the stated scope.}
\GSfield{Assessment}{literal statement: false as written; publication match: exact; verification: read the result; development: annotated; confidence: high}
\end{GSitem}
\begin{GSitem}{1973-308.7}{E}{The singular realization comparison}
\GSfield{Kind; source}{theorem; \GSpage{P000515}}
\GSfield{The notes}{For a good space, the realization of its singular complex is homotopy equivalent to the original space.}
\GSfield{Statement}{The evaluation map \(|\operatorname{Sing}X|\to X\), sending a simplex and barycentric point to its value, is a weak homotopy equivalence. If \(X\) has the homotopy type of a CW complex, it is a homotopy equivalence.}
\GSfield{Reconstruction}{``Good is replaced by the explicit CW-type condition. No homotopy-equivalence assertion is made for arbitrary spaces.}
\GSfield{Fate}{The usual comparison between realization and singular complexes underlies the constructions of classifying spaces in \citet[Chapter~IV, \S3]{weibel2013kbook}.}
\GSfield{Present status}{As of 9 November 2026: established in the stated scope.}
\GSfield{Assessment}{literal statement: correct with missing hypotheses; publication match: exact; verification: indirect; development: none; confidence: high}
\end{GSitem}
\begin{GSitem}{1973-308.8}{E}{An ordered simplicial model of a polyhedron}
\GSfield{Kind; source}{construction; \GSpage{P000515}}
\GSfield{The notes}{To reconstruct a polyhedron, the notes propose affine maps of simplices into it.}
\GSfield{Statement}{For an ordered simplicial complex \(K\), use the simplicial set whose q-simplices are weakly increasing lists of vertices lying in one simplex, allowing repetitions. Its realization is canonically homeomorphic to \(|K|\). Arbitrary affine maps with arbitrary vertex images are a larger construction and do not give this finite combinatorial description.}
\GSfield{Reconstruction}{The vertex-preserving and order-preserving restrictions are supplied.}
\GSfield{Fate}{\citet[\S4, p.~108]{segal1968classifying} stresses the need to order vertices and the natural ordering of barycentric subdivision.}
\GSfield{Present status}{As of 9 November 2026: established in the stated scope.}
\GSfield{Assessment}{literal statement: correct with missing hypotheses; publication match: exact; verification: read the result; development: none; confidence: high}
\end{GSitem}
\begin{GSitem}{1973-308.9}{E}{Realization of products and the topology convention}
\GSfield{Kind; source}{theorem; \GSpage{P000515}}
\GSfield{The notes}{\(\operatorname{real}(S\times S^{\prime})\to|S|\times|S^{\prime}|\) is asserted to be a homeomorphism.}
\GSfield{Statement}{For simplicial sets the natural map is a homeomorphism when the target product is formed in compactly generated spaces. It is also a homeomorphism with the ordinary product in the standard finite cases used here. The unqualified statement with arbitrary ordinary topological products needs this restriction.}
\GSfield{Reconstruction}{The notebook does not specify the category of spaces.}
\GSfield{Fate}{\citet[\S1, p.~106]{segal1968classifying} and \citet[\S1, realization preliminaries]{quillen1973higher} state the k-space qualification.}
\GSfield{Present status}{As of 9 November 2026: established in the stated scope.}
\GSfield{Assessment}{literal statement: correct with missing hypotheses; publication match: exact; verification: read the result; development: none; confidence: high}
\end{GSitem}
\begin{GSitem}{1973-308.10}{E}{The simplicial interval}
\GSfield{Kind; source}{construction; \GSpage{P000515}}
\GSfield{The notes}{For the ordered segment, \(S_q\) has \(q+2\) elements, given by a switch from 0 to 1.}
\GSfield{Statement}{\(\Delta[1]_q=\operatorname{Hom}_{\Delta}([q],[1])\) consists of the q+2 weakly increasing binary strings. There are two nondegenerate 0-simplices and one nondegenerate 1-simplex, and no others. Its realization is the closed interval.}
\GSfield{Fate}{The interval-category construction in \citet[Proposition~2.1]{segal1968classifying} is exactly this model.}
\GSfield{Present status}{As of 9 November 2026: established in the stated scope.}
\GSfield{Assessment}{literal statement: correct; publication match: exact; verification: read the result; development: none; confidence: high}
\end{GSitem}
\begin{GSitem}{1973-308.11}{E}{The nerve and classifying space of a small category}
\GSfield{Kind; source}{definition; \GSpage{P000516}}
\GSfield{The notes}{\((NC)_q\) is the set of composable strings of q morphisms; faces discard endpoints or compose arrows, and degeneracies insert identities.}
\GSfield{Statement}{For a small category \(C\), \((NC)_q=\operatorname{Fun}([q],C)\) and \(BC=|NC|\). For \(0<i<q\), \(d_i\) composes the arrows adjacent to \(X_i\); \(d_0,d_q\) delete the first or last object; \(s_i\) repeats \(X_i\) with its identity arrow.}
\GSfield{Fate}{\citet[\S2]{segal1968classifying}; \citet[\S1]{quillen1973higher}.}
\GSfield{Present status}{As of 9 November 2026: established in the stated scope.}
\GSfield{Assessment}{literal statement: correct; publication match: exact; verification: read the result; development: none; confidence: high}
\end{GSitem}
\begin{GSitem}{1973-308.12}{E}{Groups as one-object categories and the C2 example}
\GSfield{Kind; source}{construction; \GSpage{P000516}}
\GSfield{The notes}{A group is a one-object category, \((NG)_q=G^q\). For \(G=\mathbb Z_2\), there is one nondegenerate simplex in each dimension and the realization is \(\mathbb RP^{\infty}\).}
\GSfield{Statement}{For a discrete group \(G\), \(BG\) is a \(K(G,1)\). For \(C_2\), a string is nondegenerate exactly when none of its entries is the identity, so there is one open cell in each dimension. Its realization has the homotopy type of \(\mathbb RP^{\infty}\), the standard classifying space for \(C_2\).}
\GSfield{Reconstruction}{The robust comparison is stated up to homotopy; the notebook does not construct a specific homeomorphism of CW models.}
\GSfield{Fate}{\citet[\S3]{segal1968classifying}.}
\GSfield{Present status}{As of 9 November 2026: established in the stated scope.}
\GSfield{Assessment}{literal statement: correct; publication match: exact; verification: read the result; development: none; confidence: high}
\end{GSitem}
\begin{GSitem}{1973-308.13}{E}{The first definition of the Q-construction and higher K-groups}
\GSfield{Kind; source}{construction; \GSpage{P000513; P000516}}
\GSfield{The notes}{An abelian category or an extension-closed full subcategory \(M\) gives a category \(QM\) of spans and \(K_i(M)=\pi_{i+1}(BQM)\).}
\GSfield{Statement}{Let \(M\) be an essentially small exact category. The objects of \(QM\) are those of \(M\); a morphism \(X\to Y\) is an isomorphism class of spans \(X\leftarrow E\to Y\) whose left map is an admissible epimorphism and right map an admissible monomorphism. Composition is by pullback. With basepoint the zero object, \(K_i(M)=\pi_{i+1}(BQM,0)\), \(i\ge0\).}
\GSfield{Reconstruction}{Admissibility, middle-object isomorphism and the basepoint are made explicit; record 311 gives the detailed occurrence.}
\GSfield{Fate}{\citet[\S2, pp.~99--103]{quillen1973higher}. The manuscript was prepared in early 1973; this is an exposition of it.}
\GSfield{Present status}{As of 9 November 2026: established in the stated scope.}
\GSfield{Assessment}{literal statement: correct with missing hypotheses; publication match: exact; verification: read the result; development: none; confidence: high}
\end{GSitem}

\GStalk{309}{1973-10-08}{D. B. Fuchs}{P000516-P000522}{Categorical homotopy fibers and Mishchenko Fredholm representations}
The continuation begins on P000517 (the page says following Marchenko, without a separate mathematical block attributable to Marchenko). Fuchs checks the square as a product of simplicial intervals, derives homotopies from natural transformations, and explains Quillen Theorem B using comma categories. P000519 explicitly changes to Mishchenko and continues through P000520: Fredholm representations, family-index pairings, and a K-class on a group classifying space. Because the catalogue supplies no separate Mishchenko record, those items are retained here with a speaker override. These are largely earlier constructions, not new unsolved questions.

\begin{GSitem}{1973-309.1}{E}{The product of two simplicial intervals}
\GSfield{Kind; source}{construction; \GSpage{P000517}}
\GSfield{The notes}{The notes list the nondegenerate simplices of \(\Delta[1]\times\Delta[1]\), draw a diagonal square, and give a four-column array for the two triangles.}
\GSfield{Statement}{A q-simplex is a pair of increasing maps \([q]\to[1]\); it is degenerate exactly when both maps repeat the same adjacent pair. There are four nondegenerate vertices, five nondegenerate edges and two nondegenerate triangles, giving the usual triangulation of \([0,1]^2\). No simplex of larger dimension is nondegenerate.}
\GSfield{Reconstruction}{The four columns are read as two consecutive two-column triangles; no extra 2-simplices are inferred.}
\GSfield{Fate}{This directly instantiates the product-realization construction of \citet[\S1]{segal1968classifying}.}
\GSfield{Present status}{As of 9 November 2026: established in the stated scope.}
\GSfield{Assessment}{literal statement: correct; publication match: exact; verification: read the result; development: annotated; confidence: high}
\end{GSitem}
\begin{GSitem}{1973-309.2}{E}{The nerve revisited}
\GSfield{Kind; source}{definition; \GSpage{P000517}}
\GSfield{The notes}{The definition \(BC=\operatorname{Real}NC\) is repeated, with objects, nonidentity arrows and commutative triangles as its first cells.}
\GSfield{Statement}{For a small category, the nondegenerate q-simplices are composable q-tuples of nonidentity arrows. Composite arrows occurring as faces may be identities; this does not alter nondegeneracy of the original tuple.}
\GSfield{Fate}{\citet[\S2]{segal1968classifying}; compare the correction in item 1973-308.5.}
\GSfield{Present status}{As of 9 November 2026: established in the stated scope.}
\GSfield{Assessment}{literal statement: correct; publication match: exact; verification: read the result; development: none; confidence: high}
\end{GSitem}
\begin{GSitem}{1973-309.3}{E}{Barycentric subdivision as a nerve}
\GSfield{Kind; source}{claim; \GSpage{P000517}}
\GSfield{The notes}{Barycentric subdivision of a polyhedron of high degree is said to be the nerve of a category.}
\GSfield{Statement}{For any simplicial complex \(K\), the poset of its nonempty simplices ordered by face inclusion is a small category; its nerve realizes to the barycentric subdivision of \(K\), hence to \(|K|\). No high-dimensional or repeated-subdivision hypothesis is needed.}
\GSfield{Reconstruction}{The face-poset category identifies the unnamed category; the unclear phrase about degree is not treated as a necessary mathematical condition.}
\GSfield{Fate}{\citet[\S4, p.~108]{segal1968classifying}.}
\GSfield{Present status}{As of 9 November 2026: established in the stated scope.}
\GSfield{Assessment}{literal statement: correct; publication match: exact; verification: read the result; development: none; confidence: high}
\end{GSitem}
\begin{GSitem}{1973-309.4}{E}{Natural transformations induce homotopies}
\GSfield{Kind; source}{theorem; \GSpage{P000517-P000518}}
\GSfield{The notes}{\(B(C\times C^{\prime})=BC\times BC^{\prime}\). A natural transformation \(f\Rightarrow g\) is a functor \(C\times\{0\to1\}\to C^{\prime}\), hence a homotopy.}
\GSfield{Statement}{In compactly generated spaces, the classifying-space functor preserves finite products. A natural transformation determines \(BC\times[0,1]\to BC^{\prime}\), restricting to \(Bf\) and \(Bg\). For this interval product the usual topology already suffices.}
\GSfield{Reconstruction}{The category \(\{0\to1\}\) is distinguished from its realized interval.}
\GSfield{Fate}{\citet[Proposition~2.1 and following remark]{segal1968classifying}; \citet[\S1, Proposition~2]{quillen1973higher}.}
\GSfield{Present status}{As of 9 November 2026: established in the stated scope.}
\GSfield{Assessment}{literal statement: correct with missing hypotheses; publication match: exact; verification: read the result; development: none; confidence: high}
\end{GSitem}
\begin{GSitem}{1973-309.5}{E}{Initial and terminal objects give contractions}
\GSfield{Kind; source}{theorem; \GSpage{P000518}}
\GSfield{The notes}{If \(C\) has an initial or terminal object, \(BC\) is contractible; it is called a cone.}
\GSfield{Statement}{A small category with an initial or terminal object has contractible classifying space. The unique arrows give a natural transformation between the identity and the corresponding constant functor, hence a contraction.}
\GSfield{Reconstruction}{Contractibility is the invariant conclusion; a literal identification with a particular topological cone is unnecessary.}
\GSfield{Fate}{\citet[\S1, Corollary~2 to Proposition~2]{quillen1973higher}.}
\GSfield{Present status}{As of 9 November 2026: established in the stated scope.}
\GSfield{Assessment}{literal statement: correct; publication match: exact; verification: read the result; development: none; confidence: high}
\end{GSitem}
\begin{GSitem}{1973-309.6}{E}{Adjoints give homotopy equivalences}
\GSfield{Kind; source}{theorem; \GSpage{P000518}}
\GSfield{The notes}{A functor with a left or a right adjoint induces a homotopy equivalence.}
\GSfield{Statement}{If \(F:C\to D\) has either adjoint \(G\), the unit and counit induce homotopies showing \(BF\) and \(BG\) to be homotopy inverses.}
\GSfield{Fate}{\citet[\S1, Corollary~1 to Proposition~2]{quillen1973higher}.}
\GSfield{Present status}{As of 9 November 2026: established in the stated scope.}
\GSfield{Assessment}{literal statement: correct; publication match: exact; verification: read the result; development: none; confidence: high}
\end{GSitem}
\begin{GSitem}{1973-309.7}{E}{Opposite categories have the same classifying space}
\GSfield{Kind; source}{theorem; \GSpage{P000518}}
\GSfield{The notes}{\(BC^{\circ}=BC\).}
\GSfield{Statement}{For a small category, reversing each composable string and reversing its barycentric coordinates defines a canonical homeomorphism \(B(C^{\mathrm{op}})\cong BC\). This does not assert that the nerves are naturally isomorphic as simplicial sets.}
\GSfield{Reconstruction}{Equality is replaced by the actual canonical homeomorphism.}
\GSfield{Fate}{\citet[\S1, opposite-category construction]{quillen1973higher}.}
\GSfield{Present status}{As of 9 November 2026: established in the stated scope.}
\GSfield{Assessment}{literal statement: correct; publication match: exact; verification: read the result; development: none; confidence: high}
\end{GSitem}
\begin{GSitem}{1973-309.8}{E}{The homotopy sequence of a fibration}
\GSfield{Kind; source}{theorem; \GSpage{P000518}}
\GSfield{The notes}{The homotopy sequence of a Serre fibration is written, with an unmatched parenthesis in the fiber term.}
\GSfield{Statement}{For a pointed Serre fibration \(f:(X,x)\to(Y,y)\) with fiber \(F=f^{-1}(y)\), there is the exact sequence \(\cdots\to\pi_q(F,x)\to\pi_q(X,x)\to\pi_q(Y,y)\to\pi_{q-1}(F,x)\to\cdots\); in low degrees it is an exact sequence of pointed sets with the usual action interpretation.}
\GSfield{Reconstruction}{Basepoints and the low-degree qualification are supplied.}
\GSfield{Fate}{This is the standard homotopy-fiber convention used by \citet[\S1, Theorem~B]{quillen1973higher}.}
\GSfield{Present status}{As of 9 November 2026: established in the stated scope.}
\GSfield{Assessment}{literal statement: correct with missing hypotheses; publication match: exact; verification: read the result; development: none; confidence: high}
\end{GSitem}
\begin{GSitem}{1973-309.9}{E}{Mapping paths replace an arbitrary map by a fibration}
\GSfield{Kind; source}{construction; \GSpage{P000518}}
\GSfield{The notes}{\(\widetilde X=\{(x,s):s(0)=f(x)\}\), \(\widetilde f(x,s)=s(1)\), and \(\widetilde X\simeq X\).}
\GSfield{Statement}{For a continuous map between compactly generated spaces, give the path space the compact-open k-topology. The inclusion using constant paths \(X\to\widetilde X\) is a homotopy equivalence, and endpoint evaluation \(\widetilde f\) is a fibration. Its fiber at \(y\) is the homotopy fiber of \(f\).}
\GSfield{Reconstruction}{The path-space topology is stated explicitly.}
\GSfield{Fate}{The construction fixes the homotopy-fiber meaning in \citet[\S1, Theorem~B]{quillen1973higher}.}
\GSfield{Present status}{As of 9 November 2026: established in the stated scope.}
\GSfield{Assessment}{literal statement: correct with missing hypotheses; publication match: exact; verification: read the result; development: none; confidence: high}
\end{GSitem}
\begin{GSitem}{1973-309.10}{E}{The comma-category map to the homotopy fiber}
\GSfield{Kind; source}{construction; \GSpage{P000518-P000519}}
\GSfield{The notes}{\(Y\backslash g\) has pairs \((X,u:Y\to gX)\). On a simplex, adjoining \(Y\) produces a larger simplex in \(BC^{\prime}\), and a straight path supplies a point of the homotopy fiber.}
\GSfield{Statement}{For a functor \(g:C\to D\) of small categories and \(Y\in D\), let \((Y\downarrow g)\) have objects \((X,u:Y\to gX)\) and arrows \(a:X\to X^{\prime}\) satisfying \(g(a)u=u^{\prime}\). Forgetting \(u\) gives a map to \(BC\), and the natural arrows from constant \(Y\) to \(gX\) give a canonical map \(B(Y\downarrow g)\to\operatorname{hofib}_Y(Bg)\).}
\GSfield{Reconstruction}{The notebook uses a reversed chain convention on P000519; the natural-transformation description makes the path compatible with faces and degeneracies.}
\GSfield{Fate}{\citet[\S1, pp.~93--97]{quillen1973higher}.}
\GSfield{Present status}{As of 9 November 2026: established in the stated scope.}
\GSfield{Assessment}{literal statement: correct; publication match: exact; verification: read the result; development: none; confidence: high}
\end{GSitem}
\begin{GSitem}{1973-309.11}{E}{Quillen Theorem B}
\GSfield{Kind; source}{theorem; \GSpage{P000518-P000519}}
\GSfield{The notes}{If every target arrow \(v:Y\to Y^{\prime}\) induces a homotopy equivalence \(B(Y^{\prime}\backslash g)\to B(Y\backslash g)\), the comma category models the homotopy fiber.}
\GSfield{Statement}{Assume for every morphism \(v:Y\to Y^{\prime}\) in \(D\) that precomposition \(v^*:(Y^{\prime}\downarrow g)\to(Y\downarrow g)\) induces a homotopy equivalence on classifying spaces. Then \(B(Y\downarrow g)\to\operatorname{hofib}_Y(Bg)\) is a homotopy equivalence for every \(Y\), and yields the corresponding homotopy exact sequence.}
\GSfield{Fate}{\citet[Theorem~B, p.~97]{quillen1973higher}.}
\GSfield{Present status}{As of 9 November 2026: established in the stated scope.}
\GSfield{Assessment}{literal statement: correct; publication match: exact; verification: read the result; development: none; confidence: high}
\end{GSitem}
\begin{GSitem}{1973-309.12}{F}{A proposed limit over comma categories}
\GSfield{Kind; source}{remark; \GSpage{P000519}}
\GSfield{The notes}{Fuchs suggests that \(Y\backslash g\) should probably be defined as an inverse limit over all such categories with different \(Y\).}
\GSfield{Statement}{The notes do not specify an indexing category, the diagram, transition maps, ordinary versus homotopy limit, or the desired universal property. Theorem B already gives the precise assertion in the preceding item; this additional suggestion has no unique formal statement.}
\GSfield{Fate}{Comparison with \citet[\S1, Theorem~B and its proof]{quillen1973higher} does not identify an additional theorem with these quantifiers.}
\GSfield{Present status}{As of 9 November 2026: an exploratory comment, not a separate unresolved problem.}
\GSfield{Assessment}{literal statement: not determinable; publication match: thematic only; verification: unresolved lead; development: none; confidence: low}
\end{GSitem}
\begin{GSitem}{1973-309.13}{E}{Fredholm representations of a ring with involution}
\GSfield{Kind; source}{definition; \GSpage{P000519}; A. S. Mishchenko}
\GSfield{The notes}{\(\Lambda\) has involution, \(T_1,T_2\) act on Hilbert spaces, \(F\) is Fredholm, and \(FT_1(\lambda)-T_2(\lambda)F\) is compact for every \(\lambda\).}
\GSfield{Statement}{For a unital *-ring \(\Lambda\), a Fredholm representation consists of unital *-representations \(T_i:\Lambda\to\mathcal B(H_i)\) on separable complex Hilbert spaces and a bounded Fredholm operator \(F:H_1\to H_2\), such that \(FT_1(\lambda)-T_2(\lambda)F\) is compact for every \(\lambda\in\Lambda\). For \(\Lambda=\mathbb Z[\pi]\), this is equivalently a pair of unitary group representations with a compact-intertwining Fredholm operator.}
\GSfield{Reconstruction}{Boundedness, complex Hilbert spaces and preservation of involution are explicit.}
\GSfield{Fate}{\citet[Definition~1.1]{mishchenko1974signatures}; the paper was submitted December 1972 and its central construction was announced in 1973. The notebook generalizes the definition from group rings to *-rings. The public pre-talk announcement is \citet[\S\S1--3, p.~239]{mishchenko1973infinite}, published in the March--April 1973 issue and received 10 October 1972.}
\GSfield{Present status}{As of 9 November 2026: established in the stated scope.}
\GSfield{Assessment}{literal statement: correct with missing hypotheses; publication match: exact; verification: read the result; development: none; confidence: high}
\end{GSitem}
\begin{GSitem}{1973-309.14}{E}{Twisting a Fredholm cycle by a vector bundle}
\GSfield{Kind; source}{construction; \GSpage{P000519-P000520}; A. S. Mishchenko}
\GSfield{The notes}{For \(\Lambda=C(M)\), the notes write \(F\otimes1_\eta\) on \(H_i\otimes_{C(M)}\Gamma(\eta)\), then take the index.}
\GSfield{Statement}{For compact \(M\), represent a vector bundle \(\eta\) by a projection \(e\in M_n(C(M))\). The compressed operator \(T_2(e)F^{(n)}T_1(e):T_1(e)H_1^n\to T_2(e)H_2^n\) is Fredholm. Its index defines an additive homomorphism \(K^0(M)\to\mathbb Z\). Because F intertwines only modulo compacts, the naive balanced tensor operator need not be literally well defined.}
\GSfield{Reconstruction}{Projection compression makes the shorthand tensor construction precise and avoids treating a compactly intertwining operator as an exact module homomorphism.}
\GSfield{Fate}{The Fredholm-cycle pairing is recalled in \citet[p.~702]{douglas1980extensions}; the family-index formalism is in \citet[\S2]{atiyahsinger1971families}.}
\GSfield{Present status}{As of 9 November 2026: established in the stated scope.}
\GSfield{Assessment}{literal statement: correct with missing hypotheses; publication match: exact; verification: read the result; development: annotated; confidence: high}
\end{GSitem}
\begin{GSitem}{1973-309.15}{E}{Integer-valued pairings lose torsion information}
\GSfield{Kind; source}{claim; \GSpage{P000520}; A. S. Mishchenko}
\GSfield{The notes}{The first pairing is called an element of \(K_0(M)\), then corrected as too crude because \(K^0(M)\) can have torsion.}
\GSfield{Statement}{An integer-valued homomorphism annihilates torsion in \(K^0(M)\), and does not by itself specify every analytic K-homology class. The Fredholm cycle is the richer object; its numerical pairing is only one invariant.}
\GSfield{Reconstruction}{The initial identification is not retained as a definition of K-homology; it is qualified by the notebook itself.}
\GSfield{Fate}{The distinction between Fredholm cycles and their integer pairing is explicit in \citet[p.~702]{douglas1980extensions}.}
\GSfield{Present status}{As of 9 November 2026: established in the stated scope.}
\GSfield{Assessment}{literal statement: correct with missing hypotheses; publication match: exact; verification: read the result; development: none; confidence: high}
\end{GSitem}
\begin{GSitem}{1973-309.16}{E}{Families refine the numerical index}
\GSfield{Kind; source}{construction; \GSpage{P000520}; A. S. Mishchenko}
\GSfield{The notes}{For every \(Y\), \(T\) defines \(\operatorname{index}_T:K^0(M\times Y)\to K^0(Y)\), and this is said to be an element of \(K_0(M)\).}
\GSfield{Statement}{For compact finite CW complexes \(M,Y\), projection compression in a continuous family gives \(\operatorname{index}_T:K^0(M\times Y)\to K^0(Y)\), natural in Y. The Fredholm cycle defines an even analytic K-homology class, whose family pairings retain more structure than the numerical map at \(Y=\mathrm{pt}\).}
\GSfield{Reconstruction}{Compact parameter spaces and the index-bundle construction are supplied. A full characterization of all natural transformations is not asserted.}
\GSfield{Fate}{\citet[\S2, family-index construction]{atiyahsinger1971families}; \citet[p.~702]{douglas1980extensions} gives the cycle interpretation.}
\GSfield{Present status}{As of 9 November 2026: established in the stated scope.}
\GSfield{Assessment}{literal statement: correct with missing hypotheses; publication match: exact; verification: read the result; development: none; confidence: high}
\end{GSitem}
\begin{GSitem}{1973-309.17}{E}{Elliptic symbols and analytic K-homology}
\GSfield{Kind; source}{claim; \GSpage{P000520}; A. S. Mishchenko}
\GSfield{The notes}{An elliptic operator has symbol in \(K^0(T^*M)\), compact commutators with functions, and K-homology class dual to its symbol.}
\GSfield{Statement}{For a closed smooth manifold M and a bounded order-zero elliptic pseudodifferential operator between complex vector bundles, commutators with multiplication by smooth functions are of negative order and compact on L2. It therefore determines an even Fredholm cycle. The symbol lies in compactly supported \(K_c^0(T^*M)\); under the symbol/analytic K-homology identification, its image is this cycle. A positive-order differential operator first requires a bounded transform or Sobolev normalization.}
\GSfield{Reconstruction}{Compact supports, a closed manifold and bounded normalization are essential; the statement is not ordinary untwisted K-cohomological Poincare duality on an arbitrary manifold.}
\GSfield{Fate}{The symbol and family-index conventions are checked in \citet[Introduction and \S2]{atiyahsinger1971families}; the Fredholm-cycle framework is in \citet[p.~702]{douglas1980extensions}. The full duality identification is used here as the modern formulation of the notebook\textquoteright s Poincare-duality shorthand.}
\GSfield{Present status}{As of 9 November 2026: established in the stated scope.}
\GSfield{Assessment}{literal statement: correct with missing hypotheses; publication match: partial; verification: indirect; development: none; confidence: high}
\end{GSitem}
\begin{GSitem}{1973-309.18}{E}{The K-class of a Fredholm representation of a discrete group}
\GSfield{Kind; source}{construction; \GSpage{P000519-P000520}; A. S. Mishchenko}
\GSfield{The notes}{For \(\Lambda=\mathbb Z[\pi]\), a Fredholm representation determines \(\xi_T\in K^0(B\pi)\). The notebook gives the equivariant construction on the universal cover.}
\GSfield{Statement}{For a discrete group \(\pi\) and a CW model of \(B\pi\), choose on each orbit of vertices of \(E\pi\) the operator F and extend equivariantly by \(\widehat F(gx)=T_2(g)\widehat F(x)T_1(g)^{-1}\). All values differ by compacts; barycentric interpolation on simplices stays Fredholm and gives an equivariant norm-continuous family. It descends between Hilbert bundles over \(B\pi\). Trivializing them gives a class in representable \(K^0(B\pi)\), independent of the choices.}
\GSfield{Reconstruction}{For an infinite CW complex the K-group is representable K-theory. The stability of the Fredholm set under compact perturbations justifies each interpolation.}
\GSfield{Fate}{\citet[\S2, Lemma~2.1 and pp.~86--87]{mishchenko1974signatures}. The construction itself works beyond the finitely presented groups emphasized in that paper; the 1973 announcement predates this intervention. The public pre-talk announcement is \citet[\S\S1--3, p.~239]{mishchenko1973infinite}, published in the March--April 1973 issue and received 10 October 1972.}
\GSfield{Present status}{As of 9 November 2026: established in the stated scope.}
\GSfield{Assessment}{literal statement: correct with missing hypotheses; publication match: exact; verification: read the result; development: none; confidence: high}
\end{GSitem}
\begin{GSitem}{1973-309.19}{E}{The Fredholm space has one component for each index}
\GSfield{Kind; source}{claim; \GSpage{P000520}; A. S. Mishchenko}
\GSfield{The notes}{The notes use \(K(B\pi)=[B\pi,BU]\) and \(BU\simeq\mathcal F\), where \(\mathcal F\) is all Fredholm operators.}
\GSfield{Statement}{For infinite-dimensional separable complex H with the norm topology, \(\mathcal F(H)\simeq\mathbb Z\times BU\). The zero-index component alone has homotopy type BU. Consequently full K0 is represented by \(\mathbb Z\times BU\); reduced or fixed-index conventions must be stated to use BU alone.}
\GSfield{Reconstruction}{The missing integer component is an emendation. A general F in the notebook is not assumed to have index zero.}
\GSfield{Fate}{The family-index description in \citet[Introduction and \S2]{atiyahsinger1971families} and the Fredholm-family use in \citet[\S2]{mishchenko1974signatures} require the index components.}
\GSfield{Present status}{As of 9 November 2026: established in the stated scope.}
\GSfield{Assessment}{literal statement: false as written; publication match: exact; verification: read the result; development: annotated; confidence: high}
\end{GSitem}
\begin{GSitem}{1973-309.20}{E}{Generalized signatures of Hermitian forms}
\GSfield{Kind; source}{construction; \GSpage{P000519}; A. S. Mishchenko}
\GSfield{The notes}{For an arbitrary ring with involution and an invertible Hermitian matrix A, the notes announce an integer \(\operatorname{sign}_T A\).}
\GSfield{Statement}{Apply the two *-representations to an invertible Hermitian matrix \(A\in M_n(\Lambda)\). Their images are invertible self-adjoint operators; their positive and negative spectral subspaces are intertwined by \(F^{(n)}\) modulo compacts. The two diagonal compressions are Fredholm, and define \(\operatorname{sign}_T(A)=\operatorname{ind}F_+-\operatorname{ind}F_-\). For group rings this descends to the corresponding even Wall group in the stated signature degree.}
\GSfield{Reconstruction}{The notebook gives no formula; the positive/negative spectral-compression formula is supplied from the paper. No analogous Wall-group descent for every unspecified *-ring is inferred.}
\GSfield{Fate}{\citet[\S3, Definition~3.4 and Theorem~3.5]{mishchenko1974signatures} supplies the construction and its group-ring invariance. The public pre-talk announcement is \citet[\S\S1--3, p.~239]{mishchenko1973infinite}, published in the March--April 1973 issue and received 10 October 1972.}
\GSfield{Present status}{As of 9 November 2026: established in the stated scope.}
\GSfield{Assessment}{literal statement: correct with missing hypotheses; publication match: exact; verification: read the result; development: none; confidence: high}
\end{GSitem}

\GStalk{310}{1973-10-15}{A. N. Rudakov}{P000520-P000523}{Discrete representations of formal vector fields}
The dated talk is on P000521; the Kirillov intervention at the beginning of P000522 is kept with it and explicitly attributed. Rudakov introduces discrete continuous representations of formal vector fields, induction from finite jet quotients, and the minimal-height submodule. The high-height induction theorem and the differential-form exceptions match \citet[\S\S2--3,6,13]{rudakov1974irreducible}, received 9 October 1973 and published in 1974. Kirillov then suggests a diffeomorphism-group analogue, explicitly labelled unproved in the notebook. This is not evidence of a complete classification of arbitrary representations of the full diffeomorphism group.

\begin{GSitem}{1973-310.1}{E}{The formal vector-field algebra and its filtration}
\GSfield{Kind; source}{definition; \GSpage{P000521}}
\GSfield{The notes}{\(W_n=\sum_i k[[x_1,\ldots,x_n]]\partial_{x_i}\), with continuous inserted before derivations; \(\mathcal L_j\) is defined by degrees at least \(j+1\).}
\GSfield{Statement}{Over a field k of characteristic zero, let \(\mathfrak m=(x_1,\ldots,x_n)\) and \(W_n=\operatorname{Der}_{k,\rm cont} k[[x_1,\ldots,x_n]]\). For \(j\ge-1\), \(\mathcal L_j=\sum_i\mathfrak m^{j+1}\partial_{x_i}\). The order-of-vanishing filtration defines the linear topology; it is not the highest polynomial degree of a formal series.}
\GSfield{Reconstruction}{Characteristic zero and the meaning of degree for formal power series are supplied.}
\GSfield{Fate}{\citet[\S1, pp.~836--838]{rudakov1974irreducible}.}
\GSfield{Present status}{As of 9 November 2026: established in the stated scope.}
\GSfield{Assessment}{literal statement: correct with missing hypotheses; publication match: exact; verification: read the result; development: none; confidence: high}
\end{GSitem}
\begin{GSitem}{1973-310.2}{E}{Continuity into a discrete countable-dimensional module}
\GSfield{Kind; source}{definition; \GSpage{P000521}}
\GSfield{The notes}{The notes write \(\rho:\mathcal L\to\operatorname{Aut}E\), and require a neighborhood V with \(\rho(v)x=0\).}
\GSfield{Statement}{A continuous Lie-algebra representation is a homomorphism \(\rho:\mathcal L\to\operatorname{End}_k(E)\) for a discrete, countable-dimensional vector space E such that every \(x\in E\) is killed by some \(\mathcal L_j\). Equivalently every orbit map is locally constant at zero.}
\GSfield{Reconstruction}{Automorphisms is the wrong target for a Lie-algebra action: even the zero Lie element acts by the zero endomorphism, not an invertible map. The neighborhood quantifier means every v in V.}
\GSfield{Fate}{\citet[Lemma~2.1, p.~839]{rudakov1974irreducible}.}
\GSfield{Present status}{As of 9 November 2026: established in the stated scope.}
\GSfield{Assessment}{literal statement: false as written; publication match: exact; verification: read the result; development: none; confidence: high}
\end{GSitem}
\begin{GSitem}{1973-310.3}{E}{Inflation from a finite jet quotient followed by induction}
\GSfield{Kind; source}{construction; \GSpage{P000521}}
\GSfield{The notes}{A representation of \(M/M^{\prime}\), with \(M^{\prime}\) of finite codimension, is inflated to M and induced to \(\mathcal L\); later M is \(\mathcal L_0\).}
\GSfield{Statement}{For an open Lie subalgebra \(\mathcal L_0\), a continuous module V that factors through \(\mathcal L_0/\mathcal L_r\) inflates to \(\mathcal L_0\) and gives \(U(\mathcal L)\otimes_{U(\mathcal L_0)}V\). Under the stated countable-dimensional discrete hypotheses the induced representation is continuous. A finite-dimensional quotient algebra does not require V itself to be finite-dimensional.}
\GSfield{Reconstruction}{The abstract M in the opening programme is specialized to the open filtered subalgebra used in the theorem. No finite-dimensional-module hypothesis is added.}
\GSfield{Fate}{\citet[Definitions~2.2--2.3 and following observation, pp.~839--840]{rudakov1974irreducible}.}
\GSfield{Present status}{As of 9 November 2026: established in the stated scope.}
\GSfield{Assessment}{literal statement: correct with missing hypotheses; publication match: exact; verification: read the result; development: none; confidence: high}
\end{GSitem}
\begin{GSitem}{1973-310.4}{A}{Minimal height and the kernel submodule}
\GSfield{Kind; source}{construction; \GSpage{P000521}}
\GSfield{The notes}{\(p(x)=\min\{r:\mathcal L_rx=0\}\), \(\bar p(E)=\min_{x\ne0}p(x)\), and \(E_0=\{x:p(x)=\bar p(E)\}\) is asserted to be a submodule.}
\GSfield{Statement}{For nonzero x define p(x) by its annihilator in the filtration. Let h be the minimum over nonzero vectors. Then \(E_0=\{0\}\cup\{x\ne0:p(x)=h\}\) is a \(\mathcal L_0\)-submodule. Indeed p(x+y) is at most the maximum of the two heights, and \([\mathcal L_0,\mathcal L_h]\subseteq\mathcal L_h\).}
\GSfield{Reconstruction}{The printed paper explicitly adjoins zero; the notebook set omits it. The short argument here is a critical verification of the definition, not a new classification result.}
\GSfield{Fate}{\citet[Definition~3.1, Proposition~3.2 and Corollary~3.3, p.~840]{rudakov1974irreducible}.}
\GSfield{Present status}{As of 9 November 2026: established in the stated scope.}
\GSfield{Assessment}{literal statement: false as written; publication match: exact; verification: read the result; development: annotated; confidence: high}
\end{GSitem}
\begin{GSitem}{1973-310.5}{A}{Irreducibility and reconstruction above height one}
\GSfield{Kind; source}{theorem; \GSpage{P000521}}
\GSfield{The notes}{If \(\bar p(E)>1\), the representation is induced: \(E=\operatorname{Ind}(E_0)\); the handwritten bound is retraced.}
\GSfield{Statement}{For a Cartan-type algebra \(\mathcal L\) in Rudakov\textquoteright s filtration, every irreducible discrete continuous module E of height greater than one has irreducible \(\mathcal L_0\)-kernel \(E_0\), and \(U(\mathcal L)\otimes_{U(\mathcal L_0)}E_0\to E\) is an isomorphism. Conversely induction of an irreducible \(\mathcal L_0\)-module of height greater than one is irreducible.}
\GSfield{Reconstruction}{The primary paper resolves the provisional bound as 1. Continuity, irreducibility, topology and filtration are inherited from the preceding definitions; the assertion is not about all abstract Wn-modules.}
\GSfield{Fate}{\citet[Theorem~6.1 and Corollaries~6.2--6.3, pp.~843--844]{rudakov1974irreducible}; the manuscript was received six days before this talk and appeared in 1974.}
\GSfield{Present status}{As of 9 November 2026: established in the stated scope.}
\GSfield{Assessment}{literal statement: correct with missing hypotheses; publication match: exact; verification: read the result; development: none; confidence: high}
\end{GSitem}
\begin{GSitem}{1973-310.6}{A}{The differential-form exceptions for W\_n}
\GSfield{Kind; source}{theorem; \GSpage{P000521}}
\GSfield{The notes}{The opening says all irreducibles are induced with finitely many exceptions; the closing identifies \((d\Omega^p)^*\).}
\GSfield{Statement}{For \(W_n\), the exceptional height-one inducing \(\mathcal L_0\)-modules are the n modules inflated from the \(\mathfrak{gl}_n\)-modules with highest weights \((0,\ldots,0,-1,\ldots,-1)\), having p entries equal to \(-1\), \(1\le p\le n\). Each exceptional induced module has a unique maximal proper submodule; its irreducible quotient is the continuous dual of \(d\Omega^{p-1}\), the formal exact p-forms. Other irreducible height-one inducing modules give irreducible induced modules.}
\GSfield{Reconstruction}{The exact height-one statement is given instead of extending the shorthand to all unrestricted module categories. The dual is the continuous dual of the linearly compact formal-form space, hence discrete; the trivial representation is separately understood. Closed positive-degree formal forms are exact in characteristic zero.}
\GSfield{Fate}{\citet[Definition~13.6, Theorem~13.7 and Corollaries~13.8--13.9, pp.~865--866]{rudakov1974irreducible}.}
\GSfield{Present status}{As of 9 November 2026: established in the stated scope.}
\GSfield{Assessment}{literal statement: correct with missing hypotheses; publication match: exact; verification: read the result; development: none; confidence: high}
\end{GSitem}
\begin{GSitem}{1973-310.7}{A}{The corresponding Hamiltonian result}
\GSfield{Kind; source}{claim; \GSpage{P000521}}
\GSfield{The notes}{The last line says there are corresponding results for Hamiltonian ones as well.}
\GSfield{Statement}{For \(H_{2m}\), using its induced formal filtration, the same high-height induction and irreducibility statements hold. The notebook does not state a list of height-one Hamiltonian exceptions.}
\GSfield{Reconstruction}{Only the precise high-height assertion supported by the paper is identified. The Wn exceptional list is not transferred to the Hamiltonian algebra.}
\GSfield{Fate}{\citet[\S1 and Theorem~6.1, Corollaries~6.2--6.3]{rudakov1974irreducible} explicitly includes H alongside W, S and K. The subsequent height-one classification for Hamiltonian algebras is \citet[Definition~5.5, Corollary~5.9 and Theorem~5.10, pp.~477--480]{rudakov1975types}, over an algebraically closed field of characteristic zero; received 7 October 1974 and published in 1975. This supplies later fate, without attributing that complete exceptional list to the one-line 1973 remark.}
\GSfield{Present status}{As of 9 November 2026: established in the stated scope.}
\GSfield{Assessment}{literal statement: correct with missing hypotheses; publication match: partial; verification: read the result; development: none; confidence: high}
\end{GSitem}
\begin{GSitem}{1973-310.8}{F}{The parenthetical bounded-representation question}
\GSfield{Kind; source}{question; \GSpage{P000521}}
\GSfield{The notes}{\ru{\(\exists\) ли огр. непр. предст-я?} appears inside the discussion of height.}
\GSfield{Statement}{The abbreviation may ask about a bounded height, bounded operators under a norm, or another boundedness condition. No norm, common height bound or quantified class is specified.}
\GSfield{Fate}{\citet[\S\S2--3]{rudakov1974irreducible} defines continuity and height but does not determine the missing predicate of this parenthetical question.}
\GSfield{Present status}{As of 9 November 2026: not sufficiently definite to receive a research-problem identifier.}
\GSfield{Assessment}{literal statement: not determinable; publication match: thematic only; verification: unresolved lead; development: none; confidence: low}
\end{GSitem}
\begin{GSitem}{1973-310.9}{F}{Kirillov on finite-functional-dimension diffeomorphism representations}
\GSfield{Kind; source}{question; \GSpage{P000522}; A. A. Kirillov}
\GSfield{The notes}{Kirillov asks for representations of \(\operatorname{Diff}M\) of finite functional dimension, induced from a finite-point stabilizer modulo a finite-jet stabilizer; the note says unproved but well known.}
\GSfield{Statement}{Given finitely many points of a smooth manifold, their stabilizer has finite jet quotient groups from which one can form induced representations. The note does not specify the topology on Diff, the representation category, the definition of functional dimension, the allowed jet order or an equivalence relation; it therefore does not determine an exhaustive classification theorem.}
\GSfield{Fate}{\citet[\S1.1, pp.~5--7]{vershikgelfandgraev1975diffeomorphisms} gives relevant examples of finite functional dimension and cites Kirillov\textquoteright s 1974 preprint. This is a thematic match, not a proof of the unqualified classification sentence.}
\GSfield{Present status}{As of 9 November 2026: the construction has a documented later context, but the claimed universal classification is under-specified.}
\GSfield{Assessment}{literal statement: not determinable; publication match: thematic only; verification: read the result; development: none; confidence: medium}
\end{GSitem}

\GStalk{311}{1973-10-15}{D. B. Fuchs}{P000521-P000526}{Quillen Theorem A, the Q construction and the plus comparison}
The continuation starts below Kirillov on P000522 and ends on P000524, before the new date on P000525. It gives the dual comma category, Theorem A, exact categories, admissible base change, and Quillen\textquoteright s span category Q. It then compares the Q definition with the plus construction, explicitly retracting a guessed equality. The recorded plus formula has an erroneous one-degree shift. The direct finite-rank maps to BQ are legitimate but do not themselves supply the comparison to the loop space. \citet[\S\S1--2]{quillen1973higher} and \citet[IV.7]{weibel2013kbook} distinguish these points.

\begin{GSitem}{1973-311.1}{E}{The dual comma category and homotopy fibres}
\GSfield{Kind; source}{construction; \GSpage{P000522}}
\GSfield{The notes}{Alongside \(Y\backslash f\), the page defines \(f/Y\) with arrows \(f(X)\to Y\), and says the dual theorem holds.}
\GSfield{Statement}{For a functor f between small categories, the comma category f/Y has objects \((X,f(X)\to Y)\). A morphism \(Y\to Y^{\prime}\) induces \(f/Y\to f/Y^{\prime}\). If every such induced classifying-space map is a homotopy equivalence, the canonical map \(B(f/Y)\to\operatorname{hofib}_Y(Bf)\) is a homotopy equivalence. For \(Y\backslash f\) the variance is reversed.}
\GSfield{Reconstruction}{Theorem B\textquoteright s arrow-invariance hypothesis remains necessary; a comma category is not automatically a homotopy fibre for every functor.}
\GSfield{Fate}{\citet[Theorem~B and its dual, p.~97]{quillen1973higher}.}
\GSfield{Present status}{As of 9 November 2026: established in the stated scope.}
\GSfield{Assessment}{literal statement: correct with missing hypotheses; publication match: exact; verification: read the result; development: none; confidence: high}
\end{GSitem}
\begin{GSitem}{1973-311.2}{E}{Quillen Theorem A}
\GSfield{Kind; source}{theorem; \GSpage{P000522}}
\GSfield{The notes}{If every \(B(Y\backslash f)\) is contractible, \(Bf\) induces isomorphisms of homotopy groups.}
\GSfield{Statement}{If \(B(Y\backslash f)\) is contractible for every object Y of the small category \(C^{\prime}\), then \(Bf:BC\to BC^{\prime}\) is a homotopy equivalence. The dual criterion with f/Y is equally valid.}
\GSfield{Reconstruction}{Classifying spaces have CW type, so the weak-equivalence conclusion in the notes gives the stated homotopy equivalence, including the component condition.}
\GSfield{Fate}{\citet[Theorem~A, p.~93]{quillen1973higher}.}
\GSfield{Present status}{As of 9 November 2026: established in the stated scope.}
\GSfield{Assessment}{literal statement: correct; publication match: exact; verification: read the result; development: none; confidence: high}
\end{GSitem}
\begin{GSitem}{1973-311.3}{E}{An exact category inside an abelian category}
\GSfield{Kind; source}{definition; \GSpage{P000522}}
\GSfield{The notes}{M is a full extension-closed subcategory of an abelian category A; monomorphisms and epimorphisms in its short exact sequences are admissible.}
\GSfield{Statement}{Let M be an essentially small, full additive subcategory of an abelian category A, closed under extensions. Give M the exact structure consisting of short exact sequences in A with all three terms in M. An admissible monomorphism has cokernel in M, and an admissible epimorphism has kernel in M.}
\GSfield{Reconstruction}{Additivity and the requirement that every term belong to M are stated explicitly. Not every ambient monomorphism between objects of M is admissible.}
\GSfield{Fate}{\citet[\S2, pp.~99--100]{quillen1973higher}.}
\GSfield{Present status}{As of 9 November 2026: established in the stated scope.}
\GSfield{Assessment}{literal statement: correct with missing hypotheses; publication match: exact; verification: read the result; development: none; confidence: high}
\end{GSitem}
\begin{GSitem}{1973-311.4}{E}{Admissible epimorphisms survive pullback}
\GSfield{Kind; source}{theorem; \GSpage{P000522}}
\GSfield{The notes}{For an admissible epimorphism \(Y\to X\) and arbitrary \(Z\to X\), the pullback lies in M and its arrow to Z is admissible.}
\GSfield{Statement}{If \(0\to K\to Y\to X\to0\) is exact in M and \(Z\to X\) is any M-morphism, the ambient pullback P satisfies \(0\to K\to P\to Z\to0\). Extension closure puts P in M and makes \(P\to Z\) an admissible epimorphism.}
\GSfield{Reconstruction}{This proof supplies exactly the closure needed for the following span composition.}
\GSfield{Fate}{\citet[\S2, exact-category axiom (b), pp.~99--100]{quillen1973higher}.}
\GSfield{Present status}{As of 9 November 2026: established in the stated scope.}
\GSfield{Assessment}{literal statement: correct; publication match: exact; verification: read the result; development: annotated; confidence: high}
\end{GSitem}
\begin{GSitem}{1973-311.5}{E}{Morphisms in Q are admissible subquotient spans}
\GSfield{Kind; source}{construction; \GSpage{P000523}}
\GSfield{The notes}{Objects of QM equal those of M, and morphisms are \(X_1\leftarrow X\to X_2\), with admissible epi and mono, modulo the obvious equivalence.}
\GSfield{Statement}{A morphism \(X\to Y\) of QM is an isomorphism class of diagrams \(X\xleftarrow{p}E\xrightarrow{i}Y\), where p is admissible epic and i admissible monic. Isomorphisms of middle objects must commute with both arrows. Thus X is an admissible quotient of an admissible subobject of Y.}
\GSfield{Reconstruction}{The equivalence relation is stated explicitly; reversing either admissibility label would define a different construction.}
\GSfield{Fate}{\citet[\S2, definition of Q, p.~101]{quillen1973higher}.}
\GSfield{Present status}{As of 9 November 2026: established in the stated scope.}
\GSfield{Assessment}{literal statement: correct; publication match: exact; verification: read the result; development: none; confidence: high}
\end{GSitem}
\begin{GSitem}{1973-311.6}{E}{Composition of Q spans by fibre product}
\GSfield{Kind; source}{construction; \GSpage{P000523}}
\GSfield{The notes}{The diagram defines \(X_4=X\times_{X_2}X^{\prime}\) and uses it as the middle object of the composite.}
\GSfield{Statement}{Compose \(X_1\leftarrow E\rightarrowtail X_2\) and \(X_2\twoheadleftarrow E^{\prime}\rightarrowtail X_3\) through \(E\times_{X_2}E^{\prime}\). The projection to E is admissible epic; the projection to \(E^{\prime}\) is admissible monic, being the kernel of the composite to the admissible cokernel of \(E\rightarrowtail X_2\). Associativity follows from canonical pullback isomorphisms.}
\GSfield{Reconstruction}{An admissible epi pullback stays in M by the preceding item; the kernel argument supplies the other admissibility. The corrected middle-object label X4 in the scan is respected.}
\GSfield{Fate}{\citet[\S2, definition of composition in Q, p.~101]{quillen1973higher}.}
\GSfield{Present status}{As of 9 November 2026: established in the stated scope.}
\GSfield{Assessment}{literal statement: correct; publication match: exact; verification: read the result; development: none; confidence: high}
\end{GSitem}
\begin{GSitem}{1973-311.7}{E}{Higher K groups as shifted homotopy groups of BQ}
\GSfield{Kind; source}{definition; \GSpage{P000523}}
\GSfield{The notes}{\(K_i(M)=\pi_{i+1}(BQM,0)\); isomorphism classes of objects must form a set.}
\GSfield{Statement}{Choose a small skeleton of the exact category M and a zero object as basepoint. For \(i\ge0\), define \(K_i(M)=\pi_{i+1}(BQM,0)\). In degree zero this is canonically the Grothendieck group on objects with exact-sequence relations.}
\GSfield{Reconstruction}{Essential smallness justifies the classifying-space construction; zero objects give equivalent choices of basepoint.}
\GSfield{Fate}{\citet[\S2, definition and Theorem~1, pp.~101--103]{quillen1973higher}.}
\GSfield{Present status}{As of 9 November 2026: established in the stated scope.}
\GSfield{Assessment}{literal statement: correct with missing hypotheses; publication match: exact; verification: read the result; development: none; confidence: high}
\end{GSitem}
\begin{GSitem}{1973-311.8}{E}{The exact category of finitely generated projective modules}
\GSfield{Kind; source}{definition; \GSpage{P000523}}
\GSfield{The notes}{For an arbitrary ring A, PA consists of finitely generated projective left modules, and \(K_i(A)=K_i(PA)\).}
\GSfield{Statement}{For a unital ring A take P(A), a small skeleton of finitely generated projective unital left A-modules with its split exact structure. Its Q construction defines the connective algebraic K groups of A.}
\GSfield{Reconstruction}{The identity convention on rings and modules is explicit. All these short exact sequences split, but their span category and K groups are still nontrivial.}
\GSfield{Fate}{\citet[Introduction and \S2]{quillen1973higher}.}
\GSfield{Present status}{As of 9 November 2026: established in the stated scope.}
\GSfield{Assessment}{literal statement: correct with missing hypotheses; publication match: exact; verification: read the result; development: none; confidence: high}
\end{GSitem}
\begin{GSitem}{1973-311.9}{E}{The plus construction and its corrected degree}
\GSfield{Kind; source}{definition; \GSpage{P000523}}
\GSfield{The notes}{\(BGL(A)\to B_+GL(A)\) is an acyclic map, surjective on fundamental groups, with H-space target. The page then calls \(\pi_{i+1}(B_+GL(A))\) the old definition of Ki.}
\GSfield{Statement}{For stable \(GL(A)\), kill its perfect elementary subgroup by the plus construction. The map \(BGL(A)\to BGL(A)^+\) is acyclic, has fundamental group \(GL(A)/E(A)=K_1(A)\), and its target has an H-space structure. The correct convention is \(K_i(A)=\pi_i(BGL(A)^+)\) for \(i\ge1\). Degree zero is supplied separately by K0(A).}
\GSfield{Reconstruction}{The source transcription intentionally preserves i+1. That index is correct for BQ and incorrect for BGL-plus. The map and H-space requirements alone do not specify the plus construction without its perfect subgroup.}
\GSfield{Fate}{\citet[Introduction, pp.~85--86]{quillen1973higher}; precise conventions in \citet[Definition~IV.1.1 and IV.1.1.1]{weibel2013kbook}.}
\GSfield{Present status}{As of 9 November 2026: established in the stated scope.}
\GSfield{Assessment}{literal statement: false as written; publication match: exact; verification: read the result; development: none; confidence: high}
\end{GSitem}
\begin{GSitem}{1973-311.10}{E}{The plus/Q comparison needs a loop space}
\GSfield{Kind; source}{theorem; \GSpage{P000523}}
\GSfield{The notes}{The guessed equality \(BQPA=B_+GL(A)\) is immediately followed by this is false, but there is a map.}
\GSfield{Statement}{The correct comparison is a homotopy equivalence \(\Omega BQP(A)\simeq K_0(A)\times BGL(A)^+\). Thus the identity component of the loop space corresponds to the connected plus space. This yields agreement of the two definitions in positive K degrees.}
\GSfield{Reconstruction}{The notebook itself retracts the unlooped equality. Both the loop and the discrete K0 factor matter; a mere map to BQ is insufficient to establish the comparison.}
\GSfield{Fate}{Announced in \citet[Introduction, p.~86]{quillen1973higher}; a full proof and exact statement are in \citet[Theorem~IV.7.1 and Corollary~IV.7.2]{weibel2013kbook}.}
\GSfield{Present status}{As of 9 November 2026: established in the stated scope.}
\GSfield{Assessment}{literal statement: false as written; publication match: exact; verification: read the result; development: none; confidence: high}
\end{GSitem}
\begin{GSitem}{1973-311.11}{E}{Finite-rank automorphisms give nullhomotopic maps to BQ}
\GSfield{Kind; source}{construction; \GSpage{P000523-P000524}}
\GSfield{The notes}{Viewing GL(n,A) as a one-object category, its action on \(S_n=A^n\) gives a functor to QP(A), hence \(BGL(n,A)\to BQP(A)\).}
\GSfield{Statement}{The stated functor exists: an automorphism g gives the span \(A^n\xleftarrow{\rm id}A^n\xrightarrow{g}A^n\). Its classifying-space map is nullhomotopic, because \(0\xleftarrow{\rm id}0\rightarrowtail A^n\) is invariant under these automorphisms and gives a natural transformation from the constant zero functor.}
\GSfield{Reconstruction}{\ours{For every g the two composites of the displayed zero span agree, so its components form a natural transformation. Realizing its nerve gives the nullhomotopy. This works on the automorphism subcategory; it does not make the whole Q category contractible.}}
\GSfield{Fate}{Functoriality of nerves and the natural-transformation homotopy are \citet[\S1, Proposition~2]{quillen1973higher}. The specific nullhomotopy is our direct check of the notebook construction, not a quoted plus/Q comparison theorem.}
\GSfield{Present status}{As of 9 November 2026: established in the stated scope.}
\GSfield{Assessment}{literal statement: correct; publication match: exact; verification: read the result; development: annotated; confidence: high}
\end{GSitem}
\begin{GSitem}{1973-311.12}{E}{Adding a free summand is homotopic to the identity on BQ}
\GSfield{Kind; source}{construction; \GSpage{P000524}}
\GSfield{The notes}{\(X\mapsto X\oplus A\) is called a homotopy equivalence and used to align the finite-rank stabilization square.}
\GSfield{Statement}{Direct sum with A acts on a Q-span by summing A on each object and the identity on each arrow. The admissible inclusions \(X\to X\oplus A\), represented as Q-spans, are a natural transformation from the identity functor. Hence the induced map on BQP(A) is homotopic to the identity and is a homotopy equivalence. The finite-rank square commutes before realization with this right-hand functor.}
\GSfield{Reconstruction}{\ours{For a span \(X\xleftarrow{p}E\xrightarrow{i}Y\), pull back \(p\oplus1:E\oplus A\to X\oplus A\) along \(X\hookrightarrow X\oplus A\). The pullback is E, and both naturality composites are \(X\leftarrow E\to Y\oplus A\). Thus stabilization is homotopic to the identity. These unlooped maps do not by themselves identify the K-theory comparison.}}
\GSfield{Fate}{\citet[\S1, Proposition~2]{quillen1973higher} supplies the homotopy mechanism. The direct check below establishes the exact stabilization sentence; the separate comparison theorem remains \citet[Corollary~IV.7.2]{weibel2013kbook}.}
\GSfield{Present status}{As of 9 November 2026: established in the stated scope.}
\GSfield{Assessment}{literal statement: correct; publication match: exact; verification: read the result; development: annotated; confidence: high}
\end{GSitem}

\GStalk{313}{1973-10-22}{I. M. Gelfand}{P000525-P000526}{Nonlinear scattering and a 2-adic logarithm example}
The first block on P000525 advertises Strauss's nonlinear scattering theory through comparison of a nonlinear flow with a simpler linear flow. Klein--Gordon and Korteweg--de Vries equations are examples, without a stated dimension, data class, or completeness theorem. The explicitly attributed Cartier example ``as presented by Gelfand'' occurs inside the ensuing discussion of \(p\)-adic measures on P000526 and is indexed here as an intervention, not a second uninterrupted talk. The archive's overlapping page ranges do not establish a different speaker for the unlabelled surrounding measure theory.

\begin{GSitem}{1973-313.1}{E}{Wave operators and the nonlinear scattering map}
\GSfield{Kind; source}{definition; \GSpage{P000525}}
\GSfield{The notes}{The notes write \(U(t)U(s)=U(t+s)\), conservation of \(E\), and \(E(U(t)f-U_0(t)f_\pm)\to0\) as \(t\to\pm\infty\). The relation between \(f_-\) and \(f_+\) is denoted \(S\).}
\GSfield{Statement}{Let \(U(t)\) and \(U_0(t)\) be global flows on specified phase spaces embedded in a common normed space \(X\). A scattering state for \(f\) is an \(f_\pm\) with \(\|U(t)f-U_0(t)f_\pm\|_X\to0\). If these states exist uniquely on specified incoming and outgoing sets, the scattering map is \(S:f_-\mapsto f_+\). Existence, uniqueness, and surjectivity are separate assertions; conservation of an energy alone implies none of them.}
\GSfield{Reconstruction}{The normed phase space, domains of the maps, and uniqueness assumptions make the definition precise. Interpreting the notebook's \(E\) on a difference as a norm or a coercive error functional is a reconstruction.}
\GSfield{Fate}{This is exposition of an existing framework. Strauss's 1973 Denver lectures were published in 1974; the abstract and opening formulation of \citet[pp.~53--78]{strauss1974scattering} explicitly study asymptotic comparison with linear equations. The publication is not a theorem verifying all examples listed here.}
\GSfield{Present status}{As of 9 November 2026: The framework is standard; no universal existence or completeness assertion can be extracted from this paragraph.}
\GSfield{Assessment}{literal statement: correct with missing hypotheses; publication match: exact; verification: publisher metadata checked; development: annotated; confidence: high}
\end{GSitem}
\begin{GSitem}{1973-313.2}{E}{The cubic Klein--Gordon equation as a scattering example}
\GSfield{Kind; source}{remark; \GSpage{P000525}}
\GSfield{The notes}{The displayed example is \(u_{tt}-\Delta u+u+|u|^2u=0\), called a wave equation. The Laplacian sign is retraced in the manuscript.}
\GSfield{Statement}{For a sufficiently regular real or complex solution on \(\mathbb R^d\), with decay justifying integration by parts, the defocusing cubic Klein--Gordon energy is \[E(u,u_t)=\int_{\mathbb R^d}\left(\tfrac12|u_t|^2+\tfrac12|\nabla u|^2+\tfrac12|u|^2+\tfrac14|u|^4\right)\,dx.\] It is conserved. A scattering theorem additionally fixes \(d\), a solution class, and a norm; the notebook does not supply them.}
\GSfield{Reconstruction}{The energy and the identification as Klein--Gordon are our standard reconstruction of the displayed equation, not a recorded theorem of global scattering.}
\GSfield{Fate}{The example belongs to the nonlinear-wave scattering program advertised in \citet[abstract and pp.~53--78]{strauss1974scattering}. Only the program and bibliographic chronology were verified there; the notebook does not identify a particular theorem with fixed dimension and data.}
\GSfield{Present status}{As of 9 November 2026: The equation and its conserved energy are identifiable. A claim of all-data scattering in an arbitrary dimension would go beyond the evidence.}
\GSfield{Assessment}{literal statement: correct with missing hypotheses; publication match: partial; verification: publisher metadata checked; development: annotated; confidence: high}
\end{GSitem}
\begin{GSitem}{1973-313.3}{E}{KdV and the distinction between dispersive and soliton asymptotics}
\GSfield{Kind; source}{remark; \GSpage{P000525}}
\GSfield{The notes}{The second example is \(u_\tau+u_{xxx}+uu_x=0\), labelled the de Vries equation.}
\GSfield{Statement}{This is the Korteweg--de Vries equation after a normalization of the dependent variable. Its nonzero travelling solitons show why the preceding comparison with the free Airy flow cannot be asserted for every decaying solution without specifying the asymptotic sector. The equation itself is unambiguous; no scattering theorem is stated.}
\GSfield{Reconstruction}{The warning about solitons interprets the example; it is not attributed to Gelfand. The manuscript's \(\tau\) is retained as a time variable.}
\GSfield{Fate}{\citet[pp.~240--243]{zabuskykruskal1965solitons} had already connected KdV solitary-wave interactions with recurrence in 1965. The publisher record confirms the paper and date; this supports chronology, not an exact identification of the unspecified scattering statement.}
\GSfield{Present status}{As of 9 November 2026: An expository example of nonlinear dispersive asymptotics, with no precise new problem in the notes.}
\GSfield{Assessment}{literal statement: correct; publication match: thematic only; verification: publisher metadata checked; development: annotated; confidence: high}
\end{GSitem}
\begin{GSitem}{1973-313.4}{E}{Cartier's 2-adic logarithm example}
\GSfield{Kind; source}{theorem; \GSpage{P000526}; I. M. Gelfand (after Cartier)}
\GSfield{The notes}{``\ru{(Cartier) в изложении Гельфанда}'' precedes \(\ln(1-a)=a+a^2/2+\cdots\), followed by \(a=2\), \(\ln(-1)=0\), and \(\sum_{n\ge1}2^n/n=0\). The partial sums begin \(2,4,20/3,32/3\).}
\GSfield{Statement}{In \(\mathbb Q_2\), \(v_2(2^n/n)=n-v_2(n)\to\infty\), so the series converges. The logarithm on \(1+2\mathbb Z_2\) is \(\log(1-a)=-\sum_{n\ge1}a^n/n\). Its homomorphism law gives \(2\log(-1)=\log(1)=0\), hence \(\sum_{n\ge1}2^n/n=0\). Thus for every \(k\), all sufficiently late partial sums belong to \(2^k\mathbb Z_2\).}
\GSfield{Reconstruction}{We insert the missing minus sign and justify convergence; the final vanishing assertion survives this correction. The source gives an oral attribution, not a Cartier publication.}
\GSfield{Fate}{An elementary consequence of the pre-existing \(p\)-adic logarithm. No exact Cartier publication was located; the mathematical verification here is direct and does not turn an oral attribution into a bibliography entry.}
\GSfield{Present status}{As of 9 November 2026: Correct after repairing the logarithm identity; the attribution remains an oral report in the notebook.}
\GSfield{Assessment}{literal statement: false as written; publication match: none found; verification: indirect; development: annotated; confidence: high}
\end{GSitem}

\GStalk{312}{1973-10-22}{A. A. Kirillov}{P000525-P000527}{Metaplectic spinors and bounded p-adic measures}
After Gelfand's scattering announcement, Kirillov presents symplectic spinors through the metaplectic cover and the canonical commutation relations. The numbered second block treats \(p\)-adic measures, the bounded power-series transform, and the Newton--Mahler basis. Its speaker is not restated; assigning it to the continuing Kirillov block is an editorial reconstruction. The Cartier example explicitly credited to Gelfand is separated as item 1973-313.4. P000527 begins with the end of the Mahler argument before the new Fuchs heading.

\begin{GSitem}{1973-312.1}{E}{The metaplectic obstruction and symplectic spinors}
\GSfield{Kind; source}{construction; \GSpage{P000525}}
\GSfield{The notes}{For symplectic \(M\), the frame group is \(\operatorname{Sp}(2n,\mathbb R)\), topologically \(U(n)\); a two-sheeted cover \(M(2n,\mathbb R)\) permits a spinor bundle when \(C_1\) is divisible by two.}
\GSfield{Statement}{Let \((M^{2n},\omega)\) be a smooth symplectic manifold and \(J\) an \(\omega\)-compatible almost complex structure. Its symplectic frame bundle lifts to \(\operatorname{Mp}(2n,\mathbb R)\) exactly when \(w_2(TM)=0\), equivalently \(c_1(TM,J)\in2H^2(M;\mathbb Z)\). A chosen lift and the oscillator representation define an associated infinite-dimensional symplectic spinor bundle. The homotopy equivalence with \(U(n)\) is not an isomorphism of groups.}
\GSfield{Reconstruction}{The standard name \(\operatorname{Mp}\), the choice of compatible \(J\), and the equivalence with \(w_2=0\) explicate the shorthand. A choice of lift is necessary, not canonical.}
\GSfield{Fate}{The notes explicitly report Kostant's work. \citet[\S3, Theorem~5 and Definition~6]{krysl2017spinors} states this obstruction and associated-bundle construction, attributing it to Kostant's 1973 Rome contribution published in 1974. The underlying metaplectic representation predates the seminar; publication in 1974 does not make this Kirillov's new result.}
\GSfield{Present status}{As of 9 November 2026: An established construction, with the topological obstruction and choice of metaplectic structure made explicit.}
\GSfield{Assessment}{literal statement: correct; publication match: exact; verification: read the result; development: annotated; confidence: high}
\end{GSitem}
\begin{GSitem}{1973-312.2}{E}{Canonical commutation relations and metaplectic implementation}
\GSfield{Kind; source}{construction; \GSpage{P000525}}
\GSfield{The notes}{The algebra has \(p_iq_j-q_jp_i=\delta_{ij}\), represented by \(p_i=\partial_{x_i}\), \(q_i=x_i\) in \(L^2(\mathbb R^n)\). Symplectic changes of generators are said to induce automorphisms of bounded operators, implemented up to scalar, whose factors reduce to \(\pm1\).}
\GSfield{Statement}{On the invariant Schwartz domain, \([\partial_{x_i},x_j]=\delta_{ij}\), together with \([p_i,p_j]=[q_i,q_j]=0\). For a unitary formulation use \(P_i=-i\partial_{x_i}\), \(Q_i=x_i\) and the Weyl relations of the Heisenberg group with a fixed nontrivial central character. Uniqueness of its irreducible unitary representation gives projective implementers for \(\operatorname{Sp}(2n,\mathbb R)\); their canonical lift is the oscillator representation of its double cover.}
\GSfield{Reconstruction}{The domain and the Weyl form of the relations repair the unsupported passage from unbounded polynomial generators to bounded-operator automorphisms. The double cover requires the metaplectic construction; innerness alone does not prove it.}
\GSfield{Fate}{\citet[\S\S2.1--2.3]{krysl2017spinors} reviews the precise oscillator representation and Weyl algebra. Kostant's earlier quantization paper was published in 1970, independently of the website's later digitization date \citep[pp.~87--208]{kostant1970quantization}.}
\GSfield{Present status}{As of 9 November 2026: The representation-theoretic statement is established. The notebook proof needs domain and continuity hypotheses and a genuine metaplectic step.}
\GSfield{Assessment}{literal statement: correct with missing hypotheses; publication match: partial; verification: read the result; development: annotated; confidence: high}
\end{GSitem}
\begin{GSitem}{1973-312.3}{E}{The Clifford-algebra analogy}
\GSfield{Kind; source}{remark; \GSpage{P000525}}
\GSfield{The notes}{The parenthesis says that replacing the minus sign by a plus gives the spin representation of the orthogonal group.}
\GSfield{Statement}{For a finite-dimensional quadratic space \((V,q)\) over a field of characteristic different from two, the Clifford algebra is generated by \(V\) with \(uv+vu=2b(u,v)\), where \(b\) polarizes \(q\). Its spin modules carry representations of \(\operatorname{Spin}(V,q)\), hence generally projective representations of \(\operatorname{SO}(V,q)\). Merely changing the displayed mixed commutator is insufficient: all anticommutation relations and the quadratic form must be specified.}
\GSfield{Reconstruction}{This is an analogy reconstructed as a Clifford-algebra statement; the notebook contains no choice of signature or spin module.}
\GSfield{Fate}{The orthogonal/symplectic analogy is reviewed in \citet[\S2 and introduction]{krysl2017spinors}. It is background, not an independent 1973 discovery.}
\GSfield{Present status}{As of 9 November 2026: Standard after specifying the full Clifford relations and the spin double cover.}
\GSfield{Assessment}{literal statement: correct with missing hypotheses; publication match: partial; verification: read the result; development: annotated; confidence: high}
\end{GSitem}
\begin{GSitem}{1973-312.4}{E}{Bounded K-valued measures as continuous duals}
\GSfield{Kind; source}{definition; \GSpage{P000526}}
\GSfield{The notes}{For a finite extension \(K/\mathbb Q_p\), the notes write \(\operatorname{Hom}(C(\mathbb Z_p,K),\mathbb Q_p)=C(\mathbb Z_p,K)'=M(\mathbb Z_p,K)\).}
\GSfield{Statement}{Give \(C(\mathbb Z_p,K)\) its supremum norm. The Banach space of bounded \(K\)-valued measures is \(M(\mathbb Z_p,K)=\operatorname{Hom}_{K,\mathrm{cont}}(C(\mathbb Z_p,K),K)\), with operator norm. Equivalently these are finitely additive \(K\)-valued functions on compact open sets with uniformly bounded values, extended to continuous functions by uniform approximation.}
\GSfield{Reconstruction}{We specify continuity and \(K\)-linearity and change the codomain to \(K\). A \(\mathbb Q_p\)-linear dual is a different object, though trace pairings can relate it to a \(K\)-dual after extra choices.}
\GSfield{Fate}{\citet[\S2]{amice1965mesures} explicitly defines the measure space as the continuous dual; the lecture is dated 17 May 1965. The note's codomain mismatch is not part of that result.}
\GSfield{Present status}{As of 9 November 2026: Established Banach-space duality; the coefficient field and boundedness cannot be suppressed.}
\GSfield{Assessment}{literal statement: false as written; publication match: exact; verification: read the result; development: annotated; confidence: high}
\end{GSitem}
\begin{GSitem}{1973-312.5}{E}{Convolution of p-adic measures}
\GSfield{Kind; source}{construction; \GSpage{P000526}}
\GSfield{The notes}{Convolution is defined by the iterated integral \(\int f\,d(\mu*\nu)=\int\int f(x+y)\,d\mu(x)\,d\nu(y)\).}
\GSfield{Statement}{For bounded \(K\)-valued measures on the compact additive group \(\mathbb Z_p\), this formula defines a bounded measure. Uniform continuity makes the inner integral continuous in \(y\). It gives a commutative unital Banach algebra with unit \(\delta_0\) and \(\|\mu*\nu\|\le\|\mu\|\|\nu\|\).}
\GSfield{Reconstruction}{The compact group, boundedness, and unit are stated explicitly to justify the integral construction.}
\GSfield{Fate}{The same construction appears in \citet[\S2, ``Produit de convolution'']{amice1965mesures}, eight years before this seminar.}
\GSfield{Present status}{As of 9 November 2026: An established convolution algebra; no unbounded-measure extension follows automatically.}
\GSfield{Assessment}{literal statement: correct; publication match: exact; verification: read the result; development: annotated; confidence: high}
\end{GSitem}
\begin{GSitem}{1973-312.6}{E}{The bounded Amice transform}
\GSfield{Kind; source}{theorem; \GSpage{P000526}}
\GSfield{The notes}{``Theorem 1'' identifies \(M(\mathbb Z_p,K)\) with \(K[[T]]\), immediately restricted to series of finite coefficient supremum. It uses \(a_k=\int\binom{x}{k}\,d\mu\) and \(\int(1+T)^x\,d\mu\).}
\GSfield{Statement}{Write \[K[[T]]_{\mathrm b}=\{\sum a_kT^k:\sup_k|a_k|_p<\infty\}.\] The map \(\mathcal A\mu(T)=\sum_{k\ge0}(\int\binom{x}{k}\,d\mu)T^k\) is an isometric \(K\)-algebra isomorphism \(M(\mathbb Z_p,K)\simeq K[[T]]_{\mathrm b}\). In particular \(\mathcal A(\mu*\nu)=\mathcal A\mu\,\mathcal A\nu\). The integral character formula converges for \(|T|_p<1\).}
\GSfield{Reconstruction}{The subscript distinguishes the bounded-coefficient algebra from all formal series; the notebook's finite-norm parenthesis already supplies that restriction.}
\GSfield{Fate}{\citet[\S3]{amice1965mesures} identifies convolution measures with bounded-coefficient analytic transforms and records the isometry. The one-variable identification in the notes is its standard special case.}
\GSfield{Present status}{As of 9 November 2026: Established; the equality with unrestricted formal power series would be false.}
\GSfield{Assessment}{literal statement: correct with missing hypotheses; publication match: exact; verification: read the result; development: annotated; confidence: high}
\end{GSitem}
\begin{GSitem}{1973-312.7}{F}{Unbounded distributions are not automatically arbitrary moment series}
\GSfield{Kind; source}{claim; \GSpage{P000526}}
\GSfield{The notes}{The notes say that formal series are obtained from unbounded measures, then define such measures as finitely additive functions on elementary sets, dual to locally constant functions.}
\GSfield{Statement}{The algebraic dual of locally constant functions on \(\mathbb Z_p\) consists of compatible values on residue classes, without a uniform bound. It does not by that definition act on the non-locally-constant polynomials \(\binom{x}{k}\). Thus the displayed moment transform is not defined for an arbitrary such distribution. A distribution space with additional regularity and a growth-controlled transform must be specified before any power-series classification is meaningful.}
\GSfield{Reconstruction}{This is a diagnosis of the missing domain, not a replacement theorem asserting a particular distribution topology.}
\GSfield{Fate}{The bounded theorem in \citet[\S\S2--3]{amice1965mesures} does not establish this unbounded extension. Searches for Amice's duals and Fourier transforms identify several inequivalent distribution theories; none supplies the absent test-function space in this sentence.}
\GSfield{Present status}{As of 9 November 2026: The literal inference is unsupported. Its intended unbounded version cannot be identified from the notes; this is not evidence for an open problem.}
\GSfield{Assessment}{literal statement: not determinable; publication match: none found; verification: unresolved lead; development: annotated; confidence: medium}
\end{GSitem}
\begin{GSitem}{1973-312.8}{E}{Measures on the p-adic unit group}
\GSfield{Kind; source}{theorem; \GSpage{P000526}}
\GSfield{The notes}{``Theorem 2'' writes \(M(\mathbb Z_p^*,K)=(p-1)K[[T]]\).}
\GSfield{Statement}{For odd \(p\), \(\mathbb Z_p^\times\simeq\mu_{p-1}\times(1+p\mathbb Z_p)\), and a choice of a topological generator identifies the second factor with \(\mathbb Z_p\). Since \(K\supset\mathbb Q_p\) contains \(\mu_{p-1}\), character idempotents give \(M(\mathbb Z_p^\times,K)\simeq\prod_{\chi:\mu_{p-1}\to K^\times}K[[T]]_{\mathrm b}\). For \(p=2\), use \(\{\pm1\}\times(1+4\mathbb Z_2)\), giving two components, not \(p-1=1\).}
\GSfield{Reconstruction}{We interpret the coefficient \(p-1\) as a product decomposition of convolution algebras, add the odd-prime restriction, and make the noncanonical generator choice explicit.}
\GSfield{Fate}{This is a direct finite-character decomposition combined with the already established one-variable transform of \citet[\S3]{amice1965mesures}; the decomposition is reconstructed here, not quoted as that paper's numbered theorem.}
\GSfield{Present status}{As of 9 November 2026: Correct for odd primes in the stated form; the dyadic case requires a different finite factor.}
\GSfield{Assessment}{literal statement: correct with missing hypotheses; publication match: partial; verification: read the result; development: annotated; confidence: high}
\end{GSitem}
\begin{GSitem}{1973-312.9}{E}{Exponential coordinates and the range of the transform}
\GSfield{Kind; source}{question; \GSpage{P000526}}
\GSfield{The notes}{Putting \(1+T=e^\tau\), the notes write \(\widetilde\mu(\tau)=\int e^{\tau x}\,d\mu\), invite one to ``rewrite the analysis'', and ask what functions arise, with integral coefficients after inverse substitution.}
\GSfield{Statement}{For the already defined bounded measures, \(\widetilde\mu(\tau)=\mathcal A\mu(e^\tau-1)\) converges for \(|\tau|_p<p^{-1/(p-1)}\). Formally, \(G(\tau)\in K[[\tau]]\) is in this image exactly when \(G(\log(1+T))\in K[[T]]_{\mathrm b}\); the integral-measure unit ball corresponds to coefficients in \(\mathcal O_K\). Analytic continuation outside this exponential disk is a separate question and is not specified in the notes.}
\GSfield{Reconstruction}{This exact formal characterization is our consequence of the preceding isomorphism. We do not replace the informal request for a range description by an unrecorded continuation problem.}
\GSfield{Fate}{The character transform is already developed in \citet[\S3]{amice1965mesures}; the exponential-coordinate formulation here follows by a change of formal variable. No distinct publication or unresolved analytic-domain question is identifiable from the notebook wording.}
\GSfield{Present status}{As of 9 November 2026: The bounded formal range is determined by the transform theorem. A broader analytic range question lacks a specified domain and topology.}
\GSfield{Assessment}{literal statement: correct with missing hypotheses; publication match: partial; verification: read the result; development: annotated; confidence: medium}
\end{GSitem}
\begin{GSitem}{1973-312.10}{E}{The Newton--Mahler expansion and norm identity}
\GSfield{Kind; source}{theorem; \GSpage{P000526-P000527}}
\GSfield{The notes}{The proof uses \(P_k(x)=x(x-1)\cdots(x-k+1)/k!\), \(c_k=\Delta^kf(0)\), and claims uniform convergence for every continuous \(f\), with \(\max|f|=\max|a_k|\). The series is called Taylor and then Maclaurin.}
\GSfield{Statement}{Every \(f\in C(\mathbb Z_p,K)\) has the unique uniformly convergent expansion \(f(x)=\sum_{k\ge0}c_k\binom{x}{k}\), where \(c_k=\sum_{j=0}^k(-1)^{k-j}\binom{k}{j}f(j)\to0\), and \(\|f\|_\infty=\sup_k|c_k|\). Duality gives \(\|\mu\|=\sup_k|\int\binom{x}{k}\,d\mu|\). These are finite differences and falling factorials, not derivatives and ordinary powers.}
\GSfield{Reconstruction}{The vanishing of the coefficients and uniqueness are supplied. The \(a_k\) in the last norm formula refers to function-expansion coefficients there; elsewhere it denotes measure moments.}
\GSfield{Fate}{This is Mahler's 1958 interpolation theorem \citep[Theorem~1]{mahler1958interpolation}, extended in \citet[Chapter~II]{amice1964interpolation}. Mahler's 1961 correction concerns the differentiability theorem, not this continuous expansion \citep[opening paragraph]{mahler1961correction}.}
\GSfield{Present status}{As of 9 November 2026: Established; the notebook name ``Maclaurin must not be read as a claim of ordinary analyticity for every continuous p-adic function.}
\GSfield{Assessment}{literal statement: correct; publication match: exact; verification: read the result; development: annotated; confidence: high}
\end{GSitem}

\GStalk{314}{1973-10-22}{D. B. Fuchs}{P000527}{The Q-construction and the Grothendieck group}
The continuation on P000527 gives Quillen's exact-category \(Q\)-construction and explains why its fundamental group is the Grothendieck group: objects determine loops and short exact sequences determine their additive relations. A separate line compares \(BQM\) with \(B\operatorname{GL}(\infty,A)\), omitting the loop-space shift. This material continues record 311, rather than introducing an independent K-theory construction.

\begin{GSitem}{1973-314.1}{E}{Quillen's Q-category of an exact category}
\GSfield{Kind; source}{construction; \GSpage{P000527}}
\GSfield{The notes}{\(M\) is an additive full subcategory of an abelian category, closed under extensions. Objects of \(QM\) are those of \(M\); a morphism is drawn \(X\leftarrow X''\to X'\), with an admissible epimorphism and an admissible monomorphism.}
\GSfield{Statement}{Let \(\mathcal M\) be an essentially small extension-closed additive full subcategory of an abelian category, with the induced exact structure. A morphism \(X\to Y\) in \(Q\mathcal M\) is an isomorphism class of spans \(X\xleftarrow{p}Z\xrightarrow{i}Y\), with \(p\) an admissible epimorphism and \(i\) an admissible monomorphism. Composition uses pullback; the exact-category axioms ensure admissibility.}
\GSfield{Reconstruction}{We insert smallness, the induced exact structure, equivalence of spans, and the composition law. These are necessary to turn the diagram into a category.}
\GSfield{Fate}{\citet[\S2, construction of \(Q\mathcal M\)]{quillen1973higher} is the exact source, originating in the 1972 Seattle lectures and published in 1973. This continues the exposition in record 311.}
\GSfield{Present status}{As of 9 November 2026: Established; there is no separate Fuchs conjecture here.}
\GSfield{Assessment}{literal statement: correct with missing hypotheses; publication match: exact; verification: read the result; development: annotated; confidence: high}
\end{GSitem}
\begin{GSitem}{1973-314.2}{E}{The missing loop in comparison with the plus construction}
\GSfield{Kind; source}{claim; \GSpage{P000527}}
\GSfield{The notes}{The notebook states that \(BQM\) is homologically equivalent to \(B\operatorname{GL}(\infty,A)\).}
\GSfield{Statement}{For a unital ring \(A\), let \(\mathcal P(A)\) be its exact category of finitely generated projective modules. Then \(\Omega BQ\mathcal P(A)\simeq K_0(A)\times B\operatorname{GL}(A)^+\), after choosing compatible base components. Thus the component of the constant loop has the homotopy type of \(B\operatorname{GL}(A)^+\), and the plus map from \(B\operatorname{GL}(A)\) is an integral homology equivalence. The statement does not concern \(BQ\mathcal P(A)\) itself.}
\GSfield{Reconstruction}{The ring, category, loop-space shift, and component are supplied. For example \(H_1(BQ\mathcal P(\mathbb F_2))\cong K_0(\mathbb F_2)=\mathbb Z\), whereas \(H_1(B\operatorname{GL}(\mathbb F_2))=0\), so the literal line cannot be retained.}
\GSfield{Fate}{The comparison is announced in \citet[introduction, pp.~86--87]{quillen1973higher}, which expressly defers its proof; see the later treatment in \citet[Chapter~IV, Q-construction]{weibel2013kbook}. The corrected theorem was already available in the 1972--1973 work.}
\GSfield{Present status}{As of 9 November 2026: Established with the loop and component restored. The displayed notebook comparison is false as written.}
\GSfield{Assessment}{literal statement: false as written; publication match: partial; verification: read the result; development: annotated; confidence: high}
\end{GSitem}
\begin{GSitem}{1973-314.3}{E}{K-zero as the fundamental group of BQ}
\GSfield{Kind; source}{theorem; \GSpage{P000527}}
\GSfield{The notes}{The Grothendieck group is generated by objects with relations \(X=X'+X''\) from short exact sequences; the notes identify it with \(K_0(M)=\pi_1(BQM,0)\).}
\GSfield{Statement}{For an essentially small exact category \(\mathcal M\), let \(K_0(\mathcal M)\) be the abelian group on isomorphism classes \([X]\), modulo \([X]=[X']+[X'']\) for every admissible short exact sequence. There is a canonical isomorphism \(K_0(\mathcal M)\cong\pi_1(BQ\mathcal M,0)\).}
\GSfield{Reconstruction}{Isomorphism classes and admissibility are explicit. The split exact sequences explain commutativity but do not replace the nonsplit relations.}
\GSfield{Fate}{This is the basic \(K_0\)-identification in \citet[\S2]{quillen1973higher}, circulated in the 1972 Seattle lectures and published in 1973.}
\GSfield{Present status}{As of 9 November 2026: A standard theorem; the indexing of homotopy groups is correct in this line.}
\GSfield{Assessment}{literal statement: correct; publication match: exact; verification: read the result; development: annotated; confidence: high}
\end{GSitem}
\begin{GSitem}{1973-314.4}{E}{Object loops and additivity through two-simplices}
\GSfield{Kind; source}{method; \GSpage{P000527}}
\GSfield{The notes}{For each \(X\), the sketch defines a loop \(\zeta_X\) through \(0\) and \(X\). The triangular picture for \(0\to X'\to X\to X''\to0\) yields \(\zeta_X\sim\zeta_{X'}\zeta_{X''}\), invoking two-simplices.}
\GSfield{Statement}{There are two canonical \(Q\)-arrows \(0\to X\): the spans \(0\leftarrow0\to X\) and \(0\leftarrow X\xrightarrow{1}X\). Traversing the first and then the reverse of the second gives a based loop, with a choice of global sign convention. The nerve's two-simplices associated with an admissible exact sequence identify its class with the product of the corresponding classes for \(X'\) and \(X''\).}
\GSfield{Reconstruction}{The backwards path is a path in the realization, not an inverse morphism in \(Q\mathcal M\). We spell out the spans to disambiguate the diagram.}
\GSfield{Fate}{This is the elementary nerve proof underlying \citet[\S2, calculation of \(\pi_1\)]{quillen1973higher}, not a new additivity theorem.}
\GSfield{Present status}{As of 9 November 2026: An identifiable proof of the preceding theorem; the diagram is meaningful after distinguishing arrows from reversed paths.}
\GSfield{Assessment}{literal statement: correct with missing hypotheses; publication match: exact; verification: read the result; development: annotated; confidence: high}
\end{GSitem}

\GStalk{315}{1973-10-29}{A. A. Kirillov}{P000528}{Bundles associated with the metaplectic representation}
Kirillov resumes the symplectic-spinor discussion: local symplectic coordinates, lifted transition functions, and the Schwartz--Hilbert--tempered-distribution triple of associated bundles. A Heisenberg-algebra bundle is identified with \(TM\oplus\mathbb R\), but the accompanying degree bound for polynomials is one degree too large. Lagrangian planes are introduced at the end, before Schnirelman's heading.

\begin{GSitem}{1973-315.1}{E}{Darboux coordinates and symplectic transition functions}
\GSfield{Kind; source}{theorem; \GSpage{P000528}}
\GSfield{The notes}{Locally \(\omega=\sum dp_k\wedge dq_k\), and tangent-bundle transition maps lie in \(\operatorname{Sp}(2n,\mathbb R)\). They are assumed lifted to the two-sheeted cover.}
\GSfield{Statement}{Every smooth symplectic \(2n\)-manifold has Darboux charts in which \(\omega\) has the displayed form. Derivatives of coordinate changes preserve \(\omega\) and therefore are symplectic linear maps pointwise. A metaplectic structure is a compatible cocycle of lifts, not merely separate lifts of individual transition functions.}
\GSfield{Reconstruction}{We distinguish nonlinear coordinate changes from their derivatives, and impose the cocycle condition on the lifts.}
\GSfield{Fate}{This is the local and bundle background used in \citet[\S3]{krysl2017spinors}; the topological obstruction is item 1973-312.1. No new local normal-form theorem is proposed.}
\GSfield{Present status}{As of 9 November 2026: Standard Darboux and principal-bundle theory.}
\GSfield{Assessment}{literal statement: correct with missing hypotheses; publication match: exact; verification: read the result; development: annotated; confidence: high}
\end{GSitem}
\begin{GSitem}{1973-315.2}{E}{The Schwartz--Hilbert--distribution triple of spinor bundles}
\GSfield{Kind; source}{construction; \GSpage{P000528}}
\GSfield{The notes}{The spinor representation induces \(HM\), with \(SM\subset HM\subset S'M\); \(SM\) is labelled rapidly decreasing.}
\GSfield{Statement}{Given a metaplectic principal bundle \(P\to M\), the invariant triple \(\mathcal S(\mathbb R^n)\subset L^2(\mathbb R^n)\subset\mathcal S'(\mathbb R^n)\) yields associated bundles \(P\times_{\operatorname{Mp}}\mathcal S\), \(P\times_{\operatorname{Mp}}L^2\), and \(P\times_{\operatorname{Mp}}\mathcal S'\). The Hilbert embedding in tempered distributions uses the usual pairing convention. These are infinite-dimensional bundles with their respective locally convex topologies.}
\GSfield{Reconstruction}{The fibre models, topology, and chosen metaplectic structure are supplied; \(S'M\) is a distribution bundle, not the algebraic dual of all sections.}
\GSfield{Fate}{The spinor bundle and oscillator realization are described in \citet[\S\S2.2--3.1]{krysl2017spinors}; \citet[\S3, Definition~2]{krysl2006spinor} also constructs the associated bundle. This records Kostant's existing construction.}
\GSfield{Present status}{As of 9 November 2026: An established construction. The representation preserves Schwartz space, so the three fibres glue compatibly.}
\GSfield{Assessment}{literal statement: correct; publication match: exact; verification: read the result; development: annotated; confidence: high}
\end{GSitem}
\begin{GSitem}{1973-315.3}{E}{The Heisenberg bundle and the polynomial degree error}
\GSfield{Kind; source}{claim; \GSpage{P000528}}
\GSfield{The notes}{\(GM=TM\oplus(\mathbb R\times M)\) is called the Heisenberg Lie algebra bundle and is labelled polynomials of degree at most two in \(p,q\). It acts on the spinor bundle.}
\GSfield{Statement}{The fibre \(T_xM\oplus\mathbb R\) with \([(v,a),(w,b)]=(0,\omega_x(v,w))\) is the \((2n+1)\)-dimensional Heisenberg Lie algebra. Under Poisson brackets it corresponds to affine-linear polynomials, of degree at most one. Polynomials of degree at most two form the larger semidirect extension by \(\mathfrak{sp}(2n,\mathbb R)\). Its infinitesimal Schrödinger action is defined on the Schwartz fibres, not as bounded operators on every Hilbert fibre.}
\GSfield{Reconstruction}{We repair the degree and specify the invariant domain of the action. The central normalization may be rescaled consistently.}
\GSfield{Fate}{The Heisenberg and Weyl-algebra realization is standard background in \citet[\S\S2.1--2.3]{krysl2017spinors}. A dimension count already separates the two polynomial algebras.}
\GSfield{Present status}{As of 9 November 2026: The named Heisenberg construction is established; the degree-two label is a faithful notebook error.}
\GSfield{Assessment}{literal statement: false as written; publication match: partial; verification: read the result; development: annotated; confidence: high}
\end{GSitem}
\begin{GSitem}{1973-315.4}{E}{Lagrangian tangent planes}
\GSfield{Kind; source}{definition; \GSpage{P000528}}
\GSfield{The notes}{The final line introduces \(RM\) and \(R_x\), Lagrangian subspaces of \(T_xM\) of maximal dimension.}
\GSfield{Statement}{A subspace \(L\subset T_xM\) is Lagrangian if \(\omega_x|_L=0\) and \(\dim L=n\), equivalently \(L=L^{\perp_\omega}\). All such planes form the fibre of the Lagrangian Grassmannian bundle over \(M\). A global section is additional structure and is not asserted to exist.}
\GSfield{Reconstruction}{The equality of maximal isotropic dimension with half the tangent dimension is made explicit. The notation does not identify a chosen polarization.}
\GSfield{Fate}{This is standard symplectic linear algebra in the setting reviewed by \citet[\S3]{krysl2017spinors}; Fuchs's short intervention on P000529 continues the bundle description.}
\GSfield{Present status}{As of 9 November 2026: An established definition, with no existence theorem for a Lagrangian distribution in the notes.}
\GSfield{Assessment}{literal statement: correct; publication match: exact; verification: read the result; development: annotated; confidence: high}
\end{GSitem}

\GStalk{316}{1973-10-29}{A. I. Schnirelman}{P000528-P000529}{Fermi--Pasta--Ulam recurrence and a periodic-wave heuristic}
The talk reports numerical recurrence of energy initially placed in the first Fourier mode of a weakly nonlinear chain. Zakharov suggests approximation by an integrable differential system. The discussion passes to a string and to perturbations of the wave equation on the circle. Parameters, boundary conditions, numerical units, and the averaged first-order drift system are not fully recorded, so the numerical sketch and the asserted periodic drift cannot be promoted to general theorems.

\begin{GSitem}{1973-316.1}{E}{The nonlinear chain equation and its Hamiltonian sign}
\GSfield{Kind; source}{claim; \GSpage{P000528}}
\GSfield{The notes}{The displayed chain has \(\ddot x_k=x_{k-1}+x_{k+1}-2x_k+\varepsilon(x_{k-1}-x_k)^2+\varepsilon(x_{k+1}-x_k)^2\). Both nonlinear signs are positive in the scan.}
\GSfield{Statement}{For a nearest-neighbour Hamiltonian chain with \(r_k=x_{k+1}-x_k\), potential \(V(r)=r^2/2+\varepsilon r^3/3\), and compatible fixed or periodic boundary conditions, \(\ddot x_k=V'(r_k)-V'(r_{k-1})\). Hence the quadratic terms occur as \(\varepsilon[(x_{k+1}-x_k)^2-(x_k-x_{k-1})^2]\). The notebook equation is a different dynamical system unless an omitted convention changes the model.}
\GSfield{Reconstruction}{The Hamiltonian and boundary-condition alternatives specify the standard alpha-FPU reconstruction. We do not silently edit the source or claim that its plus-plus equation has the same conservation law.}
\GSfield{Fate}{The FPU numerical experiment was reported in \citet[LA-1940]{fermipastaulam1955studies}, written in May 1955. The primary report record was checked; the Hamiltonian differentiation above supplies the sign check independently.}
\GSfield{Present status}{As of 9 November 2026: The standard FPU model is identifiable, but its equation is not faithfully represented by the two nonlinear plus signs.}
\GSfield{Assessment}{literal statement: false as written; publication match: partial; verification: publisher metadata checked; development: annotated; confidence: high}
\end{GSitem}
\begin{GSitem}{1973-316.2}{E}{Numerical recurrence of low-mode energy}
\GSfield{Kind; source}{remark; \GSpage{P000528}}
\GSfield{The notes}{For a chain labelled 16, the initial displacement is sinusoidal, with all energy in the first harmonic. The sketch shows transfer to other modes and near return. A cubic nonlinearity is said to give the same qualitative result.}
\GSfield{Statement}{The identifiable claim is an observation for particular finite nonlinear chains and initial conditions: modal energies may show pronounced near recurrence rather than rapid equipartition. It is not exact periodicity, not a statement for every initial condition, and not a bound uniform in chain length or energy.}
\GSfield{Reconstruction}{We distinguish a numerical observation from a quantified theorem. The notebook does not specify amplitudes, coupling strength, endpoint convention, or the formula used to define modal energy.}
\GSfield{Fate}{The phenomenon is the already reported FPU experiment \citep[LA-1940]{fermipastaulam1955studies}; the connection with soliton recurrence was studied by \citet[pp.~240--243]{zabuskykruskal1965solitons}. Both precede the talk.}
\GSfield{Present status}{As of 9 November 2026: The historical phenomenon is established, but the particular plotted experiment is not reproducible from these notes alone.}
\GSfield{Assessment}{literal statement: correct with missing hypotheses; publication match: partial; verification: publisher metadata checked; development: annotated; confidence: high}
\end{GSitem}
\begin{GSitem}{1973-316.3}{F}{Unidentified recurrence times and integration-step comments}
\GSfield{Kind; source}{claim; \GSpage{P000528}}
\GSfield{The notes}{The notes contain \(t=10^{-35}\), a caption approximately 50 million periods, and comments about ten steps and a one-third period. No physical units or numerical algorithm are given.}
\GSfield{Statement}{No definite numerical estimate can be reconstructed: it is not stated which time scale, error tolerance, integration scheme, or physical system the numbers refer to. The figure supports only the qualitative recurrence item 1973-316.2.}
\GSfield{Reconstruction}{We retain the negative exponent seen in the scan rather than changing it to match a plausible timescale. The disconnected numbers are not used as data for a computation.}
\GSfield{Fate}{Comparison with the identifiable FPU report \citep[LA-1940]{fermipastaulam1955studies} does not determine these units or the stated step counts. No exact numerical publication match was found.}
\GSfield{Present status}{As of 9 November 2026: Unidentified numerical details, not an open recurrence-time conjecture.}
\GSfield{Assessment}{literal statement: not determinable; publication match: none found; verification: unresolved lead; development: annotated; confidence: low}
\end{GSitem}
\begin{GSitem}{1973-316.4}{E}{Zakharov's integrable approximation explanation}
\GSfield{Kind; source}{method; \GSpage{P000529}; V. E. Zakharov}
\GSfield{The notes}{Zakharov proposes approximation by a differential system possessing first integrals as a possible cause of recurrence. A sinusoidally excited string is mentioned, followed by the remark that other initial data can lead to equipartition.}
\GSfield{Statement}{In a long-wavelength, weak-nonlinearity scaling, an FPU chain may be approximated over a specified time regime by a dispersive PDE such as KdV. Integrability can explain recurrent or coherent behaviour in that regime; it does not imply exact integrability of the original chain or exclude later energy transfer. The recorded equipartition remark has no universal quantifiers or rate.}
\GSfield{Reconstruction}{The scaling restrictions and limitation of the approximation are supplied. The notes propose an explanation, not a convergence theorem with an error bound.}
\GSfield{Fate}{The KdV and recurrence connection was already explored by \citet[pp.~240--243]{zabuskykruskal1965solitons} in 1965. The notebook's attribution to Zakharov concerns this seminar intervention, not priority for the general connection.}
\GSfield{Present status}{As of 9 November 2026: A recognizable established mechanism, with no precise new approximation or equipartition problem in the notes.}
\GSfield{Assessment}{literal statement: correct with missing hypotheses; publication match: partial; verification: publisher metadata checked; development: annotated; confidence: high}
\end{GSitem}
\begin{GSitem}{1973-316.5}{E}{Periodic trajectories of the wave equation on a circle}
\GSfield{Kind; source}{claim; \GSpage{P000529}}
\GSfield{The notes}{For \(u_{tt}=u_{xx}\) on \(S^1\), the notes say all trajectories are closed.}
\GSfield{Statement}{On a circle of length \(2\pi\), an energy solution has Fourier modes \(a_n\cos(nt)+b_n\sin(nt)\) for \(n\ne0\) and constant mode \(a_0+b_0t\). The solution is \(2\pi\)-periodic in the energy phase space if and only if the mean initial velocity \(b_0\) is zero. The spatially constant solution \(u(t,x)=t\) refutes the unrestricted assertion.}
\GSfield{Reconstruction}{The circle length and phase space are fixed for precision. The missing zero-mode hypothesis is diagnosed directly by Fourier series, not inferred as something the speaker necessarily said.}
\GSfield{Fate}{Elementary Fourier analysis gives both the corrected theorem and the explicit counterexample. This is old background within the discussion, not a later resolution of a research problem.}
\GSfield{Present status}{As of 9 November 2026: Correct on the zero-mean-velocity subspace; false for the full phase space.}
\GSfield{Assessment}{literal statement: false as written; publication match: none found; verification: indirect; development: annotated; confidence: high}
\end{GSitem}
\begin{GSitem}{1973-316.6}{F}{Periodic first-order drift for a cubic perturbation}
\GSfield{Kind; source}{claim; \GSpage{P000529}}
\GSfield{The notes}{The equations \(u_{tt}=u_{xx}+\varepsilon f'(u)\) and \(u_{tt}=u_{xx}-\varepsilon u^3\) are followed by the claim that the first-order drift of trajectories is also periodic.}
\GSfield{Statement}{The recorded statement lacks the averaged vector field, its domain, treatment of the zero mode, and the meaning of periodicity. Consequently neither a universal periodicity theorem for an averaged flow nor periodicity of every solution of the nonlinear PDE is a supported precise formulation.}
\GSfield{Reconstruction}{We preserve the claim without supplying an unrecorded normal-form or averaging construction. The nonlinear equation alone does not determine which drift system is intended.}
\GSfield{Fate}{Searches combining Schnirelman/Shnirelman, periodic wave equations, first-order averaging, and FPU recurrence did not identify an exact theorem with these quantifiers. The familiar integrable-approximation literature is only thematic and cannot fill the gap.}
\GSfield{Present status}{As of 9 November 2026: Not identifiable as a definite theorem or open problem from the available notes.}
\GSfield{Assessment}{literal statement: not determinable; publication match: none found; verification: unresolved lead; development: annotated; confidence: low}
\end{GSitem}

\GStalk{317}{1973-10-29}{D. B. Fuchs}{P000529}{The Lagrangian Grassmannian bundle}
A boxed intervention labels the space \(\Lambda\) of Lagrangian planes in \(\mathbb R^{2n}\) and the bundle \(RM\to M\). There is no further theorem, Maslov-class calculation, or quantization assertion in this block. It links Kirillov's notation on P000528 to the ensuing discussion without supplying an independent long talk.

\begin{GSitem}{1973-317.1}{E}{The bundle of Lagrangian planes}
\GSfield{Kind; source}{construction; \GSpage{P000529}}
\GSfield{The notes}{The box headed Fuchs defines \(\Lambda\) as the set of Lagrangian planes in \(\mathbb R^{2n}\) and depicts \(RM\to M\), a bundle of Lagrangian planes on a symplectic manifold.}
\GSfield{Statement}{For the symplectic frame bundle \(\operatorname{Fr}_{\operatorname{Sp}}(TM)\), form \[\Lambda(TM)=\operatorname{Fr}_{\operatorname{Sp}}(TM)\times_{\operatorname{Sp}(2n,\mathbb R)}\Lambda(n)\longrightarrow M,\qquad \Lambda(n)\cong U(n)/O(n).\] Its fibre consists of all Lagrangian subspaces of \(T_xM\). A section is a Lagrangian distribution, which need not be integrable or globally exist.}
\GSfield{Reconstruction}{The associated-bundle formula and homogeneous-space description explicate the picture. No Maslov class or global section is attributed to the notebook.}
\GSfield{Fate}{This is standard symplectic bundle geometry in the framework reviewed by \citet[\S3]{krysl2017spinors}; the brief source gives only the construction, not a distinct research claim.}
\GSfield{Present status}{As of 9 November 2026: Established background, continuing item 1973-315.4.}
\GSfield{Assessment}{literal statement: correct; publication match: exact; verification: read the result; development: annotated; confidence: high}
\end{GSitem}

\GStalk{318}{1973-10-29}{A. S. Mishchenko}{P000529-P000530}{Fredholm representations and higher signatures}
Mishchenko defines a Fredholm representation, applies it to nonsingular Hermitian forms over a ring with involution, and constructs a signature by the difference of two Fredholm indices. The culmination is the characteristic-number formula for a closed oriented manifold. The exact constructions and formula were publicly announced in March--April 1973, from a submission received in October 1972. The fuller 1974 paper supplies proofs and a more general ring formulation. These are therefore expository E items, despite their close match to a later full publication.

\begin{GSitem}{1973-318.1}{E}{Nonsingular Hermitian forms over a ring with involution}
\GSfield{Kind; source}{definition; \GSpage{P000529}}
\GSfield{The notes}{A matrix \(\lambda\) over \(\Lambda\) is called a quadratic form when \(\lambda^*=\lambda\) and \(\mu\lambda=1\) for some \(\mu\).}
\GSfield{Statement}{Let \(\Lambda\) be a unital ring with involution and \(\lambda\in M_n(\Lambda)\) satisfy \(\lambda^*=\lambda\) and \(\mu\lambda=1\). Taking adjoints gives \(\lambda\mu^*=1\); uniqueness of a left and right inverse yields \(\mu=\mu^*\), so \(\lambda\) is invertible. It defines a nonsingular Hermitian form on the free module. Over rings where two is not invertible, this is not the same data as a quadratic refinement in surgery theory.}
\GSfield{Reconstruction}{The term Hermitian is used for the actual displayed matrix conditions. The source's one-sided inverse is sufficient because self-adjointness provides the other side.}
\GSfield{Fate}{Hermitian forms are older background. The self-adjoint invertible matrix setting over a group ring is explicit in \citet[\S3]{mishchenko1973infinite}, published before the seminar; the general ring exposition is \citet[\S3]{mishchenko1974signatures}.}
\GSfield{Present status}{As of 9 November 2026: Established; the matrix definition is precise after distinguishing Hermitian forms from quadratic refinements.}
\GSfield{Assessment}{literal statement: correct with missing hypotheses; publication match: exact; verification: read the result; development: annotated; confidence: high}
\end{GSitem}
\begin{GSitem}{1973-318.2}{E}{Fredholm representations}
\GSfield{Kind; source}{definition; \GSpage{P000529}}
\GSfield{The notes}{\(T=(T_1,F,T_2)\), where the \(T_i\) are involution-preserving representations and \(F\) is Fredholm with \(FT_1(x)-T_2(x)F\) compact for every \(x\in\Lambda\).}
\GSfield{Statement}{Take unital star representations \(T_i:\Lambda\to\mathcal B(H_i)\) and a bounded Fredholm operator \(F:H_1\to H_2\) satisfying the stated compact-intertwining condition. For \(\Lambda=\mathbb Z[\pi]\), unitary representations of \(\pi\) extend to this data. Different Hilbert spaces are allowed.}
\GSfield{Reconstruction}{The two Hilbert spaces, boundedness, and unit-preserving condition are explicit. Compactness is pointwise for each ring element; no unrecorded uniform bound on the whole ring is imposed.}
\GSfield{Fate}{The group version already appears in \citet[\S1]{mishchenko1973infinite}, received 10 October 1972 and published in the March--April 1973 issue. The general ring definition is \citet[Definition~1.1]{mishchenko1974signatures}.}
\GSfield{Present status}{As of 9 November 2026: An established precursor of Fredholm-module formulations of K-homology.}
\GSfield{Assessment}{literal statement: correct; publication match: exact; verification: read the result; development: annotated; confidence: high}
\end{GSitem}
\begin{GSitem}{1973-318.3}{E}{Spectral splitting and the compact off-diagonal blocks}
\GSfield{Kind; source}{theorem; \GSpage{P000529}}
\GSfield{The notes}{The represented forms are bounded, self-adjoint and invertible. In positive/negative decompositions, \(F\otimes1\) has compact off-diagonal entries and Fredholm diagonal entries. The lower right entry is first written \(F_3\), then called \(F_4\).}
\GSfield{Statement}{Let \(\lambda_i=T_i^{(n)}(\lambda)\), and split \(H_i^n=H_i^+\oplus H_i^-\) by their spectral projections. The relation \(F^{(n)}\lambda_1-\lambda_2F^{(n)}\in\mathcal K\) implies that the two off-diagonal blocks of \(F^{(n)}\) are compact. Since \(F^{(n)}\) is Fredholm, its diagonal blocks \(F_+:H_1^+\to H_2^+\) and \(F_-:H_1^-\to H_2^-\) are Fredholm.}
\GSfield{Reconstruction}{The decomposition is spectral, not a finite-dimensional diagonalization by an eigenbasis. Functional calculus modulo compact operators justifies preservation of spectral signs because zero is separated from the spectra.}
\GSfield{Fate}{The splitting lemma was stated in \citet[\S3]{mishchenko1973infinite}; \citet[Lemmas~3.2--3.3]{mishchenko1974signatures} gives the fuller proof. The notebook has an index typo.}
\GSfield{Present status}{As of 9 November 2026: Established for the bounded invertible represented forms; the corrected lower right block is the negative-to-negative block.}
\GSfield{Assessment}{literal statement: false as written; publication match: partial; verification: read the result; development: annotated; confidence: high}
\end{GSitem}
\begin{GSitem}{1973-318.4}{E}{The T-signature of a form}
\GSfield{Kind; source}{construction; \GSpage{P000530}}
\GSfield{The notes}{\(\operatorname{sign}_T\lambda=\operatorname{index}F_1-\operatorname{index}F_4\), independent of the splitting and of the representative of the form class.}
\GSfield{Statement}{With the spectral blocks of item 1973-318.3, define \(\operatorname{Sign}_T(\lambda)=\operatorname{ind}F_+-\operatorname{ind}F_-\). It is invariant under isometry of nonsingular forms, additive under orthogonal sum, and zero on hyperbolic forms, hence descends to the corresponding Witt/L-group under the conventions of the construction.}
\GSfield{Reconstruction}{We identify precisely the equivalences required for descent. This is not a claim of invariance under an unspecified relation on arbitrary singular matrices.}
\GSfield{Fate}{The index difference and its well-defined Wall-group homomorphism are explicit in \citet[\S3]{mishchenko1973infinite}. The fuller account is \citet[Definition~3.4 and Theorem~3.5]{mishchenko1974signatures}.}
\GSfield{Present status}{As of 9 November 2026: Established; the index difference, not the sum, is the signature invariant.}
\GSfield{Assessment}{literal statement: correct; publication match: exact; verification: read the result; development: annotated; confidence: high}
\end{GSitem}
\begin{GSitem}{1973-318.5}{E}{The Fredholm-representation higher-signature formula}
\GSfield{Kind; source}{theorem; \GSpage{P000530}}
\GSfield{The notes}{For compact \(M\) of dimension \(4K\), \(\pi_1(M)=\pi\), and \(\Lambda=\mathbb Z[\pi]\), the notes state \(\operatorname{sign}_T\sigma(M)=\langle\operatorname{ch}\xi_T\,L(M),[M]\rangle\). They call \(\xi_T\) a bundle and \(\sigma(M)\) a quadratic form.}
\GSfield{Statement}{Let \(M\) be a closed oriented smooth \(4k\)-manifold, \(f:M\to B\pi\) the classifying map of its universal cover, and \(T\) a Fredholm representation of \(\pi\). With Mishchenko's Wall/Poincare-complex class \(\sigma(M)\) and associated K-theory class \(\xi_T\) pulled back along \(f\), \[\operatorname{Sign}_T\sigma(M)=\langle L(TM)\smile\operatorname{ch}(f^*\xi_T),[M]\rangle.\] The coefficient class is virtual in general; constructing \(\sigma(M)\) uses the algebraic Poincare complex, not an unjustified free middle-homology module.}
\GSfield{Reconstruction}{Closedness and orientation, the classifying map, the virtual K-class, and the surgery-theoretic meaning of \(\sigma(M)\) are supplied.}
\GSfield{Fate}{The exact formula is \citet[Theorem~1]{mishchenko1973infinite}, publicly printed before this talk. \citet[Theorem~5.1]{mishchenko1974signatures} gives the full treatment; announcement and full result were read.}
\GSfield{Present status}{As of 9 November 2026: The formula is established for the classes produced by Fredholm representations. It is not a proof of the Novikov conjecture for every group and every cohomology class.}
\GSfield{Assessment}{literal statement: correct with missing hypotheses; publication match: exact; verification: read the result; development: annotated; confidence: high}
\end{GSitem}
\begin{GSitem}{1973-318.6}{E}{Why finite-dimensional representations do not suffice}
\GSfield{Kind; source}{remark; \GSpage{P000530}}
\GSfield{The notes}{The talk ends ``\ru{Конечномерных предст. недостаточно}''; the ordinary Hirzebruch signature is indicated by \(\operatorname{ch}\xi_T=1\).}
\GSfield{Statement}{A complex vector bundle associated to a finite-dimensional unitary representation of \(\pi\) is flat. Its positive-degree rational Chern character vanishes, so the corresponding characteristic number sees only its rank times the ordinary signature. Fredholm representation pairs can produce nontrivial K-classes with higher-degree Chern character. This explains the need for the infinite-dimensional construction, without asserting that all K-classes arise this way for every group.}
\GSfield{Reconstruction}{The rational qualification and the limitation on representability of arbitrary K-classes are essential to the interpretation.}
\GSfield{Fate}{Flat-bundle rational Chern classes are older background. The pre-talk \citet[Theorem~3 and corollary]{mishchenko1973infinite} exhibits the reach of infinite-dimensional families under geometric assumptions; \citet[introduction and \S6]{mishchenko1974signatures} explains the motivation. These stronger geometric conclusions are not imported into the notebook.}
\GSfield{Present status}{As of 9 November 2026: An established motivation, not a universal representability theorem.}
\GSfield{Assessment}{literal statement: correct with missing hypotheses; publication match: partial; verification: read the result; development: annotated; confidence: high}
\end{GSitem}

\GStalk{319}{1973-11-12}{V. A. Ponomarev}{P000531-P000532}{Free modular lattices, subspace representations and perfect elements}
Ponomarev passes from lattice identities to representations of the free modular lattice by finite families of subspaces. He defines perfect elements and gives recursive lattice words and a Boolean family generated by sums omitting one generator. These definitions and constructions match the 1974 Gelfand--Ponomarev paper. The separate conjecture that every modular lattice is a submodule lattice is false: the Arguesian condition is an older additional restriction. A question about classification must distinguish the tame four-subspace problem from the matrix-pair complexity already present for five subspaces.

\begin{GSitem}{1973-319.1}{E}{Lattice operations, order and distributivity}
\GSfield{Kind; source}{definition; \GSpage{P000531}}
\GSfield{The notes}{Two associative, commutative, idempotent operations \(\cap,+\) satisfy absorption. The order is \(a\le b\iff a\cap b=a\). Distributive lattices are described as systems of subsets.}
\GSfield{Statement}{A lattice has meet \(a\wedge b\) and join \(a\vee b\) satisfying those identities; they determine the stated partial order. Distributivity is either, equivalently both, distributive identities. Every distributive lattice embeds into a lattice of subsets with intersection and union (with the usual prime-ideal separation principle); it need not be the entire powerset, nor be closed under complements.}
\GSfield{Reconstruction}{We interpret ``systems of subsets'' as a sublattice representation, not a Boolean-algebra assertion. The conventional separation principle is explicit.}
\GSfield{Fate}{These are the introductory definitions, also reproduced in \citet[\S1.1]{gelfandponomarev1974modular}; they were standard before 1973.}
\GSfield{Present status}{As of 9 November 2026: Established elementary lattice theory.}
\GSfield{Assessment}{literal statement: correct with missing hypotheses; publication match: exact; verification: read the result; development: annotated; confidence: high}
\end{GSitem}
\begin{GSitem}{1973-319.2}{E}{The modular law for submodule lattices}
\GSfield{Kind; source}{theorem; \GSpage{P000531}}
\GSfield{The notes}{For \(a\subseteq b\), \(b\cap(a+c)=a+(b\cap c)\), called Dedekind's law. Subspace and submodule lattices are modular but not generally distributive.}
\GSfield{Statement}{For submodules \(A\subseteq B\) and \(C\) of a module over any unital ring, \(B\cap(A+C)=A+(B\cap C)\). Thus its submodule lattice is modular. A vector space of dimension at least two has a nondistributive subspace lattice: three distinct lines in a two-dimensional subspace violate distributivity. Dimensions zero and one are exceptions to the notebook's unqualified phrase.}
\GSfield{Reconstruction}{The dimension restriction is supplied; modularity needs no finite-dimensionality assumption.}
\GSfield{Fate}{This is standard background, stated in \citet[\S1.1]{gelfandponomarev1974modular}.}
\GSfield{Present status}{As of 9 November 2026: Established, with the low-dimensional distributive cases separated.}
\GSfield{Assessment}{literal statement: correct with missing hypotheses; publication match: exact; verification: read the result; development: annotated; confidence: high}
\end{GSitem}
\begin{GSitem}{1973-319.3}{E}{Are all modular lattices submodule lattices?}
\GSfield{Kind; source}{conjecture; \GSpage{P000531}}
\GSfield{The notes}{The notes ask whether there are further identities and explicitly conjecture that every modular lattice is the lattice of submodules of a module; they immediately note that subspace lattices have further identities.}
\GSfield{Statement}{The literal conjecture is false. Submodule lattices satisfy the Arguesian (lattice-theoretic Desargues) identity, whereas a non-Desarguesian projective plane gives a modular lattice of height three that fails it. Such a lattice is neither isomorphic to a full submodule lattice nor embeddable as a sublattice of one. Modularity alone is therefore insufficient.}
\GSfield{Reconstruction}{The exact conjecture concerns submodules, not merely vector subspaces. We use the older Arguesian obstruction to answer that literal formulation; we do not replace it by an unspecified coordinatization problem.}
\GSfield{Fate}{The relevant extra condition was studied in \citet[pp.~294--313]{jonsson1954modular}, published in 1954. The submodule necessity is also stated explicitly in \citet[\S6]{herrmanntakach2005subgroup}. The conjecture was thus already false in the established theory, despite its appearance as a question in the notebook.}
\GSfield{Present status}{As of 9 November 2026: Disproved; the non-Desarguesian-plane obstruction answers the recorded universal assertion.}
\GSfield{Assessment}{literal statement: false as written; publication match: partial; verification: read the result; development: annotated; confidence: high; problem \hyperref[prob:GP-1973-30]{GP-1973-30}}
\end{GSitem}
\begin{GSitem}{1973-319.4}{A}{Free modular lattices and linear equivalence of terms}
\GSfield{Kind; source}{construction; \GSpage{P000531}}
\GSfield{The notes}{\(D^r\) is generated by \(e_1,\ldots,e_r\) modulo modular-lattice relations. Terms \(P,Q\) are linearly equivalent if they agree after substituting subspaces.}
\GSfield{Statement}{The free modular lattice \(D^r\) has the universal property for assignments of its generators to any modular lattice. For this notebook's finite-dimensional representation setting, define \(a\sim b\) if \(\rho(a)=\rho(b)\) for every lattice homomorphism \(\rho:D^r\to\mathcal L(V)\), every field \(k\), and every finite-dimensional \(k\)-space \(V\). This congruence can identify distinct abstract lattice elements.}
\GSfield{Reconstruction}{The quantifiers over fields and finite dimensions follow the subsequent definition of perfect elements. Abstract equality and equality in every linear representation remain separate.}
\GSfield{Fate}{\citet[\S1.2]{gelfandponomarev1974modular} gives precisely this notion and notes that distinct linearly equivalent elements occur already for four generators. The paper was received 10 June 1974 and published later that year.}
\GSfield{Present status}{As of 9 November 2026: An established distinction; identities valid in all subspace representations do not automatically hold in the free modular lattice.}
\GSfield{Assessment}{literal statement: correct; publication match: exact; verification: read the result; development: annotated; confidence: high}
\end{GSitem}
\begin{GSitem}{1973-319.5}{E}{Subspace systems and direct-sum decomposition}
\GSfield{Kind; source}{construction; \GSpage{P000531}}
\GSfield{The notes}{A representation sends meet to intersection and join to sum. A representation of \(D^r\) is exactly a finite-dimensional space \(V\) with subspaces \(E_1^0,\ldots,E_r^0\). Direct sums and indecomposability are defined.}
\GSfield{Statement}{An ordered \(r\)-subspace system is \((V;E_1,\ldots,E_r)\); its morphisms are linear maps carrying every \(E_i\) into the corresponding target subspace. By freeness this is equivalent to a representation of \(D^r\). A direct decomposition is \(V=V_1\oplus V_2\) with \(E_i=(E_i\cap V_1)\oplus(E_i\cap V_2)\) for every \(i\). Finite-dimensional systems decompose into indecomposables, uniquely up to isomorphism and permutation.}
\GSfield{Reconstruction}{The morphisms and simultaneous compatibility condition make the category explicit; arbitrary vector-space complements need not give decompositions of the system.}
\GSfield{Fate}{The subspace-system viewpoint is already present in \citet[\S\S1--2]{gelfandponomarev1971quadruples} and \citet[introduction]{bernsteingelfandponomarev1973coxeter}. Earlier seminar occurrences include 1970-198.1 and 1970-198.13.}
\GSfield{Present status}{As of 9 November 2026: Established representation theory; the number of subspaces determines the complexity of classification.}
\GSfield{Assessment}{literal statement: correct; publication match: exact; verification: read the result; development: annotated; confidence: high}
\end{GSitem}
\begin{GSitem}{1973-319.6}{E}{Classifying all representations of D-r}
\GSfield{Kind; source}{problem; \GSpage{P000531}}
\GSfield{The notes}{The explicit problem is ``describe all representations of \(D^r\)''. The ground field is arbitrary and the representations are finite-dimensional.}
\GSfield{Statement}{For fixed \(r\) and field \(k\), classify ordered \(r\)-subspace systems up to a simultaneous linear isomorphism. Over an algebraically closed field, \(r\le3\) has finite indecomposable type; \(r=4\) has a tame classification with one-parameter families. For \(r\ge5\), the problem contains simultaneous similarity of arbitrary pairs \((A,B)\): use \(W\oplus W\) with the two coordinate subspaces, the diagonal, and the graphs of \(A\) and \(B\). An isomorphism preserving the first three is \(g\oplus g\), so the remaining two encode simultaneous conjugacy.}
\GSfield{Reconstruction}{The explicit five-subspace reduction is supplied here. Wildness is a complexity result, not a claim that no mathematically useful invariants or restricted classifications can exist. The four-subspace result quoted is restricted to the field hypotheses of the sources.}
\GSfield{Fate}{The four-subspace classification circulated in the 1970 Tihany contribution and the 1971 announcement \citep{gelfandponomarev1971quadruples}. \citet[pp.~30--32]{bernsteingelfandponomarev1973coxeter} already distinguishes affine families from matrix-pair complexity. Compare earlier occurrences 1972-277.5--1972-277.7. GP-1970-23 concerns a specific deformation example and is not this general problem.}
\GSfield{Present status}{As of 9 November 2026: The finite/tame/wild distinction is established. A request for unrestricted normal forms beyond four subspaces cannot be treated as an isolated unsolved seminar conjecture.}
\GSfield{Assessment}{literal statement: correct with missing hypotheses; publication match: partial; verification: read the result; development: annotated; confidence: high; problem \hyperref[prob:GP-1973-31]{GP-1973-31}}
\end{GSitem}
\begin{GSitem}{1973-319.7}{A}{Perfect elements and compatible complements}
\GSfield{Kind; source}{definition; \GSpage{P000531}}
\GSfield{The notes}{An element is perfect if its image in every finite-dimensional representation over every field admits a complement splitting the representation. Equivalently, its image in each indecomposable is zero or the whole space. Linearly equivalent elements are perfect together.}
\GSfield{Statement}{For a modular lattice \(L\), call \(a\in L\) perfect if each finite-dimensional representation \(\rho\) splits with \(\rho(a)\) as one summand. Equivalently \(\rho(a)\in\{0,V\}\) for every indecomposable \(\rho:L\to\mathcal L(V)\): decompose a finite-dimensional representation into indecomposables and collect those on which \(a\) acts as the whole space. Perfection therefore depends only on the linear-equivalence class.}
\GSfield{Reconstruction}{We explain the equivalence using finite-dimensional decomposition. An arbitrary complement of a subspace is insufficient; it must split every \(\rho(b)\).}
\GSfield{Fate}{This exactly matches the definition in \citet[\S1.2]{gelfandponomarev1974modular}; the proof of the cubicle perfection theorem occupies later sections. Publication followed in 1974.}
\GSfield{Present status}{As of 9 November 2026: An established definition and equivalence. The notebook does not assert that it has classified all perfect elements.}
\GSfield{Assessment}{literal statement: correct; publication match: exact; verification: read the result; development: annotated; confidence: high}
\end{GSitem}
\begin{GSitem}{1973-319.8}{A}{Recursive lattice words indexed by adjacent-distinct sequences}
\GSfield{Kind; source}{construction; \GSpage{P000531-P000532}}
\GSfield{The notes}{The first word is \(e_{i_1i_2}=e_{i_1}\,[\text{overwritten sign}]\,\sum_{k\ne i_1,i_2}e_k\). The next page gives \(e_\alpha=e_{i_1}\sum_{\beta\in\Gamma(\alpha)}e_\beta\), but the conditions defining \(\Gamma\) end in an ellipsis.}
\GSfield{Statement}{The matching published construction is as follows. Let \(\alpha=(i_1,\ldots,i_\ell)\) have adjacent distinct entries in \(\{1,\ldots,r\}\). Let \(\Gamma(\alpha)\) consist of \(\beta=(k_1,\ldots,k_{\ell-1})\) with adjacent distinct entries and \(k_j\notin\{i_j,i_{j+1}\}\) for every \(j\). Put \(e_{(i)}=e_i\) and recursively \(e_\alpha=e_{i_1}\wedge\bigvee_{\beta\in\Gamma(\alpha)}e_\beta\). Empty joins use the least element of the generated bounded lattice.}
\GSfield{Reconstruction}{The meet sign and the completed indexing conditions come from the published formula, not from deciphering the overwritten mark. They are explicitly a literature-based reconstruction. We do not assert that each individual word is perfect.}
\GSfield{Fate}{\citet[\S1.3, equation~(1.1)]{gelfandponomarev1974modular} gives the matching recurrence; the full text was opened and this definition read. The later paper resolves the incomplete source notation without legitimizing silent transcription repair.}
\GSfield{Present status}{As of 9 November 2026: The construction is identified with high confidence; its complete formula is supplied from the 1974 publication.}
\GSfield{Assessment}{literal statement: not determinable; publication match: partial; verification: read the result; development: annotated; confidence: high}
\end{GSitem}
\begin{GSitem}{1973-319.9}{A}{The first Boolean cubicle of perfect elements}
\GSfield{Kind; source}{theorem; \GSpage{P000532}}
\GSfield{The notes}{The ``trivial perfect elements'' are \(m_t=\sum_{i\ne t}e_i\); the notes say that they generate a Boolean algebra with \(r\) generators.}
\GSfield{Statement}{In the free modular lattice on \(r\ge2\) generators, the sublattice generated by \(m_t=\bigvee_{i\ne t}e_i\) is the first upper cubicle \(B^+(1)\). It is a Boolean lattice with \(2^r\) elements, and every element is perfect. The \(m_t\) are its distinguished coatoms. This is not the free Boolean algebra on \(r\) independent Boolean generators, which would have \(2^{2^r}\) elements.}
\GSfield{Reconstruction}{We identify the cubicle and distinguish lattice generators from freely independent Boolean generators. The minimum and maximum are those of this sublattice, not necessarily those of every ambient representation.}
\GSfield{Fate}{\citet[Lemma~3.1 and Corollaries~4.1--4.2]{gelfandponomarev1974modular} proves the Boolean structure, perfection, and cardinality. These statements give an exact fate for the notebook claim.}
\GSfield{Present status}{As of 9 November 2026: Established in the 1974 paper. No completeness conjecture for all perfect elements is stated in the notebook and none is added to the problem index.}
\GSfield{Assessment}{literal statement: correct with missing hypotheses; publication match: exact; verification: read the result; development: annotated; confidence: high}
\end{GSitem}

\GStalk{321}{1973-11-12}{S. I. Gelfand}{P000532}{Affine and spherical buildings}
After the end of Ponomarev's lattice discussion, S. I. Gelfand describes the Bruhat--Tits complex as a discrete analogue of a symmetric space. Maximal compact subgroups and Iwahori subgroups encode vertices and chambers. The rank-one picture is a tree; a finite-dimensional flag-incidence example over a finite field illustrates a spherical building. In source order this block precedes Kazhdan's curve-counting talk, although the archive numbers it later.

\begin{GSitem}{1973-321.1}{E}{A p-adic homogeneous space as a vertex set}
\GSfield{Kind; source}{construction; \GSpage{P000532}}
\GSfield{The notes}{For a \(p\)-adic group \(G\), a maximal compact \(K\) is open, so \(G/K\) is discrete. The example is \(G=\operatorname{GL}_2(\mathbb Q_p)/C\), with \(K=\operatorname{GL}_2(\mathcal O)\).}
\GSfield{Statement}{For \(G=\operatorname{PGL}_n(\mathbb Q_p)\), let \(K\) be the image of \(\operatorname{GL}_n(\mathbb Z_p)\). Then \(K\) is compact open and \(G/K\) identifies with homothety classes of full \(\mathbb Z_p\)-lattices in \(\mathbb Q_p^n\). These are vertices of the affine building. For a general reductive group, one quotient by a chosen parahoric need not contain all vertex types.}
\GSfield{Reconstruction}{The quotient image of \(K\), the lattice model, and the restriction to type A are supplied. Maximal compact openness is not asserted for every conceivable topological group.}
\GSfield{Fate}{The affine-building framework is developed in \citet[\S\S2 and 6]{bruhattits1972groupes}, published in 1972. The notebook is an introductory exposition.}
\GSfield{Present status}{As of 9 November 2026: Established; the quotient and vertex-type qualifications prevent an overgeneralized construction.}
\GSfield{Assessment}{literal statement: correct with missing hypotheses; publication match: exact; verification: read the result; development: annotated; confidence: high}
\end{GSitem}
\begin{GSitem}{1973-321.2}{E}{The standard Iwahori subgroup}
\GSfield{Kind; source}{definition; \GSpage{P000532}}
\GSfield{The notes}{The displayed subgroup has integral upper-triangular entries, entries divisible by \(p\) below the diagonal, and unit diagonal entries; it is called Iwahori.}
\GSfield{Statement}{The standard Iwahori subgroup of \(\operatorname{GL}_n(\mathbb Q_p)\) is \(I=\{g\in\operatorname{GL}_n(\mathbb Z_p):g\bmod p\text{ is upper triangular}\}\). Its image in \(\operatorname{PGL}_n\) fixes the corresponding alcove in the affine building. A complete lattice chain determines the same stabilizer.}
\GSfield{Reconstruction}{We make the reduction map and passage to the projective group explicit; the notebook's letter \(B\) denotes this compact-open subgroup, not the algebraic Borel of \(G\).}
\GSfield{Fate}{This is standard Iwahori/parahoric structure within \citet[\S\S4 and 6]{bruhattits1972groupes}.}
\GSfield{Present status}{As of 9 November 2026: Established; the matrix display is correct with the group and reduction convention specified.}
\GSfield{Assessment}{literal statement: correct; publication match: exact; verification: read the result; development: annotated; confidence: high}
\end{GSitem}
\begin{GSitem}{1973-321.3}{E}{Simplices from common Iwahori stabilizers}
\GSfield{Kind; source}{construction; \GSpage{P000532}}
\GSfield{The notes}{Points \(x_i\in G/K\) are identified with conjugates \(K_i\); they span a simplex if a conjugate of the Iwahori subgroup lies in \(\bigcap K_i\).}
\GSfield{Statement}{For \(\operatorname{PGL}_n(\mathbb Q_p)\), distinct lattice-homothety vertices form a simplex when representatives can be ordered in a chain \(pL\subsetneq L_r\subsetneq\cdots\subsetneq L_1\subsetneq L\). Equivalently they are a subset of the vertices of one alcove, whose Iwahori fixes them. Vertex stabilizers are the corresponding conjugates of \(K\). Repeated vertices do not enlarge a simplex.}
\GSfield{Reconstruction}{The chain formulation specifies the construction in the example where the notebook's single vertex orbit works. General buildings require all relevant parahoric types.}
\GSfield{Fate}{The simplex/parahoric correspondence belongs to the affine-building construction in \citet[\S\S2 and 6]{bruhattits1972groupes}.}
\GSfield{Present status}{As of 9 November 2026: Established for the stated type-A model; a literal construction from an arbitrary single G/K is too broad.}
\GSfield{Assessment}{literal statement: correct with missing hypotheses; publication match: partial; verification: read the result; development: annotated; confidence: high}
\end{GSitem}
\begin{GSitem}{1973-321.4}{E}{The rank-one Bruhat--Tits tree}
\GSfield{Kind; source}{theorem; \GSpage{P000532}}
\GSfield{The notes}{For \(n=2\) the complex is declared contractible, with a tree drawn below.}
\GSfield{Statement}{The affine building of \(\operatorname{PGL}_2(\mathbb Q_p)\) is the tree whose vertices are homothety classes of full lattices in \(\mathbb Q_p^2\), with adjacency represented by \(pL\subsetneq L'\subsetneq L\). Its neighbours at \([L]\) correspond to lines in \(L/pL\cong\mathbb F_p^2\), so the valency is \(p+1\). A tree is contractible.}
\GSfield{Reconstruction}{The lattice description and valency explain the drawing. The manuscript does not label it as a finite tree or attach a numerical branching count.}
\GSfield{Fate}{The rank-one instance is contained in the pre-existing affine-building theory of \citet[\S2]{bruhattits1972groupes}.}
\GSfield{Present status}{As of 9 November 2026: Established; the contractibility claim is correct in this example.}
\GSfield{Assessment}{literal statement: correct; publication match: exact; verification: read the result; development: annotated; confidence: high}
\end{GSitem}
\begin{GSitem}{1973-321.5}{E}{The finite flag-incidence building in dimension three}
\GSfield{Kind; source}{construction; \GSpage{P000532}}
\GSfield{The notes}{In \(V_3(\mathbb F_q)\), vertices of type one are lines and of type two are planes; incidence spans a simplex.}
\GSfield{Statement}{The spherical building of \(\operatorname{PGL}_3(\mathbb F_q)\) has as vertices the nonzero proper subspaces of \(\mathbb F_q^3\); edges are flags \(L\subset P\) with dimensions one and two. It is a bipartite incidence graph, not the contractible affine tree of the preceding item.}
\GSfield{Reconstruction}{We distinguish this finite spherical example from the affine example. No claim that every building in the block is contractible is imputed to the notes.}
\GSfield{Fate}{This is the flag-complex model underlying the building theory used in \citet[\S1, combinatorial background]{bruhattits1972groupes}.}
\GSfield{Present status}{As of 9 November 2026: Established incidence geometry; its type and topology differ from the preceding rank-one affine example.}
\GSfield{Assessment}{literal statement: correct; publication match: exact; verification: read the result; development: annotated; confidence: high}
\end{GSitem}

\GStalk{320}{1973-11-12}{D. A. Kazhdan}{P000532-P000533}{The Stepanov--Bombieri method for counting points on curves}
Kazhdan explains a Riemann--Roch and Frobenius construction after Bombieri's June 1973 Bourbaki lecture on Stepanov's method. The notes omit the constant one in the projective Weil bound and compress two distinct multiplication maps into a single diagram. They also abbreviate the deduction of the full Weil estimate from a coarse upper bound, omitting the uniformity and covering argument emphasized by Bombieri. These are identifiable proof-compression defects, not new unsolved point-counting questions.

\begin{GSitem}{1973-320.1}{F}{Cancelled rational-function and exponential-sum fragment}
\GSfield{Kind; source}{remark; \GSpage{P000532}}
\GSfield{The notes}{Before the curve-counting discussion, crossed-out lines contain \(ax+b/x\) with \(a,b\in\mathbb Z\), a partly illegible superscript, and an exponential-sum fragment.}
\GSfield{Statement}{No summation range, finite field, character, or estimate survives. The fragment cannot identify a definite Kloosterman-sum claim or any other precise exponential-sum theorem.}
\GSfield{Reconstruction}{The visible rational expression is retained as cancelled source evidence. We do not infer the missing summand from a familiar formula.}
\GSfield{Fate}{No exact publication identification is warranted by the cancelled fragment. The subsequent Bombieri discussion is independently identifiable and does not restore this text.}
\GSfield{Present status}{As of 9 November 2026: Unidentifiable cancelled material; it has no problem-index entry.}
\GSfield{Assessment}{literal statement: not determinable; publication match: none found; verification: unresolved lead; development: annotated; confidence: low}
\end{GSitem}
\begin{GSitem}{1973-320.2}{E}{The Weil bound and the missing constant one}
\GSfield{Kind; source}{theorem; \GSpage{P000532}}
\GSfield{The notes}{For a curve \(X/\mathbb F_q\) of genus \(g\), the notes write \(|[X_q]-q|\le2g\sqrt q\) and identify it with the Weil hypothesis.}
\GSfield{Statement}{For a smooth projective geometrically connected curve \(X/\mathbb F_q\) of genus \(g\), \[|\#X(\mathbb F_q)-(q+1)|\le2g\sqrt q.\] Equivalently the reciprocal roots of the numerator of its zeta function have complex absolute value \(\sqrt q\); the equivalence uses the point counts over all finite extensions. An affine curve requires adjustment for omitted points.}
\GSfield{Reconstruction}{Smoothness, projectivity, geometric connectedness, and the extension-field formulation are supplied. For \(\mathbb P^1\), the notebook expression gives \(1\le0\), proving that the missing one matters.}
\GSfield{Fate}{This is the established curve Riemann hypothesis reviewed in \citet[\S I]{bombieri1974counting}. Bombieri's lecture was delivered in June 1973 and printed in 1974; it is not the first proof of Weil's theorem.}
\GSfield{Present status}{As of 9 November 2026: Established after restoring q+1 and the curve hypotheses; the literal displayed bound is false.}
\GSfield{Assessment}{literal statement: false as written; publication match: exact; verification: read the result; development: annotated; confidence: high}
\end{GSitem}
\begin{GSitem}{1973-320.3}{E}{A coarse upper bound and the deduction of the full Weil estimate}
\GSfield{Kind; source}{method; \GSpage{P000532}}
\GSfield{The notes}{The notes attribute \([X_q]\le q+4g\sqrt q\) to Stepanov through Bombieri and say that any constant independent of \(q\) would formally suffice for the sharp result.}
\GSfield{Statement}{The Stepanov--Bombieri strategy constructs auxiliary functions to obtain a square-root error estimate and then uses extension fields and suitable covers or twists to recover the zeta-function conclusion. A one-sided upper estimate for one curve over one field, or even only its extension sequence without the necessary companion information, is not the complete deduction. The printed \(q+4g\sqrt q\) also needs an endpoint correction in genus zero.}
\GSfield{Reconstruction}{We retain the method and do not manufacture an exact theorem from the abbreviated constant. Bombieri's hypotheses for the intermediate estimate and the extra covering step are part of the literature, not present in the notebook.}
\GSfield{Fate}{\citet[\S\S~II--III]{bombieri1974counting} is the direct source. Section III explicitly warns that an upper bound alone does not establish the Riemann hypothesis, and supplies the additional argument. The June 1973 circulation predates this talk.}
\GSfield{Present status}{As of 9 November 2026: The method is established; the compressed inference is incomplete and the displayed constant is not quoted as an exact published theorem.}
\GSfield{Assessment}{literal statement: correct with missing hypotheses; publication match: partial; verification: read the result; development: annotated; confidence: high}
\end{GSitem}
\begin{GSitem}{1973-320.4}{E}{Riemann--Roch dimension for a large divisor}
\GSfield{Kind; source}{theorem; \GSpage{P000532}}
\GSfield{The notes}{For a divisor \(D\), the notes call \(L(D)\) the number of functions divisible by \(D\), and write \(L(D)=\deg D-g+1\) for \(\deg D>2g\).}
\GSfield{Statement}{For a divisor \(D\) on a smooth projective geometrically connected curve, put \(L(D)=H^0(X,\mathcal O_X(D))=\{f\in k(X):\operatorname{div}(f)+D\ge0\}\cup\{0\}\) and \(\ell(D)=\dim_k L(D)\). Riemann--Roch gives \(\ell(D)=\deg D-g+1\) when \(\deg D>2g-2\). The notebook's stronger threshold is sufficient. Over \(\mathbb F_q\), the actual number of elements is \(q^{\ell(D)}\).}
\GSfield{Reconstruction}{We distinguish the vector space, its dimension, and its cardinality, and fix the divisor sign convention.}
\GSfield{Fate}{This is the standard Riemann--Roch input used in \citet[\S~II]{bombieri1974counting}.}
\GSfield{Present status}{As of 9 November 2026: Established; the notebook formula concerns dimension, not the cardinality of a space of functions.}
\GSfield{Assessment}{literal statement: correct with missing hypotheses; publication match: exact; verification: read the result; development: annotated; confidence: high}
\end{GSitem}
\begin{GSitem}{1973-320.5}{E}{Frobenius powers of a Riemann--Roch space}
\GSfield{Kind; source}{construction; \GSpage{P000533}}
\GSfield{The notes}{A point \(x_0\) is fixed, \(L_e\) consists of functions with a pole of order at most \(e\) there, and \(L_e^q\) has the same number of elements as \(L_e\).}
\GSfield{Statement}{Assume \(x_0\in X(\mathbb F_q)\) and \(L_e=L(ex_0)\). Then \(f\mapsto f^q\) is an injective \(\mathbb F_q\)-linear map with image \(L_e^q\subset L(qex_0)\), so \(\dim L_e^q=\dim L_e\). It is not the assertion that \(L(qex_0)=L_e^q\). Over an algebraic closure the map is semilinear, with the same dimension conclusion.}
\GSfield{Reconstruction}{The rationality of the base point, the image-space notation, and the increased pole order are supplied.}
\GSfield{Fate}{This Frobenius device is part of \citet[\S~II]{bombieri1974counting}. The notebook retains the essential equal-dimension observation.}
\GSfield{Present status}{As of 9 November 2026: Established with the image and pole order distinguished.}
\GSfield{Assessment}{literal statement: correct with missing hypotheses; publication match: exact; verification: read the result; development: annotated; confidence: high}
\end{GSitem}
\begin{GSitem}{1973-320.6}{E}{The two multiplication maps in the auxiliary-function construction}
\GSfield{Kind; source}{method; \GSpage{P000533}}
\GSfield{The notes}{A kernel is drawn in \(L_{q-1}\times L_e^q\to L_{q-1}L_e\subset L_{q+e-1}\), with ``isom. tensor prod.'' below. The dimension comparison is \((q-g)(e-g+1)>q+e-g\); the resulting functions vanish on rational points.}
\GSfield{Statement}{Distinguish \(m:L_{q-1}\otimes_{\mathbb F_q}L_e\to L_{q+e-1}\), \(a\otimes b\mapsto ab\), from \(m_q:L_{q-1}\otimes L_e\to L_{q-1+qe}\), \(a\otimes b\mapsto ab^q\). If the recorded dimension inequality applies, \(\ker m\ne0\). For \(v=\sum a_i\otimes b_i\in\ker m\), \(m_q(v)\) vanishes at every rational point except the pole point because \(b_i(P)^q=b_i(P)\). To count zeros one must additionally prove \(m_q(v)\ne0\) and control vanishing multiplicities. The notebook omits these essential steps.}
\GSfield{Reconstruction}{The tensor product and the separation into two maps are our reconstruction from the source method. We do not claim that this bare dimension count by itself proves the square-root bound.}
\GSfield{Fate}{The auxiliary-function and Frobenius-independence argument is in \citet[\S~II]{bombieri1974counting}; the notes compress different spaces and maps into one line. The publication is an exact methodological match, not a validation of the literal diagram.}
\GSfield{Present status}{As of 9 November 2026: The established method works with additional nonvanishing and multiplicity arguments; the notebook diagram has a genuine pole-order mismatch.}
\GSfield{Assessment}{literal statement: false as written; publication match: partial; verification: read the result; development: annotated; confidence: high}
\end{GSitem}

\GStalk{322}{1973-11-19}{L. N. Vasserstein}{P000534-P000535}{Stable rank, Serre's problem and a conic reduction question}
The opening announcement reports a three-variable case of Serre's problem. Vasserstein then defines stable rank, connects the two-variable case to \(K_2(k)\), and exhibits the real degree obstruction. An incompletely defined winding symbol and symplectic \(K_1\) lead to a concrete conic question for polynomial, entire and continuous functions. The 1974 announcement matches both genuine results and an erroneous field dichotomy; the 1976 full paper separates the remaining field cases. The conic wording needs its coefficient field and the meaning of ``common divisor'' restored before it has a definite answer.

\begin{GSitem}{1973-322.1}{A}{The announced three-variable case of Serre's problem}
\GSfield{Kind; source}{claim; \GSpage{P000534}}
\GSfield{The notes}{Before the speaker heading the notes say: ``\ru{Решена пробл. Серра о мног. от 3-х переменных}''. No coefficient field or characteristic restriction is recorded.}
\GSfield{Statement}{The nearby 1974 result states that every finitely generated projective \(k[x_1,x_2,x_3]\)-module is free when \(k\) is perfect or \(\operatorname{char}k\ne2\). The modern unrestricted theorem holds for every field \(k\) and every finite number of variables. These are different scopes: the notebook alone does not establish the omitted imperfect characteristic-two case in November 1973.}
\GSfield{Reconstruction}{The restriction is supplied from the publication, not read into the handwritten announcement. The all-variable result is its later fate, not a theorem attributed to this meeting.}
\GSfield{Fate}{\citet[Theorem~1 and following remarks]{vasersteinsuslin1974serre}, received 26 November 1973, records work begun in January 1973 and the stated three-variable scope. The unrestricted result followed in 1976: \citet[\S III, Theorem~2]{suslin1976projective} and independently \citet[pp.~167--171]{quillen1976projective}. Their submission dates were 26 February and 1 March 1976 respectively.}
\GSfield{Present status}{As of 9 November 2026: Serre's polynomial-ring problem over fields is solved. The 1973 announcement should retain its contemporary scope.}
\GSfield{Assessment}{literal statement: correct with missing hypotheses; publication match: partial; verification: read the result; development: annotated; confidence: high; problem \hyperref[prob:GP-1973-33]{GP-1973-33}}
\end{GSitem}
\begin{GSitem}{1973-322.2}{E}{Bass stable rank at most two}
\GSfield{Kind; source}{definition; \GSpage{P000534}}
\GSfield{The notes}{For every unimodular column \((a,b,c)^t\), there should exist \(\alpha,\beta\in A\) making \((a+\alpha c,b+\alpha c)^t\) unimodular. The scan repeats \(\alpha\) in the second component, despite quantifying \(\beta\). Unimodular is glossed as relatively prime.}
\GSfield{Statement}{For a commutative unital ring \(A\), write \(\operatorname{Um}_r(A)=\{(a_1,\ldots,a_r):\sum Aa_i=A\}\). Then \(\operatorname{sr}(A)\le2\) means that for every \((a,b,c)\in\operatorname{Um}_3(A)\) some \(\alpha,\beta\in A\) give \((a+\alpha c,b+\beta c)\in\operatorname{Um}_2(A)\). Having no common nonunit factor is weaker than generating the unit ideal in a multivariable polynomial ring.}
\GSfield{Reconstruction}{We repair the second coefficient and use the unit-ideal definition. The discussion specializes to commutative rings; noncommutative left/right conventions are not reconstructed unnecessarily.}
\GSfield{Fate}{Bass introduced the stable-range framework \citep[Chapter~I, \S4]{bass1964stable}; \citet[p.~17 and Theorem~1]{vaserstein1971stable} gives the precise unimodular-column formulation before this seminar.}
\GSfield{Present status}{As of 9 November 2026: Standard definition, with two independent coefficients and genuine unimodularity.}
\GSfield{Assessment}{literal statement: false as written; publication match: exact; verification: read the result; development: annotated; confidence: high}
\end{GSitem}
\begin{GSitem}{1973-322.3}{E}{Nonzero K2 forces stable rank three in two variables}
\GSfield{Kind; source}{theorem; \GSpage{P000534}}
\GSfield{The notes}{The first branch of the displayed theorem is \(\operatorname{sr}k[x_1,x_2]=3\) when \(K_2k\ne0\).}
\GSfield{Statement}{For any field \(k\), if its Milnor group \(K_2^M(k)\) is nonzero, then \(\operatorname{sr}(k[x_1,x_2])=3\). For fields the usual algebraic \(K_2\) agrees with this Milnor group. Bass's dimension bound gives the upper bound three; the nontrivial content is nonreducibility of some unimodular triple.}
\GSfield{Reconstruction}{We specify the field and the \(K_2\) convention. This theorem concerns existence of an obstructed triple, not every triple, and therefore does not contradict the reducibility of the particular conic triple below.}
\GSfield{Fate}{\citet[\S19, p.~994]{vasersteinsuslin1976serre} explicitly attributes this result to \citet{krusemeyer1973fundamental}, published in March 1973 from a submission of June 1972. The precise consequence was read in the former; the original article's chronology was publisher-verified.}
\GSfield{Present status}{As of 9 November 2026: Established and already available before the talk. The later symplectic formulation is not evidence of a new 1973 theorem by Vasserstein.}
\GSfield{Assessment}{literal statement: correct; publication match: exact; verification: read the result; development: annotated; confidence: high}
\end{GSitem}
\begin{GSitem}{1973-322.4}{A}{The unrestricted zero-K2 branch of the stable-rank assertion}
\GSfield{Kind; source}{claim; \GSpage{P000534}}
\GSfield{The notes}{The second branch says \(K_2k=0\Longrightarrow\operatorname{sr}k[x_1,x_2]=2\), with no qualification on \(k\).}
\GSfield{Statement}{The verified theorem is \(\operatorname{sr}(k[x_1,x_2])=2\) for fields \(k\) algebraic over a finite field. The notebook's stronger implication for every field with \(K_2(k)=0\) is a definite assertion, but the sources checked here do not establish it. Infinite algebraic extensions of global fields with vanishing \(K_2\) are not excluded by the hypothesis.}
\GSfield{Reconstruction}{We separate the proved algebraic-over-finite-field case from the stronger recorded implication. The absence of a verified modern resolution of the latter in this bounded search is not asserted to mean it is open today.}
\GSfield{Fate}{\citet[Theorem~4 and following remark]{vasersteinsuslin1974serre} prints the same conclusion through the false field implication in item 1973-322.5. The fuller \citet[\S19, p.~994]{vasersteinsuslin1976serre} explicitly leaves the additional field cases unresolved at that time; \citet[introduction]{suslin1980reciprocity} retains the distinction. Class A records publication of the assertion, not proof of its unrestricted validity.}
\GSfield{Present status}{As of 9 November 2026: Verified for fields algebraic over finite fields. The unrestricted zero-K2 implication is not settled by the literature verified in this chapter; no present-day C or D claim is made.}
\GSfield{Assessment}{literal statement: not determinable; publication match: exact; verification: read the result; development: annotated; confidence: medium}
\end{GSitem}
\begin{GSitem}{1973-322.5}{A}{The erroneous classification of fields with vanishing K2}
\GSfield{Kind; source}{claim; \GSpage{P000534}}
\GSfield{The notes}{The parenthesis asserts \(K_2k=0\) for algebraic extensions of finite fields and \(K_2k\ne0\) in every other case.}
\GSfield{Statement}{The forward assertion for fields algebraic over a finite field is correct, but the claimed converse is false: \(K_2(\overline{\mathbb Q})=0\), while \(\overline{\mathbb Q}\) has characteristic zero and is not algebraic over any finite field.}
\GSfield{Reconstruction}{The explicit counterexample answers the literal field assertion. It does not by itself decide the stable rank of \(\overline{\mathbb Q}[x_1,x_2]\).}
\GSfield{Fate}{The erroneous implication is also printed after Theorem~4 in \citet{vasersteinsuslin1974serre}. It is replaced by a more careful statement in \citet[\S19]{vasersteinsuslin1976serre}. The concrete counterexample is recorded in \citet[III, Exercise~6.5, p.~243]{weibel2013kbook}.}
\GSfield{Present status}{As of 9 November 2026: False as written; publication of the same claim does not validate it.}
\GSfield{Assessment}{literal statement: false as written; publication match: exact; verification: read the result; development: annotated; confidence: high}
\end{GSitem}
\begin{GSitem}{1973-322.6}{E}{The real degree obstruction to reducing the sphere row}
\GSfield{Kind; source}{theorem; \GSpage{P000535}}
\GSfield{The notes}{Under ``example over \(\mathbb R\)'', the column \((x_1,x_2,1-x_1^2-x_2^2)^t\) cannot be reduced; the reason is ``\ru{из степени}''. The coefficients are written \(\alpha,\beta\in\mathbb R\), without explicitly saying polynomial functions.}
\GSfield{Statement}{The row \((x,y,1-x^2-y^2)\) is unimodular over \(\mathbb R[x,y]\), because \(x^2+y^2+(1-x^2-y^2)=1\). There are no continuous real-valued \(f_1,f_2\) on the closed unit disk for which \((x+(1-x^2-y^2)f_1,y+(1-x^2-y^2)f_2)\) is everywhere nonzero. Normalization would extend the identity map of the boundary circle to a disk-to-circle map, contradicting its degree. In particular polynomial reduction is impossible over \(\mathbb R\).}
\GSfield{Reconstruction}{The stronger continuous statement and the disk formulation make the recorded degree argument precise. It applies to real-valued functions and cannot be transferred unchanged to complex-valued functions.}
\GSfield{Fate}{This is the obstruction used in \citet[Theorem~8 and its proof, p.~27]{vaserstein1971stable}, received in April 1970 and published in 1971. Our appendix supplies the elementary disk proof for the exact two-variable example.}
\GSfield{Present status}{As of 9 November 2026: Established. The coefficient field is decisive for the subsequent conic discussion.}
\GSfield{Assessment}{literal statement: correct with missing hypotheses; publication match: exact; verification: read the result; development: annotated; confidence: high}
\end{GSitem}
\begin{GSitem}{1973-322.7}{F}{An incompletely specified algebraic winding symbol}
\GSfield{Kind; source}{method; \GSpage{P000535}}
\GSfield{The notes}{For a curve \(C(x_1,x_2)=0\) and polynomials \(a,b,c\), the notes mention a winding index and a symbol read uncertainly as \(c(a,b)\) or \(e(a,b)\). For \(a=x_1,b=x_2\), it is called the winding number around the origin.}
\GSfield{Statement}{A classical integer winding number requires an oriented real one-cycle avoiding \((a,b)=(0,0)\), together with a map to \(\mathbb R^2\setminus\{0\}\). A symplectic Mennicke symbol instead starts from a unimodular pair in a specified quotient ring and takes values in a specified \(K_1\)-group. The notebook does not say which cycle, field, quotient, or equivalence defines its proposed symbol, so no single exact invariant can be recovered.}
\GSfield{Reconstruction}{The two possible constructions explain the missing data without identifying them. The presence of the next \(K_1\operatorname{Sp}\) definition suggests an algebraic direction but does not prove equivalence with a topological degree.}
\GSfield{Fate}{\citet[\S4, pp.~948--949]{vasersteinsuslin1976serre} defines symplectic Mennicke symbols precisely and is a thematic lead. Searches for the exact conic/index wording did not determine the notebook's symbol or its claimed vanishing criterion.}
\GSfield{Present status}{As of 9 November 2026: The intended invariant is unidentified. The separately displayed function-existence question is concrete enough to analyze without inventing that invariant.}
\GSfield{Assessment}{literal statement: not determinable; publication match: thematic only; verification: unresolved lead; development: annotated; confidence: low}
\end{GSitem}
\begin{GSitem}{1973-322.8}{E}{Stable symplectic K1}
\GSfield{Kind; source}{definition; \GSpage{P000535}}
\GSfield{The notes}{\(K_1\mathrm{Sp}\,B\) is described as symplectic matrices of order \(2n\) modulo elementary matrices.}
\GSfield{Statement}{For a commutative unital ring \(B\), let \(\operatorname{Sp}(B)=\varinjlim_n\operatorname{Sp}_{2n}(B)\) under hyperbolic stabilization and let \(\operatorname{ESp}(B)\) be its elementary symplectic subgroup. Then \(K_1^{\operatorname{Sp}}(B)=\operatorname{Sp}(B)/\operatorname{ESp}(B)\). Stabilization makes the subgroup normal; an arbitrary fixed-rank quotient needs separate stability and normality hypotheses.}
\GSfield{Reconstruction}{The stable direct limit is inserted. The source's unspecified \(2n\) is not treated as an assertion that every finite rank already computes stable symplectic \(K_1\).}
\GSfield{Fate}{The pre-existing symplectic \(K_1\)-functor is recalled in \citet[\S4, Theorem~4.1 and following definition]{vasersteinsuslin1976serre}; the 1974 announcement also explicitly takes the stable union and quotient.}
\GSfield{Present status}{As of 9 November 2026: Standard stable definition. It alone does not identify the preceding winding symbol.}
\GSfield{Assessment}{literal statement: correct with missing hypotheses; publication match: exact; verification: read the result; development: annotated; confidence: high}
\end{GSitem}
\begin{GSitem}{1973-322.9}{B}{Reducing the conic triple by entire functions}
\GSfield{Kind; source}{question; \GSpage{P000535}}
\GSfield{The notes}{For \(z_1^2+z_2^2=1\), the notes ask whether the index is zero and, equivalently, whether two functions \(f_1,f_2\), entire or of exponential growth, can make \((z_1^2+z_2^2-1)f_1+z_1\) and \((z_1^2+z_2^2-1)f_2+z_2\) have ``\ru{не имеют общ. делит.}''. They add ``\ru{в полиномах нет, в непр. ф-ях есть}''. The coefficient field is not explicitly restated after the earlier real example.}
\GSfield{Statement}{For the complex formulation put \(c=x^2+y^2-1\). There exist even \(f_1,f_2,A,B\in\mathbb C[x,y]\) with \[A(x+cf_1)+B(y+cf_2)=1.\] In fact take \(f_1=1/2\), \(f_2=(i(x+2)-y)/2\), \(a=x+c/2\), \(s=x+1-iy\), \(A=s(a+1)-a-2\), \(B=is^2/2\). This also answers the complex entire and finite-exponential-type versions. With real-valued continuous coefficients the corresponding nowhere-zero pair is impossible by item 1973-322.6.}
\GSfield{Reconstruction}{We replace the ambiguous ``no common divisor'' by the stronger, stable-rank-relevant unit-ideal condition, and explicitly choose complex coefficients. Literal absence of a common polynomial factor is already satisfied by \(f_1=f_2=0\). The polynomial parenthesis is false for unrestricted complex coefficients, but a lost real-coefficient restriction cannot be ruled out from the notes. The exact relation with the unidentified index is not asserted.}
\GSfield{Fate}{The general entire-function reduction is subsumed by \citet[Theorem~1.1]{brudnyi2019stein}: \(\operatorname{sr}\mathcal O(X)=\lfloor\dim_{\mathbb C}X/2\rfloor+1\), so \(\operatorname{sr}\mathcal O(\mathbb C^2)=2\). The preprint circulated in 2018, publication followed in 2019. It imposes no growth bound. The explicit polynomial certificate is verified here as a critical emendation, without a priority claim or a new C/D designation.}
\GSfield{Present status}{As of 9 November 2026: The stated complex existence problem is settled, including finite exponential type. The historical coefficient restriction and the unidentified index remain source ambiguities, not proved open mathematical problems.}
\GSfield{Assessment}{literal statement: not determinable; publication match: partial; verification: read the result; development: annotated; confidence: medium; problem \hyperref[prob:GP-1973-32]{GP-1973-32}}
\end{GSitem}

\GStalk{323}{1973-11-19}{D. A. Kazhdan}{P000535-P000536}{Derived categories and total derived Hom}
Kazhdan develops the derived category from ordinary Ext: projective resolutions, comparison up to homotopy, complexes of maps, and localization at quasi-isomorphisms. He then describes total derived functors and their universal transformation. This is established Grothendieck--Verdier homological algebra, connected to the earlier 1971 derived-category discussions. The essential repairs are boundedness for the free-complex comparison theorem, variance for a contravariant functor, and the direction in which the universal transformation induces a map on natural transformations.

\begin{GSitem}{1973-323.1}{E}{Ext from a free resolution}
\GSfield{Kind; source}{definition; \GSpage{P000535}}
\GSfield{The notes}{For \(A\)-modules \(M,N\), a free resolution \(\cdots\to M_1\to M_0\to M\to0\) gives \(C^i=\operatorname{Hom}(M_i,N)\), and its cohomology is \(\operatorname{Ext}^i(M,N)\). In degree zero the notes recover \(\operatorname{Hom}(M,N)\).}
\GSfield{Statement}{Let \(A\) be a unital ring and \(M,N\) left \(A\)-modules. Choose a projective resolution \(P_\bullet\to M\), concentrated in nonnegative homological degrees. Set \(C^i=\operatorname{Hom}_A(P_i,N)\) with differential induced by \(P_{i+1}\to P_i\). For \(i\ge0\), \(\operatorname{Ext}_A^i(M,N)=H^i(C^\bullet)\), and \(H^0(C^\bullet)=\operatorname{Hom}_A(M,N)\). Free resolutions exist and suffice.}
\GSfield{Reconstruction}{We specify module handedness, exactness of the augmented resolution, and the grading. Hom and Ext are abelian groups unless extra central coefficient structure is specified.}
\GSfield{Fate}{This is classical homological algebra preceding the seminar. The projective-resolution conventions and existence are checked in \citet[Tag~0643, Definition~13.19.1 and Lemma~13.19.3]{stacks2026foundations}; the total derived interpretation is reviewed in \citet[\S10.7]{weibel1994homological}.}
\GSfield{Present status}{As of 9 November 2026: Established definition, independent of the chosen projective resolution.}
\GSfield{Assessment}{literal statement: correct; publication match: exact; verification: read the result; development: annotated; confidence: high}
\end{GSitem}
\begin{GSitem}{1973-323.2}{E}{Comparison of resolutions and the Hom-complex invariant}
\GSfield{Kind; source}{theorem; \GSpage{P000535}}
\GSfield{The notes}{``Any two resolutions are homotopic''; the invariant is said to be the complex defined up to homotopy, not merely its cohomology groups.}
\GSfield{Statement}{Two projective resolutions of the same module admit chain maps lifting its identity, unique up to chain homotopy, and these maps are homotopy inverses. Applying \(\operatorname{Hom}_A(-,N)\) gives homotopy-equivalent cochain complexes. Consequently the Hom complex is determined up to chain homotopy by \(M,N\) and the projective-resolution construction.}
\GSfield{Reconstruction}{The phrase ``resolutions are homotopic'' is interpreted as chain-homotopy equivalence of projective resolutions. An arbitrary acyclic augmentation by nonprojective modules does not satisfy the same comparison theorem.}
\GSfield{Fate}{The comparison result is an older part of the theory. A modern precise form is \citet[Tag~0643, Lemmas~13.19.6--13.19.8]{stacks2026foundations}. Its bounded-above projectivity hypothesis also explains the later qualification in item 1973-323.6.}
\GSfield{Present status}{As of 9 November 2026: Established for projective resolutions; the stronger chain-homotopy invariant is correctly identified.}
\GSfield{Assessment}{literal statement: correct with missing hypotheses; publication match: exact; verification: read the result; development: annotated; confidence: high}
\end{GSitem}
\begin{GSitem}{1973-323.3}{E}{Cohomology groups do not determine chain-homotopy type}
\GSfield{Kind; source}{remark; \GSpage{P000535}}
\GSfield{The notes}{The notes warn that complexes can have isomorphic homology groups without being homotopy equivalent, with an analogy to topology.}
\GSfield{Statement}{For complexes of modules, equality of graded cohomology does not determine chain-homotopy type. For example the acyclic complex \(0\to\mathbb Z\xrightarrow{2}\mathbb Z\to\mathbb Z/2\to0\) is not contractible: a contracting homotopy would split the surjection \(\mathbb Z\to\mathbb Z/2\), which cannot split. It is quasi-isomorphic to zero but not homotopy equivalent to zero.}
\GSfield{Reconstruction}{The example is supplied as an elementary verification of the recorded distinction; it is not a recovered example of the speaker. This does not conflict with the special comparison theorem for bounded-above projective complexes.}
\GSfield{Fate}{The distinction is built into the older localization construction of the derived category \citep[Chapter~I]{hartshorne1966residues}; \citet[Tag~05RR]{stacks2026foundations} specifies which morphisms are inverted. The particular nonsplit example is checked directly here.}
\GSfield{Present status}{As of 9 November 2026: Correct. Merely matching cohomology groups is even weaker than having a specified quasi-isomorphism.}
\GSfield{Assessment}{literal statement: correct; publication match: exact; verification: read the result; development: annotated; confidence: high}
\end{GSitem}
\begin{GSitem}{1973-323.4}{E}{The complex of graded homomorphisms}
\GSfield{Kind; source}{construction; \GSpage{P000535-P000536}}
\GSfield{The notes}{For complexes \(\dot M,\dot N\), the notes write \(\operatorname{Hom}(\dot M,\dot N)\) and describe \(\operatorname{Hom}^i(R,N)\) as maps with shift \(i\), with source and target arrows pointing in opposite directions in the sketch.}
\GSfield{Statement}{For cochain complexes \(P^\bullet,N^\bullet\) of left \(A\)-modules define \[\operatorname{Hom}_A^n(P,N)=\prod_j\operatorname{Hom}_A(P^j,N^{j+n}),\qquad d\phi=d_N\phi-(-1)^n\phi d_P.\] Then \(d^2=0\). Its degree-zero cycles are chain maps, and degree-zero boundaries are null-homotopic maps. For bounded source and target the product is finite in every degree.}
\GSfield{Reconstruction}{The totalization, product convention and differential signs are supplied; the notebook gives only the shifted-map description and a schematic diagram.}
\GSfield{Fate}{The construction is standard and is given exactly in \citet[Tag~0A8H, \S15.73]{stacks2026foundations}. It predates the seminar as part of the homological-algebra framework.}
\GSfield{Present status}{As of 9 November 2026: Established. The differential is essential: the graded collection of shifted maps alone is not yet the Hom complex.}
\GSfield{Assessment}{literal statement: correct with missing hypotheses; publication match: exact; verification: read the result; development: annotated; confidence: high}
\end{GSitem}
\begin{GSitem}{1973-323.5}{E}{A resolution as a quasi-isomorphism to a complex}
\GSfield{Kind; source}{definition; \GSpage{P000535-P000536}}
\GSfield{The notes}{A module is regarded as the complex concentrated in one degree. Replacing it by a free complex inducing the same cohomology is called a resolution; the underlined term is quasi-isomorphism.}
\GSfield{Statement}{A chain map \(f:C^\bullet\to D^\bullet\) is a quasi-isomorphism when every \(H^n(f)\) is an isomorphism. A projective resolution of \(M[0]\) is a quasi-isomorphism \(P^\bullet\to M[0]\) with \(P^j=0\) for \(j>0\) and every \(P^j\) projective. More generally a bounded-above complex of modules has a bounded-above projective replacement.}
\GSfield{Reconstruction}{We distinguish a map inducing cohomology isomorphisms from two complexes with merely abstractly isomorphic cohomology. The boundedness is the version justified by the notebook's ordinary resolutions.}
\GSfield{Fate}{See \citet[Tag~0643, Definition~13.19.1 and Lemma~13.19.3]{stacks2026foundations}. These definitions belong to the pre-1973 derived-category theory documented by \citet[Chapter~I]{hartshorne1966residues}.}
\GSfield{Present status}{As of 9 November 2026: Established; arbitrary unbounded replacements require a further resolution theory.}
\GSfield{Assessment}{literal statement: correct with missing hypotheses; publication match: exact; verification: read the result; development: annotated; confidence: high}
\end{GSitem}
\begin{GSitem}{1973-323.6}{E}{When a quasi-isomorphism of free complexes is a homotopy equivalence}
\GSfield{Kind; source}{theorem; \GSpage{P000536}}
\GSfield{The notes}{The boxed assertion says: two complexes of free modules are quasi-isomorphic if and only if they are homotopy equivalent. No boundedness restriction is repeated.}
\GSfield{Statement}{Every quasi-isomorphism between bounded-above cochain complexes of projective modules is a chain-homotopy equivalence. The unrestricted statement is false. Over \(A=k[\varepsilon]/(\varepsilon^2)\), the doubly infinite complex with \(P^n=A\) and \(d^n=\varepsilon\) in every degree is acyclic, but not contractible: each possible expression \(dh+hd\) has image in \(\varepsilon A\) and cannot equal the identity. Thus \(P\to0\) is an unbounded counterexample even with rank-one free terms.}
\GSfield{Reconstruction}{The bounded-above convention is a reconstruction supported by the surrounding resolution diagrams. The counterexample checks the need for that convention; it is not attributed to Kazhdan.}
\GSfield{Fate}{\citet[Tag~0643, Lemma~13.19.4 and Remark~13.19.5]{stacks2026foundations} gives the correct bounded-above statement. The notebook's omission does not create a new theorem or question about arbitrary unbounded complexes.}
\GSfield{Present status}{As of 9 November 2026: Correct for bounded-above projectives and, more generally, suitable K-projective complexes; false without such a hypothesis.}
\GSfield{Assessment}{literal statement: correct with missing hypotheses; publication match: partial; verification: read the result; development: annotated; confidence: high}
\end{GSitem}
\begin{GSitem}{1973-323.7}{E}{The homotopy category and localization at quasi-isomorphisms}
\GSfield{Kind; source}{construction; \GSpage{P000536}}
\GSfield{The notes}{\(K(A)\) has complexes as objects and homotopy classes of maps as morphisms. \(D(A)\) has the same objects; its maps are described as strings of forward maps and backward quasi-isomorphisms.}
\GSfield{Statement}{Let \(K(A)\) be the chain-homotopy category of complexes of left \(A\)-modules. The derived category is the localization \(D(A)=K(A)[S^{-1}]\) at quasi-isomorphisms. Zigzags represent morphisms only after imposing the localization relations; concatenated raw strings are not distinct morphisms merely because their presentations differ. For suitable calculus of fractions one may use roofs \(M\leftarrow P\to N\), with the left arrow a quasi-isomorphism.}
\GSfield{Reconstruction}{The equivalence relation on zigzags and the universal localization property complete the shorthand. Neither a triangulated structure nor later enhancements are claimed to have been fully developed in these two pages.}
\GSfield{Fate}{This is the Grothendieck--Verdier construction already in \citet[Chapter~I, pp.~19--84]{hartshorne1966residues}, based on the Harvard seminar of 1963--1964. The precise localization definition was reread in \citet[Tag~05RR, Definition~13.11.3]{stacks2026foundations}. Related earlier occurrences are 1971-266.3 and 1971-267.1.}
\GSfield{Present status}{As of 9 November 2026: Established. The quotient relations are part of the definition, and the earlier discussions are expository links rather than recurring open problems.}
\GSfield{Assessment}{literal statement: correct with missing hypotheses; publication match: exact; verification: read the result; development: annotated; confidence: high}
\end{GSitem}
\begin{GSitem}{1973-323.8}{E}{Derived Hom and graded morphisms in the derived category}
\GSfield{Kind; source}{construction; \GSpage{P000536}}
\GSfield{The notes}{The stated goal is to make Hom into a complex; the source displays \(\operatorname{Hom}^i_{D(A)}(\dot M,\dot N)\) after constructing Hom from a free replacement.}
\GSfield{Statement}{For \(M\in D^-(A)\), choose a bounded-above projective replacement \(P\to M\). For a complex \(N\), the object \(R\operatorname{Hom}_A(M,N)\) of \(D(\mathbb Z)\) is represented by \(\operatorname{Hom}_A^\bullet(P,N)\). Its cohomology is \[H^iR\operatorname{Hom}_A(M,N)\simeq\operatorname{Hom}_{D(A)}(M,N[i]).\] For modules in degree zero and \(i\ge0\), this is \(\operatorname{Ext}_A^i(M,N)\). For a commutative coefficient ring the usual additional module structure is available.}
\GSfield{Reconstruction}{A derived Hom object is distinguished from the single morphism group in the ordinary category \(D(A)\). The notebook's superscript is interpreted as the graded shifted morphism group, not an intrinsic differential on every ordinary Hom set.}
\GSfield{Fate}{The formulas are checked in \citet[Tags~0A8H and 0A5W, equation~15.75.0.2]{stacks2026foundations}, with projective computation justified by Tag~0643. They are standard derived-category constructions, not a new enrichment theorem of this seminar.}
\GSfield{Present status}{As of 9 November 2026: Established with the specified replacement and grading conventions.}
\GSfield{Assessment}{literal statement: correct with missing hypotheses; publication match: exact; verification: read the result; development: annotated; confidence: high}
\end{GSitem}
\begin{GSitem}{1973-323.9}{E}{Deriving a contravariant functor by projective replacements}
\GSfield{Kind; source}{method; \GSpage{P000536}}
\GSfield{The notes}{The functor \(F\) is called contravariant, but the diagrams show \(F:M(A)\to M(B)\) and \(\dot F:D(A)\to D(B)\). One chooses free representatives and applies \(F\); cohomology gives ordinary derived functors.}
\GSfield{Statement}{For an additive contravariant left-exact functor \(F:(A\text{-}\mathrm{Mod})^{\mathrm{op}}\to B\text{-}\mathrm{Mod}\), applying \(F\) to a bounded-above projective replacement, with degrees reversed, constructs \(RF:D^-(A)^{\mathrm{op}}\to D^+(B)\). Its cohomology on a module recovers the right derived functors of \(F\). Additivity preserves chain homotopies; bounded-above projective comparison proves independence. A general covariant functor has a different left/right resolution convention.}
\GSfield{Reconstruction}{We restore the opposite category, degree reversal, additivity, left exactness and bounded ranges. Without those data the notebook's covariance diagram and the asserted recovery in degree zero are not a precise theorem.}
\GSfield{Fate}{The bounded derived-functor construction is part of the theory already presented in \citet[Chapter~I]{hartshorne1966residues}. A checked formulation and the role of exactness are in \citet[Theorem~10.5.6 and Remark~10.5.8]{weibel1994homological}, with passage to opposite categories for the contravariant case.}
\GSfield{Present status}{As of 9 November 2026: Established after repairing variance and specifying the resolution range; arbitrary unbounded free replacement is not sufficient.}
\GSfield{Assessment}{literal statement: false as written; publication match: partial; verification: read the result; development: annotated; confidence: high}
\end{GSitem}
\begin{GSitem}{1973-323.10}{E}{The universal transformation for a total right derived functor}
\GSfield{Kind; source}{definition; \GSpage{P000536}}
\GSfield{The notes}{The last lines require \(QF\to\dot FQ\) and display an isomorphism \[\operatorname{Hom}(QF,GQ)\to\operatorname{Hom}(\dot FQ,GQ)\] for every \(G\), with the arrow pointing right.}
\GSfield{Statement}{Let \(Q_A:\mathcal K_A\to\mathcal D_A\) and \(Q_B:\mathcal K_B\to\mathcal D_B\) be the chosen localizations and let \(F:\mathcal K_A\to\mathcal K_B\) be an additive functor compatible with the relevant triangulated structures. A total right derived functor is \(RF:\mathcal D_A\to\mathcal D_B\) with \(\eta:Q_BF\to RFQ_A\), universal in the sense that composition with \(\eta\) induces the bijection \[\operatorname{Nat}(RF,G)\longrightarrow\operatorname{Nat}(Q_BF,GQ_A),\quad\theta\longmapsto(\theta Q_A)\eta.\] Existence is an additional requirement. For contravariant \(F\), first use the appropriate opposite category.}
\GSfield{Reconstruction}{This states the covariant categorical version consistently. The displayed notebook direction is the inverse factorization map if its universal property has already been established; direct precomposition with \(\eta\) goes the other way. Thus the arrow need not be declared a false theorem, but its intended meaning must be specified.}
\GSfield{Fate}{\citet[Definition~10.5.1, p.~391]{weibel1994homological} states exactly this universal property. The source theory predates 1973; the existence hypotheses are developed in \citet[Chapter~I, \S5]{hartshorne1966residues}. The Stacks formulation also treats existence separately \citep[Tag~05S7]{stacks2026foundations}.}
\GSfield{Present status}{As of 9 November 2026: Established universal property. The notebook has a direction ambiguity, not evidence for an additional universal construction beyond standard derived functors.}
\GSfield{Assessment}{literal statement: correct with missing hypotheses; publication match: partial; verification: read the result; development: annotated; confidence: high}
\end{GSitem}

\GStalk{324}{}{}{P000536-P000538}{Calendar heading: 1974}
The catalogue supplies a calendar-year divider, not a dated talk. Its overlapping page range contains the end of Kazhdan's previous entry, Springer's short entry, and the opening of Novikov's talk. No additional mathematical assertion is assigned to this heading. The academic-year prefix remains 1973.


\GStalk{325}{1974-02-11}{T. A. Springer}{P000537}{Green polynomials}
Only the title and the distinction between a finite general linear group and its algebraic group over the algebraic closure survive. Neither a Green-polynomial formula nor a theorem of Springer can be recovered from this page.

\begin{GSitem}{1973-325.1}{E}{Finite and algebraic general linear groups}
\GSfield{Kind; source}{definition; \GSpage{P000537}}
\GSfield{The notes}{The page writes \(G=GL_n(k)\), \(|k|=q\), and \(\underline G=GL(\bar k)\).}
\GSfield{Statement}{For a finite field \(k=\mathbb F_q\) and an integer \(n\ge1\), let \(\mathbf G=GL_{n,\bar k}\). Then \(G=\mathbf G(\bar k)^F=GL_n(k)\), where \(F\) raises entries to the \(q\)-th power.}
\GSfield{Reconstruction}{The subscript \(n\) on the algebraic group and the Frobenius description are explanatory additions.}
\GSfield{Fate}{This is the standard setting of \citet[opening definitions]{green1955characters} (1955), long before the talk. The finite/algebraic distinction is identifiable without reconstructing any missing character formula. The AMS PDF was located but direct opening returned 403, so its entry remains unverified.}
\GSfield{Present status}{As of 9 November 2026: Established in the stated scope; no unresolved claim is inferred from this occurrence.}
\GSfield{Assessment}{literal statement: correct with missing hypotheses; publication match: exact; verification: indirect; development: annotated; confidence: high}
\end{GSitem}
\begin{GSitem}{1973-325.2}{F}{The announced Green-polynomial subject}
\GSfield{Kind; source}{remark; \GSpage{P000537}}
\GSfield{The notes}{The handwritten heading is in English: Green polynomials. No polynomial, formula, or subsequent argument is recorded.}
\GSfield{Statement}{The surviving text identifies the subject of Green polynomials associated with finite general linear groups, but does not specify any particular identity, realization, or theorem.}
\GSfield{Fate}{\citet{green1955characters} supplies a historical subject match only. Assigning Springer resolutions or a cohomological character formula to this page would exceed the evidence. Earlier occurrences 1971-231.2 and 1971-231.5 concern Springer geometry, not an exact repeated assertion here.}
\GSfield{Present status}{As of 9 November 2026: The mathematical content beyond the heading is not recoverable from the supplied page.}
\GSfield{Assessment}{literal statement: not determinable; publication match: thematic only; verification: unresolved lead; development: annotated; confidence: high}
\end{GSitem}

\GStalk{326}{1974-02-25}{S. P. Novikov}{P000538-P000541}{Periodic KdV, Floquet theory, and finite-gap potentials}
Novikov's continued talk interleaves Gelfand's explanation of Floquet decomposition with stationary KdV and its spectral curve. Distinct questions concern monotonicity inside bands, missing inverse-spectral data, and density of finite-gap potentials. Novikov's 1974 paper was received on 18 February, one week before this recorded session; its later journal appearance is distinguished from this prior manuscript date.

\begin{GSitem}{1973-326.1}{E}{The first two KdV Lax operators}
\GSfield{Kind; source}{construction; \GSpage{P000538}}
\GSfield{The notes}{\(L=-D^2+u\), \(A_0=D\), \(A_1=4D^3-3(uD+Du)\).}
\GSfield{Statement}{For smooth real \(u\), as formal differential operators, \([A_0,L]=u_x\) and \([A_1,L]=u_{xxx}-6uu_x\). Thus \(L_t=[A_1,L]\) gives \(u_t=u_{xxx}-6uu_x\); reversing the commutator reverses time.}
\GSfield{Reconstruction}{Operator composition, commutator order, and the resulting sign are made explicit; the sign of the later \(P_1\) on P000541 uses the opposite convention.}
\GSfield{Fate}{\citet[opening equations and \S2]{novikov1974periodic} uses the Lax hierarchy and cites its established antecedents. This is the same general line as 1970-219.1 and 1970-219.6, not a new problem.}
\GSfield{Present status}{As of 9 November 2026: Established in the stated scope; no unresolved claim is inferred from this occurrence.}
\GSfield{Assessment}{literal statement: correct; publication match: exact; verification: read the result; development: annotated; confidence: high}
\end{GSitem}
\begin{GSitem}{1973-326.2}{E}{The real periodic Schrödinger operator is self-adjoint}
\GSfield{Kind; source}{claim; \GSpage{P000538}; I. M. Gelfand}
\GSfield{The notes}{Gelfand writes \(y^{\prime\prime}+uy=\lambda y\), real periodic \(u\), and says the closure from \(C_c^\infty\) is self-adjoint.}
\GSfield{Statement}{For real smooth \(T\)-periodic \(u\), the operator \(-D^2+u\) on \(C_c^\infty(\mathbb R)\) is essentially self-adjoint in \(L^2(\mathbb R)\); its closure has domain \(H^2(\mathbb R)\).}
\GSfield{Reconstruction}{The uniform convention \(-D^2+u\), \(T>0\), and a sufficient smoothness hypothesis are supplied. Gelfand's displayed \(+D^2\) has the opposite principal sign, so it cannot share all later spectral sign conventions unchanged.}
\GSfield{Fate}{\citet[\S1, pp.~54--58 of the Russian original]{novikov1974periodic} recalls this classical Floquet theory; the paper was received on 18 February 1974 and published later in 1974.}
\GSfield{Present status}{As of 9 November 2026: Established in the stated scope; no unresolved claim is inferred from this occurrence.}
\GSfield{Assessment}{literal statement: correct with missing hypotheses; publication match: exact; verification: read the result; development: annotated; confidence: high}
\end{GSitem}
\begin{GSitem}{1973-326.3}{E}{Periodic and antiperiodic interlacing}
\GSfield{Kind; source}{theorem; \GSpage{P000538}; I. M. Gelfand}
\GSfield{The notes}{Both boundary conditions carry the same \(\pm\), and the labels occur in pairs: periodic, antiperiodic, antiperiodic, periodic, periodic, and so on.}
\GSfield{Statement}{For real smooth \(T\)-periodic \(u\), impose either \(y(T)=y(0),y_x(T)=y_x(0)\), or both minus signs. The combined eigenvalues satisfy \(E_0<E_1\le E_2<E_3\le E_4<\cdots\); \(E_0\) is periodic and the subsequent boundary types alternate in equal pairs. Equalities mean closed gaps.}
\GSfield{Fate}{\citet[\S1, pp.~54--58 of the Russian original]{novikov1974periodic} recalls this classical Floquet theory; the paper was received on 18 February 1974 and published later in 1974.}
\GSfield{Present status}{As of 9 November 2026: Established in the stated scope; no unresolved claim is inferred from this occurrence.}
\GSfield{Assessment}{literal statement: correct with missing hypotheses; publication match: exact; verification: read the result; development: annotated; confidence: high}
\end{GSitem}
\begin{GSitem}{1973-326.4}{E}{Self-adjoint Bloch fibers}
\GSfield{Kind; source}{construction; \GSpage{P000539}; I. M. Gelfand}
\GSfield{The notes}{For every phase \(t\), \(y(x+T)=e^{it}y(x)\) and \(y_x(x+T)=e^{it}y_x(x)\) define a self-adjoint discrete-spectrum problem.}
\GSfield{Statement}{For \(\theta\in\mathbb R/2\pi\mathbb Z\), let \(L_\theta=-D^2+u\) on \(L^2(0,T)\), with domain \(\{f\in H^2(0,T):f(T)=e^{i\theta}f(0),f_x(T)=e^{i\theta}f_x(0)\}\). It is self-adjoint with compact resolvent.}
\GSfield{Reconstruction}{The phase is renamed \(\theta\) to avoid confusion with evolution time.}
\GSfield{Fate}{\citet[\S1, pp.~54--58 of the Russian original]{novikov1974periodic} recalls this classical Floquet theory; the paper was received on 18 February 1974 and published later in 1974.}
\GSfield{Present status}{As of 9 November 2026: Established in the stated scope; no unresolved claim is inferred from this occurrence.}
\GSfield{Assessment}{literal statement: correct with missing hypotheses; publication match: exact; verification: read the result; development: annotated; confidence: high}
\end{GSitem}
\begin{GSitem}{1973-326.5}{E}{Gelfand transform and its normalization}
\GSfield{Kind; source}{theorem; \GSpage{P000539}; I. M. Gelfand}
\GSfield{The notes}{The lemma writes an integral reconstruction, a quasiperiodicity condition, and an \(L^2\) identity, but mixes period 1 and \(2\pi\), integrates over \(dx\) instead of phase, and drops one subscript.}
\GSfield{Statement}{For \(f\in L^2(\mathbb R)\), define \(Uf(\theta,x)=\sum_{m\in\mathbb Z}e^{-im\theta}f(x+mT)\) first for compactly supported smooth \(f\), then by completion. It satisfies \(Uf(\theta,x+T)=e^{i\theta}Uf(\theta,x)\), \(f(x)=(2\pi)^{-1}\int_0^{2\pi}Uf(\theta,x)\,d\theta\), and \(\|f\|_2^2=(2\pi)^{-1}\int_0^{2\pi}\int_0^T|Uf|^2\,dx\,d\theta\). The identities hold in the corresponding \(L^2\) spaces.}
\GSfield{Reconstruction}{The transform itself and a consistent normalization are supplied; the manuscript's three defective displayed identities are retained in the corrections table.}
\GSfield{Fate}{The notes explicitly credit Gelfand's circa-1950 Doklady paper. The concrete normalization here follows Fourier-series Parseval; it is an editorial reconstruction of a classical unitary direct-integral transform. \citet[\S1]{novikov1974periodic} uses the resulting band decomposition.}
\GSfield{Present status}{As of 9 November 2026: Established in the stated scope; no unresolved claim is inferred from this occurrence.}
\GSfield{Assessment}{literal statement: false as written; publication match: exact; verification: read the result; development: annotated; confidence: high}
\end{GSitem}
\begin{GSitem}{1973-326.6}{E}{Bloch eigenfunction expansion and band derivatives}
\GSfield{Kind; source}{construction; \GSpage{P000539}; I. M. Gelfand}
\GSfield{The notes}{Each \(f_t\) is expanded in \(y_{t,n}\), with coefficients written without a complex conjugate; the notes then mention computing \(d\lambda/dt\).}
\GSfield{Statement}{Choose an orthonormal eigenbasis \(\phi_{n,\theta}\) of each fiber. Then \(Uf(\theta,\cdot)=\sum_n\langle Uf,\phi_{n,\theta}\rangle\phi_{n,\theta}\) in \(L^2(0,T)\), with the complex conjugate in the inner product. If \(\lambda_n(\theta)\) is simple, gauge to periodic \(v=e^{-i\theta x/T}\phi\); then \(\lambda^{\prime}_n(\theta)=\langle v,[-2iT^{-1}(D+i\theta/T)]v\rangle\) for normalized \(v\).}
\GSfield{Reconstruction}{The derivative formula is a labeled reconstruction by differentiating the gauged operator; no such formula survives in the notes.}
\GSfield{Fate}{\citet[\S1, pp.~54--58 of the Russian original]{novikov1974periodic} recalls this classical Floquet theory; the paper was received on 18 February 1974 and published later in 1974.}
\GSfield{Present status}{As of 9 November 2026: Established in the stated scope; no unresolved claim is inferred from this occurrence.}
\GSfield{Assessment}{literal statement: false as written; publication match: exact; verification: read the result; development: annotated; confidence: high}
\end{GSitem}
\begin{GSitem}{1973-326.7}{E}{Monodromy in a conjugate solution basis}
\GSfield{Kind; source}{construction; \GSpage{P000539}}
\GSfield{The notes}{With \(\lambda=k^2\), initial data \((y,y_x)=(1,ik)\), the pair \(y,\bar y\) gives \(T=\left(\begin{smallmatrix}a&b\\\bar b&\bar a\end{smallmatrix}\right)\), \(|a|^2-|b|^2=1\).}
\GSfield{Statement}{For real \(k\ne0\) and real \(u\), those two solutions are independent because their Wronskian is nonzero. Transfer by one period has determinant one and the conjugation symmetry stated in the notes. Its half trace \(d(\lambda)=\tfrac12\operatorname{tr}T=\operatorname{Re}a\) is independent of the basepoint and extends as an entire function of \(\lambda\).}
\GSfield{Reconstruction}{The basis restriction \(k\in\mathbb R\setminus\{0\}\) is essential; for other spectral parameters one uses a real initial-data fundamental matrix, not the claimed conjugate basis.}
\GSfield{Fate}{\citet[\S1, pp.~54--58 of the Russian original]{novikov1974periodic} recalls this classical Floquet theory; the paper was received on 18 February 1974 and published later in 1974.}
\GSfield{Present status}{As of 9 November 2026: Established in the stated scope; no unresolved claim is inferred from this occurrence.}
\GSfield{Assessment}{literal statement: correct with missing hypotheses; publication match: exact; verification: read the result; development: annotated; confidence: high}
\end{GSitem}
\begin{GSitem}{1973-326.8}{E}{Floquet multipliers and spectral bands}
\GSfield{Kind; source}{theorem; \GSpage{P000539-P000540}}
\GSfield{The notes}{The multipliers are \(a_R\pm\sqrt{a_R^2-1}\); bounded solutions are equated to \(|a_R|\le1\), and outside the strip the notes write \(e^{\pm t}\).}
\GSfield{Statement}{For real \(\lambda\), the multipliers solve \(\rho^2-2d(\lambda)\rho+1=0\), and \(\sigma(L)=\{\lambda:|d(\lambda)|\le1\}\). Inside an open band all solutions are bounded. At a band edge at least one nonzero bounded periodic or antiperiodic solution exists, but the second solution can grow linearly. In a gap the real multipliers are \(\pm e^{\pm\kappa}\), with a common minus sign when \(d<-1\).}
\GSfield{Reconstruction}{The edge quantifier and the omitted negative-multiplier case are restored.}
\GSfield{Fate}{\citet[\S1, pp.~54--58 of the Russian original]{novikov1974periodic} recalls this classical Floquet theory; the paper was received on 18 February 1974 and published later in 1974.}
\GSfield{Present status}{As of 9 November 2026: Established in the stated scope; no unresolved claim is inferred from this occurrence.}
\GSfield{Assessment}{literal statement: correct with missing hypotheses; publication match: exact; verification: read the result; development: annotated; confidence: high}
\end{GSitem}
\begin{GSitem}{1973-326.9}{E}{Monotonicity of the discriminant in each band}
\GSfield{Kind; source}{question; \GSpage{P000540}; I. M. Gelfand}
\GSfield{The notes}{Novikov says monotonicity was proved for small potentials; Gelfand announces a \ru{конкурс на доказательство} for the next meeting.}
\GSfield{Statement}{For any real smooth periodic potential, is \(d_{\lambda}(\lambda)\ne0\) whenever \(|d(\lambda)|<1\)? Equivalently, does each noncollapsed band interior map monotonically onto \((-1,1)\)? The answer is yes; critical values of the real Hill discriminant lie outside the open strip.}
\GSfield{Fate}{\citet[\S1, pp.~54--58 of the Russian original]{novikov1974periodic} recalls this classical Floquet theory; the paper was received on 18 February 1974 and published later in 1974. The challenge asks for a proof of a classical scalar Sturm--Liouville fact, not a still-open small-potential extension.}
\GSfield{Present status}{As of 9 November 2026: Affirmatively settled in scalar periodic Sturm--Liouville theory, with no smallness assumption.}
\GSfield{Assessment}{literal statement: correct with missing hypotheses; publication match: exact; verification: read the result; development: annotated; confidence: high; problem \hyperref[prob:GP-1973-50]{GP-1973-50}}
\end{GSitem}
\begin{GSitem}{1973-326.10}{E}{Constant potentials and the free monodromy}
\GSfield{Kind; source}{claim; \GSpage{P000540}; I. M. Gelfand}
\GSfield{The notes}{For constant \(u\), Gelfand says the spectrum is positive. The later free matrix has diagonal entries \(e^{\pm iT}\) and \(a_R=\cos T\), omitting spectral dependence.}
\GSfield{Statement}{For \(u\equiv c\in\mathbb R\), \(\sigma(-D^2+c)=[c,\infty)\) and \(d(\lambda)=\cos(T\sqrt{\lambda-c})\), analytically continued below \(c\). For \(c=0\), the plane-wave basis at \(\lambda=k^2>0\) gives \(\operatorname{diag}(e^{ikT},e^{-ikT})\). All finite gaps are closed.}
\GSfield{Reconstruction}{The general constant \(c\), the threshold, and \(k\)-dependence are restored.}
\GSfield{Fate}{\citet[\S1, pp.~54--58 of the Russian original]{novikov1974periodic} recalls this classical Floquet theory; the paper was received on 18 February 1974 and published later in 1974. The corrections also follow directly from Fourier transformation.}
\GSfield{Present status}{As of 9 November 2026: Established in the stated scope; no unresolved claim is inferred from this occurrence.}
\GSfield{Assessment}{literal statement: false as written; publication match: exact; verification: read the result; development: annotated; confidence: high}
\end{GSitem}
\begin{GSitem}{1973-326.11}{F}{Mathieu stability sketch}
\GSfield{Kind; source}{remark; \GSpage{P000540}; I. M. Gelfand}
\GSfield{The notes}{The equation is \(-y^{\prime\prime}+(a+b\cos2x)y=0\); the first stability band is associated with \(a+b\cos2x<0\), followed by a contrast with constant coefficients.}
\GSfield{Statement}{For each real \((a,b)\), stability of the Mathieu equation is governed by the trace of its period-\(\pi\) monodromy: at least one bounded solution exists precisely when its half trace has absolute value at most one. The inequality in the notes does not specify whether it is pointwise, averaged, or a statement along a parameter curve, and does not identify an exact first-band boundary.}
\GSfield{Reconstruction}{Only the well-defined monodromy criterion is supplied as context; the sketch is not promoted to a new theorem.}
\GSfield{Fate}{Floquet stability is classical; \citet[\S1]{novikov1974periodic} supports the trace criterion. No exact publication match is claimed for the ambiguous inequality or sketch.}
\GSfield{Present status}{As of 9 November 2026: The general criterion is established, but the intended parameter-region assertion remains unidentified.}
\GSfield{Assessment}{literal statement: not determinable; publication match: thematic only; verification: unresolved lead; development: annotated; confidence: medium}
\end{GSitem}
\begin{GSitem}{1973-326.12}{E}{Closed gaps and double boundary eigenvalues}
\GSfield{Kind; source}{theorem; \GSpage{P000540}}
\GSfield{The notes}{Periodic and antiperiodic points have \(|a_R|=1\); a double degeneracy has \(b(x_0,\lambda_n)\equiv0\).}
\GSfield{Statement}{For real periodic scalar \(L\), a closed finite gap is a double eigenvalue of the relevant periodic or antiperiodic problem. Its monodromy is \(M=\pm I\); in any admissible conjugate basis \(b=0\). At an open-gap edge the boundary eigenvalue is simple and \(M\) is a nontrivial Jordan matrix with eigenvalue \(\pm1\).}
\GSfield{Reconstruction}{The claim \(b\ne0\) is restricted to nonzero real \(k\) where the notebook basis exists; the basis-independent formulation uses \(M\).}
\GSfield{Fate}{\citet[\S1, pp.~54--58 of the Russian original]{novikov1974periodic} recalls this classical Floquet theory; the paper was received on 18 February 1974 and published later in 1974.}
\GSfield{Present status}{As of 9 November 2026: Established in the stated scope; no unresolved claim is inferred from this occurrence.}
\GSfield{Assessment}{literal statement: correct with missing hypotheses; publication match: exact; verification: read the result; development: annotated; confidence: high}
\end{GSitem}
\begin{GSitem}{1973-326.13}{B}{Which supplementary data determine a periodic potential?}
\GSfield{Kind; source}{question; \GSpage{P000540}}
\GSfield{The notes}{The half trace does not determine \(u\). The notes ask which scattering data beyond the spectrum do, mentioning Dubrovin and Faddeev's students.}
\GSfield{Statement}{At a fixed basepoint, supplement the band spectrum of a real smooth finite-gap Schrödinger operator with its Dirichlet divisor: one point on each real gap oval of the hyperelliptic spectral curve, including its sheet, with the usual endpoint identification. These data determine the potential. Band edges alone leave an isospectral family and cannot even distinguish translations.}
\GSfield{Reconstruction}{The finite-gap restriction and sheet-labelled divisor are added explicitly; no unrestricted theorem is inferred solely from the finite-gap reference.}
\GSfield{Fate}{\citet[\S\S2--3]{novikov1974periodic} develops the stationary spectral curve. \citet[Chapter~2, \S3, pp.~84--87 of the Russian original, especially Eqs.~(2.3.11)--(2.3.12)]{dubrovinmatveevnovikov1976nonlinear} gives one sheet-labelled pole on each gap oval, its evolution, and the trace formula reconstructing the potential. These pages were read. The latter paper was received 2 June 1975 and published in 1976. This is an exact answer in the finite-gap scope, a partial match to the notes' unrestricted periodic wording.}
\GSfield{Present status}{As of 9 November 2026: The finite-gap inverse problem is established. General periodic inverse spectral theory also requires supplementary data; the finite-gap reference here is not represented as proving the general theorem.}
\GSfield{Assessment}{literal statement: correct with missing hypotheses; publication match: partial; verification: read the result; development: annotated; confidence: high; problem \hyperref[prob:GP-1973-51]{GP-1973-51}}
\end{GSitem}
\begin{GSitem}{1973-326.14}{E}{Finite-gap periodic potentials}
\GSfield{Kind; source}{definition; \GSpage{P000540}}
\GSfield{The notes}{Eventually all boundary eigenvalues become double and the forbidden intervals shrink to points.}
\GSfield{Statement}{A real smooth periodic potential is finite-gap if only finitely many bounded complementary intervals of \(\sigma(-D^2+u)\) are nonempty. Equivalently, all sufficiently high periodic/antiperiodic pairs coincide. The half-line below the bottom of the spectrum is not counted as a finite gap.}
\GSfield{Fate}{\citet[\S2]{novikov1974periodic} uses this precise spectral class.}
\GSfield{Present status}{As of 9 November 2026: Established in the stated scope; no unresolved claim is inferred from this occurrence.}
\GSfield{Assessment}{literal statement: correct; publication match: exact; verification: read the result; development: annotated; confidence: high}
\end{GSitem}
\begin{GSitem}{1973-326.15}{A}{Stationary KdV and finite-gap spectra}
\GSfield{Kind; source}{theorem; \GSpage{P000540-P000541}}
\GSfield{The notes}{If \([A,L]=P(u,\ldots)\) is multiplication and \(P=0\) has a periodic solution, the potential is finite-gap; the converse is also asserted. The next page specifies one operator at every odd order and their linear combinations.}
\GSfield{Statement}{For real smooth periodic \(u\), a nontrivial odd-order differential operator \(A\) commuting with \(L=-D^2+u\), with nonzero constant leading coefficient, forces finitely many open gaps. Conversely a smooth real periodic finite-gap potential satisfies a stationary equation of the KdV hierarchy. A trivial commuting operator \(A=p(L)\) is not a sufficient hypothesis.}
\GSfield{Reconstruction}{The indispensable odd-order/nontrivial condition is obtained from the following page, not from the isolated theorem sentence. The converse is separated from the forward theorem.}
\GSfield{Fate}{\citet[Theorem~2.1 and its stationary construction, pp.~59--62]{novikov1974periodic} is the near-contemporary primary publication. The finite-gap inverse and converse framework is developed in \citet{dubrovinmatveevnovikov1976nonlinear}. The converse is explicitly proved in Chapter~2, \S3, statement3, p.~86 of the Russian original of \citet{dubrovinmatveevnovikov1976nonlinear}, which was also read.}
\GSfield{Present status}{As of 9 November 2026: Established in the stated scope; no unresolved claim is inferred from this occurrence.}
\GSfield{Assessment}{literal statement: correct with missing hypotheses; publication match: partial; verification: read the result; development: annotated; confidence: high}
\end{GSitem}
\begin{GSitem}{1973-326.16}{A}{Finite-dimensional Hamiltonian stationary flows}
\GSfield{Kind; source}{theorem; \GSpage{P000541}}
\GSfield{The notes}{The stationary hierarchy equations are Hamiltonian systems with \(n\) degrees of freedom and a maximal set of commuting integrals.}
\GSfield{Statement}{Fix a nondegenerate stationary KdV equation of order \(2n+1\) and integrate its translation-invariant first integral once. On the resulting regular \(2n\)-dimensional phase space, the induced finite-dimensional flow has a Hamiltonian structure and \(n\) independent commuting integrals on a regular open set. Singular levels and loss of independence are not covered by the Liouville-torus conclusion.}
\GSfield{Reconstruction}{The regular-level and independence qualifications distinguish Liouville integrability from a claim about every singular point.}
\GSfield{Fate}{\citet[Theorem~2.1, pp.~59--61]{novikov1974periodic} supplies the stationary Hamiltonian systems and commuting integrals; manuscript received 18 February 1974.}
\GSfield{Present status}{As of 9 November 2026: Established in the stated scope; no unresolved claim is inferred from this occurrence.}
\GSfield{Assessment}{literal statement: correct with missing hypotheses; publication match: exact; verification: read the result; development: annotated; confidence: high}
\end{GSitem}
\begin{GSitem}{1973-326.17}{B}{Density of finite-gap potentials}
\GSfield{Kind; source}{question; \GSpage{P000541}}
\GSfield{The notes}{Novikov asks \ru{любой ли потенциал аппроксим. конечнозонными?}; neither function space nor topology is stated.}
\GSfield{Statement}{For every real \(q\in H^m(\mathbb R/T\mathbb Z)\), \(m\ge0\), and every \(\varepsilon>0\), does a real periodic finite-gap \(q_*\) exist with \(\|q-q_*\|_{H^m}<\varepsilon\)? Yes. In particular smooth periodic potentials can be approximated in any prescribed finite collection of Sobolev seminorms.}
\GSfield{Reconstruction}{The periodic setting comes from the talk; the real Sobolev topology and quantifiers are an explicit modern formulation, not written in the notebook.}
\GSfield{Fate}{\citet[Theorem~1.2 and its proof in \S5]{grebertkappelerposchel2004gaps} proves finite-gap density in \(H^m\), \(m\ge0\), attributes the question to Novikov, and credits the first proof to Marchenko--Ostrovskii's 1975 characterization of Hill spectra. The author-hosted full PDF was downloaded and these passages read. The manuscript was received 15 September 2003, revised in 2004, and published in 2004. First-proof credit is distinguished from the later proof actually inspected.}
\GSfield{Present status}{As of 9 November 2026: Affirmatively settled in every real Sobolev space \(H^m\), \(m\ge0\); it is not a C/D candidate.}
\GSfield{Assessment}{literal statement: correct with missing hypotheses; publication match: exact; verification: read the result; development: annotated; confidence: high; problem \hyperref[prob:GP-1973-52]{GP-1973-52}}
\end{GSitem}
\begin{GSitem}{1973-326.18}{A}{A stationary matrix equation and its spectral polynomial}
\GSfield{Kind; source}{construction; \GSpage{P000541}}
\GSfield{The notes}{A traceless \(2\times2\) matrix \(\Lambda\) satisfies \(\Lambda^{\prime}=[\Lambda,Q]\). Its eigenvalues are written \(\pm\sqrt{\det\Lambda}\), followed by the words the word polynomial.}
\GSfield{Statement}{For any differentiable traceless matrix solution \(\Lambda_x=[\Lambda,Q]\), its characteristic polynomial is \(\mu^2+\det\Lambda\); hence \(\mu=\pm\sqrt{-\det\Lambda}\) and \(\det\Lambda\) is independent of \(x\). In the stationary KdV construction a suitable gauge and spectral parameter yield polynomial coefficients of the spectral curve. The square roots themselves are generally algebraic functions, not polynomials.}
\GSfield{Reconstruction}{The notes mix \(\delta\), \(\delta^2\), and a \(\delta^{-1}\) gauge term; no exact matrix-degree assertion is recovered without fixing that gauge. The polynomial claim is safely assigned to the determinant/spectral polynomial.}
\GSfield{Fate}{\citet[\S2, pp.~61--62]{novikov1974periodic} constructs the matrix and spectral polynomial. The sign follows independently from the characteristic polynomial of a traceless \(2\times2\) matrix.}
\GSfield{Present status}{As of 9 November 2026: Established in the stated scope; no unresolved claim is inferred from this occurrence.}
\GSfield{Assessment}{literal statement: false as written; publication match: exact; verification: read the result; development: annotated; confidence: high}
\end{GSitem}
\begin{GSitem}{1973-326.19}{E}{Riccati expansion and conservation-law densities}
\GSfield{Kind; source}{method; \GSpage{P000541}}
\GSfield{The notes}{\(\psi_k=e^{ikx+\int_{-\infty}^x\chi}\), \(\chi=\sum_{n\ge1}\chi_nk^{-n}\), and \(I_n=\int_{\mathbb R}\chi_n\) are conserved.}
\GSfield{Statement}{For a nonvanishing local solution of \(-\psi_{xx}+u\psi=k^2\psi\), put \(\psi^{\prime}/\psi=ik+\chi\). Then \(\chi_x+2ik\chi+\chi^2=u\), which recursively determines a formal Laurent series with differential-polynomial coefficients. Under smooth rapid decay its convergent spatial density integrals are KdV conservation laws; for periodic \(u\) use one-period integrals of the densities. The formal expansion need not converge in \(k^{-1}\).}
\GSfield{Reconstruction}{The notes temporarily use the decaying-line normalization inside a periodic discussion. We separate the two regimes and do not integrate a nondecaying periodic density over the whole line.}
\GSfield{Fate}{\citet[trace identities and conservation laws]{zakharovfaddeev1971kdv} already develops this connection in 1971; it also matches the earlier occurrence 1970-219.6. \citet[\S2]{novikov1974periodic} adapts the hierarchy to periodic potentials.}
\GSfield{Present status}{As of 9 November 2026: Established in the stated scope; no unresolved claim is inferred from this occurrence.}
\GSfield{Assessment}{literal statement: correct with missing hypotheses; publication match: exact; verification: publisher metadata checked; development: annotated; confidence: high}
\end{GSitem}
\begin{GSitem}{1973-326.20}{A}{Evolution of a normalized eigenfunction basis}
\GSfield{Kind; source}{construction; \GSpage{P000541}}
\GSfield{The notes}{\(\dot\psi=A\psi+\lambda\psi+\mu\bar\psi\), with \(\Lambda\) formed from the last two coefficients.}
\GSfield{Statement}{If \(L_t=[A,L]\) and \((L-k^2)\psi=0\), then \((L-k^2)(\psi_t-A\psi)=0\). In a two-dimensional fundamental solution space, the difference is a linear combination of basis solutions; the coefficients depend on time and spectral parameter but not on the spatial variable in that fixed basis. They enforce the chosen initial-data normalization.}
\GSfield{Reconstruction}{The notebook's \(\lambda\) here is a gauge coefficient, not the spectral value used on preceding pages; this collision of notation is resolved.}
\GSfield{Fate}{\citet[Proposition~2.2 and its derivation, p.~61]{novikov1974periodic} gives the normalization correction and monodromy evolution.}
\GSfield{Present status}{As of 9 November 2026: Established in the stated scope; no unresolved claim is inferred from this occurrence.}
\GSfield{Assessment}{literal statement: correct with missing hypotheses; publication match: exact; verification: read the result; development: annotated; confidence: high}
\end{GSitem}
\begin{GSitem}{1973-326.21}{A}{Monodromy conjugation under time and basepoint changes}
\GSfield{Kind; source}{theorem; \GSpage{P000541}}
\GSfield{The notes}{\(\dot T=[\Lambda,T]\), \(T_{x_0}=[Q,T]\), followed by the stationary equation \(\Lambda^{\prime}=[\Lambda,Q]\).}
\GSfield{Statement}{For a smooth periodic Lax system with compatible auxiliary operator, the period monodromy evolves by conjugation under time evolution and under moving the basepoint. With a consistently chosen fundamental matrix the equations are commutator equations, so trace, determinant, and multipliers are invariant. The sign of the displayed commutator depends on whether the basis or coordinate vectors are propagated.}
\GSfield{Reconstruction}{The common conjugation statement is independent of basis conventions; the notes use both row-like and column-like transformations, so individual commutator signs are not silently unified.}
\GSfield{Fate}{\citet[Proposition~2.2 and Corollary~2.3, pp.~61--62]{novikov1974periodic} is an exact primary match to this argument.}
\GSfield{Present status}{As of 9 November 2026: Established in the stated scope; no unresolved claim is inferred from this occurrence.}
\GSfield{Assessment}{literal statement: correct with missing hypotheses; publication match: exact; verification: read the result; development: annotated; confidence: high}
\end{GSitem}
\begin{GSitem}{1973-326.22}{A}{Finite spectral curve and band edges}
\GSfield{Kind; source}{claim; \GSpage{P000541}}
\GSfield{The notes}{For stationary solutions, ``band edges are the roots of a polynomial.}
\GSfield{Statement}{For a regular finite-gap potential with \(g\) open gaps, the reduced hyperelliptic curve can be written \(\mu^2=\prod_{j=0}^{2g}(\lambda-E_j)\). Its simple branch points are the bottom endpoint and the endpoints of the open gaps. Infinitely many closed-gap double points of the full discriminant are not distinct roots of this finite reduced polynomial.}
\GSfield{Reconstruction}{The reduction by double factors is essential: a finite polynomial cannot list every high periodic/antiperiodic double eigenvalue.}
\GSfield{Fate}{\citet[\S2, pp.~61--62]{novikov1974periodic} obtains the algebraic spectral polynomial. The reduced-curve terminology is a modern clarification.}
\GSfield{Present status}{As of 9 November 2026: Established in the stated scope; no unresolved claim is inferred from this occurrence.}
\GSfield{Assessment}{literal statement: correct with missing hypotheses; publication match: exact; verification: read the result; development: annotated; confidence: high}
\end{GSitem}

\GStalk{327}{1974-03-04}{V. E. Zakharov}{P000542-P000545}{Dressing, compatible operators, and integrable wave equations}
The talk explains the Zakharov--Shabat dressing scheme through a Gelfand--Levitan equation, then obtains scalar and matrix nonlinear equations from compatibility. It includes KP, three-wave interaction, sine-Gordon, and nonlinear Schrödinger examples, and an unquantified Hamiltonian conjecture. The manuscript writes 1973 in the heading; catalogue chronology and adjacent dates require 1974. The general joint paper was received ten days after the talk, whereas Shabat's 1973 scalar precursor already existed.

\begin{GSitem}{1973-327.1}{A}{Constructing the odd KdV hierarchy}
\GSfield{Kind; source}{question; \GSpage{P000542}; I. M. Gelfand}
\GSfield{The notes}{Gelfand asks how to find \(A_n=D^{2n+1}+\cdots\) such that \([A_n,-D^2+u]\) is multiplication, how the resulting \(P_n\) are related, and how this connects with trace formulas and other equations.}
\GSfield{Statement}{In the formal algebra of differential polynomials in a scalar smooth potential, odd-order differential operators can be chosen recursively so that \([A_n,L]\) has order zero. They are determined up to lower hierarchy flows and scalar polynomials in \(L\), after a leading normalization is fixed. The multiplication coefficients are the KdV hierarchy; the Riccati recursion supplies its conserved densities. The question about other equations is answered here only for the displayed dressing examples.}
\GSfield{Reconstruction}{The formal differential-polynomial category and normalization freedom are specified. The notebook's derivative order \(u^{(n)}\) is not correct literally for an order-\(2n+1\) flow; the leading nonlinear equation contains \(u^{(2n+1)}\).}
\GSfield{Fate}{\citet[\S2]{novikov1974periodic} and \citet[\S\S1--2]{zakharovshabat1974scheme} provide the near-contemporary constructions. The general KdV/scattering connection also recurs from 1970-219.6 and 1970-219.9, without a prior problem ID.}
\GSfield{Present status}{As of 9 November 2026: Established in the stated scope; no unresolved claim is inferred from this occurrence.}
\GSfield{Assessment}{literal statement: correct with missing hypotheses; publication match: partial; verification: read the result; development: annotated; confidence: high; problem \hyperref[prob:GP-1973-53]{GP-1973-53}}
\end{GSitem}
\begin{GSitem}{1973-327.2}{E}{Gelfand--Levitan equation and Volterra factorization}
\GSfield{Kind; source}{construction; \GSpage{P000542}}
\GSfield{The notes}{\(K+F+\int_x^\infty K(x,s)F(s,y)\,ds=0\), and \(1+\hat F=(1+\hat K_+)^{-1}(1+\hat K_-)\), with opposite triangular supports.}
\GSfield{Statement}{For sufficiently smooth, decaying scalar or matrix kernels, suppose each restricted integral equation on \([x,\infty)\) is invertible and the corresponding Volterra operators are well-defined. Solving the equation for \(y\ge x\) gives the upper triangular factor \(K_+\); multiplication determines a lower triangular remainder \(K_-\). Conversely such an invertible factorization gives that kernel equation. An unsigned integral bound on \(F\) alone does not guarantee invertibility.}
\GSfield{Reconstruction}{Absolute-integrability/operator hypotheses and invertibility are added; the source gives only an insufficient unsigned bound. The construction is already present in Shabat's 1973 scalar setting, so this item is E, with the matrix version a 1974 extension.}
\GSfield{Fate}{\citet[\S1, equations (4)--(10), pp.~43--45 of the Russian original]{zakharovshabat1974scheme} is the primary match; received 14 March 1974 and published later in 1974. Its scalar dressing antecedent is \citet[Lemma~1 and Theorem~1]{shabat1973kdv}, received 13 December 1972 and published in 1973.}
\GSfield{Present status}{As of 9 November 2026: Established in the stated scope; no unresolved claim is inferred from this occurrence.}
\GSfield{Assessment}{literal statement: correct with missing hypotheses; publication match: partial; verification: read the result; development: annotated; confidence: high}
\end{GSitem}
\begin{GSitem}{1973-327.3}{E}{Parameter differentiation of the integral equation}
\GSfield{Kind; source}{method; \GSpage{P000542}}
\GSfield{The notes}{Differentiating in \(t\) gives \(K_t+F_t+K_t*F+K*F_t=0\).}
\GSfield{Statement}{If the kernels and their \(t\)-derivatives are integrable uniformly in a parameter neighborhood, and the inverses of the restricted operators vary continuously, differentiation under the integral is valid and the displayed linear equation uniquely determines \(K_t\). Here \((K*F)(x,y)=\int_x^\infty K(x,s)F(s,y)\,ds\).}
\GSfield{Reconstruction}{Uniform domination and differentiability of the inverse are sufficient analytic assumptions, omitted in the notebook.}
\GSfield{Fate}{\citet[\S1, equations (4)--(10), pp.~43--45 of the Russian original]{zakharovshabat1974scheme} is the primary match; received 14 March 1974 and published later in 1974. Its scalar dressing antecedent is \citet[Lemma~1 and Theorem~1]{shabat1973kdv}, received 13 December 1972 and published in 1973.}
\GSfield{Present status}{As of 9 November 2026: Established in the stated scope; no unresolved claim is inferred from this occurrence.}
\GSfield{Assessment}{literal statement: correct with missing hypotheses; publication match: exact; verification: read the result; development: annotated; confidence: high}
\end{GSitem}
\begin{GSitem}{1973-327.4}{E}{Lifting kernel derivations through dressing}
\GSfield{Kind; source}{method; \GSpage{P000542}}
\GSfield{The notes}{The target identity is \(\widetilde DK+DF+\widetilde DK*F+K*DF=0\); for \(D=\partial_x+\partial_y\) the same operator works on \(K\).}
\GSfield{Statement}{For smooth decaying kernels, \(D=\partial_x+\partial_y\) obeys the Leibniz identity for the truncated convolution: integration by parts cancels the lower-limit boundary term. For higher free commutator derivations, the boundary terms require lower-order corrections, encoded by a dressed differential operator \(L_n\).}
\GSfield{Reconstruction}{The cancellation of the diagonal boundary contribution, essential to the first-order assertion, is stated explicitly.}
\GSfield{Fate}{\citet[\S1, equations (4)--(10), pp.~43--45 of the Russian original]{zakharovshabat1974scheme} is the primary match; received 14 March 1974 and published later in 1974. Its scalar dressing antecedent is \citet[Lemma~1 and Theorem~1]{shabat1973kdv}, received 13 December 1972 and published in 1973.}
\GSfield{Present status}{As of 9 November 2026: Established in the stated scope; no unresolved claim is inferred from this occurrence.}
\GSfield{Assessment}{literal statement: correct with missing hypotheses; publication match: exact; verification: read the result; development: annotated; confidence: high}
\end{GSitem}
\begin{GSitem}{1973-327.5}{E}{Kernel commutators and formal adjoints}
\GSfield{Kind; source}{construction; \GSpage{P000542-P000543}}
\GSfield{The notes}{Equation (4) has an obliterated sign; equation (5) prints a plus sign. The bottom line identifies \(K_{xx}-K_{yy}\) with \([\hat K,D^2]\). Later formulas use an adjoint superscript inconsistently.}
\GSfield{Statement}{For a full-line smooth rapidly decaying kernel, the kernel of \([D^n,F]\) is \(\partial_x^nF-(-1)^n\partial_y^nF\). More generally \([A,F]\) has kernel \(A_xF-A_y^{\mathrm t}F\), with the formal transpose acting in the second variable. For triangular kernels there are additional diagonal boundary terms, absorbed in the dressed coefficients. Thus the displayed second-order difference corresponds to \([D^2,\hat K]\), not the reversed commutator.}
\GSfield{Reconstruction}{The missing sign is reconstructed by integration by parts, not deciphered from overwritten ink. Formal transpose is distinguished from Hermitian adjoint for complex/matrix kernels.}
\GSfield{Fate}{\citet[\S1, equations (4)--(10), pp.~43--45 of the Russian original]{zakharovshabat1974scheme} is the primary match; received 14 March 1974 and published later in 1974. Its scalar dressing antecedent is \citet[Lemma~1 and Theorem~1]{shabat1973kdv}, received 13 December 1972 and published in 1973.}
\GSfield{Present status}{As of 9 November 2026: Established in the stated scope; no unresolved claim is inferred from this occurrence.}
\GSfield{Assessment}{literal statement: false as written; publication match: exact; verification: read the result; development: annotated; confidence: high}
\end{GSitem}
\begin{GSitem}{1973-327.6}{E}{Diagonal jets and dressed differential coefficients}
\GSfield{Kind; source}{construction; \GSpage{P000543}}
\GSfield{The notes}{\(L_n=D^n+n(dK(x,x)/dx)D^{n-2}+\cdots\), with all lower coefficients formed from \(\xi_j=(\partial_x-\partial_y)^jK|_{y=x}\).}
\GSfield{Statement}{When the dressing intertwining equation is solved recursively for a smooth Volterra kernel, each coefficient of the monic dressed operator is a universal differential polynomial in finitely many diagonal jets of \(K\), with products in their prescribed order for matrix kernels. The leading correction is \(n\partial_xK(x,x)D^{n-2}\) in the notebook's convention; the absent \(D^{n-1}\) term is a normalization.}
\GSfield{Reconstruction}{The words 'nonlinear operator in \(K\)' mean nonlinear dependence of coefficients on \(K\); \(L_n\) remains linear on the function to which it is applied.}
\GSfield{Fate}{\citet[\S1, recurrence (10)]{zakharovshabat1974scheme} gives the coefficient recurrence; \citet[p.~1311]{shabat1973kdv} already prints the first scalar dressed operators.}
\GSfield{Present status}{As of 9 November 2026: Established in the stated scope; no unresolved claim is inferred from this occurrence.}
\GSfield{Assessment}{literal statement: correct with missing hypotheses; publication match: exact; verification: read the result; development: annotated; confidence: high}
\end{GSitem}
\begin{GSitem}{1973-327.7}{F}{A larger Lie algebra of admissible derivations}
\GSfield{Kind; source}{remark; \GSpage{P000543}}
\GSfield{The notes}{Variable-coefficient operators are also possible; all extendible \(D\) form a Lie algebra, described by A. Shabat.}
\GSfield{Statement}{Commutators of admissible differentiations preserve a common kernel algebra whenever the differentiations preserve the domain and the dressing map is invertible. The notes do not specify the domain, all allowed variable coefficients, or the claimed complete classification.}
\GSfield{Reconstruction}{A closure property is supplied as context; the unspecified completeness claim is left unidentified.}
\GSfield{Fate}{\citet[p.~1310, formulas (2)--(3)]{shabat1973kdv} explicitly studies a subalgebra generated by multiplication differences and differential commutators. That is a partial identifiable antecedent, not verification that all extendible operators have been classified.}
\GSfield{Present status}{As of 9 November 2026: The asserted exhaustive classification cannot be identified from the notebook.}
\GSfield{Assessment}{literal statement: not determinable; publication match: partial; verification: read the result; development: annotated; confidence: medium}
\end{GSitem}
\begin{GSitem}{1973-327.8}{E}{Linear constraints become dressed nonlinear constraints}
\GSfield{Kind; source}{theorem; \GSpage{P000543}}
\GSfield{The notes}{Unique solvability of the Gelfand--Levitan equation and \(DF=0\) imply \(\widetilde DK=0\); the next line strengthens this to an equivalence.}
\GSfield{Statement}{Under the differentiated integral identity and invertibility of \(I+F_x\), \(DF=0\) implies \(\widetilde DK=0\). The reverse implication follows when the dressing correspondence is also invertible in the other direction on the chosen class, so that the resulting homogeneous equation for \(DF\) has only the zero solution.}
\GSfield{Reconstruction}{Unique solvability only for \(K\) does not, by itself, justify both arrows; the additional invertibility assumption is isolated.}
\GSfield{Fate}{\citet[Lemma~1]{shabat1973kdv} proves the equivalence for regular differentiable kernels with its explicit assumptions. \citet[Theorem~1]{zakharovshabat1974scheme} gives the general operator formulation.}
\GSfield{Present status}{As of 9 November 2026: Established in the stated scope; no unresolved claim is inferred from this occurrence.}
\GSfield{Assessment}{literal statement: correct with missing hypotheses; publication match: exact; verification: read the result; development: annotated; confidence: high}
\end{GSitem}
\begin{GSitem}{1973-327.9}{A}{Two compatible dressed operators}
\GSfield{Kind; source}{theorem; \GSpage{P000543}}
\GSfield{The notes}{If \(D_1F=D_2F=0\), then \(\beta L^{(1)}_z-\alpha L^{(2)}_t=[L^{(1)},L^{(2)}]\).}
\GSfield{Statement}{Let \(M_{10}=\alpha\partial_t+L_{10}\) and \(M_{20}=\beta\partial_z+L_{20}\) commute, with constant coefficient differential parts. Suppose an invertible Volterra dressing operator intertwines both with \(M_1=\alpha\partial_t+L_1\) and \(M_2=\beta\partial_z+L_2\). Then \([M_1,M_2]=0\), equivalently \(\beta L_{1,z}-\alpha L_{2,t}=[L_1,L_2]\). Matrix coefficients require the free operators to commute, not merely to be individually constant.}
\GSfield{Reconstruction}{Commutativity of the free pair and the shared invertible dressing are indispensable hypotheses, made explicit.}
\GSfield{Fate}{\citet[Theorem~2, pp.~45--46]{zakharovshabat1974scheme} is an exact primary match. This general two-operator scheme was published after the 4 March 1974 talk.}
\GSfield{Present status}{As of 9 November 2026: Established in the stated scope; no unresolved claim is inferred from this occurrence.}
\GSfield{Assessment}{literal statement: correct with missing hypotheses; publication match: exact; verification: read the result; development: annotated; confidence: high}
\end{GSitem}
\begin{GSitem}{1973-327.10}{A}{Generating solutions from linear kernel equations}
\GSfield{Kind; source}{method; \GSpage{P000544}}
\GSfield{The notes}{Choose a pair of free operators, solve \(D_1F=D_2F=0\), then recover \(K\) and its diagonal jets to obtain a nonlinear solution.}
\GSfield{Statement}{For kernels meeting the dressing hypotheses, the sequence \(F\mapsto K\mapsto(\xi_0,\xi_1,\ldots)\) maps solutions of the compatible linear kernel constraints to solutions of the coefficient equations obtained from zero curvature. This is a solution construction on the invertible dressing class; it does not alone prove that every solution or every Cauchy datum belongs to that class.}
\GSfield{Reconstruction}{Surjectivity, a global inverse map, and well-posedness for unrestricted data are not added to the notebook claim.}
\GSfield{Fate}{\citet[Theorems~1--2 and \S2]{zakharovshabat1974scheme} states the scheme; the corresponding scalar construction is already in \citet[Theorem~1]{shabat1973kdv}.}
\GSfield{Present status}{As of 9 November 2026: Established in the stated scope; no unresolved claim is inferred from this occurrence.}
\GSfield{Assessment}{literal statement: correct with missing hypotheses; publication match: exact; verification: read the result; development: annotated; confidence: high}
\end{GSitem}
\begin{GSitem}{1973-327.11}{A}{The second/third-order KP pair}
\GSfield{Kind; source}{construction; \GSpage{P000544}}
\GSfield{The notes}{\(P=D^2+u\), \(Q=D^3+\frac34(uD+Du)+\lambda D+w\), with \(u=2\partial_xK(x,x)\); \(w\) is expressed using \(\xi_1+\xi_0^2\).}
\GSfield{Statement}{For scalar smooth \(u,w\), the displayed operators satisfy \([P,Q]=2w_xD+w_{xx}-\tfrac14u_{xxx}-\tfrac32uu_x-\lambda u_x\). Assign the transverse derivative to \(\beta\partial_z+P\) and the evolution derivative to \(\alpha\partial_t+Q\). Their compatibility gives \(w_x=-\tfrac34\beta u_z\) and \(\beta w_z=\alpha u_t+\tfrac14u_{xxx}+\tfrac32uu_x+\lambda u_x\).}
\GSfield{Reconstruction}{The notebook switches the labels of evolution and transverse variables and writes \(w_y\) in one line. The modern statement fixes an explicit assignment instead of pretending the printed system is internally consistent.}
\GSfield{Fate}{\citet[\S2, second/third-order example]{zakharovshabat1974scheme} gives the KP reduction. The exact displayed-coefficient commutator in this statement is an editorial calculation reproduced in the included verification code.}
\GSfield{Present status}{As of 9 November 2026: Established in the stated scope; no unresolved claim is inferred from this occurrence.}
\GSfield{Assessment}{literal statement: false as written; publication match: partial; verification: read the result; development: annotated; confidence: high}
\end{GSitem}
\begin{GSitem}{1973-327.12}{A}{Scalar KP equation and its KdV reduction}
\GSfield{Kind; source}{theorem; \GSpage{P000544}}
\GSfield{The notes}{The notes print \(\partial_x\{\alpha u_t+6uu_x+u_{xxx}+\lambda u_x=(3/4)\beta^2u_{zz}\), and say that independence of \(z\) gives KdV.}
\GSfield{Statement}{With the literal operator pair and assignment just specified, elimination of \(w\) gives \((4\alpha u_t+u_{xxx}+6uu_x+4\lambda u_x)_x+3\beta^2u_{zz}=0\). For real nonzero \(\beta\), real rescalings preserve the relative sign of the transverse term and the fourth \(x\)-derivative, giving one real KP form. The other real KP form requires an imaginary choice of \(\beta\) with a compatible reality reduction, or an explicitly stated complex continuation; it does not follow from real rescaling alone. If \(u_z=0\), decay at spatial infinity, or an appropriate fixed integration constant, gives the corresponding KdV equation after integrating once in \(x\). Without such a condition the integrated right side is a function of time.}
\GSfield{Reconstruction}{The closing brace, coefficients, transverse sign, and integration hypothesis are supplied explicitly. Parameters cannot be reused unchanged across the inconsistent displayed formulas. The two real KP forms require separate reality conditions; the positive factor \(\beta^2\) for real \(\beta\) cannot be changed by a real coordinate rescaling.}
\GSfield{Fate}{\citet[\S2]{zakharovshabat1974scheme} identifies the equation and reduction. This is a corrected normalization of the notebook's calculation, not a new result.}
\GSfield{Present status}{As of 9 November 2026: Established in the stated scope; no unresolved claim is inferred from this occurrence.}
\GSfield{Assessment}{literal statement: false as written; publication match: partial; verification: read the result; development: annotated; confidence: high}
\end{GSitem}
\begin{GSitem}{1973-327.13}{A}{Triangular matrix solutions encode linearization}
\GSfield{Kind; source}{construction; \GSpage{P000544}}
\GSfield{The notes}{Putting \(u=\left(\begin{smallmatrix}u_0&u_1\\0&u_0\end{smallmatrix}\right)\) yields a linearized equation for \(u_1\), followed by higher linearizations.}
\GSfield{Statement}{For a differential-polynomial equation \(E(u)=0\), substitute \(u=u_0I+u_1N\), where \(N^2=0\). Then \(E(u)=E(u_0)I+(DE)_{u_0}[u_1]N\). Thus the diagonal part solves the nonlinear scalar equation and the upper entry solves its linearization. Larger nilpotent algebras encode higher variation equations, including their inhomogeneous lower-variation terms.}
\GSfield{Reconstruction}{The square-zero calculation is a reconstruction; 'higher linearizations' does not mean independent homogeneous copies at every order.}
\GSfield{Fate}{The matrix extension belongs to \citet[\S\S1--2]{zakharovshabat1974scheme}; the square-zero substitution is an elementary algebraic explanation of the notebook statement.}
\GSfield{Present status}{As of 9 November 2026: Established in the stated scope; no unresolved claim is inferred from this occurrence.}
\GSfield{Assessment}{literal statement: correct with missing hypotheses; publication match: partial; verification: read the result; development: annotated; confidence: high}
\end{GSitem}
\begin{GSitem}{1973-327.14}{E}{Physical interpretations and real forms of KP}
\GSfield{Kind; source}{remark; \GSpage{P000544}}
\GSfield{The notes}{Real or imaginary \(\beta\), the drift parameter \(\lambda\), thin-wing flow, nonlinear strings, sound dispersion, and shallow water are mentioned; Karpman's book is recommended.}
\GSfield{Statement}{In a normalized real KP equation the sign of the transverse second derivative distinguishes two real forms. A constant longitudinal drift can be removed by a moving coordinate. The notes identify physical settings, but do not state the scaling limits, boundary conditions, or a theorem validating any particular thin-wing or shallow-water approximation.}
\GSfield{Reconstruction}{Only the elementary effects of sign and drift are stated precisely; no physical error estimate is inferred.}
\GSfield{Fate}{\citet[\S2]{zakharovshabat1974scheme} is an equation-level match, not a proof of the physical approximations. The incomplete book reference is recorded as a bibliographic pointer rather than an independently recovered theorem.}
\GSfield{Present status}{As of 9 November 2026: The equation-level interpretations are standard; the physical approximation hypotheses are absent.}
\GSfield{Assessment}{literal statement: not determinable; publication match: thematic only; verification: read the result; development: annotated; confidence: medium}
\end{GSitem}
\begin{GSitem}{1973-327.15}{A}{First-order matrix dressing and the N-wave system}
\GSfield{Kind; source}{construction; \GSpage{P000544-P000545}}
\GSfield{The notes}{For commuting diagonal \(\ell_1,\ell_2\), \(L_j=\ell_jD+[\ell_j,\xi_0]\); the first nontrivial Hermitian case is \(3\times3\).}
\GSfield{Statement}{Let \(A,B\) be commuting constant diagonal matrices and \(X=X(x,t,z)\) smooth. Compatibility of \(\partial_t+AD+[A,X]\) and \(\partial_z+BD+[B,X]\) is the matrix PDE \([B,X]_t-[A,X]_z+A[B,X]_x-B[A,X]_x+[[A,X],[B,X]]=0\). It closes on off-diagonal entries after choosing a diagonal gauge; reality reductions require consistent Hermitian/sign conventions.}
\GSfield{Reconstruction}{The notebook's first free-operator line mixes \(x\) and \(y\). The pair is reconstructed from its two separately displayed dressed operators, with the compatibility equation fully specified.}
\GSfield{Fate}{\citet[\S3, matrix first-order operators]{zakharovshabat1974scheme} develops this matrix construction.}
\GSfield{Present status}{As of 9 November 2026: Established in the stated scope; no unresolved claim is inferred from this occurrence.}
\GSfield{Assessment}{literal statement: correct with missing hypotheses; publication match: partial; verification: read the result; development: annotated; confidence: high}
\end{GSitem}
\begin{GSitem}{1973-327.16}{A}{Resonant three-wave interaction}
\GSfield{Kind; source}{claim; \GSpage{P000545}}
\GSfield{The notes}{For \(u_1=\xi_{13}\), \(u_2=\xi_{12}\), \(u_3=\xi_{23}\), the notes print \((\partial_t+V_1\cdot\nabla)u_1=i\lambda u_2u_3\), \((\partial_t+V_2\cdot\nabla)u_2=i\lambda u_1\bar u_3\), and \((\partial_t+V_3\cdot\nabla)u_3=i\lambda u_1\bar u_2\).}
\GSfield{Statement}{The displayed system is a three-complex-amplitude resonant interaction model in two spatial variables, with constant velocities and coupling. Its connection with the preceding integrable matrix system requires a realizable diagonal-matrix parametrization, a compatible reality reduction, and nondegenerate rescalings of the three components. Arbitrary degeneracies of the velocity parameters are not justified by the notes.}
\GSfield{Reconstruction}{The velocity/coupling compatibility qualifications are supplied. The explicit three-wave equations, rather than an entire classification of models, are the identifiable item.}
\GSfield{Fate}{\citet[\S3]{zakharovshabat1974scheme} supplies the integrable first-order matrix family. This is a partial match to the notebook's unrestricted phrase 'arbitrary vectors'; no global existence or unrestricted Cauchy theorem is asserted.}
\GSfield{Present status}{As of 9 November 2026: Established in the stated scope; no unresolved claim is inferred from this occurrence.}
\GSfield{Assessment}{literal statement: correct with missing hypotheses; publication match: partial; verification: read the result; development: annotated; confidence: medium}
\end{GSitem}
\begin{GSitem}{1973-327.17}{F}{The prohibition on adding dissipation}
\GSfield{Kind; source}{remark; \GSpage{P000545}}
\GSfield{The notes}{\ru{Диссипацию ввести нельзя (никогда)}: dissipation cannot be introduced, ever.}
\GSfield{Statement}{No precise theorem is specified: neither the allowed modification of the equations nor the meaning of preserving integrability is defined. The remark can be read as a restriction of the conservative isospectral scheme being discussed; it cannot establish a universal impossibility theorem for dissipative solvable equations.}
\GSfield{Fate}{The isospectral framework in \citet[\S1]{zakharovshabat1974scheme} gives context only. The absolute wording has no identifiable quantified publication match.}
\GSfield{Present status}{As of 9 November 2026: An informal restriction of the discussion, not a well-posed open problem.}
\GSfield{Assessment}{literal statement: not determinable; publication match: thematic only; verification: unresolved lead; development: annotated; confidence: high}
\end{GSitem}
\begin{GSitem}{1973-327.18}{F}{The unquantified Hamiltonian conjecture}
\GSfield{Kind; source}{conjecture; \GSpage{P000545}}
\GSfield{The notes}{\ru{Гипотеза}: all these systems are Hamiltonian; one \(L\) gives the symplectic form, the other the Hamiltonian.}
\GSfield{Statement}{The strongest literal content is a proposed Hamiltonian organization of the displayed dressing equations. No phase space, reduction, boundary conditions, degeneracy quotient, or quantifier over allowable operator pairs is stated, so there is no unique precise universal conjecture to classify as solved or open.}
\GSfield{Fate}{Hamiltonian KdV was already treated in \citet{zakharovfaddeev1971kdv}. The two-operator method of \citet{zakharovshabat1974scheme} does not by itself verify the notebook's universal Hamiltonian claim. Those publications are partial examples and context, not a proof of every possible interpretation.}
\GSfield{Present status}{As of 9 November 2026: No definite unresolved conjecture is extracted; choosing a modern phase space and reduction would create a new problem.}
\GSfield{Assessment}{literal statement: not determinable; publication match: partial; verification: unresolved lead; development: annotated; confidence: high}
\end{GSitem}
\begin{GSitem}{1973-327.19}{E}{Sine-Gordon as an integrable reduction}
\GSfield{Kind; source}{claim; \GSpage{P000545}}
\GSfield{The notes}{\(u_{tt}-u_{xx}=\sin u\) is said to arise from two displayed diagonal \(4\times4\) matrices.}
\GSfield{Statement}{After a sign/coordinate convention is fixed, sine-Gordon admits an auxiliary linear system whose compatibility is the nonlinear equation and an inverse-scattering construction for an appropriate decaying or finite-energy scattering class. The two free diagonal matrices alone do not specify the needed potential matrix, reality constraint, or reduction.}
\GSfield{Reconstruction}{No full four-dimensional Lax matrix is supplied from the missing entries. The primary publisher abstract was read; the exact representation was not verified.}
\GSfield{Fate}{\citet[abstract and pp.~1262--1264]{ablowitzkaupnewellsegur1973sine} was received 7 March 1973 and published 18 June 1973, before this talk. Its inverse-scattering result is an exact subject/result match; the specific undocumented \(4\times4\) realization remains only partial.}
\GSfield{Present status}{As of 9 November 2026: Established in the stated scope; no unresolved claim is inferred from this occurrence.}
\GSfield{Assessment}{literal statement: correct with missing hypotheses; publication match: partial; verification: publisher metadata checked; development: annotated; confidence: high}
\end{GSitem}
\begin{GSitem}{1973-327.20}{F}{Laser-pulse and two-level Bloch equations}
\GSfield{Kind; source}{remark; \GSpage{P000545}}
\GSfield{The notes}{An equation for propagation of a laser pulse is mentioned with an ellipsis; the next line says Bloch equations for density matrices in a two-level quantum system.}
\GSfield{Statement}{The subject suggests coupled optical-field and two-level density-matrix equations, but the notes supply no system, coupling, relaxation terms, initial data, or exact assertion. A particular Maxwell--Bloch system cannot be assigned as though it were recorded.}
\GSfield{Fate}{The general matrix-integrability discussion in \citet{zakharovshabat1974scheme} supplies only thematic context. No exact equation-level identification is claimed.}
\GSfield{Present status}{As of 9 November 2026: The omitted model is unidentifiable from the surviving text.}
\GSfield{Assessment}{literal statement: not determinable; publication match: thematic only; verification: unresolved lead; development: annotated; confidence: high}
\end{GSitem}
\begin{GSitem}{1973-327.21}{E}{Nonlinear Schrödinger inverse scattering}
\GSfield{Kind; source}{claim; \GSpage{P000545}}
\GSfield{The notes}{A first-order diagonal free operator paired with a second-order one yields \(i\xi_t+\xi_{xx}\pm|\xi|^2\xi=0\); there is a scheme for the Cauchy problem.}
\GSfield{Statement}{For smooth scattering data in the appropriate one-dimensional focusing or defocusing NLS class, the Zakharov--Shabat direct transform turns the evolution into explicit time factors; an inverse integral/scattering problem reconstructs the solution while its solvability conditions hold. The decaying focusing problem and the nonzero-background defocusing soliton problem require different spectral and boundary hypotheses.}
\GSfield{Reconstruction}{Boundary classes and the distinction between a scattering construction and unrestricted well-posedness are added.}
\GSfield{Fate}{\citet[\S\S1--3]{zakharovshabat1972nls} was submitted 22 December 1970, published in Russian in July 1971 and in English in January 1972. The stable-medium extension is \citet{zakharovshabat1973stable}. Both predate the seminar; the notebook's unrestricted phrase 'Cauchy problem' is not a modern all-data well-posedness theorem.}
\GSfield{Present status}{As of 9 November 2026: Established in the stated scope; no unresolved claim is inferred from this occurrence.}
\GSfield{Assessment}{literal statement: correct with missing hypotheses; publication match: partial; verification: read the result; development: annotated; confidence: high}
\end{GSitem}

\GStalk{328}{1974-03-11}{A. Ya. Povzner}{P000546}{Whitham modulation of nonlinear waves}
Povzner presents Whitham's slowly modulated traveling-wave method for nonlinear Klein--Gordon equations. The one-page outline mixes fast and slow variables and gives an incompletely specified action relation. The identifiable reference is Whitham's 1970 article, although the notebook writes 1972.

\begin{GSitem}{1973-328.1}{E}{Slow variables for a nonlinear Klein--Gordon wave}
\GSfield{Kind; source}{method; \GSpage{P000546}}
\GSfield{The notes}{\(u_{tt}-u_{xx}+T'(u)=0\); far-field traveling waves are studied by \(x\mapsto x/\varepsilon\), \(t\mapsto t/\varepsilon\).}
\GSfield{Statement}{Let \(V\) be a smooth potential. If \(v(X,S)=u(X/\varepsilon,S/\varepsilon)\), then the exact rescaled equation is \(\varepsilon^2(v_{SS}-v_{XX})+V'(v)=0\). Slow modulation concerns families whose phase varies on this long scale; the scaling identity alone gives no approximation estimate.}
\GSfield{Reconstruction}{The potential is renamed \(V\) to avoid a collision with period notation. Fast/slow coordinates are distinguished, and the notebook's \(v\)/\(u\) inconsistency is corrected.}
\GSfield{Fate}{\citet[abstract and nonlinear-wave variational formulation]{whitham1970twotiming} describes two-timing and averaging. The article was received 10 January 1970 and published in 1970, before the 1974 exposition; the Caltech primary repository metadata and abstract were opened.}
\GSfield{Present status}{As of 9 November 2026: Established in the stated scope; no unresolved claim is inferred from this occurrence.}
\GSfield{Assessment}{literal statement: correct with missing hypotheses; publication match: exact; verification: publisher metadata checked; development: annotated; confidence: high}
\end{GSitem}
\begin{GSitem}{1973-328.2}{E}{Rapid phase and the leading oscillator equation}
\GSfield{Kind; source}{construction; \GSpage{P000546}}
\GSfield{The notes}{The ansatz is \(u=f(\Theta_0/\varepsilon;x,t)\), with \(\omega_2=\Theta_{0t}\), \(\omega_1=\Theta_{0x}\). The purported limiting equation still contains \(\varepsilon^2\).}
\GSfield{Statement}{For \(v(X,S)=f(\theta;X,S)\), \(\theta=\Theta(X,S)/\varepsilon\), write \(\omega=\Theta_S\), \(k=\Theta_X\). Substitution into the scaled equation gives the leading equation \((\omega^2-k^2)f_{\theta\theta}+V'(f)=0\). Slow derivatives give terms of order \(\varepsilon\) and \(\varepsilon^2\), provided the required derivatives of \(f,\Theta\) are bounded.}
\GSfield{Reconstruction}{This chain-rule calculation is supplied explicitly; the persisting \(\varepsilon^2\) and the extra overwritten factor after \(T'(u)\) are not mathematically consistent with this fast phase.}
\GSfield{Fate}{\citet[abstract and nonlinear-wave variational formulation]{whitham1970twotiming} describes two-timing and averaging. The article was received 10 January 1970 and published in 1970, before the 1974 exposition; the Caltech primary repository metadata and abstract were opened.}
\GSfield{Present status}{As of 9 November 2026: Established in the stated scope; no unresolved claim is inferred from this occurrence.}
\GSfield{Assessment}{literal statement: false as written; publication match: exact; verification: publisher metadata checked; development: annotated; confidence: high}
\end{GSitem}
\begin{GSitem}{1973-328.3}{E}{Traveling-wave quadrature and integration constants}
\GSfield{Kind; source}{construction; \GSpage{P000546}}
\GSfield{The notes}{The leading second-order ODE has constants \(E_1,E_2\), and the notes integrate \(\sqrt{E-T(u)}\) between turning points.}
\GSfield{Statement}{Freeze \(X,S\), suppose \(a=\omega^2-k^2>0\), and assume an energy \(E\) has two simple turning points \(u_-<u_+\) with \(V(u)<E\) in between. The periodic oscillator obeys \(\tfrac12a f_\theta^2+V(f)=E\), with phase shift as the second integration constant and period \(\sqrt{2a}\int_{u_-}^{u_+}(E-V(s))^{-1/2}\,ds\).}
\GSfield{Reconstruction}{A nondegenerate periodic orbit, the sign of \(a\), and the distinction between energy and phase shift are supplied; arbitrary integration constants do not always produce periodic real waves.}
\GSfield{Fate}{\citet[abstract and nonlinear-wave variational formulation]{whitham1970twotiming} describes two-timing and averaging. The article was received 10 January 1970 and published in 1970, before the 1974 exposition; the Caltech primary repository metadata and abstract were opened.}
\GSfield{Present status}{As of 9 November 2026: Established in the stated scope; no unresolved claim is inferred from this occurrence.}
\GSfield{Assessment}{literal statement: correct with missing hypotheses; publication match: exact; verification: publisher metadata checked; development: annotated; confidence: high}
\end{GSitem}
\begin{GSitem}{1973-328.4}{E}{Averaged variational equations}
\GSfield{Kind; source}{method; \GSpage{P000546}}
\GSfield{The notes}{An additional equation for the phase is sought after inserting the local oscillatory family; Whitham's variational work is cited.}
\GSfield{Statement}{For a smooth family of periodic waves normalized to a fixed phase period, average \(\mathcal L=\tfrac12(\omega^2-k^2)f_\theta^2-V(f)\) over that period. Formal variation of the slow phase gives \(\partial_S\overline{\mathcal L}_{\omega}+\partial_X\overline{\mathcal L}_k=0\), together with \(k_S=\omega_X\). Additional wave parameters satisfy their averaged stationarity equations. This is a formal leading modulation system, not an all-time error theorem.}
\GSfield{Reconstruction}{The averaged Euler--Lagrange and compatibility equations are a modern reconstruction of the cited method; their precise formulas are absent from the notebook.}
\GSfield{Fate}{\citet[abstract and nonlinear-wave variational formulation]{whitham1970twotiming} describes two-timing and averaging. The article was received 10 January 1970 and published in 1970, before the 1974 exposition; the Caltech primary repository metadata and abstract were opened.}
\GSfield{Present status}{As of 9 November 2026: Established in the stated scope; no unresolved claim is inferred from this occurrence.}
\GSfield{Assessment}{literal statement: correct with missing hypotheses; publication match: partial; verification: publisher metadata checked; development: annotated; confidence: high}
\end{GSitem}
\begin{GSitem}{1973-328.5}{F}{The unexplained constant action relation}
\GSfield{Kind; source}{claim; \GSpage{P000546}}
\GSfield{The notes}{\(\sqrt{\omega_2^2-\omega_1^2}\,\varphi(E_1,E_2)=\mathrm{const}\), where \(\varphi=\int_{u_0}^{u_1}\sqrt{E-T(u)}\,du\).}
\GSfield{Statement}{For the periodic oscillator just specified, \(\sqrt{\omega^2-k^2}\int_{u_-}^{u_+}\sqrt{E-V(s)}\,ds\) is proportional to its action. General two-variable modulation imposes a conservation law for an action density and flux, not an unspecified constant throughout space-time. A constant action can arise only after a further reduction or normalization that is not recorded.}
\GSfield{Reconstruction}{The proportionality to oscillator action is identified; the missing reduction is not invented.}
\GSfield{Fate}{\citet[abstract and nonlinear-wave variational formulation]{whitham1970twotiming} describes two-timing and averaging. The article was received 10 January 1970 and published in 1970, before the 1974 exposition; the Caltech primary repository metadata and abstract were opened. The primary article is a thematic match for action averaging; no exact reduction producing the notebook's particular constant relation has been verified.}
\GSfield{Present status}{As of 9 November 2026: The intended constant-action reduction remains unidentified; the standard formal modulation method is established.}
\GSfield{Assessment}{literal statement: not determinable; publication match: thematic only; verification: unresolved lead; development: annotated; confidence: medium}
\end{GSitem}
\begin{GSitem}{1973-328.6}{E}{Scope of the Whitham exposition}
\GSfield{Kind; source}{remark; \GSpage{P000546}}
\GSfield{The notes}{The closing references mention KdV, heavy fluids, radiophysics, and lasers, then J. Fluid Mechanics, 1972, 'Two-timing ...'.}
\GSfield{Statement}{The recognizable bibliographic reference is G. B. Whitham, 'Two-timing, variational principles and waves', Journal of Fluid Mechanics 44(2), 373--395 (1970). The listed applications are pointers, with no separate PDE or theorem stated on this page.}
\GSfield{Reconstruction}{The publication year and full title are reconstructed from the repository record; the application list is not expanded into unrecorded mathematical claims.}
\GSfield{Fate}{\citet[abstract and nonlinear-wave variational formulation]{whitham1970twotiming} describes two-timing and averaging. The article was received 10 January 1970 and published in 1970, before the 1974 exposition; the Caltech primary repository metadata and abstract were opened.}
\GSfield{Present status}{As of 9 November 2026: Established in the stated scope; no unresolved claim is inferred from this occurrence.}
\GSfield{Assessment}{literal statement: false as written; publication match: exact; verification: publisher metadata checked; development: annotated; confidence: high}
\end{GSitem}

\GStalk{329}{1974-03-11}{V. N. Tulovskii}{P000547-P000548}{Hörmander wavefront sets and oscillatory integrals}
Tulovskii explicitly presents Hörmander's work: conic Fourier decay, kernel relations, distribution products, symbol estimates, and cutoff oscillatory integrals. The content matches the 1971 Acta paper, whose precursors include 1969 lectures. Gelfand contributes geometric explanations. The notebook's word 'continuous' must be read in the original sequential sense; its sign convention for kernel relations also needs repair.

\begin{GSitem}{1973-329.1}{E}{The smooth wavefront set}
\GSfield{Kind; source}{definition; \GSpage{P000547}}
\GSfield{The notes}{\((x_0,\xi_0)\notin WF(u)\) if a localization \(\varphi(x_0)\ne0\) has Fourier transform decaying faster than every power in a cone around \(\xi_0\).}
\GSfield{Statement}{For \(u\in\mathcal D'(\Omega)\), \(\xi_0\ne0\), choose \(\varphi\in C_c^\infty(\Omega)\) with \(\varphi(x_0)\ne0\) and an open conic neighborhood \(V\) of \(\xi_0\). The point is outside \(WF(u)\) when for every \(N\ge0\) a constant \(C_N\) satisfies \(|\widehat{\varphi u}(\xi)|\le C_N(1+|\xi|)^{-N}\) for all \(\xi\in V\). The same localization and cone serve for all \(N\).}
\GSfield{Fate}{\citet[Proposition~2.5.5 and its concluding equivalence, pp.~122--123]{hormander1971fourier} is exact. The paper appeared in 1971 and developed material circulated in the 1969 Nordic summer-school lectures.}
\GSfield{Present status}{As of 9 November 2026: Established in the stated scope; no unresolved claim is inferred from this occurrence.}
\GSfield{Assessment}{literal statement: correct; publication match: exact; verification: read the result; development: annotated; confidence: high}
\end{GSitem}
\begin{GSitem}{1973-329.2}{E}{Step functions, boundary normals, and smoothing}
\GSfield{Kind; source}{construction; \GSpage{P000547}}
\GSfield{The notes}{A step across a line and the characteristic function of a disk are drawn, with normal directions; convolution with a delta-like sequence is suggested.}
\GSfield{Statement}{If \(\Omega_0\) has a smooth boundary and \(u=1_{\Omega_0}\), then locally at each smooth boundary point \(WF(u)=N^*(\partial\Omega_0)\setminus0\). A smooth nonvanishing jump across a hypersurface has the same conormal set. Mollification approximates the distribution; conic rapid decay is uniform only away from the indicated wavefront directions, with the relevant microlocal convergence conditions.}
\GSfield{Reconstruction}{Smoothness of the boundary and a nonzero jump are supplied; corners would add other directions.}
\GSfield{Fate}{\citet[\S2.4 and pp.~125--126]{hormander1971fourier} develops conormal distributions and smoothing in the wavefront pseudotopology. This is explanatory classical content, not a new singularity theorem by Tulovskii.}
\GSfield{Present status}{As of 9 November 2026: Established in the stated scope; no unresolved claim is inferred from this occurrence.}
\GSfield{Assessment}{literal statement: correct with missing hypotheses; publication match: exact; verification: read the result; development: annotated; confidence: high}
\end{GSitem}
\begin{GSitem}{1973-329.3}{E}{Gelfand's slicing explanation of directional singularities}
\GSfield{Kind; source}{remark; \GSpage{P000547}; I. M. Gelfand}
\GSfield{The notes}{Gelfand reduces a planar Fourier integral to a one-dimensional oscillatory integral in a sliced function \(\psi(t)\); drawings contrast tangent and transverse cuts of a circle.}
\GSfield{Statement}{For a smooth compactly supported density on a domain with smooth boundary, integrating over level sets of a linear coordinate produces a smooth density near transverse intersections. At critical values where the level hyperplane is tangent to the boundary, this pushforward can be singular. Such critical directions are conormal directions; stationary-phase decay depends on the local geometry and degeneracy.}
\GSfield{Reconstruction}{The original \(\varphi\) in the planar integral cannot be a globally smooth compactly supported function if a boundary jump is intended. The characteristic-domain factor and the transversality qualification are supplied.}
\GSfield{Fate}{\citet[\S1.2 and Proposition~2.5.7]{hormander1971fourier} gives the oscillatory framework. The precise scalar slice density is an explanatory reconstruction of the drawings, not an exact formula recoverable from them.}
\GSfield{Present status}{As of 9 November 2026: Established in the stated scope; no unresolved claim is inferred from this occurrence.}
\GSfield{Assessment}{literal statement: correct with missing hypotheses; publication match: partial; verification: read the result; development: annotated; confidence: medium}
\end{GSitem}
\begin{GSitem}{1973-329.4}{E}{Wavefront projection equals singular support}
\GSfield{Kind; source}{theorem; \GSpage{P000547}}
\GSfield{The notes}{\(\pi WF(u)=\operatorname{sing\,supp}u\), \(\pi(x,\xi)=x\).}
\GSfield{Statement}{For every distribution on a smooth manifold, the projection of its smooth wavefront set to the base manifold is its singular support. In particular \(WF(u)=\varnothing\) if and only if \(u\) is smooth.}
\GSfield{Fate}{\citet[\S2.5, definition and Fourier characterization]{hormander1971fourier} establishes this fundamental property.}
\GSfield{Present status}{As of 9 November 2026: Established in the stated scope; no unresolved claim is inferred from this occurrence.}
\GSfield{Assessment}{literal statement: correct; publication match: exact; verification: read the result; development: annotated; confidence: high}
\end{GSitem}
\begin{GSitem}{1973-329.5}{E}{Schwartz kernels and the required continuity}
\GSfield{Kind; source}{theorem; \GSpage{P000547}}
\GSfield{The notes}{Every \(P:C_c^\infty(\Omega)\to C^\infty(\Omega_1)\) is assigned \(K_P\in\mathcal D'(\Omega\times\Omega_1)\).}
\GSfield{Statement}{If \(P:C_c^\infty(X)\to\mathcal D'(Y)\) is continuous in the standard test-function and distribution topologies, there is a unique distribution kernel \(K\in\mathcal D'(Y\times X)\) satisfying \(\langle Pu,v\rangle=\langle K,v\otimes u\rangle\). Mere linearity without continuity is insufficient.}
\GSfield{Reconstruction}{Continuity is supplied, and output-first kernel order is fixed for all subsequent modern formulas.}
\GSfield{Fate}{\citet[\S2.5, discussion preceding Theorem~2.5.12, p.~129]{hormander1971fourier} applies the classical kernel theorem in this setting.}
\GSfield{Present status}{As of 9 November 2026: Established in the stated scope; no unresolved claim is inferred from this occurrence.}
\GSfield{Assessment}{literal statement: correct with missing hypotheses; publication match: exact; verification: read the result; development: annotated; confidence: high}
\end{GSitem}
\begin{GSitem}{1973-329.6}{E}{Extension to compactly supported distributions}
\GSfield{Kind; source}{theorem; \GSpage{P000547}}
\GSfield{The notes}{If the kernel wavefront contains no point with either component covector zero, \(P\) extends from test functions to \(\mathcal E'(\Omega)\to\mathcal D'(\Omega_1)\).}
\GSfield{Statement}{For \(K\in\mathcal D'(Y\times X)\), assume \((y,0;x,\xi)\notin WF(K)\) for every nonzero \(\xi\). Then the kernel pairs with arbitrary compactly supported distributions in \(X\), defining \(P:\mathcal E'(X)\to\mathcal D'(Y)\). Excluding also \((y,\eta;x,0)\), \(\eta\ne0\), removes source-independent output singularities. No global support condition is needed for a compactly supported input and compactly supported output test function.}
\GSfield{Reconstruction}{The two zero-component exclusions have different roles; their joint sufficient condition in the notes is retained.}
\GSfield{Fate}{\citet[p.~130, discussion preceding Theorem~2.5.14]{hormander1971fourier} gives the kernel action on compactly supported distributions and precisely the zero-output-covector obstruction.}
\GSfield{Present status}{As of 9 November 2026: Established in the stated scope; no unresolved claim is inferred from this occurrence.}
\GSfield{Assessment}{literal statement: correct with missing hypotheses; publication match: exact; verification: read the result; development: annotated; confidence: high}
\end{GSitem}
\begin{GSitem}{1973-329.7}{E}{The primed canonical relation and its sign}
\GSfield{Kind; source}{definition; \GSpage{P000547}}
\GSfield{The notes}{The notebook orders the kernel coordinates as input \(x\), output \(y\), but reverses the output covector \(\eta\) in its definition of \(WF'(P)\).}
\GSfield{Statement}{With output-first kernel \(K(y,x)\), define \[C_P=\{(y,\eta;x,\xi):(y,\eta;x,-\xi)\in WF(K)\}.\] Using input-first ordering, the same relation is written \[\{(x,\xi;y,\eta):(x,-\xi;y,\eta)\in WF(K)\}.\] Thus it is the input covector that must change sign when the relation acts on \(WF(u)\) without another sign reversal.}
\GSfield{Reconstruction}{Both coordinate order and the sign convention are written explicitly; this is more than simply reversing the order of factors in a composition formula.}
\GSfield{Fate}{\citet[Eq.~(2.5.9) and Theorem~2.5.14, p.~130]{hormander1971fourier} fixes the input sign in the operator relation. A simultaneous sign reversal can conceal the discrepancy for real kernels, but not for arbitrary complex kernels.}
\GSfield{Present status}{As of 9 November 2026: Established in the stated scope; no unresolved claim is inferred from this occurrence.}
\GSfield{Assessment}{literal statement: false as written; publication match: exact; verification: read the result; development: annotated; confidence: high}
\end{GSitem}
\begin{GSitem}{1973-329.8}{E}{Wavefront propagation through a kernel}
\GSfield{Kind; source}{theorem; \GSpage{P000547}}
\GSfield{The notes}{\(WF(Pu)\subset WF'(P)\circ WF(u)\) under the zero-component exclusions.}
\GSfield{Statement}{Under both exclusions of the preceding extension item, for compactly supported \(u\), \(WF(Pu)\subset C_P\circ WF(u)\), where \(C_P\) uses the corrected input-sign convention. Without the exclusion of kernel points with input covector zero, an additional output-wavefront term is necessary.}
\GSfield{Reconstruction}{The notebook's stronger two-sided exclusion permits its short formula once the canonical-relation sign is repaired.}
\GSfield{Fate}{\citet[Theorem~2.5.14, pp.~130--131]{hormander1971fourier} proves the estimate under the stated zero-component exclusions. The preceding Theorem~2.5.12 gives the extra output term for smooth inputs when that exclusion is absent.}
\GSfield{Present status}{As of 9 November 2026: Established in the stated scope; no unresolved claim is inferred from this occurrence.}
\GSfield{Assessment}{literal statement: correct with missing hypotheses; publication match: exact; verification: read the result; development: annotated; confidence: high}
\end{GSitem}
\begin{GSitem}{1973-329.9}{E}{Composition of kernel relations}
\GSfield{Kind; source}{theorem; \GSpage{P000547}}
\GSfield{The notes}{For \(Q\circ P\), the notes write \(WF'(P)\circ WF'(Q)\) and explicitly define composition in input-first order.}
\GSfield{Statement}{Suppose \(P:X\to Y\) and \(Q:Y\to Z\) have kernels with the zero-component exclusions and enough proper support in the intermediate variable for the composition to be defined. Then, in output-first convention, \(C_{QP}\subset C_Q\circ C_P\). The notebook's reversed written order is consistent with its explicitly defined input-first convention and is not itself an error.}
\GSfield{Reconstruction}{Proper support and the convention conversion are stated. Abstractly mapping \(\mathcal E'\) into \(\mathcal E'\) does not replace all kernel support hypotheses without continuity and support control.}
\GSfield{Fate}{\citet[Theorem~2.5.15 and its concluding remark, pp.~132--133]{hormander1971fourier} gives kernel composition with support and wavefront hypotheses.}
\GSfield{Present status}{As of 9 November 2026: Established in the stated scope; no unresolved claim is inferred from this occurrence.}
\GSfield{Assessment}{literal statement: correct with missing hypotheses; publication match: exact; verification: read the result; development: annotated; confidence: high}
\end{GSitem}
\begin{GSitem}{1973-329.10}{E}{Distribution spaces with controlled wavefront}
\GSfield{Kind; source}{definition; \GSpage{P000548}}
\GSfield{The notes}{For a closed conic \(\Gamma\subset T^*\Omega\setminus0\), the notes write \(\mathcal D_\Gamma=\{u\in\mathcal D':WF(u)\subset\Gamma\}\), without a prime on the left.}
\GSfield{Statement}{Write \(\mathcal D'_\Gamma(\Omega)=\{u\in\mathcal D'(\Omega):WF(u)\subset\Gamma\}\). Hörmander's original convergence means weak distribution convergence together with rapid Fourier decay estimates converging to zero on every compactly localized closed cone disjoint from \(\Gamma\).}
\GSfield{Reconstruction}{The prime is a modern notational clarification; the defining set in the source already makes the intended space unambiguous.}
\GSfield{Fate}{\citet[p.~125]{hormander1971fourier} defines this sequential pseudotopology and proves smooth sequential approximation.}
\GSfield{Present status}{As of 9 November 2026: Established in the stated scope; no unresolved claim is inferred from this occurrence.}
\GSfield{Assessment}{literal statement: correct; publication match: exact; verification: read the result; development: annotated; confidence: high}
\end{GSitem}
\begin{GSitem}{1973-329.11}{E}{Multiplication under the no-opposite-covectors condition}
\GSfield{Kind; source}{theorem; \GSpage{P000548}}
\GSfield{The notes}{If \(\Gamma_1+\Gamma_2\) avoids zero covectors, multiplication extends uniquely and continuously into the cone union \(\Gamma_1\cup\Gamma_2\cup(\Gamma_1+\Gamma_2)\).}
\GSfield{Statement}{For closed conic sets with \(\Gamma_1\cap(-\Gamma_2)=\varnothing\), the product \(uv\) is canonically defined for \(u\in\mathcal D'_{\Gamma_1}\), \(v\in\mathcal D'_{\Gamma_2}\), and \(WF(uv)\subset\Gamma_1\cup\Gamma_2\cup(\Gamma_1+\Gamma_2)\). The extension agrees with smooth multiplication and is sequentially continuous for the original Hörmander convergence, which gives uniqueness by smoothing. 'Continuous' is not an unrestricted assertion of joint continuity for an unspecified locally convex topology.}
\GSfield{Reconstruction}{The continuity qualification is essential historical interpretation, not a newly discovered open defect in multiplication.}
\GSfield{Fate}{\citet[Theorem~2.5.10, pp.~126--127]{hormander1971fourier} is an exact 1971 match. \citet[Introduction and multiplication theorem]{brouderdanghelein2016continuity}, first circulated on arXiv in September 2014 and published in 2016, distinguishes sequential continuity from topological continuity and proves hypocontinuity in the normal topology.}
\GSfield{Present status}{As of 9 November 2026: Established in the stated scope; no unresolved claim is inferred from this occurrence.}
\GSfield{Assessment}{literal statement: correct with missing hypotheses; publication match: exact; verification: read the result; development: annotated; confidence: high}
\end{GSitem}
\begin{GSitem}{1973-329.12}{E}{Homogeneous phases and symbol classes}
\GSfield{Kind; source}{definition; \GSpage{P000548}}
\GSfield{The notes}{A phase is real, homogeneous of degree one in \(\Theta\), with total differential nonzero. Symbols obey \(|\partial_x^\alpha\partial_\Theta^\beta a|\le C_{\alpha\beta}(1+|\Theta|)^{m-\rho|\beta|+\delta|\alpha|}\).}
\GSfield{Statement}{Let \(\varphi\) be smooth on a conic subset of \(X\times(\mathbb R^N\setminus0)\), real, positively homogeneous of degree one, with \(d_{x,\Theta}\varphi\ne0\). For a sufficient standard class take \(0\le\delta<\rho\le1\); a smooth amplitude is in \(S^m_{\rho,\delta}\) when the displayed estimates hold uniformly on each compact \(x\)-set for all multiindices. Nondegeneracy of the critical manifold is an additional assumption needed for the usual nondegenerate FIO calculus; it is not part of the recorded phase definition.}
\GSfield{Reconstruction}{Smoothness, compact-set quantifiers, an admissible parameter range, and the distinction between a phase and a nondegenerate phase are supplied.}
\GSfield{Fate}{\citet[\S1.1, \S1.2, and phase definition]{hormander1971fourier} is the direct source.}
\GSfield{Present status}{As of 9 November 2026: Established in the stated scope; no unresolved claim is inferred from this occurrence.}
\GSfield{Assessment}{literal statement: correct with missing hypotheses; publication match: exact; verification: read the result; development: annotated; confidence: high}
\end{GSitem}
\begin{GSitem}{1973-329.13}{E}{Cutoff definition of an oscillatory distribution}
\GSfield{Kind; source}{construction; \GSpage{P000548}}
\GSfield{The notes}{\(\langle A,u\rangle=\lim_{\varepsilon\to0}\int e^{i\varphi}a\chi(\varepsilon\Theta)u\,dx\,d\Theta\), with compactly supported \(\chi=1\) near zero.}
\GSfield{Statement}{With a phase and symbol as above, local compact support in \(x\) after pairing with a test function, and the amplitude cut off near \(\Theta=0\) where necessary, integration by parts defines an oscillatory distribution. For \(\rho>0\) and \(\delta<1\) in the standard symbol setup, the cutoff limit exists distributionally and is independent of \(\chi\). To define an operator, use a kernel phase in both output and input variables; the formula displayed on the page itself defines a distribution.}
\GSfield{Reconstruction}{The distribution/operator distinction and low-frequency cutoff are supplied. No unrecorded canonical-relation theorem is added to this last formula.}
\GSfield{Fate}{\citet[Proposition~1.2.2, pp.~89--90]{hormander1971fourier} proves cutoff independence. The broader oscillatory-integral line continues 1965-84.20, without repeating its specific characteristic-set problem.}
\GSfield{Present status}{As of 9 November 2026: Established in the stated scope; no unresolved claim is inferred from this occurrence.}
\GSfield{Assessment}{literal statement: correct with missing hypotheses; publication match: exact; verification: read the result; development: annotated; confidence: high}
\end{GSitem}

\GStalk{330}{1974-03-18}{A. B. Shabat}{P000549-P000550}{Soliton collisions, rational interpolation, and the Miura transformation}
Shabat discusses one- and two-soliton waves, collision amplitude and shifts, and an interpolation construction of polynomial-exponential Schrödinger solutions. Several signs, a nonlinear coefficient, and the interpolation normalization are inconsistent in the notebook; an auxiliary operator is literally left empty. Shabat's 2014 paper explicitly recalls this 1974 Gelfand seminar, making it an unusually strong provenance match, while its precise theorem also exposes the missing assumptions.

\begin{GSitem}{1973-330.1}{E}{The one-soliton wave and its normalization}
\GSfield{Kind; source}{construction; \GSpage{P000549}}
\GSfield{The notes}{The equation is \(4f_t+f_{xxx}+ff_x=0\), followed by \(f=2\lambda^2\operatorname{sech}^2[\lambda(x-\lambda^2t-h)]\).}
\GSfield{Statement}{For \(\kappa>0\) and \(h\in\mathbb R\), \(q(t,x)=2\kappa^2\operatorname{sech}^2[\kappa(x-\kappa^2t-h)]\) solves \(4q_t+q_{xxx}+6qq_x=0\), has amplitude \(2\kappa^2\), and speed \(\kappa^2\). It does not solve the notebook's coefficient-one equation except in the zero solution limit.}
\GSfield{Reconstruction}{The coefficient 6 is an emendation of a faithful source error, not a newly supplied soliton theorem.}
\GSfield{Fate}{\citet[\S1, Eq.~(1.8)]{lax1968integrals} gives the classical solitary wave in another amplitude/time normalization. The Russian translation and its derivation were read. Direct substitution here independently identifies the missing factor 6.}
\GSfield{Present status}{As of 9 November 2026: Established in the stated scope; no unresolved claim is inferred from this occurrence.}
\GSfield{Assessment}{literal statement: false as written; publication match: exact; verification: read the result; development: annotated; confidence: high}
\end{GSitem}
\begin{GSitem}{1973-330.2}{E}{Two-soliton asymptotic separation}
\GSfield{Kind; source}{theorem; \GSpage{P000549}}
\GSfield{The notes}{A two-soliton solution is written schematically as two pulses plus \(O(e^{-\varkappa t})\) as \(t\to+\infty\).}
\GSfield{Statement}{For the reflectionless KdV two-soliton solution with distinct \(\kappa_1,\kappa_2>0\) and positive norming constants, the field separates as \(t\to\pm\infty\) into two solitary waves with speeds \(\kappa_j^2\), amplitudes \(2\kappa_j^2\), and asymptotic centers \(h_j^\pm\), with an exponentially small remainder for this finite pure-soliton family. This does not assert exponential decay of radiation for arbitrary decaying initial data.}
\GSfield{Reconstruction}{The finite reflectionless class, distinct positive parameters, and both time limits are explicit reconstruction; the notebook displays only the positive limit.}
\GSfield{Fate}{\citet[\S2]{lax1968integrals} studies two solitary waves before the seminar. The inverse-scattering setting also appears in \citet{shabat1973kdv}; that paper's general radiative asymptotics are not the exponential estimate claimed for this reflectionless special case.}
\GSfield{Present status}{As of 9 November 2026: Established in the stated scope; no unresolved claim is inferred from this occurrence.}
\GSfield{Assessment}{literal statement: correct with missing hypotheses; publication match: exact; verification: read the result; development: annotated; confidence: high}
\end{GSitem}
\begin{GSitem}{1973-330.3}{E}{Uniqueness of a prescribed pure two-soliton solution}
\GSfield{Kind; source}{claim; \GSpage{P000549}}
\GSfield{The notes}{The notes assert uniqueness: \ru{можно док., что оно единственно}.}
\GSfield{Statement}{Within the regular real reflectionless class, the two distinct discrete eigenvalues and their positive norming constants uniquely determine the potential and its KdV evolution. Equivalently, specifying the complete outgoing two-soliton scattering data determines this member of the pure-soliton family. Merely specifying two speeds, without positions/norming constants or excluding radiation, is insufficient.}
\GSfield{Reconstruction}{The data and solution class are supplied because the pronoun 'it' has no fully specified uniqueness quantifier in the notes.}
\GSfield{Fate}{\citet[\S3, reflectionless specialization]{shabat2014interpolation} revisits the finite-data construction and cites the earlier inverse-scattering framework. The underlying classical reflectionless inverse problem also occurred at 1970-219.3; no new general multi-soliton uniqueness theorem is attributed to this elliptical sentence.}
\GSfield{Present status}{As of 9 November 2026: Established in the stated scope; no unresolved claim is inferred from this occurrence.}
\GSfield{Assessment}{literal statement: correct with missing hypotheses; publication match: partial; verification: read the result; development: annotated; confidence: high}
\end{GSitem}
\begin{GSitem}{1973-330.4}{F}{The damaged collision-amplitude bound}
\GSfield{Kind; source}{problem; \GSpage{P000549}}
\GSfield{The notes}{Problem 1 asks for total amplitude at collision and writes \(|f(t,x)|\le2\sum_{j=0}^{n}\lambda_j(\lambda_{j+1}-\lambda_j)\), \(\lambda_0=0\), with an uncertain upper summation limit.}
\GSfield{Statement}{The general task is to bound a pure multi-soliton profile during interaction in terms of its discrete spectral parameters. The displayed proposed bound cannot be made a definite theorem from the surviving ink: the upper index is uncertain, parameter ordering is absent, and the indicated upper limit requires an undefined \(\lambda_{n+1}\).}
\GSfield{Reconstruction}{No plausible replacement sum is substituted for the damaged one, and no new problem-index entry is created without a definite inequality.}
\GSfield{Fate}{The primary papers \citet{lax1968integrals,shabat1973kdv,shabat2014interpolation} were examined for the collision and reflectionless context. None was verified to contain this precise bound with recoverable indexing. It remains a historical formula-identification gap, not evidence of an open mathematical problem.}
\GSfield{Present status}{As of 9 November 2026: The exact proposed estimate is unidentified; the notes alone do not establish that a corresponding modern extremal problem is open.}
\GSfield{Assessment}{literal statement: not determinable; publication match: none found; verification: unresolved lead; development: annotated; confidence: medium}
\end{GSitem}
\begin{GSitem}{1973-330.5}{E}{Collision-induced shifts of soliton centers}
\GSfield{Kind; source}{question; \GSpage{P000549}}
\GSfield{The notes}{Problem 2: the centers shift; calculate the displacement. The answer is described as standard scattering.}
\GSfield{Statement}{For a regular real two-soliton solution with \(\kappa_1>\kappa_2>0\), compare the intercepts of the incoming and outgoing center trajectories. They satisfy \(\Delta h_1=\kappa_1^{-1}\log[(\kappa_1+\kappa_2)/(\kappa_1-\kappa_2)]\) and \(\Delta h_2=-\kappa_2^{-1}\log[(\kappa_1+\kappa_2)/(\kappa_1-\kappa_2)]\). Thus the larger soliton advances and the smaller one is delayed. In the regular finite \(N\)-soliton family, shifts add pairwise.}
\GSfield{Reconstruction}{The notebook gives no formula. The reflectionless hypothesis, positive ordered wave parameters, and the definition outgoing minus incoming are supplied.}
\GSfield{Fate}{\citet[\S2]{lax1968integrals} is the pre-seminar two-wave source. The precise normalized pairwise formula is independently visible in \citet[Eqs.~(2.4)--(2.5)]{bonnemaindoyonel2022hydrodynamics}; that modern restatement is not claimed as first publication.}
\GSfield{Present status}{As of 9 November 2026: The classical inverse-scattering phase shifts answer the recorded question.}
\GSfield{Assessment}{literal statement: correct with missing hypotheses; publication match: exact; verification: read the result; development: annotated; confidence: high; problem \hyperref[prob:GP-1973-54]{GP-1973-54}}
\end{GSitem}
\begin{GSitem}{1973-330.6}{E}{Bargmann potentials and polynomial-exponential solutions}
\GSfield{Kind; source}{definition; \GSpage{P000549}}
\GSfield{The notes}{A Bargmann potential is characterized by two solutions of \(y_{xx}+(q+\zeta^2)y=0\) of the form \(P_{1,2}(x,\zeta)e^{\mp i\zeta x}\), polynomial in \(\zeta\).}
\GSfield{Statement}{For a regular real decaying finite reflectionless potential, suitable spectral normalizations of the Jost solutions give polynomial-exponential functions of \(\zeta\); they are independent at generic spectral parameters. They cease to be independent at the discrete spectral points, and real regularity is a separate restriction on the data. The notebook's algebraic property may also describe complex or singular potentials if those restrictions are omitted.}
\GSfield{Reconstruction}{Generic spectral independence, permissible spectral rescaling, and the distinction between formal algebraic and regular real scattering potentials are supplied.}
\GSfield{Fate}{\citet[\S3, especially the Schrödinger specialization]{shabat2014interpolation} identifies this construction and references Shabat's 1970 work on zero-reflection potentials. The general classical reflectionless line is already present at 1970-219.3. The 1970 original was not opened and is not treated as a verified exact interpolation source.}
\GSfield{Present status}{As of 9 November 2026: Established in the stated scope; no unresolved claim is inferred from this occurrence.}
\GSfield{Assessment}{literal statement: correct with missing hypotheses; publication match: partial; verification: read the result; development: annotated; confidence: high}
\end{GSitem}
\begin{GSitem}{1973-330.7}{F}{Rational interpolation for the auxiliary solutions}
\GSfield{Kind; source}{construction; \GSpage{P000550}}
\GSfield{The notes}{\(P/Q=c_j e^{2i\zeta_jx}\) at the interpolation nodes and \(P/Q=1\) at infinity. The number of nodes is written \(2n\,?\,1\), with an overwritten sign. The notes claim arbitrary parameters produce a fundamental system of a scalar Sturm--Liouville equation.}
\GSfield{Statement}{A precise related theorem uses \(P,Q\) of degrees \(m,n\), \(N=m+n\), and \(N+1\) distinct interpolation nodes \(k_j\), with \(P(k_j)/Q(k_j)=c_j e^{-2k_jx}\). On an interval where the interpolation system is nondegenerate, normalize the leading coefficients by \(p_0q_0=1\). Then \(e^{kx}P\) and \(e^{-kx}Q\) solve a common equation \(\psi_{xx}=(k^2+ku_1+u_2)\psi\). Obtaining an ordinary Schrödinger potential with \(u_1=0\), and a globally regular real potential, requires additional reduction and positivity/reality conditions.}
\GSfield{Reconstruction}{This is a separately labeled modern reconstruction from the 2014 theorem, with \(k=i\zeta\) up to swapping the two solutions; it is not a decipherment of the obscured node-count sign.}
\GSfield{Fate}{\citet[Theorem~1, \S3, pp.~313--315 of the Russian original]{shabat2014interpolation} is a close mathematical match and explicitly mentions the 1974 Gelfand seminar on p.~305. It appeared forty years later, not 'soon after'. The source's unrestricted scalar assertion is therefore classified F: the missing reduction prevents an exact theorem identification, despite unusually strong provenance.}
\GSfield{Present status}{As of 9 November 2026: The nondegenerate rational-interpolation construction is published; the notebook omits restrictions needed for its stronger scalar and arbitrary-parameter wording.}
\GSfield{Assessment}{literal statement: false as written; publication match: partial; verification: read the result; development: annotated; confidence: high}
\end{GSitem}
\begin{GSitem}{1973-330.8}{E}{Polynomial Wronskian and the missing normalization}
\GSfield{Kind; source}{claim; \GSpage{P000550}}
\GSfield{The notes}{\(P_xQ-Q_xP-2i\zeta PQ=\prod(\zeta-\zeta_j)\). The preceding line has \(y_1=Pe^{+i\zeta x}\), \(y_2=Qe^{-i\zeta x}\).}
\GSfield{Statement}{If the ratio condition has exponent \(+2i\zeta_jx\), the polynomial \(W_0=P_xQ-Q_xP-2i\zeta PQ\) vanishes at each node. Under the matching degree and leading-coefficient hypotheses, \(W_0=-2i\,p_0q_0\prod_j(\zeta-\zeta_j)\). It is the Wronskian \(y_1'y_2-y_1y_2'\) for \(y_1=Pe^{-i\zeta x}\), \(y_2=Qe^{+i\zeta x}\), with these corrected exponents. In the 2014 normalization the analogous constant is \(-2\).}
\GSfield{Reconstruction}{The degree restriction and leading factor are restored. No independence is asserted at zeros of this Wronskian.}
\GSfield{Fate}{\citet[Theorem~1 and its proof]{shabat2014interpolation} proves the polynomial identity with its leading coefficient. The vanishing follows directly by differentiating the interpolation data; the notebook omitted the normalization and reversed the compatible exponential assignment.}
\GSfield{Present status}{As of 9 November 2026: Established in the stated scope; no unresolved claim is inferred from this occurrence.}
\GSfield{Assessment}{literal statement: false as written; publication match: partial; verification: read the result; development: annotated; confidence: high}
\end{GSitem}
\begin{GSitem}{1973-330.9}{E}{Cubic time evolution of reflectionless data}
\GSfield{Kind; source}{method; \GSpage{P000550}}
\GSfield{The notes}{\(c_j(t)=c_j(0)e^{2i\zeta^3t}\), substituted into the same interpolation relation.}
\GSfield{Statement}{At each fixed node, take \(c_j(t)=c_j(0)e^{2i\zeta_j^3t}\). In a compatible regular reflectionless reduction, the cubic time factor is the KdV scattering evolution after the notebook's time rescaling. For \(\zeta_j=i\kappa_j\) the combined factor is \(e^{-2\kappa_jx+2\kappa_j^3t}\), giving center speed \(\kappa_j^2\). Arbitrary unconstrained complex interpolation parameters need not give a globally nonsingular real solution.}
\GSfield{Reconstruction}{The node subscript is restored on the exponent, and the real regular reduction is stated separately.}
\GSfield{Fate}{\citet[\S3]{shabat2014interpolation} revisits interpolation and soliton dynamics, while \citet{shabat1973kdv} gives the earlier inverse-scattering evolution. This is the known cubic-flow mechanism, not a proof of the notebook's unrestricted construction.}
\GSfield{Present status}{As of 9 November 2026: Established in the stated scope; no unresolved claim is inferred from this occurrence.}
\GSfield{Assessment}{literal statement: correct with missing hypotheses; publication match: partial; verification: read the result; development: annotated; confidence: high}
\end{GSitem}
\begin{GSitem}{1973-330.10}{F}{The incomplete auxiliary time operator}
\GSfield{Kind; source}{claim; \GSpage{P000550}}
\GSfield{The notes}{\((4\partial_t+\partial_x^3-3q\partial_x)P=(\qquad)Q=0\); the second operator is left empty.}
\GSfield{Statement}{No unique pair of auxiliary equations for the polynomial factors \(P,Q\) can be recovered. Conjugating a differential time operator on full eigenfunctions by \(e^{\pm i\zeta x}\) generally introduces \(\zeta\)-dependent terms; these cannot simply be suppressed when the same letter denotes only the polynomial factor.}
\GSfield{Reconstruction}{The blank is preserved, and no fabricated equation is attributed to Shabat.}
\GSfield{Fate}{The compatible Lax and interpolation frameworks in \citet{lax1968integrals,shabat2014interpolation} are thematic guides, but do not authorize filling the literal empty parentheses with a guessed operator.}
\GSfield{Present status}{As of 9 November 2026: The exact missing operator and gauge are not identifiable from the surviving source.}
\GSfield{Assessment}{literal statement: not determinable; publication match: thematic only; verification: unresolved lead; development: annotated; confidence: high}
\end{GSitem}
\begin{GSitem}{1973-330.11}{E}{Riccati logarithmic derivatives and the Miura map}
\GSfield{Kind; source}{construction; \GSpage{P000550}}
\GSfield{The notes}{\(f=\partial_x\log(Pe^{?i\zeta x})\), \(4f_t+f_{xxx}+3(f^2+\zeta^2)f_x=0\), and \(f_x+f^2=q\) are said to give \(4q_t+q_{xxx}+6qq_x=0\).}
\GSfield{Statement}{For the recorded Schrödinger equation, any nonvanishing local solution has logarithmic derivative \(v=y_x/y\) satisfying \(v_x+v^2=-q-\zeta^2\). Independently, a consistent Miura normalization is: if \(v_t-6v^2v_x+v_{xxx}=0\), then \(u=v_x+v^2\) satisfies \(u_t-6uu_x+u_{xxx}=0\). Indeed its residual equals \((\partial_x+2v)(v_t-6v^2v_x+v_{xxx})\). Time rescaling and a constant spectral shift give equivalent conventions, but do not make all displayed notebook signs and coefficients mutually consistent.}
\GSfield{Reconstruction}{The consistent Miura pair is reconstructed in new symbols \(u,v\), not passed off as recovered ink. The original logarithm's exponent sign remains undeciphered.}
\GSfield{Fate}{\citet{miura1968transformation} is the 1968 historical source; its primary PDF could not be opened, so the bibliographic entry is explicitly unverified. The algebraic identity is proved directly by expansion and checked independently in the supplied code. \citet[\S1]{lax1968integrals} supplies the earlier Riccati/conservation-law context.}
\GSfield{Present status}{As of 9 November 2026: Established in the stated scope; no unresolved claim is inferred from this occurrence.}
\GSfield{Assessment}{literal statement: false as written; publication match: partial; verification: indirect; development: annotated; confidence: high}
\end{GSitem}

\GStalk{331}{1974-03-25}{A. A. Kirillov}{P000551}{Infinite-dimensional groups and symmetric Fock space}
This source-recovered supplementary entry has its own Kirillov heading on P000551. The official MCCME archive identifies record 331 but incorrectly dates it 25 March 1975; the folio clearly reads 25/III-74 between 18 March and 1 April 1974. Kirillov names diffeomorphism/current groups and bundle automorphisms, then sketches symmetric Fock space. No particular representation theorem survives.

\begin{GSitem}{1973-331.1}{F}{Representations of diffeomorphism and current groups}
\GSfield{Kind; source}{remark; \GSpage{P000551}}
\GSfield{The notes}{Kirillov announces descriptions of representations of \(G_1=\operatorname{Diff}M\) and \(G_2=H^M=C^\infty(M,H)\), with \(H\) a group.}
\GSfield{Statement}{For a smooth manifold \(M\) and a finite-dimensional Lie group \(H\), the smooth current group is \(C^\infty(M,H)\) with pointwise multiplication. A representation-classification problem requires a topology, representation category, regularity, and equivalence notion. None is specified on the page, and no classification assertion can be recovered.}
\GSfield{Reconstruction}{A Lie group and smooth-category hypotheses make the displayed current-group notation precise; no intended representation class is invented.}
\GSfield{Fate}{\citet[\S2]{araki1970factorizable} is a verified earlier source for continuous tensor-product and current-group representations, but its measurable-current framework is not an exact match for an absent Kirillov theorem. Related earlier material occurs at 1970-228.5 and 1972-291.6--291.9.}
\GSfield{Present status}{As of 9 November 2026: The announced representation description is unspecified; it does not define a new open classification problem.}
\GSfield{Assessment}{literal statement: not determinable; publication match: thematic only; verification: unresolved lead; development: annotated; confidence: high}
\end{GSitem}
\begin{GSitem}{1973-331.2}{E}{The bundle-automorphism group extension}
\GSfield{Kind; source}{claim; \GSpage{P000551}}
\GSfield{The notes}{For a principal bundle \(E\to M\) with structure group \(H\), its automorphism group projects onto \(G_1\) with fiber \(G_2\).}
\GSfield{Statement}{For a smooth principal \(H\)-bundle \(E\to M\), the exact group sequence is \(1\to\operatorname{Gau}(E)\to\operatorname{Aut}_H(E)\to\operatorname{Diff}(M)_{[E]}\to1\). The image consists of diffeomorphisms whose pullback bundle is isomorphic to \(E\). The gauge group is the smooth sections of \(E\times_{\operatorname{Ad}}H\), not generally \(C^\infty(M,H)\). For a chosen trivialization of a trivial bundle, all diffeomorphisms lift and the extension is a semidirect product with the ordinary current group.}
\GSfield{Reconstruction}{The subgroup preserving the bundle isomorphism class and the adjoint-bundle gauge group replace the unrestricted base and fiber. The Lie-group topology theorem is not imported into the 1974 claim.}
\GSfield{Fate}{\citet[introduction and bundle-automorphism extension theorem]{wockel2007symmetry} states the precise modern extension, with its additional Lie-group hypotheses. The elementary group-level distinction is a correction of the notebook's omitted trivial-bundle assumption, not an attribution of Wockel's 2007 theorem to Kirillov.}
\GSfield{Present status}{As of 9 November 2026: Established in the stated scope; no unresolved claim is inferred from this occurrence.}
\GSfield{Assessment}{literal statement: false as written; publication match: partial; verification: read the result; development: annotated; confidence: high}
\end{GSitem}
\begin{GSitem}{1973-331.3}{E}{Symmetric Fock space and the vacuum}
\GSfield{Kind; source}{construction; \GSpage{P000551}}
\GSfield{The notes}{\(e^H=H_0\oplus H\oplus S^2(H)\oplus\cdots\), with one-dimensional \(H_0\).}
\GSfield{Statement}{For a complex Hilbert space \(\mathcal H\), symmetric Fock space is the Hilbert direct sum \(\Gamma_s(\mathcal H)=\bigoplus_{n=0}^\infty\mathcal H^{\widehat\otimes_s n}\), where the zero-particle summand is \(\mathbb C\). Its vacuum is \((1,0,0,\ldots)\). The letter \(\mathcal H\) is separate from the structure group denoted \(H\) earlier on the page.}
\GSfield{Reconstruction}{Hilbert completion and the different meanings of \(H\) are clarified.}
\GSfield{Fate}{\citet[\S4, pp.~370--373]{araki1970factorizable} gives the construction in Fock space and exponential vectors before the seminar. The Hilbert direct sum and its normalization are standard explanatory background. Related constructions occur at 1969-180.7--180.8 and 1970-228.5; these are conceptual continuations, not repeated problem IDs.}
\GSfield{Present status}{As of 9 November 2026: Established in the stated scope; no unresolved claim is inferred from this occurrence.}
\GSfield{Assessment}{literal statement: correct with missing hypotheses; publication match: partial; verification: read the result; development: annotated; confidence: high}
\end{GSitem}
\begin{GSitem}{1973-331.4}{E}{Exponential vectors and factorial normalization}
\GSfield{Kind; source}{construction; \GSpage{P000551}}
\GSfield{The notes}{The row of components begins \(1,\xi,\xi\otimes\xi/2!,\ldots\), next to the norm formula \(|\xi|^2=\sum|\xi_n|^2/n!\).}
\GSfield{Statement}{With the ordinary Hilbert norm on each symmetric tensor power, the exponential vector is \(\operatorname{Exp}(\xi)=\bigoplus_{n\ge0}\xi^{\otimes n}/\sqrt{n!}\), and its squared norm is \(e^{\|\xi\|^2}\). Alternatively, unscaled components \(\xi^{\otimes n}\) may be given weight \(1/n!\), or components \(\xi^{\otimes n}/n!\) weight \(n!\). Mixing these conventions without identifying the components changes the norm.}
\GSfield{Reconstruction}{One consistent standard normalization is explicitly chosen. The source's component symbol is not silently redefined.}
\GSfield{Fate}{\citet[\S4, exponential vectors and their unitary action, pp.~371--373]{araki1970factorizable} is the earlier primary framework. The notebook's two rows require a component convention; their proximity does not establish which convention was intended. The norm of the separately chosen standard exponential vector follows directly from the exponential series.}
\GSfield{Present status}{As of 9 November 2026: Established in the stated scope; no unresolved claim is inferred from this occurrence.}
\GSfield{Assessment}{literal statement: not determinable; publication match: partial; verification: read the result; development: annotated; confidence: high}
\end{GSitem}

\GStalk{332}{1974-03-25}{B. M. Gurevich}{P000551}{KdV by Whitham modulation}
A second, separate heading on the same 25 March 1974 folio reads Gurevich, on the KdV equation by Whitham's method. The official archive identifies record 332 and B. M. Gurevich but shares record 331's erroneous 1975 date. No mathematical development follows the title. The famous A. V. Gurevich--Pitaevskii paper is a different author identity and supplies only thematic context.

\begin{GSitem}{1973-332.1}{F}{The unrecorded Whitham--KdV exposition}
\GSfield{Kind; source}{remark; \GSpage{P000551}}
\GSfield{The notes}{\ru{Гуревич. Об ур-} [overwritten ending] \ru{К.-д.-Ф. методом Уидема}.}
\GSfield{Statement}{The source identifies Whitham averaging applied to KdV but supplies no modulation system, solution class, initial data, dispersive-shock boundary conditions, or theorem. No exact assertion can be reconstructed.}
\GSfield{Reconstruction}{The official record identity is supplied; no famous paper is substituted for the absent lecture content.}
\GSfield{Fate}{\citet{whitham1970twotiming} is established background. \citet{gurevichpitaevskii1974shock} was submitted 3 April 1973, published in Russian in August 1973 and English in February 1974; it is relevant to the subject but authored by A. V. Gurevich, whereas record 332 identifies B. M. Gurevich. Neither authorship nor an exact result is conflated.}
\GSfield{Present status}{As of 9 November 2026: The mathematical content beyond the title remains unrecoverable.}
\GSfield{Assessment}{literal statement: not determinable; publication match: thematic only; verification: unresolved lead; development: annotated; confidence: high}
\end{GSitem}

\GStalk{333}{1974-04-01}{Yu. Koblenz}{P000552-P000553}{Heat invariants, closed trajectories, and asymptotic spectral multiplicity}
Koblenz reports heat-trace coefficients through the fourth term and their curvature interpretation. Gelfand inserts an unresolved lattice-point analogy; Shnirelman introduces closed billiard paths. The discussion passes from heat to wave traces and to near-degenerate eigenvalue pairs. The all-indices two-torus statement has an exact primary match in Shnirelman's 1975 note, received 7 May 1974. Its square-root spectral parameter must be distinguished from the Laplacian eigenvalues at the start of the talk.

\begin{GSitem}{1973-333.1}{E}{The Laplace operator and its sign convention}
\GSfield{Kind; source}{definition; \GSpage{P000552}}
\GSfield{The notes}{On a closed Riemannian manifold, \(0=\lambda_0\le\lambda_1\le\cdots\), \(Z(s)=\operatorname{Tr}e^{-s\Delta}\), but the displayed coordinate formula is \(+\operatorname{div}\operatorname{grad}\).}
\GSfield{Statement}{Let \(M\) be a smooth compact Riemannian manifold without boundary. Put \(L=-\operatorname{div}_g\operatorname{grad}_g\ge0\), the self-adjoint scalar Laplacian on \(L^2(M,dV_g)\). Its eigenvalues \(0=\mu_0\le\mu_1\le\cdots\) are listed with multiplicities and tend to infinity. The heat trace is \(\operatorname{Tr}e^{-tL}\) for \(t>0\). With the opposite differential sign one must instead use \(e^{t\Delta}\).}
\GSfield{Reconstruction}{Nonnegative \(L\), eigenvalues \(\mu_j\), and later frequencies \(\omega_j=\sqrt{\mu_j}\) are distinguished. A minus possibly hidden inside the overwritten block is not claimed as recovered ink.}
\GSfield{Fate}{\citet[opening heat expansion and coefficient discussion]{mckeansinger1967curvature} is a pre-seminar source; its 1969 Russian translation, pp.~138--141, was read. It uses a nonpositive Laplacian, so its signs are translated consistently here.}
\GSfield{Present status}{As of 9 November 2026: Established in the stated scope; no unresolved claim is inferred from this occurrence.}
\GSfield{Assessment}{literal statement: false as written; publication match: exact; verification: read the result; development: annotated; confidence: high}
\end{GSitem}
\begin{GSitem}{1973-333.2}{E}{The heat kernel eigenfunction expansion and trace}
\GSfield{Kind; source}{construction; \GSpage{P000552}}
\GSfield{The notes}{\(E(m,m',t)=\sum e^{-\lambda_i t}\varphi_i(m)\varphi_i(m')\), and \(Z(t)=\int_M E(m,m,t)\,dm\).}
\GSfield{Statement}{For an orthonormal eigenbasis of \(L\), \(K_t(x,y)=\sum_j e^{-t\mu_j}\varphi_j(x)\overline{\varphi_j(y)}\). The series is smooth for \(t>0\); integrating its diagonal gives \(\operatorname{Tr}e^{-tL}=\sum_j e^{-t\mu_j}\). The conjugate can be omitted if a real orthonormal eigenbasis is chosen.}
\GSfield{Reconstruction}{The Riemannian measure, orthonormality, complex-conjugation convention, and positive time are supplied.}
\GSfield{Fate}{\citet[opening heat expansion and coefficient discussion]{mckeansinger1967curvature} is a pre-seminar source; its 1969 Russian translation, pp.~138--141, was read. It uses a nonpositive Laplacian, so its signs are translated consistently here.}
\GSfield{Present status}{As of 9 November 2026: Established in the stated scope; no unresolved claim is inferred from this occurrence.}
\GSfield{Assessment}{literal statement: correct with missing hypotheses; publication match: exact; verification: read the result; development: annotated; confidence: high}
\end{GSitem}
\begin{GSitem}{1973-333.3}{E}{The small-time heat-trace expansion}
\GSfield{Kind; source}{theorem; \GSpage{P000552}}
\GSfield{The notes}{Minakshisundaram, 1949: \(Z(t)=(4\pi t)^{-n/2}(a_0+a_1t+a_2t^2+\cdots)\), \(a_0=\operatorname{vol}M\), with \(\operatorname{Re}t>0\).}
\GSfield{Statement}{For the scalar Laplacian on a smooth closed \(n\)-manifold, as real \(t\downarrow0\), \(\operatorname{Tr}e^{-tL}\sim(4\pi t)^{-n/2}\sum_{r=0}^{\infty}a_rt^r\), with \(a_0=\operatorname{vol}_g(M)\). This is an asymptotic series, not generally a convergent power series. Boundary problems require additional boundary coefficients, including half-integer powers in the usual smooth case.}
\GSfield{Reconstruction}{The asymptotic meaning and positive-real approach are made explicit; no uniform assertion over the whole complex right half-plane is inferred.}
\GSfield{Fate}{\citet[opening heat expansion and coefficient discussion]{mckeansinger1967curvature} is a pre-seminar source; its 1969 Russian translation, pp.~138--141, was read. It uses a nonpositive Laplacian, so its signs are translated consistently here. The historical Minakshisundaram--Pleijel construction predates the talk; the notebook's abbreviated attribution is not treated as an assertion that only one author contributed.}
\GSfield{Present status}{As of 9 November 2026: Established in the stated scope; no unresolved claim is inferred from this occurrence.}
\GSfield{Assessment}{literal statement: correct with missing hypotheses; publication match: exact; verification: read the result; development: annotated; confidence: high}
\end{GSitem}
\begin{GSitem}{1973-333.4}{E}{Curvature contractions in the heat coefficients}
\GSfield{Kind; source}{definition; \GSpage{P000552}}
\GSfield{The notes}{The page defines scalar curvature, Ricci curvature, and \(|R|^2\), with overwritten indices; it says the sum of squared components is in normal coordinates.}
\GSfield{Statement}{Use the curvature convention in which the round sphere has positive scalar curvature \(R\). Define \(\operatorname{Ric}\) and \(R\) by metric contraction of the Riemann tensor. The invariant squared norms are \(|\operatorname{Riem}|^2=R_{ijkl}R^{ijkl}\) and \(|\operatorname{Ric}|^2=R_{ij}R^{ij}\); at a point in an orthonormal frame these equal sums of squares of components. The coordinate Euclidean sums are valid at the origin of normal coordinates, not throughout a normal chart without the metric contractions.}
\GSfield{Reconstruction}{The badly overwritten scalar/Ricci index placements are replaced by invariant contractions without purporting to decipher them.}
\GSfield{Fate}{\citet[opening heat expansion and coefficient discussion]{mckeansinger1967curvature} is a pre-seminar source; its 1969 Russian translation, pp.~138--141, was read. It uses a nonpositive Laplacian, so its signs are translated consistently here.}
\GSfield{Present status}{As of 9 November 2026: Established in the stated scope; no unresolved claim is inferred from this occurrence.}
\GSfield{Assessment}{literal statement: correct with missing hypotheses; publication match: exact; verification: read the result; development: annotated; confidence: high}
\end{GSitem}
\begin{GSitem}{1973-333.5}{E}{The integrated scalar-curvature coefficient}
\GSfield{Kind; source}{theorem; \GSpage{P000552}}
\GSfield{The notes}{\(a_1=\int_M\tau\,dV/6\).}
\GSfield{Statement}{For the preceding nonnegative scalar Laplacian on a smooth closed manifold, \(a_1=\frac16\int_M R\,dV_g\). The operator and curvature sign conventions are those already fixed.}
\GSfield{Fate}{\citet[opening heat expansion and coefficient discussion]{mckeansinger1967curvature} is a pre-seminar source; its 1969 Russian translation, pp.~138--141, was read. It uses a nonpositive Laplacian, so its signs are translated consistently here.}
\GSfield{Present status}{As of 9 November 2026: Established in the stated scope; no unresolved claim is inferred from this occurrence.}
\GSfield{Assessment}{literal statement: correct; publication match: exact; verification: read the result; development: annotated; confidence: high}
\end{GSitem}
\begin{GSitem}{1973-333.6}{E}{The quadratic curvature heat coefficient}
\GSfield{Kind; source}{theorem; \GSpage{P000552}}
\GSfield{The notes}{\(a_2=\frac1{360}\int(2|R|^2-2|\rho|^2+5\tau^2)\), without an explicit measure in this display.}
\GSfield{Statement}{For the smooth closed scalar case, \(a_2=\frac1{360}\int_M(2|\operatorname{Riem}|^2-2|\operatorname{Ric}|^2+5R^2)\,dV_g\). A local diagonal coefficient may also contain a Laplacian of \(R\); its integral vanishes on a closed manifold. Thus the displayed formula is an integrated heat invariant, not the complete pointwise coefficient.}
\GSfield{Reconstruction}{The measure and distinction between local and integrated coefficients are supplied.}
\GSfield{Fate}{\citet[opening heat expansion and coefficient discussion]{mckeansinger1967curvature} is a pre-seminar source; its 1969 Russian translation, pp.~138--141, was read. It uses a nonpositive Laplacian, so its signs are translated consistently here.}
\GSfield{Present status}{As of 9 November 2026: Established in the stated scope; no unresolved claim is inferred from this occurrence.}
\GSfield{Assessment}{literal statement: correct with missing hypotheses; publication match: exact; verification: read the result; development: annotated; confidence: high}
\end{GSitem}
\begin{GSitem}{1973-333.7}{E}{The off-diagonal heat parametrix}
\GSfield{Kind; source}{construction; \GSpage{P000552}}
\GSfield{The notes}{Near the diagonal, \(E(t,m,m')=(4\pi t)^{-n/2}(u_0+u_1t+\cdots)\), with no Gaussian factor.}
\GSfield{Statement}{In a geodesically convex neighborhood of the diagonal, \(K_t(x,y)\sim(4\pi t)^{-n/2}e^{-d_g(x,y)^2/(4t)}\sum_{r\ge0}u_r(x,y)t^r\), with smooth transport coefficients and \(u_0(x,x)=1\). On the diagonal the Gaussian is one. A formula on a neighborhood of the diagonal must include it; global continuation through cut/conjugate loci is a separate issue.}
\GSfield{Reconstruction}{The Gaussian and local geometric scope are restored; the recorded display is valid as a diagonal schematic expansion only.}
\GSfield{Fate}{\citet[opening heat expansion and coefficient discussion]{mckeansinger1967curvature} is a pre-seminar source; its 1969 Russian translation, pp.~138--141, was read. It uses a nonpositive Laplacian, so its signs are translated consistently here.}
\GSfield{Present status}{As of 9 November 2026: Established in the stated scope; no unresolved claim is inferred from this occurrence.}
\GSfield{Assessment}{literal statement: false as written; publication match: exact; verification: read the result; development: annotated; confidence: high}
\end{GSitem}
\begin{GSitem}{1973-333.8}{E}{Sakai and the fourth heat coefficient}
\GSfield{Kind; source}{claim; \GSpage{P000552}}
\GSfield{The notes}{\(a_3\) is attributed to Sakai and the speaker independently; no formula is written.}
\GSfield{Statement}{The coefficient of \(t^3\) inside the normalized heat-trace series is a universal integrated combination of curvature and its covariant derivatives. The notes provide no particular expression to compare term by term. Sakai's computation is identifiable, but an exact formula is not reconstructed from the attribution alone.}
\GSfield{Reconstruction}{The index \(a_3\) is interpreted in the source's normalized convention, rather than the \(a_6\) convention also common in later literature.}
\GSfield{Fate}{\citet{sakai1971eigenvalues} appeared in 1971, before this seminar. Its publisher issue record was opened; \citet[Introduction, p.~201]{gilkey1979recursion}, whose full text was read, explicitly credits Sakai with computing the fourth coefficient. This verifies the historical subject match, not an unread expression from Sakai's paper.}
\GSfield{Present status}{As of 9 November 2026: Established in the stated scope; no unresolved claim is inferred from this occurrence.}
\GSfield{Assessment}{literal statement: correct with missing hypotheses; publication match: partial; verification: publisher metadata checked; development: annotated; confidence: high}
\end{GSitem}
\begin{GSitem}{1973-333.9}{F}{Koblenz's claimed independent fourth-coefficient computation}
\GSfield{Kind; source}{remark; \GSpage{P000552}}
\GSfield{The notes}{The same line says \ru{и докладчик (независимо)} after Sakai.}
\GSfield{Statement}{The source records an independent computation by the speaker, but neither its formula nor a publication title, preprint, or full author identity is supplied. It cannot be compared with Sakai's expression or assigned a publication priority from these pages alone.}
\GSfield{Reconstruction}{The speaker's informal source name is retained without guessing an expanded identity.}
\GSfield{Fate}{Searches for Koblenz/Koblents and the Russian \ru{Кобленц} with Laplace asymptotics and heat coefficients did not identify a primary publication matching this claim. \citet{sakai1971eigenvalues,gilkey1979recursion} establish the surrounding known result, not the claimed independent attribution.}
\GSfield{Present status}{As of 9 November 2026: The independent-attribution claim remains unverified; the coefficient itself is established.}
\GSfield{Assessment}{literal statement: not determinable; publication match: none found; verification: unresolved lead; development: annotated; confidence: medium}
\end{GSitem}
\begin{GSitem}{1973-333.10}{F}{Gelfand's lattice-point and zeta-zero analogy}
\GSfield{Kind; source}{remark; \GSpage{P000552}; I. M. Gelfand}
\GSfield{The notes}{Gelfand writes a disk lattice-point count as \(\pi r^2+\sum_{\lambda_n}\tau^{1/2}\sin(\tau/\lambda_n)\), identifies \(\lambda_n\) as Riemann-zeta zeros, and adds 'mod Riemann hypothesis'.}
\GSfield{Statement}{The variables, summation prescription, powers, and spectral parameters are too incomplete to define a convergent or distributionally meaningful identity. The disk lattice-point problem has classical Poisson-summation/Bessel expansions, while explicit formulas involving Riemann-zeta zeros concern other counting functions. No exact bridge or Riemann-hypothesis consequence is identifiable in this displayed analogy.}
\GSfield{Fate}{The Russian text and scan confirm the explicit zeta-zero parenthesis; it cannot be silently replaced by dual-lattice radii to manufacture a valid circle formula. The nearby heat/spectral literature is only thematic context \citep{mckeansinger1967curvature}. No precise published match was found.}
\GSfield{Present status}{As of 9 November 2026: The intended identity remains unidentified; it is not entered as a new version of the Riemann hypothesis.}
\GSfield{Assessment}{literal statement: not determinable; publication match: none found; verification: unresolved lead; development: annotated; confidence: high}
\end{GSitem}
\begin{GSitem}{1973-333.11}{E}{Closed billiard paths and oscillations in spectral density}
\GSfield{Kind; source}{claim; \GSpage{P000552-P000553}; A. I. Shnirelman}
\GSfield{The notes}{Under an explicit Shnirelman heading, the ball is associated with lengths of closed billiard paths; the smoothed density is written \(f_0(\lambda)+\sum_kc_k\lambda^{\alpha_k}\sin(\ell_k\lambda+\omega_k)\), with Balian--Bloch numerical/experimental confirmation.}
\GSfield{Statement}{The identifiable content is a periodic-orbit expansion: after appropriate spectral smoothing, closed classical paths contribute oscillations whose phases involve their lengths times the spectral frequency. A rigorous expansion requires a specified domain, boundary condition, smoothing, and treatment of degenerate periodic families. The schematic infinite sum is not an ordinary convergent equality for every ball spectrum, and path lengths are not themselves Laplacian eigenvalues.}
\GSfield{Reconstruction}{A qualitative oscillatory-trace statement is retained; the notebook supplies too few coefficients or boundary data for an exact expansion.}
\GSfield{Fate}{\citet[abstract and sphere examples]{balianbloch1972oscillations} concerns precisely the smoothed-density oscillations and closed reflected paths, with numerical comparisons for the sphere; its primary publisher search record was located, but direct access returned 403. \citet[Introduction]{colindeverdiere1973geodesicsii}, published in 1973, supplies a rigorous related closed-manifold trace framework, not the same boundary problem.}
\GSfield{Present status}{As of 9 November 2026: Established in the stated scope; no unresolved claim is inferred from this occurrence.}
\GSfield{Assessment}{literal statement: correct with missing hypotheses; publication match: partial; verification: indirect; development: annotated; confidence: high}
\end{GSitem}
\begin{GSitem}{1973-333.12}{E}{Cauchy smoothing of the spectral measure}
\GSfield{Kind; source}{construction; \GSpage{P000552}}
\GSfield{The notes}{\(\sum_k\delta(\lambda-\lambda_k)*\varepsilon/(\lambda^2+\varepsilon^2)=f(\lambda)\).}
\GSfield{Statement}{For a discrete spectral measure \(\rho=\sum_k\delta_{\mu_k}\), the written smoothing formally gives \(\sum_k\varepsilon/[(\lambda-\mu_k)^2+\varepsilon^2]\). For Laplacian eigenvalues on a closed \(n\)-manifold, Weyl growth makes this series converge at fixed \(\lambda,\varepsilon>0\) when \(n<4\), but diverge for \(n\ge4\). If the parameters are frequencies \(\sqrt{\mu_k}\), it already diverges for \(n\ge2\). A Schwartz smoothing or explicit regularization makes a dimension-independent statement possible.}
\GSfield{Reconstruction}{The convergence test is supplied here. The missing \(1/\pi\) changes only normalization; the slow decay and eigenvalue/frequency distinction are the substantive issues.}
\GSfield{Fate}{The Weyl/heat asymptotic in \citet{mckeansinger1967curvature} supplies the counting growth. The convergence threshold is an editorial comparison with \(\sum k^{-4/n}\), or \(\sum k^{-2/n}\) for frequencies. The source's switch of spectral parameter is therefore material even for a two- or three-dimensional billiard.}
\GSfield{Present status}{As of 9 November 2026: Established in the stated scope; no unresolved claim is inferred from this occurrence.}
\GSfield{Assessment}{literal statement: correct with missing hypotheses; publication match: partial; verification: read the result; development: annotated; confidence: high}
\end{GSitem}
\begin{GSitem}{1973-333.13}{E}{Heat and wave traces must be distinguished}
\GSfield{Kind; source}{claim; \GSpage{P000553}}
\GSfield{The notes}{\(\sum_ke^{i\lambda_kz}=\operatorname{Sp}(e^{-z\Delta})\), with continuation downward from the upper half-plane.}
\GSfield{Statement}{For \(L\ge0\) with eigenvalues \(\mu_k\), the heat trace is \(H(z)=\sum_ke^{-z\mu_k}\) for \(\operatorname{Re}z>0\). The half-wave trace is \(W(z)=\sum_ke^{iz\sqrt{\mu_k}}=\operatorname{Tr}e^{iz\sqrt L}\) for \(\operatorname{Im}z>0\). Its real-axis boundary value is a distribution. These are different functions, and neither the factor \(i\) nor the square root may be dropped without changing the operator and variable.}
\GSfield{Reconstruction}{Analytic continuation across the entire real axis is not asserted; singular boundary values are exactly the feature relevant to closed trajectories.}
\GSfield{Fate}{\citet[Introduction]{colindeverdiere1973geodesicsii} studies complex-time spectral traces, while \citet{duistermaatguillemin1975spectrum} develops the positive first-order operator wave trace. The latter paper was received 20 October 1974 and published in February 1975; its publisher metadata was opened. The elementary spectral identities themselves follow directly from the functional calculus.}
\GSfield{Present status}{As of 9 November 2026: Established in the stated scope; no unresolved claim is inferred from this occurrence.}
\GSfield{Assessment}{literal statement: false as written; publication match: partial; verification: publisher metadata checked; development: annotated; confidence: high}
\end{GSitem}
\begin{GSitem}{1973-333.14}{F}{The underspecified twin-gap question}
\GSfield{Kind; source}{question; \GSpage{P000553}}
\GSfield{The notes}{\ru{Пробл. близнецов}: \(\lambda_{k+1}-\lambda_k>c\lambda_k^{-\nu}\) \ru{для общих операторов}; the next sentence mentions double eigenvalues on the circle.}
\GSfield{Statement}{For spectra counted with multiplicity, the universal interpretation---a positive polynomial lower bound on every high adjacent gap for every operator in the discussed class---is false, already for the circle. If \ru{общих} means generic rather than arbitrary, an operator class, topology, choice of spectral parameter, and multiplicity convention are missing. The source does not specify a recoverable generic-gap conjecture.}
\GSfield{Reconstruction}{The two possible meanings of 'general' are separated. No problem-index entry is created without choosing an unrecorded operator class.}
\GSfield{Fate}{\citet[theorem and opening discussion]{shnirelman1975multiplicity} gives a strong nearby obstruction: superpolynomially close neighbors on near-flat two-tori. The universal reading has an elementary negative answer, while the generic reading is not identified by the source. The notebook phrase is therefore not used to invent a new open gap problem.}
\GSfield{Present status}{As of 9 November 2026: The universal reading is false; a distinct generic formulation is not specified and is not classified as open.}
\GSfield{Assessment}{literal statement: not determinable; publication match: partial; verification: read the result; development: annotated; confidence: high}
\end{GSitem}
\begin{GSitem}{1973-333.15}{E}{The circle supplies exact spectral doublets}
\GSfield{Kind; source}{claim; \GSpage{P000553}}
\GSfield{The notes}{\ru{но на окружности двойные} (stably asymptotically twofold spectrum).}
\GSfield{Statement}{For the scalar Laplacian on a circle of length \(2\pi\), the eigenvalues are \(0,1,1,4,4,9,9,\ldots\). Thus \(\mu_{2j-1}=\mu_{2j}=j^2\), and the frequencies are similarly paired. Every Riemannian metric on a one-dimensional circle is isometric to a circle of its total length; the scalar metric Laplacian therefore retains these doublets after a metric perturbation, with an overall length scaling. A separate added potential need not retain exact degeneracy.}
\GSfield{Reconstruction}{The scalar metric-Laplacian class is essential to the exact stability statement; arbitrary operator perturbations are not included.}
\GSfield{Fate}{This is the classical Fourier-series spectrum and can be checked directly with \(e^{\pm ijx}\). \citet[opening discussion]{shnirelman1975multiplicity} places the later asymptotic multiplicity question in this context; the explicit circle computation is elementary and independently included in the verification output.}
\GSfield{Present status}{As of 9 November 2026: Established in the stated scope; no unresolved claim is inferred from this occurrence.}
\GSfield{Assessment}{literal statement: correct with missing hypotheses; publication match: exact; verification: read the result; development: annotated; confidence: high}
\end{GSitem}
\begin{GSitem}{1973-333.16}{B}{Superpolynomially close neighbors on a near-flat two-torus}
\GSfield{Kind; source}{theorem; \GSpage{P000553}}
\GSfield{The notes}{For \(T^2\), \(g_{ij}\sim\delta_{ij}\), the notes state \(\min(\lambda_{k+1}-\lambda_k,\lambda_k-\lambda_{k-1})<C_Nk^{-N}\), every \(N\), attributed parenthetically to Lazutkin and Shnirelman.}
\GSfield{Statement}{For a smooth metric on \(T^2\) sufficiently close to a flat metric in the KAM regime, let \(\omega_k\) be the ordered eigenvalues of \(\sqrt{-\Delta_g}\), counted with multiplicity. Then for every \(N>0\) there is \(C_N>0\) such that, for every \(k>1\), \(\min(\omega_{k+1}-\omega_k,\omega_k-\omega_{k-1})<C_Nk^{-N}\). The proof uses four distinct invariant KAM tori in the unit cosphere bundle, paired by time reversal, with Diophantine rotation data. The same qualitative superpolynomial statement for \(-\Delta_g\) follows after adjusting the exponent by Weyl growth.}
\GSfield{Reconstruction}{Smoothness, KAM smallness, multiplicities, all-index quantifiers, and the positive square-root spectral parameter are restored from the primary publication. No numerical universal smallness threshold is asserted.}
\GSfield{Fate}{\citet[unnumbered theorem and Lemmas~1--3, pp.~265--266]{shnirelman1975multiplicity} is an exact primary match, including the all-\(k\) quantifier. It was received by the Moscow Mathematical Society on 7 May 1974 and published in 1975. It explicitly strengthens Lazutkin's positive-density subsequence result and thanks him for discussions. Both pages and the proof sketch were read. Classification B reflects Koblenz reporting a result by others; the parenthetical attribution does not establish a separate speaker intervention.}
\GSfield{Present status}{As of 9 November 2026: The stated smooth near-flat two-torus theorem is published and established in its specified KAM scope; it is not an open-problem candidate.}
\GSfield{Assessment}{literal statement: correct with missing hypotheses; publication match: exact; verification: read the result; development: annotated; confidence: high}
\end{GSitem}
\begin{GSitem}{1973-333.17}{F}{The tentative three-torus comparison}
\GSfield{Kind; source}{remark; \GSpage{P000553}}
\GSfield{The notes}{\ru{Для \(T^3\) это уже не так (по-видимому)}.}
\GSfield{Statement}{The sentence suggests failure of the immediately preceding stable near-flat two-torus phenomenon in three dimensions, but gives no metric, topology, counterexample, quantifiers, or replacement claim. It cannot mean that every three-torus lacks close pairs: every flat three-torus has at least the \(\pm k\) pairing of nonzero Fourier modes.}
\GSfield{Reconstruction}{The scope of 'this' is not silently expanded into a universal assertion about all three-torus spectra.}
\GSfield{Fate}{\citet{shnirelman1975multiplicity} is explicitly two-dimensional. Its four separating invariant tori argument does not automatically extend to three dimensions. No exact primary counterexample matching the notebook's tentative assertion was identified; a dimensional analogy is not an open conjecture with determined scope.}
\GSfield{Present status}{As of 9 November 2026: The intended higher-dimensional failure statement remains unspecified.}
\GSfield{Assessment}{literal statement: not determinable; publication match: none found; verification: unresolved lead; development: annotated; confidence: high}
\end{GSitem}
\begin{GSitem}{1973-333.18}{E}{The half-wave propagator and its distributional trace}
\GSfield{Kind; source}{construction; \GSpage{P000553}}
\GSfield{The notes}{\(u_t=i\sqrt{-\Delta}u\), \(u|_{t=0}=\delta_x\); then \(\sum e^{i\lambda_kt}=\int_Mu(x,x,t)\,dx=f(t)\).}
\GSfield{Statement}{For \(L=-\Delta_g\ge0\) on a smooth closed manifold, define \(U(t)=e^{it\sqrt L}\). Its kernel is a distribution and solves \(\partial_tU=i\sqrt L\,U\), \(U(0)=I\). The trace is the distribution \(\sum_ke^{it\omega_k}\), \(\omega_k=\sqrt{\mu_k}\); its pairing with a smooth compactly supported function of \(t\) is defined by rapidly decaying Fourier coefficients. The diagonal integral is understood as this distributional kernel trace, not a pointwise integral of an ordinary function for every \(t\).}
\GSfield{Reconstruction}{The square-root frequencies and distributional interpretation are made explicit.}
\GSfield{Fate}{\citet[ChapterVI]{duistermaathormander1972fourierii} supplies the earlier Fourier-integral propagation framework; the later trace treatment is \citet{duistermaatguillemin1975spectrum}. This reconciles the source's final wave equation with the spectral notation introduced in the heat discussion.}
\GSfield{Present status}{As of 9 November 2026: Established in the stated scope; no unresolved claim is inferred from this occurrence.}
\GSfield{Assessment}{literal statement: correct with missing hypotheses; publication match: partial; verification: read the result; development: annotated; confidence: high}
\end{GSitem}
\begin{GSitem}{1973-333.19}{E}{Unit-speed propagation of singularities and caustics}
\GSfield{Kind; source}{theorem; \GSpage{P000553}}
\GSfield{The notes}{The kernel is said to be singular where \(|x-y|=t\); singularities move at speed 1, with bifurcations at some times.}
\GSfield{Statement}{Let \(G^s\) be the Hamilton flow of \(|\xi|_g\). With \(U(t)=e^{+it\sqrt L}\) and the output-first primed kernel convention, its canonical relation maps an input covector to \(G^{-t}\) of that covector. In Euclidean coordinates the phase \((x-y)\cdot\xi+t|\xi|\) gives \(x=y-t\xi/|\xi|\). The projected geodesics still have unit speed. Locally before the cut/conjugate locus, the wavefront lies on \(d_g(x,y)=|t|\). Globally one must use all geodesic connections of length \(|t|\), including nonminimizing ones, rather than only the metric-distance sphere. Caustics are singularities of the projection of this smooth canonical relation; a single distance-phase formula fails there.}
\GSfield{Reconstruction}{The metric-distance qualification is local; global propagation is formulated in cotangent phase space. No unrecorded caustic classification is added. The negative flow time follows from the chosen positive sign in the propagator exponent; reversing the exponent reverses the flow direction.}
\GSfield{Fate}{\citet[Theorem~6.1.1 and its proof]{duistermaathormander1972fourierii} establishes propagation along real-principal bicharacteristics in 1972. This is a pre-seminar theorem. The source's word 'bifurcations' is a geometric description, not a classification of all caustic types.}
\GSfield{Present status}{As of 9 November 2026: Established in the stated scope; no unresolved claim is inferred from this occurrence.}
\GSfield{Assessment}{literal statement: correct with missing hypotheses; publication match: exact; verification: read the result; development: annotated; confidence: high}
\end{GSitem}
\begin{GSitem}{1973-333.20}{E}{Closed-geodesic contributions to the wave trace}
\GSfield{Kind; source}{claim; \GSpage{P000553}}
\GSfield{The notes}{After the diagonal trace, a loop from a point back to itself is drawn and propagation at speed 1 is emphasized.}
\GSfield{Statement}{For the half-wave trace on a smooth closed manifold, nonzero singular times are contained in the lengths (with signs) of closed geodesics. A returning path contributes to the trace's stationary canonical relation only when both position and covector return. Under nondegenerate or clean fixed-point hypotheses, localized trace expansions encode the periodic orbits. The inclusion need not be equality when contributions cancel, and a single drawing does not specify the requisite nondegeneracy.}
\GSfield{Reconstruction}{Closed covectors, possible cancellation, and clean/nondegenerate hypotheses are supplied. The preceding billiard remark requires a separate boundary analysis.}
\GSfield{Fate}{\citet[Introduction and periodic-geodesic expansion]{colindeverdiere1973geodesicsii} is a pre-seminar rigorous precursor; \citet{duistermaatguillemin1975spectrum} is the subsequent general positive-elliptic-operator treatment. This is a partial match to the source's schematic explanation; the latter publication's metadata, not its entire proof, was checked.}
\GSfield{Present status}{As of 9 November 2026: Established in the stated scope; no unresolved claim is inferred from this occurrence.}
\GSfield{Assessment}{literal statement: correct with missing hypotheses; publication match: partial; verification: publisher metadata checked; development: annotated; confidence: high}
\end{GSitem}

\section{Developments}\label{dev:1973}
\noindent\textbf{Stage-2 selection and the scope of this appendix.}
No class-C or class-D item justified a bounded research development in this run.
The following elementary proofs support the critical edition of already identified
material. They verify source readings and reconstructions, with no historical
novelty claim. The item records retain development status \textit{annotated};
none is reclassified as an open problem to justify this appendix.


% Critical checks of source formulas, not Stage 2 research developments.
\subsection{The sign distribution and its moments}
The conventions change between P000503 and P000505. We use the latter: for an odd
prime \(p\), put \(N=p^m\) and, for \(0\le a<N\), set
\[
 \mu(a+N\mathbb Z_p)=\begin{cases}-1,&2a<N,\\+1,&2a>N.\end{cases}
\]
The following verifies the notebook formula and its normalization; it is not a
claim to a new theorem on \(p\)-adic integration. The underlying regularization is
part of \citet[\S8.4]{manin1973periods}.

\begin{proposition}[Moment correction for items 1973-303.14 and 1973-304.2--3]\ours\quad
These values define a bounded clopen distribution. With Bernoulli numbers
\(B_j\) defined by \(t/(e^t-1)=\sum_{j\ge0}B_jt^j/j!\), its polynomial moments are
\[
 I_k:=\int_{\mathbb Z_p}x^k\,d\mu
   =\frac{(4-2^{1-k})B_{k+1}}{k+1},\qquad k\ge0.
\]
Moreover \(\int_{\mathbb Z_p^\times}x^k\,d\mu=(1-p^k)I_k\).
\end{proposition}
\begin{proof}
Write \(p=2r+1\). Among the children \(a+jN\), \(0\le j<p\), the first \(r\)
lie below \(pN/2\), the last \(r\) lie above it, and the middle child lies below
exactly when \(a<N/2\). The first and last children cancel in the sum; the middle
child has the parent's sign. Thus the distribution compatibility equation holds
at every level, including level zero. Every residue-ball value has absolute value one, so every clopen value has absolute value at most one by the non-Archimedean triangle inequality.
Uniform approximation by locally constant functions therefore defines its
integral on continuous \(\mathbb Q_p\)-valued functions.

For the polynomial \(x^k\), the canonical residue-tag sum is
\[
 S_k(N)=\sum_{a=0}^{N-1}a^k-2\sum_{a=0}^{(N-1)/2}a^k,
\]
with the convention \(a^0=1\), also at \(a=0\). Let
\(B_j(X)\) be the Bernoulli polynomial. Its elementary difference identity
\(B_{k+1}(X+1)-B_{k+1}(X)=(k+1)X^k\) gives the exact rational identity
\[
 S_k(N)=\frac{B_{k+1}(N)-2B_{k+1}((N+1)/2)+B_{k+1}(0)}{k+1}.
\]
As \(m\to\infty\), \(N=p^m\to0\) in \(\mathbb Q_p\). Polynomial continuity yields
\[
 I_k=\frac{2(B_{k+1}-B_{k+1}(1/2))}{k+1}.
\]
The Bernoulli duplication identity gives
\(B_{k+1}(1/2)=(2^{-k}-1)B_{k+1}\), which proves the displayed result.
In particular \(I_0=-1\), \(I_1=1/4\), \(I_3=-1/32\), and \(I_{2r}=0\) for
\(r\ge1\). These values alone already distinguish the source's convention from
several possible Bernoulli-index corrections.

Finally \(\mu(pU)=\mu(U)\) for each residue ball \(U\): the defining inequality
is unchanged when both its numerator and modulus are multiplied by \(p\).
Taking finite sums and then limits gives
\(\int_{p\mathbb Z_p}x^k\,d\mu=p^kI_k\); subtracting from the total proves the
unit formula. The opposite convention on P000503 negates every moment.
\end{proof}

\paragraph{Independent verification of the moment calculation.}
\oursComputed\ The accompanying \texttt{code/independent\_checks.py} derives
Bernoulli numbers recursively using exact rational arithmetic. It checks the
finite-sum identity for \(p=3,5,7\), \(k=0,\ldots,9\), and \(m=1,\ldots,4\):
all 120 comparisons agree exactly. The recorded \(p\)-adic error valuations
are descriptive checks; the proof of the limit is polynomial continuity,
not extrapolation from those cases. The sign convention, duplication
identity and unit restriction were independently rederived, with the same
formula and the values \(-1,1/4,-1/32\) at \(k=0,1,3\).


\subsection{Dependence on tags for the invariant distribution}
\begin{proposition}[Item 1973-303.16]\ours\quad
The translation-invariant distribution of total mass one on \(\mathbb Z_p\)
does not have a tag-independent Riemann integral of \(x\).
\end{proposition}
\begin{proof}
On level \(m\) it gives each residue ball mass \(p^{-m}\). With canonical tags,
\[
 p^{-m}\sum_{a=0}^{p^m-1}a=(p^m-1)/2\longrightarrow-1/2.
\]
Change the tag in just the ball \(p^m\mathbb Z_p\) from \(0\) to \(p^m\). This
adds \(p^mp^{-m}=1\) at every level and changes the limit to \(1/2\). Both
partitions have mesh tending to zero; only their tags differ. This does not
contradict integration with a separately fixed representative convention.
\end{proof}

\subsection{Two checks on categorical realization}
\begin{example}[A nondegenerate simplex with a degenerate face, item 1973-308.5]\ours\quad
In the nerve of the group \(C_2=\{1,g\}\), the two-simplex \((g,g)\) is
nondegenerate, since degeneracies insert an identity arrow. Its middle face is
\(g^2=1\), the degenerate one-simplex at the unique vertex. Thus the disputed
sentence on P000515 is false even for the elementary group example on P000516.
The usual realization quotient still attaches a cell for every nondegenerate
simplex; its attaching map may identify part of its boundary with lower cells.
\end{example}

\begin{proposition}[Two Q-construction checks, items 1973-311.11--12]\ours\quad
Let \(A\) be a unital ring and \(P(A)\) its exact category of finitely generated
projective modules. The map \(BGL_n(A)\to BQP(A)\) obtained by viewing
\(A^n\) and its automorphisms in \(QP(A)\) is nullhomotopic. Direct sum with
\(A\) induces a map \(BQP(A)\to BQP(A)\) homotopic to the identity.
\end{proposition}
\begin{proof}
For the first assertion, a Q-morphism \(0\to A^n\) is represented by
\(0\xleftarrow{\mathrm{id}}0\rightarrowtail A^n\). Postcomposing it with any
automorphism of \(A^n\) gives the same Q-morphism. Hence these arrows form a
natural transformation from the constant zero functor on the one-object
category \(GL_n(A)\) to the automorphism functor. Natural transformations realize
to homotopies, by \citet[\S1, Proposition~2]{quillen1973higher}.

For the second assertion, let \(S(X)=X\oplus A\), and define S on a span by
summing the identity of A on both arrows. The admissible inclusion
\(X\hookrightarrow X\oplus A\) gives a Q-morphism represented by
\(X\xleftarrow{\mathrm{id}}X\rightarrowtail X\oplus A\). To verify naturality
for \(X\xleftarrow{p}E\xrightarrow{i}Y\), pull back
\(p\oplus1:E\oplus A\to X\oplus A\) along the inclusion of X. This pullback
is E, so either composite is represented by
\(X\xleftarrow{p}E\xrightarrow{(i,0)}Y\oplus A\). Thus these components give
\(\mathrm{id}\Rightarrow S\), proving the required homotopy.

The first natural transformation only concerns the automorphism subcategory.
It does not extend to a contraction of Q: pulling a general admissible
epimorphism back along zero can have a nonzero kernel. Neither assertion replaces
the substantive comparison theorem
\(\Omega BQP(A)\simeq K_0(A)\times BGL(A)^+\), checked in
\citet[Corollary~IV.7.2]{weibel2013kbook}.
\end{proof}

\subsection{Critical verification of the conic example (1973-322.9)}
\label{sec:conic-verification-1973}

This is a verification appendix to the settled complex formulation indexed as
GP-1973-32. It is not a Stage~2 contribution to a C or D problem. The explicit
certificates below are proved here to resolve a consequential ambiguity in the
notes; no historical priority or novelty is claimed.

\paragraph{What the source supports.}
P000535 first gives an example explicitly over \(\mathbb R\), with the
unimodular triple \((x,y,1-x^2-y^2)\) and the comment that degree obstructs its
reduction. After an undefined winding symbol and the definition of symplectic
\(K_1\), the page asks about \(z_1^2+z_2^2=1\), entire functions, and entire
functions of exponential growth. The two resulting functions are required to
have \ru{не имеют общ. делит.}, literally no common divisor. The parenthesis
\ru{в полиномах нет, в непр. ф-ях есть} says that polynomials do not suffice
but continuous functions do. These words and the equation are legible in the
scan; the coefficient field is not explicitly restated.

The strongest coherent existence question furnished by the displayed formulas
is the following, once the algebra is specified. Put
\[
 c=x^2+y^2-1,
 \qquad g_1=x+cf_1,\qquad g_2=y+cf_2.
\]
Can \(f_1,f_2\) be chosen so that \(g_1,g_2\) generate the unit ideal in the
chosen function algebra? This is stronger than the literal common-divisor
condition. In \(\mathbb C[x,y]\), taking \(f_1=f_2=0\) already leaves the
relatively prime pair \(x,y\), although that pair has a common zero. The
stable-rank context is the reason for reconstructing unimodularity. It is not
an independent reading of an unrecorded condition. The initial triple really
is unimodular, since \(x\cdot x+y\cdot y-c=1\).

\paragraph{The verified literature and its scope.}
The real obstruction is an old one: Vaserstein's 1971 Theorem~8 and its proof
use the same sphere triple to establish the full polynomial stable rank over
subfields of \(\mathbb R\) \citep[Russian p.~27]{vaserstein1971stable}.
For unrestricted holomorphic functions, Brudnyi's theorem gives
\[
 \operatorname{sr}\mathcal O(X)
   =\left\lfloor\frac{\dim_{\mathbb C}X}{2}\right\rfloor+1
\]
for finite-dimensional Stein spaces, and therefore
\(\operatorname{sr}\mathcal O(\mathbb C^2)=2\)
\citep[Theorem~1.1]{brudnyi2019stein}. Its preprint circulated in March 2018;
the journal article appeared in January 2019. This theorem applies to every
unimodular triple of entire functions on \(\mathbb C^2\), but supplies no
exponential-growth restriction. The following particular example needs neither
that general theorem nor a new open-problem designation.

\begin{proposition}[A polynomial certificate over the complex numbers]\ours\quad
Let \(c=x^2+y^2-1\) and set
\[
 f_1=\frac12,\qquad
 f_2=\frac{i(x+2)-y}{2},\qquad
 a=x+cf_1,\qquad b=y+cf_2.
\]
Then \(a,b\) generate the unit ideal in \(\mathbb C[x,y]\). More precisely,
with
\[
 s=x+1-iy,\qquad
 A=s(a+1)-a-2,\qquad B=\frac{i s^2}{2},
\]
one has \(Aa+Bb=1\) identically.
\end{proposition}

\begin{proof}
Put \(t=x+1+iy\). Direct expansion gives
\[
 ts=(x+1)^2+y^2=2(a+1),
 \qquad
 b=\frac{t^2}{2i}+ita.
\]
For completeness, the second identity follows by expanding its right side:
its coefficient of \(i\), after grouping over \(\mathbb Q[x,y]\), is
\(i((x+1)a-((x+1)^2-y^2)/2)\), and its remaining part is
\((x+1)y-ya\); substituting \(a=x+c/2\) gives
\(y+c(i(x+2)-y)/2\). Hence
\[
 Bb=\frac{s^2t^2}{4}-\frac{s^2t a}{2}
    =(a+1)^2-s(a+1)a.
\]
Combining this with \(Aa=s(a+1)a-a^2-2a\) proves
\[
 Aa+Bb=(a+1)^2-a^2-2a=1.
\]
Every coefficient belongs to \(\mathbb Q(i)\), so this is an exact polynomial
identity. In particular \(a,b\) have no simultaneous zero anywhere in
\(\mathbb C^2\).
\end{proof}

\begin{corollary}[Entire and exponential-type formulations]\ours\quad
The complex conic triple is reducible in the polynomial algebra, in
\(\mathcal O(\mathbb C^2)\), and in the algebra
\[
 \mathcal E(\mathbb C^2)=
 \{f\in\mathcal O(\mathbb C^2):
   \exists C,T>0\ \forall z\in\mathbb C^2,   |f(z)|\le C e^{T(|z_1|+|z_2|)}\}.
\]
The reduction and its Bezout coefficients can all be taken to be polynomials.
\end{corollary}

\begin{proof}
The certificate persists under either inclusion of polynomial algebras.
Every polynomial \(p\) satisfies
\(|p(z)|\le C_\epsilon e^{\epsilon(|z_1|+|z_2|)}\) for every
\(\epsilon>0\), because a fixed power of \(1+r\) times
\(e^{-\epsilon r}\) is bounded on \([0,\infty)\). The algebra
\(\mathcal E\) is closed under addition and multiplication, with the
constants and types allowed to depend on the function. Thus both generators
and the explicit Bezout coefficients lie in it.
\end{proof}

\begin{proposition}[An entire certificate keeping the first coordinate]\ours\quad
There is also a solution with \(f_1=0\). Define
\[
 h(w)=\exp\!\left(\frac{i\pi(w-1)}2\right),\qquad
 F(w)=\frac{h(w)-w}{w^2-1},
\]
where the two removable values are filled in. Then
\[
 f_2(x,y)=F(y),\qquad
 g_1=x,\qquad g_2=h(y)+x^2F(y),
\]
and
\[
 -\frac{xF(y)}{h(y)}\,g_1+\frac1{h(y)}\,g_2=1.
\]
All the displayed entire functions belong to \(\mathcal E(\mathbb C^2)\).
\end{proposition}

\begin{proof}
The values \(h(1)=1\) and \(h(-1)=-1\) show that the numerator of \(F\)
vanishes at both roots of the denominator. Both roots are simple, and the
removable values obtained by differentiation are
\[
 F(1)=\frac{i\pi}{4}-\frac12,\qquad
 F(-1)=\frac{i\pi}{4}+\frac12.
\]
Consequently \(F\) is entire and
\[
 cF(y)=x^2F(y)+h(y)-y.
\]
This is precisely the required form for \(g_2=y+cf_2\). The function \(h\)
never vanishes. Substitution proves the Bezout identity for all \((x,y)\),
including \(y=\pm1\) by removability.

Here is a uniform growth check, not only a pointwise asymptotic argument.
Write \(\alpha=\pi/2\). For every \(w\),
\(|h(w)|,|h(w)^{-1}|\le e^{\alpha|w|}\). If \(|w|\ge2\), then
\[
 |w^2-1|\ge |w|^2-1\ge\tfrac34|w|^2,
 \qquad
 |F(w)|\le\frac{4(e^{\alpha|w|}+|w|)}{3|w|^2}.
\]
On \(|w|\le2\), the removable extension is bounded by compactness. Therefore
some finite \(K\) satisfies \(|F(w)|\le K e^{\alpha|w|}\) for all \(w\),
and \(|F(w)/h(w)|\le K e^{2\alpha|w|}\). Using
\(r\le e^r\) and \(r^2\le2e^r\) for \(r\ge0\) gives, for example,
\[
 \begin{split}
 |g_2(x,y)|&\le(1+2K)e^{|x|+\alpha|y|},\\
 |xF(y)/h(y)|&\le K e^{|x|+\pi|y|},\\
 |1/h(y)|&\le e^{\alpha|y|}.
 \end{split}
\]
Together with the bounds for \(g_1,f_1,f_2\), these imply a common estimate
\(C e^{(\pi+1)(|x|+|y|)}\) after increasing \(C\). No optimal type is
claimed.
\end{proof}

\begin{proposition}[The real-valued interpretation has the opposite answer]\ours\quad
Let \(D^2=\{(x,y)\in\mathbb R^2:x^2+y^2\le1\}\). For continuous
real-valued \(f_1,f_2\) on \(D^2\), the pair
\[
 (x+(x^2+y^2-1)f_1,\ y+(x^2+y^2-1)f_2)
\]
has a zero in \(D^2\). In particular there is no such reduction over
\(\mathbb R[x,y]\).
\end{proposition}

\begin{proof}
If this pair were nowhere zero, dividing it by its Euclidean norm would give
a continuous map \(r:D^2\to S^1\). On \(S^1\) the factor
\(x^2+y^2-1\) is zero, so \(r\) restricts to the identity. The composite
\(S^1\hookrightarrow D^2\xrightarrow{r}S^1\) would induce the identity on
\(H_1(S^1;\mathbb Z)=\mathbb Z\), but it factors through
\(H_1(D^2;\mathbb Z)=0\). This contradiction proves the assertion. Changing
\(c\) to \(-c\), as in the preceding notebook example, only changes the
signs of the coefficient functions.
\end{proof}

\paragraph{Independent checks.}
\oursComputed\ The accompanying \texttt{check\_conic.py} uses exact pairs of
rational numbers for coefficients in \(\mathbb Q(i)\), expands both
intermediate identities, and verifies that every coefficient of
\(Aa+Bb-1\) is zero. Its recorded output gives degrees
\(0,1,2,3,3,2\) for \(f_1,f_2,a,b,A,B\). Independently, a second checker
uses exact Gaussian-rational evaluation on a \(6\times6\) Cartesian grid.
The difference polynomial has degree at most five in each variable, so this
grid check is an exact interpolation proof, not merely numerical evidence;
four further complex points also pass. The two algebraic identities were
also rederived independently by hand.

For the alternative entire construction, eight selected complex points,
including both removable values, give maximum floating-point residuals
\(2.220\times10^{-16}\) for the row identity and
\(4.965\times10^{-16}\) for the Bezout identity. These are only implementation
checks. The proof of removability and the uniform growth estimate above,
independently rechecked, supplies the all-point justification. The scripts and
outputs are included with the chapter.

\paragraph{What remains uncertain.}
There is no remaining existence obstruction for the complex question just
formulated. The unanswered issue is documentary: the notes do not define the
index, clarify the coefficient field after the real example, or explain what
meaning of common divisor was intended. The complex polynomial certificate
invalidates the parenthesis under unrestricted complex coefficients; the real
degree argument invalidates its continuous-function half if all functions
remain real-valued. Neither fact licenses a reconstruction of the speaker's
missing restrictions. No assertion that a new open mathematical problem has
been solved, and no recommendation for a further research run, follows from
this critical verification.

\section{Problem index}\label{problems:1973}
Earlier-year occurrence identifiers are retained as text; current-year identifiers link to their entries.

\subsection{GP-1973-01: Finite presentation over a finite-field polynomial ring}\label{prob:GP-1973-01}
For every prime power \(q\) and integer \(n\geq3\), determine finite presentability of \(\mathrm{SL}_n(\mathbb F_q[t])\).
\par\noindent\textit{Occurrences:} \hyperref[item:1973-302.4]{1973-302.4}.
\par\noindent\textit{Status (2026-11):} Resolved: no at \(n=3\); yes at every \(n\geq4\), for all finite fields.

\subsection{GP-1973-02: Finite presentation over a finite-field Laurent ring}\label{prob:GP-1973-02}
For every finite field \(\mathbb F_q\) and integer \(n\geq2\), determine finite presentability of \(\mathrm{SL}_n(\mathbb F_q[t,t^{-1}])\).
\par\noindent\textit{Occurrences:} \hyperref[item:1973-302.5]{1973-302.5}.
\par\noindent\textit{Status (2026-11):} Resolved: no at \(n=2\); yes at every \(n\geq3\), for all finite fields.

\subsection{GP-1973-03: Coherence of SL3(Z)}\label{prob:GP-1973-03}
Is every finitely generated subgroup \(H\leq\mathrm{SL}_3(\mathbb Z)\) finitely presented?
\par\noindent\textit{Occurrences:} \hyperref[item:1973-302.11]{1973-302.11}.
\par\noindent\textit{Status (2026-11):} Open in checked primary literature \citep[p.~5]{audibert2025zariski}.

\subsection{GP-1973-04: Coherence of SL2(Z[1/p])}\label{prob:GP-1973-04}
For every prime \(p\), is every finitely generated subgroup of \(\mathrm{SL}_2(\mathbb Z[1/p])\) finitely presented?
\par\noindent\textit{Occurrences:} \hyperref[item:1973-302.12]{1973-302.12}.
\par\noindent\textit{Status (2026-11):} Open unconditionally; the Greenberg--Shalom hypothesis implies a positive answer \citep[Theorem~4.2, version~3]{brody2025commensurator}.

\subsection{GP-1973-05: Ranks for coherence of GLn(Q)}\label{prob:GP-1973-05}
For which integers \(n\geq1\) is the abstract discrete group \(\mathrm{GL}_n(\mathbb Q)\) coherent?
\par\noindent\textit{Occurrences:} \hyperref[item:1973-302.13]{1973-302.13}.
\par\noindent\textit{Status (2026-11):} Coherent only at \(n=1\); rational triangular subgroups obstruct every \(n\geq2\) \citep[p.~74]{brody2022trees}.

\subsection{GP-1972-403: Ordinary p-adic interpolation of modular L-values}\label{prob:GP-1972-403}
For a normalized Hecke eigenform and a prime with unit eigenvalue, construct a bounded p-adic distribution from algebraically normalized rational-cusp periods whose primitive-character Mellin values interpolate the critical twisted L-values, with the Euler root and character conventions specified.
\par\noindent\textit{Occurrences:} 1972-299.10, 1972-299.12, 1972-299.13, 1972-299.14, \hyperref[item:1973-304.12]{1973-304.12}, \hyperref[item:1973-305.1]{1973-305.1}, \hyperref[item:1973-305.5]{1973-305.5}, \hyperref[item:1973-305.10]{1973-305.10}, \hyperref[item:1973-305.11]{1973-305.11}, \hyperref[item:1973-305.13]{1973-305.13}.
\par\noindent\textit{Status (2026-11):} As of 9 November 2026: Established in the ordinary scope; \citet[\S9.5 and Theorem~10.2]{manin1973periods}. The repeated September exposition is not a fresh open problem. No unrestricted nonordinary boundedness theorem is inferred.

\subsection{GP-1973-10: Transcendence of zeta(3)}\label{prob:GP-1973-10}
Is \(\zeta(3)\) transcendental over \(\mathbb Q\)?
\par\noindent\textit{Occurrences:} \hyperref[item:1973-305.15]{1973-305.15}.
\par\noindent\textit{Status (2026-11):} As of 9 November 2026: Ap\'ery\textquoteright s irrationality theorem resolves a weaker property; transcendence is still recorded as unknown in \citet[\S1]{nica2024irrationality}, and no accepted subsequent proof was found. \citet[Theorem~1.6]{lai2026linear} is a recent preprint reporting a weaker linear-independence result; it leaves this transcendence question untouched.

\subsection{GP-1973-11: Transcendence of Euler's constant (elliptical notebook occurrence)}\label{prob:GP-1973-11}
Is Euler\textquoteright s constant \(\gamma\) transcendental? The notebook\textquoteright s “also unproved” is read with the preceding predicate; the historical attribution is less certain than for zeta(3).
\par\noindent\textit{Occurrences:} \hyperref[item:1973-305.16]{1973-305.16}.
\par\noindent\textit{Status (2026-11):} As of 9 November 2026: Even irrationality is unresolved in the accepted literature checked; \citet[Conjectures~1.0.1--1.0.2 and \S3]{lagarias2013euler}. Purported isolated proofs in search results are not adopted as established results.

\subsection{GP-1973-30: Representability of every modular lattice by submodules}\label{prob:GP-1973-30}
Is every modular lattice isomorphic to the lattice of all submodules of some module over some ring?
\par\noindent\textit{Occurrences:} \hyperref[item:1973-319.3]{1973-319.3}.
\par\noindent\textit{Status (2026-11):} As of November 2026: false. Submodule lattices are Arguesian; non-Desarguesian projective-plane lattices are modular and not Arguesian. This obstruction predates the 1973 talk \citep{jonsson1954modular,herrmanntakach2005subgroup}.

\subsection{GP-1973-31: Classification of finite-dimensional r-subspace systems}\label{prob:GP-1973-31}
For a field \(k\) and fixed \(r\), describe all ordered systems \((V;E_1,\ldots,E_r)\) of subspaces in a finite-dimensional \(k\)-space up to simultaneous linear isomorphism, equivalently the finite-dimensional representations of \(D^r\).
\par\noindent\textit{Occurrences:} 1970-198.1, 1970-198.6, 1972-277.5, 1972-277.6, 1972-277.7, \hyperref[item:1973-319.6]{1973-319.6}.
\par\noindent\textit{Status (2026-11):} As of November 2026: the algebraically closed-field classification is finite for at most three subspaces and tame for four; five subspaces already encode arbitrary matrix pairs. Earlier occurrences are special cases of this general problem. This is distinct from the deformation question GP-1970-23 \citep{gelfandponomarev1971quadruples,bernsteingelfandponomarev1973coxeter}.

\subsection{GP-1973-32: The conic triple in polynomial and entire-function rings}\label{prob:GP-1973-32}
For \(c=z_1^2+z_2^2-1\), do there exist \(f_1,f_2\) in the specified function algebra such that \((z_1+cf_1,z_2+cf_2)\) generates its unit ideal? The coefficient field and function algebra are part of the question; the notebook's ``common divisor'' is an imprecise formulation.
\par\noindent\textit{Occurrences:} \hyperref[item:1973-322.9]{1973-322.9}.
\par\noindent\textit{Status (2026-11):} As of 9 November 2026: affirmative over \(\mathbb C[z_1,z_2]\), hence for complex entire functions and the algebra of entire functions of finite exponential type, by the certificate checked in the critical appendix. Negative for real-valued continuous coefficients on the real disk. Brudnyi's later general theorem settles unrestricted entire reduction. No unresolved C/D contribution is inferred from the notebook's missing coefficient restriction.

\subsection{GP-1973-33: Serre's problem for projective modules over a polynomial ring}\label{prob:GP-1973-33}
For a field \(k\) and finite \(n\), is every finitely generated projective \(k[x_1,\ldots,x_n]\)-module free? This occurrence announces the case \(n=3\) without recording its then-needed field restriction.
\par\noindent\textit{Occurrences:} \hyperref[item:1973-322.1]{1973-322.1}.
\par\noindent\textit{Status (2026-11):} As of 9 November 2026: solved affirmatively by Quillen and Suslin in 1976. The 1974 Vaserstein--Suslin announcement proves the three-variable case for perfect fields or characteristic different from two. The statement of the general modern theorem must not erase that historical distinction.

\subsection{GP-1973-50: Monotonicity in a Hill spectral band}\label{prob:GP-1973-50}
For real smooth periodic \(u\), prove \(d_{\lambda}(\lambda)\ne0\) on \(\{|d|<1\}\), where \(d=\frac12\operatorname{tr}M\).
\par\noindent\textit{Occurrences:} \hyperref[item:1973-326.9]{1973-326.9}.
\par\noindent\textit{Status (2026-11):} As of 9 November 2026: A classical affirmative result; the seminar requested a proof. Closed-gap endpoints are excluded from the derivative assertion.

\subsection{GP-1973-51: Supplementary data for the periodic inverse spectral problem}\label{prob:GP-1973-51}
For real periodic \(u\), determine what data in addition to the Hill discriminant reconstruct \(u\), with basepoint fixed.
\par\noindent\textit{Occurrences:} \hyperref[item:1973-326.13]{1973-326.13}.
\par\noindent\textit{Status (2026-11):} As of 9 November 2026: Affirmatively solved in particular by band edges plus a sheet-labelled Dirichlet divisor in the finite-gap class; the cited verification here is limited to that match.

\subsection{GP-1973-52: Approximation by finite-gap periodic potentials}\label{prob:GP-1973-52}
Is the set of real periodic finite-gap Hill potentials dense among real periodic potentials? The precise modern version uses the \(H^m\) norm for each \(m\ge0\).
\par\noindent\textit{Occurrences:} \hyperref[item:1973-326.17]{1973-326.17}.
\par\noindent\textit{Status (2026-11):} As of 9 November 2026: Yes; finite-gap potentials are dense in every \(H^m\), \(m\ge0\). The original notebook omitted the topology.

\subsection{GP-1973-53: Recursive construction of KdV Lax operators}\label{prob:GP-1973-53}
Given \(L=-D^2+u\), construct a normalized odd-order \(A_n\) for which \([A_n,L]\) is multiplication, and relate the resulting hierarchy to conserved-density recursions.
\par\noindent\textit{Occurrences:} \hyperref[item:1973-327.1]{1973-327.1}.
\par\noindent\textit{Status (2026-11):} As of 9 November 2026: The formal hierarchy and its dressing construction are established. The seminar's additional question about unspecified other equations is not a universal classification problem.

\subsection{GP-1973-54: Soliton collision phase shifts}\label{prob:GP-1973-54}
For a real pure KdV two-soliton solution with distinct positive wave parameters and fixed scattering data, compute the outgoing minus incoming intercept of each soliton center.
\par\noindent\textit{Occurrences:} \hyperref[item:1973-330.5]{1973-330.5}.
\par\noindent\textit{Status (2026-11):} As of 9 November 2026: The logarithmic pairwise shifts are established; the source itself calls the scattering answer standard.

\section{Corrections to the transcription and translation}\label{corr:1973}
The evidence distinguishes a transcription correction from a mathematical emendation of wording faithfully preserved in the notebook.
\subsection*{P000501 (both; faithful error): Imaginary-quadratic coefficient ring}
\noindent\textit{Transcription reads.} \(\Lambda\): \ru{чисто мнимые элем. квадр. поля}\par
\noindent\textit{Critical reading.} The standard Bianchi-group interpretation uses \(\Lambda=\mathcal O_K\), the ring of integers of an imaginary quadratic field.\par
\noindent\textit{Evidence.} The scan supports the transmitted phrase, but purely imaginary elements are not a unital ring. This is a proposed mathematical reconstruction, not a transcription repair.\par
\subsection*{P000501 (both; ambiguity): Plan item 5, uncertain surname}
\noindent\textit{Transcription reads.} \ru{Герстенцуса}\par
\noindent\textit{Critical reading.} Keep the name unresolved.\par
\noindent\textit{Evidence.} Scan and context do not authorize replacing it with Gersten or Gerstenhaber.\par
\subsection*{P000501 (both; ambiguity): Plan item 10, spectral-asymmetry attribution}
\noindent\textit{Transcription reads.} Atiyah, Bott, Patodi\par
\noindent\textit{Critical reading.} Retain the names as written. The spectral-asymmetry announcement is by Atiyah, Patodi and Singer. Atiyah, Bott and Patodi also published their heat-equation/index paper in 1973, so the agenda may conflate two related subjects.\par
\noindent\textit{Evidence.} The scan supports \ru{Ботт}; bibliographical comparison does not authorize changing the diplomatic text. Compare \citet{atiyahpatodisinger1973spectral} and \citet{atiyahbottpatodi1973heat}.\par
\subsection*{P000501-P000502 (both; provenance): Record 302 catalogue range}
\noindent\textit{Transcription reads.} Record 302: P000501--P000502\par
\noindent\textit{Critical reading.} The plan/interlude ends on P000501; its catalogue range overlaps Manin's talk.\par
\noindent\textit{Evidence.} The explicit Manin heading establishes the boundary; all record 302 items are located on P000501.\par
\subsection*{P000501 (en; translation): Manin heading: X}
\noindent\textit{Transcription reads.} compact p-adic analytic set\par
\noindent\textit{Critical reading.} compact p-adic analytic manifold\par
\noindent\textit{Evidence.} Russian transcription and scan read analytic mn-zie, the abbreviation for manifold; local ball charts and Serre theorem confirm the meaning.\par
\subsection*{P000503 (both; faithful error): projective-line example}
\noindent\textit{Transcription reads.} \(\mathbb P^1(K)=\mathbb Q_p\cup\{\infty\}\), \(pA\cup A^{-1}\)\par
\noindent\textit{Critical reading.} \(\mathbb P^1(K)=K\cup\{\infty\}\); the small chart is \(\mathfrak m\), with infinity in the inverse chart\par
\noindent\textit{Evidence.} Scan and both transcriptions preserve Qp and pA. For general finite K/Qp, pA need not be the maximal ideal. This is a scope emendation, item1973-303.9.\par
\subsection*{P000503 (both; faithful error): compact-group example}
\noindent\textit{Transcription reads.} \(N=r(X)\) and an SL2 congruence modulo p\par
\noindent\textit{Critical reading.} \(i(G)\equiv N\pmod{p-1}\); for SL2(Zp), \(N=p^2-1\)\par
\noindent\textit{Evidence.} The modulus p is visible in the scan; Serre1965 example(ii),p410 gives residue class modulo p-1. Item1973-303.10.\par
\subsection*{P000503; P000505 (both; ambiguity): two-valued distribution sign}
\noindent\textit{Transcription reads.} P503 lower half +1; P505 lower half -1\par
\noindent\textit{Critical reading.} Retain the two opposing sign conventions; all P505 moment calculations use lower half -1\par
\noindent\textit{Evidence.} Both signs are visible in the scans and are not an English translation error. Items1973-303.14 and1973-304.2.\par
\subsection*{P000503-P000506 (both; faithful error): two-valued distribution and character decomposition}
\noindent\textit{Transcription reads.} p unrestricted\par
\noindent\textit{Critical reading.} Require odd p for these formulas; at p=2 use \(\{\pm1\}\times(1+4\mathbb Z_2)\) for units\par
\noindent\textit{Evidence.} At p=2 the threshold equality is unassigned and1+p is not a generator of the torsion-free unit factor. This is a missing-hypothesis emendation.\par
\subsection*{P000505 (both; faithful error): odd moments}
\noindent\textit{Transcription reads.} \((-1)^r(1-2^{-r})B_{r-1}/(2r)\)\par
\noindent\textit{Critical reading.} For standard Bernoulli numbers and lower-half sign -1, \(\int x^k d\mu=(4-2^{1-k})B_{k+1}/(k+1)\)\par
\noindent\textit{Evidence.} The scan preserves the original formula. Exact Bernoulli-polynomial calculation in critical\_verifications.tex and root independent checks determines the normalization; k=1 gives1/4. Item1973-304.2.\par
\subsection*{P000505 (both; faithful error): restriction to units}
\noindent\textit{Transcription reads.} \(\int x^{2m-1}\) in the intermediate formula for the \(2r-1\) moment\par
\noindent\textit{Critical reading.} \(\int x^{2r-1}\) throughout, with factor \(1-p^{2r-1}\)\par
\noindent\textit{Evidence.} The two exponents coexist in the scan. Scaling the distribution by p proves the correction. Item1973-304.3.\par
\subsection*{P000508 (both; faithful error): tentative correction of the distribution}
\noindent\textit{Transcription reads.} \(\rho^{-m}P(a/p^m)-p^w\rho^{-(m+m)}P(a/p^{m+1})\)\par
\noindent\textit{Critical reading.} At level n use \(\rho^{-n}P(a/p^n)-p^w\rho^{-n-1}P(a/p^{n-1})\)\par
\noindent\textit{Evidence.} The page has a question mark. The corrected recurrence is explicitly repeated on P509, with an overall rho factor; Manin1973 Lemma9.4 and equation39. No silent replacement of the source text.\par
\subsection*{P000509 (both; faithful error): continued-fraction path sum}
\noindent\textit{Transcription reads.} first sum begins at \(k=0\) but invokes \(b_{k-1}/d_{k-1}\)\par
\noindent\textit{Critical reading.} use \(k=1,\ldots,m\) for the listed edges, or explicitly supply the omitted initial convergent and orientation\par
\noindent\textit{Evidence.} The following sum on the same scan starts at1. Item1973-305.3 explains the modern edge decomposition.\par
\subsection*{P000509 (both; faithful error): period reality}
\noindent\textit{Transcription reads.} even \(\tau_j\) imaginary, odd \(\tau_j\) real\par
\noindent\textit{Critical reading.} add that the Fourier coefficients of the cusp form are real\par
\noindent\textit{Evidence.} The real-form condition is supplied on P510; otherwise multiply a form by a nonreal scalar to disprove the literal assertion. Item1973-305.4.\par
\subsection*{P000510 (both; faithful error): rewritten Hecke formula}
\noindent\textit{Transcription reads.} minus before the residue-class average\par
\noindent\textit{Critical reading.} plus before the residue-class average\par
\noindent\textit{Evidence.} The minus is visible in the scan; the P507 formula and Manin1973 Hecke relation have plus. Item1973-305.8.\par
\subsection*{P000511 (both; faithful error): nonunit root growth}
\noindent\textit{Transcription reads.} the other root gives a growing distribution\par
\noindent\textit{Critical reading.} the unit root gives a bounded distribution; the other root has a growing upper bound, without automatic noncancellation\par
\noindent\textit{Evidence.} Scan preserves the categorical wording. The valuation argument bounds terms and does not alone prove unboundedness. Item1973-305.11.\par
\subsection*{P000512 (both; faithful error): polynomial ordering}
\noindent\textit{Transcription reads.} order of factors in the polynomial expression is immaterial\par
\noindent\textit{Critical reading.} the trace of the commutator is independent of the ordering convention\par
\noindent\textit{Evidence.} Scan and Helton--Howe1975 TheoremI and preceding paragraph; operator values can differ even when their commutator traces agree. Item1973-307.2.\par
\subsection*{P000512 (both; faithful error): trace-measure hypotheses}
\noindent\textit{Transcription reads.} self-adjoint operators commuting modulo nuclear operators\par
\noindent\textit{Critical reading.} bounded self-adjoint operators with trace-class commutator; \(\sigma\) is a finite purely imaginary measure in this normalization\par
\noindent\textit{Evidence.} Helton--Howe1975 TheoremI and Sugahara2026 Theorem3.3. Missing boundedness causes domain issues; a real measure would require an extra factor of i.\par
\subsection*{P000512 (both; provenance): record order}
\noindent\textit{Transcription reads.} catalogue orders record306 before307\par
\noindent\textit{Critical reading.} in notebook order place the unlabelled Gelfand opening307 before the explicit Kirillov heading306\par
\noindent\textit{Evidence.} Scan shows the trace discussion above the Kirillov heading; catalogue attribution to Gelfand is retained as catalogue evidence.\par
\subsection*{P000512 (both; faithful error): extension category}
\noindent\textit{Transcription reads.} extension of Banach algebras\par
\noindent\textit{Critical reading.} C*-algebra extension with *-homomorphisms, essentiality and stable equivalence in the BDF setting\par
\noindent\textit{Evidence.} The later Calkin embedding and unitary equivalence require the BDF C*-context; Douglas ICM pp699–701. Item1973-306.1.\par
\subsection*{P000513 (both; faithful error): boxed Ext/K0 relation}
\noindent\textit{Transcription reads.} \(\operatorname{Ext}_1\oplus\mathbb Z=K_0\), with overwritten operation\par
\noindent\textit{Critical reading.} for connected based finite CW X: \(\operatorname{Ext}(\Sigma X)\oplus\mathbb Z\cong K_0(X)\)\par
\noindent\textit{Evidence.} The scan confirms subscript1, but not a new decipherment of the overwritten operation. With P512 conventions the degree must be0; X=S1 disproves the unshifted formula. Item1973-306.6.\par
\subsection*{P000513 (both; faithful error): circle extension}
\noindent\textit{Transcription reads.} \(TT^*=1\bmod\mathcal K\)\par
\noindent\textit{Critical reading.} also \(T^*T=1\bmod\mathcal K\)\par
\noindent\textit{Evidence.} Both relations follow from the unitary generator under a *-homomorphism; the displayed relation alone is insufficient for Fredholmness. Item1973-306.8.\par
\subsection*{P000513 (both; ambiguity): affine group example and pair closure claim}
\noindent\textit{Transcription reads.} infinity symbol; unspecified operator pair and ambient class\par
\noindent\textit{Critical reading.} retain as unresolved notation/context, without supplying a numerical answer or a research conjecture\par
\noindent\textit{Evidence.} Scan inspected; published one-operator theorem located, pair counterexample not located. Items1973-306.7 and306.10.\par
\subsection*{P000515 (both; faithful error): faces of nondegenerate simplices}
\noindent\textit{Transcription reads.} arrows from right to left take nondegenerate simplices to nondegenerate simplices\par
\noindent\textit{Critical reading.} a face of a nondegenerate simplex can be degenerate\par
\noindent\textit{Evidence.} Scan says this explicitly. In the nerve of C2, the nondegenerate2-simplex(g,g) has degenerate middle face g*g=1. Item1973-308.5.\par
\subsection*{P000515 (both; faithful error): simplicial topological spaces}
\noindent\textit{Transcription reads.} a topological structure no longer arises\par
\noindent\textit{Critical reading.} a quotient topology still arises; an ordinary simplicial-set CW decomposition need not\par
\noindent\textit{Evidence.} Russian and scan have the same wording. Segal1968 section1 defines realization for simplicial spaces. Item1973-308.6.\par
\subsection*{P000515 (both; faithful error): product realization}
\noindent\textit{Transcription reads.} unqualified homeomorphism to the ordinary product\par
\noindent\textit{Critical reading.} homeomorphism in compactly generated spaces, or with suitable extra hypotheses for the ordinary product\par
\noindent\textit{Evidence.} Segal1968 section1 explicitly uses the category of compactly generated spaces. Item1973-308.9.\par
\subsection*{P000519-P000520 (both; provenance): speaker attribution}
\noindent\textit{Transcription reads.} no separate catalogue record for the Mishchenko block\par
\noindent\textit{Critical reading.} Retain these interventions in record 309 and attribute them to A. S. Mishchenko.\par
\noindent\textit{Evidence.} The scan has a clear Mishchenko heading on P519; eight items309.13–309.20 cover it.\par
\subsection*{P000519-P000520 (both; faithful error): tensor shorthand}
\noindent\textit{Transcription reads.} \(F\otimes1\) on a balanced \(C(M)\)-tensor product\par
\noindent\textit{Critical reading.} use projection compression \(T_2(e)F^{(n)}T_1(e)\)\par
\noindent\textit{Evidence.} F is a module map only modulo compact operators, so the literal tensor map need not descend. The compressed Fredholm operator has the intended index. Item1973-309.14.\par
\subsection*{P000520 (both; faithful error): classifying space of Fredholm operators}
\noindent\textit{Transcription reads.} \(BU\simeq\mathcal F\), \(K(B\pi)=[B\pi,BU]\)\par
\noindent\textit{Critical reading.} \(\mathcal F\simeq\mathbb Z\times BU\); BU alone denotes the zero-index component\par
\noindent\textit{Evidence.} The scan omits the integer factor and does not assume index0. Components are indexed by the Fredholm index. Item1973-309.19.\par
\subsection*{P000521 (both; faithful error): representation target}
\noindent\textit{Transcription reads.} \(\rho:\mathcal L\to\operatorname{Aut}E\)\par
\noindent\textit{Critical reading.} \(\rho:\mathcal L\to\operatorname{End}_k(E)\)\par
\noindent\textit{Evidence.} Aut is clear in the scan; Rudakov1974 section2 uses the endomorphism algebra. A Lie-algebra element need not act invertibly. Item1973-310.2.\par
\subsection*{P000521 (both; faithful error): minimal-height submodule}
\noindent\textit{Transcription reads.} \(E_0=\{x:p(x)=\bar p(E)\}\)\par
\noindent\textit{Critical reading.} \(E_0=\{0\}\cup\{x\ne0:p(x)=\bar p(E)\}\)\par
\noindent\textit{Evidence.} The source omits adjoining zero; Rudakov1974 Corollary3.3 explicitly includes it. Item1973-310.4.\par
\subsection*{P000521 (both; faithful error): filtration degree}
\noindent\textit{Transcription reads.} \(\deg a_i\ge j+1\)\par
\noindent\textit{Critical reading.} order of vanishing of \(a_i\) at least \(j+1\)\par
\noindent\textit{Evidence.} For formal series this means membership in \(\mathfrak m^{j+1}\), not polynomial degree. Rudakov1974 section1. Item1973-310.1.\par
\subsection*{P000521 (both; ambiguity): retraced induction bound}
\noindent\textit{Transcription reads.} \(\bar p(E)>1\), provisional digit\par
\noindent\textit{Critical reading.} retain1, now corroborated by the published theorem\par
\noindent\textit{Evidence.} Scan uncertainty is not erased; Rudakov1974 Theorem6.1 and Corollary6.3 explicitly use greater than1.\par
\subsection*{P000523 (both; faithful error): old plus-definition degree}
\noindent\textit{Transcription reads.} \(K_i(A)=\pi_{i+1}(B_+GL(A))\)\par
\noindent\textit{Critical reading.} \(K_i(A)=\pi_i(BGL(A)^+)\) for i at least1\par
\noindent\textit{Evidence.} Scan and transcription annotation confirm i+1 was written. The shift belongs to BQ, not BGL-plus; Weibel DefinitionIV1.1. Item1973-311.9.\par
\subsection*{P000523 (both; faithful error): retracted plus/Q equality}
\noindent\textit{Transcription reads.} \(BQPA=B_+GL(A)\), followed by this is false\par
\noindent\textit{Critical reading.} \(\Omega BQP(A)\simeq K_0(A)\times BGL(A)^+\)\par
\noindent\textit{Evidence.} The retraction is preserved. Quillen introduction and Weibel CorollaryIV7.2 give the loop and K0 factor. Item1973-311.10.\par
\subsection*{P000524 (both; faithful error): stabilization map and comparison}
\noindent\textit{Transcription reads.} \(BGL(A)\to BQPA\) from finite-rank automorphisms\par
\noindent\textit{Critical reading.} valid unlooped constructions should not be called the plus/Q comparison\par
\noindent\textit{Evidence.} The finite-rank maps are nullhomotopic, and adding A is homotopic to the identity by explicit natural transformations; see critical\_verifications.tex. The comparison targets the loop space. Items1973-311.11–311.12.\par
\subsection*{P000526 (both; faithful error): Cartier example: logarithm series}
\noindent\textit{Transcription reads.} \(\ln(1-a)=a+a^2/2+\cdots\)\par
\noindent\textit{Critical reading.} \(\ln(1-a)=-a-a^2/2-\cdots\)\par
\noindent\textit{Evidence.} Russian text and scan both have positive signs. This is a faithful mathematical error; the specialization at a=2 still sums to zero.\par
\subsection*{P000526 (both; faithful error): measure dual codomain}
\noindent\textit{Transcription reads.} \(\operatorname{Hom}(C(\mathbb Z_p,K),\mathbb Q_p)\)\par
\noindent\textit{Critical reading.} \(\operatorname{Hom}_{K,\mathrm{cont}}(C(\mathbb Z_p,K),K)\)\par
\noindent\textit{Evidence.} The scan confirms the p-adic base field; the subsequent K-algebra and K-valued moments require the continuous K-dual. This is an emendation, not a misread letter.\par
\subsection*{P000526 (both; faithful error): Theorem 2: unit group}
\noindent\textit{Transcription reads.} \((p-1)K[[T]]\)\par
\noindent\textit{Critical reading.} \(\prod_{\chi\in\widehat\mu_{p-1}}K[[T]]_{\mathrm b}\) for odd p; two components for p=2\par
\noindent\textit{Evidence.} The scan has no odd-prime hypothesis. The structure of the p-adic unit group supplies the missing restriction; the bounded-coefficient convention is retained.\par
\subsection*{P000527 (both; faithful error): first paragraph: expansion name}
\noindent\textit{Transcription reads.} Maclaurin series\par
\noindent\textit{Critical reading.} Newton--Mahler interpolation series\par
\noindent\textit{Evidence.} The displayed falling factorials and finite differences, confirmed in the scan, determine the meaning. The source itself uses the misleading name.\par
\subsection*{P000527 (both; faithful error): comparison after Q-definition}
\noindent\textit{Transcription reads.} \(BQM\) homologically equivalent to \(B\operatorname{GL}(\infty,A)\)\par
\noindent\textit{Critical reading.} \(\Omega_0 BQ\mathcal P(A)\simeq B\operatorname{GL}(A)^+\); the plus map is a homology equivalence\par
\noindent\textit{Evidence.} The scan confirms the missing loop symbol. Fundamental groups already refute the literal statement over the field with two elements. The projective-module category and base component must be specified.\par
\subsection*{P000528 (both; faithful error): Heisenberg bundle polynomial label}
\noindent\textit{Transcription reads.} polynomials \(\le2\) in p and q\par
\noindent\textit{Critical reading.} polynomials \(\le1\) for the stated Heisenberg bundle\par
\noindent\textit{Evidence.} The scan and Russian text clearly give degree 2. TM plus one scalar has dimension 2n+1; degree-two polynomials additionally contain the symplectic Lie algebra.\par
\subsection*{P000528 (both; ambiguity): metaplectic notation}
\noindent\textit{Transcription reads.} \(M(n,\mathbb R)\) and \(M(2n,\mathbb R)\)\par
\noindent\textit{Critical reading.} Use a single convention \(\operatorname{Mp}(2n,\mathbb R)\)\par
\noindent\textit{Evidence.} Both notations occur visibly on the page and refer to the same double cover of Sp(2n,R). This is notation normalization, not a second covering group.\par
\subsection*{P000528 (both; faithful error): nonlinear chain equation}
\noindent\textit{Transcription reads.} \(+\varepsilon(x_{k-1}-x_k)^2+\varepsilon(x_{k+1}-x_k)^2\)\par
\noindent\textit{Critical reading.} For the standard alpha-FPU Hamiltonian: \(\varepsilon[(x_{k+1}-x_k)^2-(x_k-x_{k-1})^2]\)\par
\noindent\textit{Evidence.} Both plus signs are confirmed by Russian text and scan. Differentiating the nearest-neighbour cubic potential gives a difference. This correction is conditional on the stated FPU identification.\par
\subsection*{P000529 (both; faithful error): linear wave claim}
\noindent\textit{Transcription reads.} all trajectories are closed\par
\noindent\textit{Critical reading.} all trajectories with zero mean initial velocity are periodic\par
\noindent\textit{Evidence.} The scan makes the universal claim. The spatially constant solution u=t is a direct counterexample; Fourier modes give the corrected condition.\par
\subsection*{P000529 (both; faithful error): matrix of F tensor identity}
\noindent\textit{Transcription reads.} \(\begin{pmatrix}F_1&F_2\\F_3&F_3\end{pmatrix}\)\par
\noindent\textit{Critical reading.} \(\begin{pmatrix}F_1&F_2\\F_3&F_4\end{pmatrix}\)\par
\noindent\textit{Evidence.} Russian text and scan repeat the lower-left block index, while the next line and P000530 use a fourth index for the second Fredholm diagonal block. This is a faithful index error.\par
\subsection*{P000531 (both; faithful error): universal submodule conjecture}
\noindent\textit{Transcription reads.} every modular lattice is a submodule lattice\par
\noindent\textit{Critical reading.} The conjecture is false; the Arguesian identity is an additional necessary condition\par
\noindent\textit{Evidence.} The conjectural wording is clear in Russian and the scan. This is a mathematical annotation, not a transcription change; the obstruction was known before 1973.\par
\subsection*{P000531-P000532 (both; ambiguity): recursive lattice words}
\noindent\textit{Transcription reads.} overwritten operation and incomplete Gamma conditions\par
\noindent\textit{Critical reading.} Published reconstruction: meet with the join over adjacent-distinct \(k_j\), each avoiding \(i_j,i_{j+1}\)\par
\noindent\textit{Evidence.} Scan cannot resolve the overwritten operation. Gelfand--Ponomarev 1974, §1.3 equation (1.1), supplies the matching complete construction; reconstruction must remain separate from the diplomatic text.\par
\subsection*{P000532 (both; ambiguity): maximal compact in projective group}
\noindent\textit{Transcription reads.} \(G=\operatorname{GL}_2(\mathbb Q_p)/C,\ K=\operatorname{GL}_2(\mathbb Z_p)\)\par
\noindent\textit{Critical reading.} \(K\) is the image of \(\operatorname{GL}_2(\mathbb Z_p)\) in \(G\)\par
\noindent\textit{Evidence.} The manuscript suppresses passage to the central quotient. The unquotiented matrix group is not literally a subgroup of the projective group.\par
\subsection*{P000532 (both; faithful error): Weil point bound}
\noindent\textit{Transcription reads.} \(|[X_q]-q|\le2g\sqrt q\)\par
\noindent\textit{Critical reading.} \(|\#X(\mathbb F_q)-(q+1)|\le2g\sqrt q\) for smooth projective geometrically connected X\par
\noindent\textit{Evidence.} The scan has q, not q+1. The genus-zero projective line gives the immediate counterexample; the affine/projective convention is otherwise unstated.\par
\subsection*{P000532 (both; ambiguity): Riemann--Roch notation}
\noindent\textit{Transcription reads.} L(D) is the number of functions\par
\noindent\textit{Critical reading.} \(\ell(D)=\dim L(D)\); over \(\mathbb F_q\) the cardinality is \(q^{\ell(D)}\)\par
\noindent\textit{Evidence.} The Russian says number of functions, while the next page explicitly uses dimensions. The correct interpretation is vector-space dimension.\par
\subsection*{P000533 (both; faithful error): kernel multiplication diagram}
\noindent\textit{Transcription reads.} \(L_{q-1}\times L_e^q\to L_{q-1}L_e\subset L_{q+e-1}\)\par
\noindent\textit{Critical reading.} Separate ordinary and Frobenius-twisted tensor multiplication; the latter lands in \(L_{q-1+qe}\)\par
\noindent\textit{Evidence.} Scan and Russian confirm both the superscript q and the untwisted pole bound. Raising to q multiplies pole orders by q; rational-point evaluation explains the connection of the two maps.\par
\subsection*{P000534 (both; faithful error): stable-rank definition: reduced column}
\noindent\textit{Transcription reads.} \((a+\alpha c,b+\alpha c)^t\)\par
\noindent\textit{Critical reading.} \((a+\alpha c,b+\beta c)^t\)\par
\noindent\textit{Evidence.} Russian and scan clearly repeat alpha; two independent coefficients were quantified and the next-page real example uses beta in the second component.\par
\subsection*{P000534 (both; ambiguity): unimodular parenthetical gloss}
\noindent\textit{Transcription reads.} relatively prime elements\par
\noindent\textit{Critical reading.} elements generating the unit ideal\par
\noindent\textit{Evidence.} In several variables gcd one is insufficient: x and y have no common factor but generate a proper ideal. The stable-rank definition requires a Bezout relation.\par
\subsection*{P000534 (both; faithful error): field dichotomy following stable-rank theorem}
\noindent\textit{Transcription reads.} \(K_2k\ne0\) unless k is algebraic over a finite field\par
\noindent\textit{Critical reading.} The converse fails; \(K_2(\overline{\mathbb Q})=0\)\par
\noindent\textit{Evidence.} Russian and scan faithfully record the claim. The same error is printed in the 1974 announcement; the 1976 full paper distinguishes additional algebraic global-field cases. Weibel III Exercise 6.5 gives the concrete counterexample.\par
\subsection*{P000534 (both; ambiguity): zero-K2 branch of stable-rank assertion}
\noindent\textit{Transcription reads.} \(K_2k=0\Rightarrow\operatorname{sr}k[x_1,x_2]=2\)\par
\noindent\textit{Critical reading.} Verified here for k algebraic over a finite field; do not infer the unrestricted branch from the false field dichotomy\par
\noindent\textit{Evidence.} The 1974 announcement repeats the notebook argument, but Vaserstein--Suslin 1976 section 19 explicitly leaves extra field cases unresolved at that time. Present unrestricted status was not established by this bounded search.\par
\subsection*{P000535 (both; ambiguity): conic question: common divisor and polynomial parenthesis}
\noindent\textit{Transcription reads.} no common divisor; impossible in polynomials, possible for continuous functions\par
\noindent\textit{Critical reading.} Specify coefficients and require a unit ideal/no simultaneous zeros; unrestricted complex polynomials already suffice\par
\noindent\textit{Evidence.} Both phrases are legible in the scan. The displayed complex polynomial Bezout certificate refutes the parenthesis under that interpretation. The preceding real example and the real degree obstruction prevent an unconditional attribution of a false complex theorem to the speaker.\par
\subsection*{P000536 (both; ambiguity): equivalence for complexes of free modules}
\noindent\textit{Transcription reads.} quasi-isomorphic if and only if homotopy equivalent\par
\noindent\textit{Critical reading.} For cochain complexes of projectives bounded above, a quasi-isomorphism is a homotopy equivalence.\par
\noindent\textit{Evidence.} The scan omits the boundedness qualifier. Surrounding ordinary-resolution diagrams support this intended scope. The doubly infinite epsilon complex over the ring of dual numbers is an explicit counterexample to the unrestricted version.\par
\subsection*{P000536 (both; ambiguity): derived category morphisms}
\noindent\textit{Transcription reads.} strings of maps and inverse quasi-isomorphisms\par
\noindent\textit{Critical reading.} equivalence classes of such strings under the localization relations\par
\noindent\textit{Evidence.} The manuscript gives only representatives. Without the quotient relation the construction is not the derived category.\par
\subsection*{P000536 (both; faithful error): contravariant functor and derived-functor diagrams}
\noindent\textit{Transcription reads.} contravariant F, but \(F:M(A)\to M(B)\) and \(\dot F:D(A)\to D(B)\)\par
\noindent\textit{Critical reading.} use the opposite source category and reverse complex degrees\par
\noindent\textit{Evidence.} Russian and scan explicitly say contravariant while drawing covariant arrows; the Hom example confirms the missing variance convention.\par
\subsection*{P000536 (both; ambiguity): last natural-transformation arrow}
\noindent\textit{Transcription reads.} \(\operatorname{Nat}(QF,GQ)\to\operatorname{Nat}(\dot FQ,GQ)\)\par
\noindent\textit{Critical reading.} the arrow can mean unique factorization; direct composition with QF to dotFQ induces the reverse direction\par
\noindent\textit{Evidence.} The rightward arrow is legible. It is not silently reversed: the two directions are inverse only after the stated universal property is imposed. Weibel Definition 10.5.1 checks the canonical composition direction.\par
\subsection*{P000537 (both; ambiguity): algebraic-group notation}
\noindent\textit{Transcription reads.} \(GL(\bar k)\)\par
\noindent\textit{Critical reading.} \(GL_n(\bar k)\) if it denotes the algebraic group corresponding to the preceding \(GL_n(k)\).\par
\noindent\textit{Evidence.} The Russian and image omit the rank; the common rank is an editorial completion, not recovered ink.\par
\subsection*{P000538 (both; ambiguity): Gelfand spectral equation}
\noindent\textit{Transcription reads.} \(y\prime\prime+uy=\lambda y\)\par
\noindent\textit{Critical reading.} Keep the plus sign diplomatically; use \(-y\prime\prime+uy=\lambda y\) only when adopting Novikov\'s \(L=-D^2+u\) convention.\par
\noindent\textit{Evidence.} Russian and scan agree on the plus sign. This is a convention conflict inside the notes.\par
\subsection*{P000539 (both; faithful error): Floquet reconstruction lemma}
\noindent\textit{Transcription reads.} \(f(x)=\int_0^{2\pi}f_t(x)dx\); \(f_t(x+2\pi)=e^{it}f(x)\); mixed period/phase Parseval bounds\par
\noindent\textit{Critical reading.} With phase \(\theta\), period \(T\): \(f=(2\pi)^{-1}\int_0^{2\pi}Uf\,d\theta\), \(Uf(\theta,x+T)=e^{i\theta}Uf(\theta,x)\), and normalized Parseval on \([0,2\pi]\times[0,T]\).\par
\noindent\textit{Evidence.} All defects are visible in the scan and Russian. Fourier-series inversion fixes the integration variable, subscript, and normalization.\par
\subsection*{P000539 (both; faithful error): Bloch coefficient integral}
\noindent\textit{Transcription reads.} \(\int f_t y_{t,n}\,dx\)\par
\noindent\textit{Critical reading.} \(\int f_t\overline{y_{t,n}}\,dx\) for an orthonormal complex fiber eigenbasis.\par
\noindent\textit{Evidence.} No conjugate is present in either transcription or image; complex quasiperiodic eigenfunctions require the Hermitian inner product.\par
\subsection*{P000539 (both; faithful error): multipliers outside a band}
\noindent\textit{Transcription reads.} \(e^{\pm t}\)\par
\noindent\textit{Critical reading.} \(e^{\pm\kappa}\) for positive trace, \(-e^{\pm\kappa}\) for trace below \(-2\).\par
\noindent\textit{Evidence.} The characteristic equation shown immediately above admits negative reciprocal multipliers too.\par
\subsection*{P000540 (both; faithful error): constant potential spectrum}
\noindent\textit{Transcription reads.} constant \(u\): spectrum \(>0\)\par
\noindent\textit{Critical reading.} For \(L=-D^2+c\), spectrum \([c,\infty)\).\par
\noindent\textit{Evidence.} The scan faithfully has >0; Fourier transformation proves the corrected threshold.\par
\subsection*{P000540 (both; faithful error): free monodromy}
\noindent\textit{Transcription reads.} \(e^{\pm iT}\), \(a_R=\cos T\)\par
\noindent\textit{Critical reading.} At \(u=0,\lambda=k^2\): \(e^{\pm ikT}\), \(a_R=\cos(kT)\).\par
\noindent\textit{Evidence.} The spectral factor is absent in the scan; direct substitution of plane waves supplies it.\par
\subsection*{P000541 (both; faithful error): eigenvalues of traceless matrix}
\noindent\textit{Transcription reads.} \(\alpha_\pm=\pm\sqrt{\det\Lambda}\), a polynomial\par
\noindent\textit{Critical reading.} \(\alpha_\pm=\pm\sqrt{-\det\Lambda}\); the determinant/spectral polynomial is polynomial, generally not its square root.\par
\noindent\textit{Evidence.} The radicand has no minus in the scan; the characteristic polynomial is \(\mu^2+\det\Lambda\).\par
\subsection*{P000542 (both; provenance): date in heading}
\noindent\textit{Transcription reads.} 4/III-73\par
\noindent\textit{Critical reading.} 4 March 1974.\par
\noindent\textit{Evidence.} The image and Russian faithfully have 73; the official catalogue and neighboring 25 February/11 March 1974 entries fix 1974.\par
\subsection*{P000542 (both; faithful error): unsigned kernel bound}
\noindent\textit{Transcription reads.} \(\int F(x,y)\,dy<C\)\par
\noindent\textit{Critical reading.} Use absolute-integrability/operator bounds and an invertibility hypothesis.\par
\noindent\textit{Evidence.} The 1974 primary paper states absolute bounds; the unsigned notebook bound alone is inadequate.\par
\subsection*{P000542 (both; ambiguity): kernel commutator signs}
\noindent\textit{Transcription reads.} Eq. (4): obliterated sign; Eq. (5): plus; bottom: \([\hat K,D^2]\)\par
\noindent\textit{Critical reading.} Under \([A,B]=AB-BA\), \([D^n,F]\) has kernel \(F_{x^n}-(-1)^nF_{y^n}\); the second-order difference is \([D^2,\hat K]\).\par
\noindent\textit{Evidence.} Russian and scans confirm the printed signs/order. Integration by parts determines the consistent convention; the obscured sign is not deciphered.\par
\subsection*{P000544 (both; faithful error): KP name}
\noindent\textit{Transcription reads.} Kadomtsev--Petriashvili (?)\par
\noindent\textit{Critical reading.} Kadomtsev--Petviashvili.\par
\noindent\textit{Evidence.} The handwritten question mark belongs to the source. The identified KP equation and primary joint paper establish the standard surname.\par
\subsection*{P000544 (both; faithful error): KP variables and coefficients}
\noindent\textit{Transcription reads.} \(w_y\) in second equation; coefficients in the printed scalar KP equation do not agree with the displayed pair.\par
\noindent\textit{Critical reading.} For \(\beta\partial_z+P\), \(\alpha\partial_t+Q\): \(w_x=-3\beta u_z/4\), \(\beta w_z=\alpha u_t+u_{xxx}/4+3uu_x/2+\lambda u_x\).\par
\noindent\textit{Evidence.} The source clearly has y. Direct commutator calculation, independently rechecked in code, yields the consistent system; no single-letter correction alone fixes all formulas.\par
\subsection*{P000546 (both; faithful error): scaled unknown}
\noindent\textit{Transcription reads.} \(\varepsilon^2(v_{tt}-v_{xx})+T'(u)=0\)\par
\noindent\textit{Critical reading.} \(\varepsilon^2(v_{SS}-v_{XX})+T'(v)=0\) under the indicated slow-coordinate scaling.\par
\noindent\textit{Evidence.} The scan has v in the derivatives and u in the potential. The chain rule requires a single rescaled dependent variable.\par
\subsection*{P000546 (both; faithful error): limiting fast-phase ODE}
\noindent\textit{Transcription reads.} \(\varepsilon^2(\omega_2^2-\omega_1^2)u_{\Theta\Theta}+T'(u)=0\) after \(\varepsilon\to0\).\par
\noindent\textit{Critical reading.} \((\omega_2^2-\omega_1^2)f_{\Theta\Theta}+T'(f)=0\) for \(\Theta=\Theta_0/\varepsilon\).\par
\noindent\textit{Evidence.} The factor epsilon squared survives in the image; two fast derivatives cancel it. The overwritten extra mark after the potential derivative is left undeciphered.\par
\subsection*{P000546 (both; faithful error): Whitham reference}
\noindent\textit{Transcription reads.} J. Fluid Mechanics 1972, Two-timing ...\par
\noindent\textit{Critical reading.} Whitham (1970), Two-timing, variational principles and waves, JFM 44(2), 373--395.\par
\noindent\textit{Evidence.} Primary Caltech repository record and DOI 10.1017/S002211207000188X confirm 1970.\par
\subsection*{P000547 (both; faithful error): sign of primed kernel relation}
\noindent\textit{Transcription reads.} Input-first kernel coordinates \((x,y;\xi,\eta)\), but \(WF'(P)\) reverses \(\eta\).\par
\noindent\textit{Critical reading.} For action on input \(WF(u)\), reverse input \(\xi\); or switch to output-first coordinates and reverse the second covector.\par
\noindent\textit{Evidence.} The source explicitly labels input x and output y. Hörmander Theorem 2.5.14 supplies the convention. Reversing the output instead is not equivalent for arbitrary complex kernels.\par
\subsection*{P000548 (both; ambiguity): product continuity}
\noindent\textit{Transcription reads.} a continuous mapping\par
\noindent\textit{Critical reading.} sequentially continuous in Hörmander's original pseudotopology; normal-topology multiplication is hypocontinuous.\par
\noindent\textit{Evidence.} Hörmander 1971 Theorem 2.5.10 explicitly says sequentially continuous; Brouder--Dang--Hélein 2016 explains the topological distinction.\par
\subsection*{P000549 (both; faithful error): nonlinear KdV coefficient}
\noindent\textit{Transcription reads.} \(4f_t+f_{xxx}+ff_x=0\)\par
\noindent\textit{Critical reading.} The displayed \(2\lambda^2\operatorname{sech}^2[\lambda(x-\lambda^2t-h)]\) solves the equation with \(6ff_x\).\par
\noindent\textit{Evidence.} The scan explicitly has no coefficient before \(ff_x\). Symbolic substitution verifies the factor 6; the error is faithful in both transcriptions.\par
\subsection*{P000549 (both; ambiguity): collision-amplitude summation bound}
\noindent\textit{Transcription reads.} Upper index read provisionally as \(n\), requiring \(\lambda_{n+1}\).\par
\noindent\textit{Critical reading.} Retain the uncertain index; no uniquely justified corrected inequality was located.\par
\noindent\textit{Evidence.} The small upper sign is entangled with the previous line. The parameters and ordering are not supplied.\par
\subsection*{P000550 (both; ambiguity): interpolation node count}
\noindent\textit{Transcription reads.} \(2n\,?\,1\) and infinity.\par
\noindent\textit{Critical reading.} Keep the sign uncertain. A modern theorem with \(\deg P=\deg Q=n\) uses \(2n+1\) finite nodes, subject to additional normalization and nondegeneracy.\par
\noindent\textit{Evidence.} The ink overwrites the sign; Shabat 2014 Theorem 1 supplies a mathematical comparator, not decisive paleographic recovery.\par
\subsection*{P000550 (both; faithful error): Wronskian and exponent assignment}
\noindent\textit{Transcription reads.} \(y_1=Pe^{+i\zeta x},y_2=Qe^{-i\zeta x}\), but \(P_xQ-Q_xP-2i\zeta PQ=\prod(\zeta-\zeta_j)\).\par
\noindent\textit{Critical reading.} The written minus-Wronskian corresponds to \(y_1=Pe^{-i\zeta x},y_2=Qe^{+i\zeta x}\); restore the leading factor \(-2ip_0q_0\).\par
\noindent\textit{Evidence.} Both signs and the omitted prefactor are visible in the scan. Differentiation and leading-degree comparison give the emendation.\par
\subsection*{P000550 (both; faithful error): Riccati and modified-KdV signs}
\noindent\textit{Transcription reads.} \(f_x+f^2=q\), with the preceding \(y_{xx}+(q+\zeta^2)y=0\); cubic coefficient 3 in the f-equation.\par
\noindent\textit{Critical reading.} For \(f=y_x/y\), \(f_x+f^2=-q-\zeta^2\). Use a separately stated consistent Miura normalization rather than changing a single sign.\par
\noindent\textit{Evidence.} The mutually inconsistent equations are faithfully transcribed; logarithmic differentiation and the independently checked Miura identity diagnose the conflict.\par
\subsection*{P000551 (both; provenance): archive roster records 331 and 332}
\noindent\textit{Transcription reads.} Both official MCCME catalogue entries are dated 25 March 1975; the supplied academic-year roster omits them.\par
\noindent\textit{Critical reading.} 25 March 1974, as written on the common source folio; include separate supplementary records 331 and 332.\par
\noindent\textit{Evidence.} Official all-links identities plus the unmistakable 25/III-74 heading and physical sequence 18 March/1 April 1974. This is an archive metadata error, not an error in either source transcription.\par
\subsection*{P000551 (both; faithful error): bundle automorphism base and fiber}
\noindent\textit{Transcription reads.} Projection onto \(\operatorname{Diff}M\) with fiber \(C^\infty(M,H)\) for a principal bundle.\par
\noindent\textit{Critical reading.} For a general bundle, use \(\operatorname{Diff}(M)_{[E]}\) and \(\Gamma(E\times_{\operatorname{Ad}}H)\); the written groups are justified after trivializing a trivial bundle.\par
\noindent\textit{Evidence.} The source does not state triviality; the exact group-level kernel and image establish the necessary restriction.\par
\subsection*{P000551 (both; ambiguity): Fock factorial convention}
\noindent\textit{Transcription reads.} Components \(1,\xi,\xi^{\otimes2}/2!,\ldots\), norm weighted by \(1/n!\).\par
\noindent\textit{Critical reading.} Specify one consistent tensor-component normalization; standard Hilbert summands use \(\xi^{\otimes n}/\sqrt{n!}\).\par
\noindent\textit{Evidence.} The written rows and weights are clear but their coordinate conventions are not. No unique paleographic repair is justified.\par
\subsection*{P000552 (both; ambiguity): Laplacian sign}
\noindent\textit{Transcription reads.} \(\Delta=+\operatorname{div}\operatorname{grad}\), nonnegative eigenvalues and \(e^{-t\Delta}\).\par
\noindent\textit{Critical reading.} Use \(L=-\operatorname{div}\operatorname{grad}\ge0\) with \(e^{-tL}\), or retain the negative operator and use \(e^{t\Delta}\).\par
\noindent\textit{Evidence.} The displayed surviving sign is plus; the immediately preceding overwritten block is not deciphered. The three visible conventions cannot hold simultaneously.\par
\subsection*{P000552 (both; faithful error): off-diagonal heat expansion}
\noindent\textit{Transcription reads.} \((4\pi t)^{-n/2}\sum u_rt^r\) in a neighborhood of the diagonal.\par
\noindent\textit{Critical reading.} Multiply by \(e^{-d_g(x,y)^2/(4t)}\) in a convex normal neighborhood; no factor is needed after setting \(x=y\).\par
\noindent\textit{Evidence.} The Gaussian is absent from both transcriptions and image. The classical local heat parametrix supplies it.\par
\subsection*{P000552 (both; ambiguity): curvature indices}
\noindent\textit{Transcription reads.} Scalar and Ricci index contractions contain overwritten or struck-out characters.\par
\noindent\textit{Critical reading.} Keep the uncertain ink unresolved; state invariant contractions separately.\par
\noindent\textit{Evidence.} The scan confirms the overwritten indices; standard tensor formulas do not uniquely recover the handwriting.\par
\subsection*{P000553 (both; faithful error): mixed heat and wave trace}
\noindent\textit{Transcription reads.} \(\sum e^{i\lambda_kz}=\operatorname{Sp}(e^{-z\Delta})\).\par
\noindent\textit{Critical reading.} Heat: \(\sum e^{-z\mu_k}=\operatorname{Tr}e^{-zL}\). Half-wave: \(\sum e^{iz\sqrt{\mu_k}}=\operatorname{Tr}e^{iz\sqrt L}\).\par
\noindent\textit{Evidence.} Both exponents are unmistakable in the scan. The final wave equation confirms the square-root frequency convention for that discussion.\par
\subsection*{P000553 (both; ambiguity): spectral parameter in the torus gap theorem}
\noindent\textit{Transcription reads.} The same \(\lambda_k\) as the earlier Laplacian spectrum, without explicit redefinition.\par
\noindent\textit{Critical reading.} Shnirelman (1975) uses eigenvalues of \(\sqrt{-\Delta}\); use \(\omega_k\) to separate them from Laplacian eigenvalues \(\mu_k\).\par
\noindent\textit{Evidence.} The exact primary theorem, received 7 May 1974, explicitly defines the positive square-root operator. Superpolynomial closeness survives squaring after exponent adjustment.\par
\subsection*{P000553 (both; faithful error): global location of kernel singularities}
\noindent\textit{Transcription reads.} Singular where \(|x-y|=t\).\par
\noindent\textit{Critical reading.} Locally \(d_g(x,y)=|t|\) before cut/conjugate points; globally use covectors connected by geodesic flow for time \(t\).\par
\noindent\textit{Evidence.} Unit-speed propagation is correct but a metric-distance sphere omits longer geodesic connections and caustics.\par
